% main.tex -- generated by build_monograph.py; edit frontmatter/*.tex, not chapters/*.tex
\documentclass[11pt,twoside,openright]{book}
\usepackage{lmodern}
\usepackage[T1]{fontenc}
\usepackage[utf8]{inputenc}
\usepackage{amsmath,amssymb,amsthm,mathtools}
\usepackage[margin=1.05in]{geometry}
\usepackage{booktabs,array,longtable}
\usepackage{graphicx}
\usepackage{microtype}
\usepackage{url}
\usepackage[hidelinks]{hyperref}
\usepackage{fancyhdr}
\raggedbottom
\emergencystretch=2em
% page numbers centred at the foot of every page (front matter, chapter openings, parts and
% text pages alike); running heads without a rule
\setlength{\headheight}{14pt}
\pagestyle{fancy}
\fancyhf{}
\fancyhead[LE]{\small\slshape\nouppercase{\leftmark}}
\fancyhead[RO]{\small\slshape\nouppercase{\rightmark}}
\fancyfoot[C]{\thepage}
\renewcommand{\headrulewidth}{0pt}
\fancypagestyle{plain}{\fancyhf{}\fancyfoot[C]{\thepage}\renewcommand{\headrulewidth}{0pt}}
\setcounter{tocdepth}{1}
\setcounter{secnumdepth}{2}

\numberwithin{equation}{section}
\theoremstyle{plain}
\newtheorem{theorem}{Theorem}[section]
\newtheorem{proposition}[theorem]{Proposition}
\newtheorem{lemma}[theorem]{Lemma}
\newtheorem{corollary}[theorem]{Corollary}
\newtheorem{conjecture}[theorem]{Conjecture}
\newtheorem{hypothesis}[theorem]{Hypothesis}
\newtheorem{problem}[theorem]{Problem}
\theoremstyle{definition}
\newtheorem{definition}[theorem]{Definition}
\newtheorem{algorithm}[theorem]{Algorithm}
\newtheorem{construction}[theorem]{Construction}
\newtheorem{convention}[theorem]{Convention}
\newtheorem{example}[theorem]{Example}
\theoremstyle{remark}
\newtheorem{remark}[theorem]{Remark}

\newenvironment{chapterabstract}{\begin{quotation}\small\noindent\textbf{Summary.}\ }{\end{quotation}\medskip}
% compatibility layer for the revtex-style applied note
\usepackage{bm}
\providecommand{\onlinecite}[1]{\cite{#1}}
\newenvironment{ruledtabular}{}{}
\providecommand{\colrule}{\midrule}
\newenvironment{acknowledgments}{\section*{Acknowledgments}}{}
% union of the papers' macros; first definition wins, chapters override locally
\newcommand{\Z}{\mathbb{Z}}
\newcommand{\Q}{\mathbb{Q}}
\newcommand{\R}{\mathbb{R}}
\newcommand{\C}{\mathbb{C}}
\newcommand{\N}{\mathbb{N}}
\newcommand{\F}{\mathbb{F}}
\newcommand{\A}{\mathbb{A}}
\newcommand{\PPone}{\mathbb{P}^1}
\newcommand{\cA}{\mathcal{A}}
\newcommand{\cE}{\mathcal{E}}
\newcommand{\cG}{\mathcal{G}}
\newcommand{\cH}{\mathcal{H}}
\newcommand{\cK}{\mathcal{K}}
\newcommand{\cL}{\mathcal{L}}
\newcommand{\cO}{\mathcal{O}}
\newcommand{\cR}{\mathcal{R}}
\newcommand{\cT}{\mathcal{T}}
\newcommand{\pp}{\mathfrak{p}}
\newcommand{\Ns}{\mathrm{N}}
\newcommand{\Frob}{\operatorname{Frob}}
\newcommand{\Gal}{\operatorname{Gal}}
\newcommand{\Sat}{\operatorname{Sat}}
\newcommand{\Prob}{\operatorname{Prob}}
\newcommand{\coker}{\operatorname{coker}}
\newcommand{\Tors}{\operatorname{Tors}}
\newcommand{\SNF}{\operatorname{SNF}}
\newcommand{\Jac}{\operatorname{Jac}}
\newcommand{\rank}{\operatorname{rank}}
\newcommand{\Spec}{\operatorname{Spec}}
\newcommand{\Ind}{\operatorname{Ind}}
\newcommand{\Irr}{\operatorname{Irr}}
\newcommand{\HS}{\mathrm{HS}}
\newcommand{\KMS}{\mathrm{KMS}}
\newcommand{\id}{\mathrm{id}}
\newcommand{\Aut}{\operatorname{Aut}}
\newcommand{\GL}{\mathrm{GL}}
\newcommand{\PGL}{\mathrm{PGL}}
\newcommand{\Sym}{\operatorname{Sym}}
\newcommand{\Tr}{\operatorname{Tr}}
\newcommand{\ind}{\operatorname{ind}}
\newcommand{\shadow}{Y^{\sharp}}
\newcommand{\cI}{\mathcal{I}}
\newcommand{\ord}{\operatorname{ord}}
\newcommand{\supp}{\operatorname{supp}}
\newcommand{\Zl}{\Z_\ell}
\newcommand{\La}{\Lambda}
\newcommand{\Tor}{\operatorname{Tor}}
\newcommand{\car}{\operatorname{char}_{\La}}
\newcommand{\Fitt}{\operatorname{Fitt}}
\newcommand{\Ql}{\Q_\ell}
\newcommand{\Hom}{\operatorname{Hom}}
\newcommand{\rk}{\operatorname{rk}}
\newcommand{\Res}{\operatorname{Res}}
\newcommand{\tr}{\operatorname{tr}}
\newcommand{\adj}{\operatorname{adj}}
\newcommand{\Reff}{R_{\mathrm{eff}}}
\newcommand{\Anb}{A^{\mathrm{nb}}}
\newcommand{\Dp}{\mathsf{D}_p}
\newcommand{\Alb}{\operatorname{Alb}}
\newcommand{\Trs}{\operatorname{Tr}_s}
\newcommand{\Var}{\operatorname{Var}}
\newcommand{\HC}{\mathit{HC}}
\newcommand{\HH}{\mathit{HH}}
\newcommand{\Zp}{\Z_p}
\newcommand{\ta}{\tau_*}
\newcommand{\taup}{\tau^*}
\newcommand{\et}{\eta_*}
\newcommand{\etp}{\eta^*}
\newcommand{\Qp}{\Q_p}
\newcommand{\cM}{\mathcal{M}}
\newcommand{\pdet}{\operatorname{pdet}}
\newcommand{\ev}{\operatorname{ev}}
\newcommand{\End}{\operatorname{End}}
\newcommand{\XX}{\mathfrak{X}_\infty}
\newcommand{\Ext}{\operatorname{Ext}}
\newcommand{\Tup}{\cT^{\,k}}
\newcommand{\Tdn}{\cT_{k}}
\newcommand{\cX}{\mathcal{X}}
\newcommand{\KK}{\mathit{KK}}
\newcommand{\sissharp}{\sharp}
\newcommand{\Ysh}[1]{Y_{#1}^{\sharp}}
\newcommand{\Yt}{\widetilde Y_\alpha}
\newcommand{\PG}{\mathrm{PG}}
\newcommand{\Atwo}{\widetilde A_2}
\newcommand{\DeltaH}{\Delta^{\mathrm{H}}}
\newcommand{\JacH}{\Jac^{\mathrm{H}}}
\newcommand{\Cthree}{C^{(3)}}
\newcommand{\spec}{\operatorname{spec}}
\newcommand{\LE}{L_{E}}
\newcommand{\LF}{L_{F}}
\newcommand{\SE}{\sim_{\Z}}
\newcommand{\LB}{L_B}
\newcommand{\im}{\operatorname{im}}
\newcommand{\Ch}{\widehat C}
\newcommand{\Chp}{\widehat C{}'}
\newcommand{\Bh}{\widehat B}
\newcommand{\Cpp}{C''}
\newcommand{\Atilde}[1]{\widetilde A_{#1}}
\newcommand{\tp}{\mathsf{T}}
\newcommand{\T}{\mathbb{T}}
\newcommand{\PP}{\mathbb{P}}
\newcommand{\III}{\mathrm{III}}
\newcommand{\M}{\mathcal{M}}
\newcommand{\K}{\mathcal{K}}
\newcommand{\Loc}{\mathrm{Loc}}
\newcommand{\ec}{\mathrm{ec}}
\newcommand{\DB}{\mathrm{DB}}
\newcommand{\Arb}{\mathrm{Arb}}
\newcommand{\Rec}{\mathrm{Rec}}
\newcommand{\Sk}{\mathrm{Sk}}
\newcommand{\mult}{\operatorname{mult}}
\newcommand{\Sig}{\Sigma}


\title{\bfseries Algorithmic Non-commutative Class Field Theory\\[6pt]
\large Hecke operators, sandpile groups and operator algebras on trees and buildings,\\ with applications}
\author{Ruqing Chen\\[4pt]\normalsize GUT Geoservice Inc., Rivi\`ere-Beaudette, Qu\'ebec, Canada\\
\normalsize\texttt{ruqing@hotmail.com}}
\date{September 2026\\[10pt]\normalsize Sources, code, data and figures:\\
\url{https://github.com/Ruqing1963/anccft}}

\begin{document}
\frontmatter
\maketitle
\chapter*{Preface}
\addcontentsline{toc}{chapter}{Preface}

This book collects, in a single logical order, the twenty papers of the series
\emph{Algorithmic non-commutative class field theory} (September 2026) and the eleven
applied notes that grew out of them. The papers were written in the order in which the
questions arose; the chapters are arranged in the order in which the mathematics is
best learned. Each chapter is a paper, lightly edited: the ``task specification'' framing
of the originals has been left in place where it records a correction, because those
corrections are part of the record, and every chapter keeps its own verification
battery, its own ``Solved / open / impossible'' remark and its own provenance
paragraph. Nothing has been re-proved and nothing has been quietly upgraded.

Three conventions govern the whole book.

\emph{Every number is machine-verified.} Each chapter ends with a table of keyed
checks, each labelled \emph{exact} (integer or symbolic arithmetic), \emph{sampled}
(a statistical test with its $p$-value) or \emph{numeric} (floating point to a stated
tolerance). The scripts that produce the tables, the data they produce, the figures, the
\LaTeX\ sources of every chapter and of this book, and the archived outputs are in the
public repository
\begin{center}
\url{https://github.com/Ruqing1963/anccft}
\end{center}
(\texttt{code/}, \texttt{data/}, \texttt{figures/}, \texttt{papers/},
\texttt{applications/}, \texttt{monograph/}); a frozen copy of each release is archived
on Zenodo under the DOI given in the repository's \texttt{README}. The chapters cite
one another as chapters, never as ``preprints'': nothing in this corpus has been
published elsewhere, and the bibliography lists only the external literature.

\emph{Proof boundaries are stated.} Every result is one of: proved here; quoted from a
named source; verified computationally on named instances; or conjectural. The status
paragraphs at the head of each chapter say which. Where an earlier chapter's claim was
later found to be wrong or imprecise, the correction is recorded in the later chapter and
flagged in the closing chapter of this book.

\emph{The negative results are results.} A large part of the programme consists of
showing that the obvious next step does not exist: no non-negative matrix realizes the
Ihara companion; no commutative Iwasawa algebra governs the congruence tower; no
separable $C^*$-algebra has the pro-module as its $K_0$; no cyclic cocycle computes the
$L$-invariant; no pairwise invariant of a chord diagram computes the arrangement group
of a multi-copy repeat. These are stated as theorems with witnesses, and the closing
chapter lists what remains genuinely open.

The concordance that follows the table of contents maps paper numbers to chapters.
Inside a chapter, ``Theorem 3.7'' means Theorem $n$.3.7 of that chapter $n$; a reference
``[Ch.~1, Thm.~3.7]'' from another chapter points to the same statement.

\bigskip
\noindent Rivi\`ere-Beaudette, September 2026

\tableofcontents
\chapter*{Concordance}
\addcontentsline{toc}{chapter}{Concordance}
The chapters are the papers of the series, re-ordered by subject. A reference of the form ``[Ch.~$n$, Thm.~$a.b$]'' points to Theorem $n.a.b$ of this book; the papers' own section-wise numbering is preserved inside each chapter. References to the applied notes (Phys~1, Bio~1--5, \dots) are by the labels below.

\begin{longtable}{llrp{8.2cm}}
\toprule
paper & folder & chapter & title\\
\midrule
\endhead
I & \texttt{Paper 1} & 2 & the Langlands--thermodynamic bridge over Bruhat--Tits trees\\
II & \texttt{Paper 2} & 3 & the trace obstruction and the K-homological bulk\\
VI & \texttt{Paper 6} & 4 & the shadow algebra --- a functorial Kirchberg realization of the critical group\\
III & \texttt{Paper 3} & 5 & the Iwahori--Hecke tower and the type III dissolution of the trace obstruction\\
IV & \texttt{Paper 4} & 6 & the limit module of abelian towers and the characteristic identity\\
IX & \texttt{Paper 9} & 7 & exceptional zeros, the graph $\cL$-invariant, and the operator-algebraic
Mazur--Tate--Teitelbaum formula\\
XIV & \texttt{Paper 14} & 8 & the graph $\cL$-invariant as spectral action, Bloch effective mass and Albanese length, and what the cyclic Chern character sees\\
V & \texttt{Paper 5} & 9 & the degeneracy structure of the Iwahori tower and the Euler correction\\
VII & \texttt{Paper 7} & 10 & the snake factorization of the transfer and the evaluated characteristic form\\
XIII & \texttt{Paper 13} & 11 & universal localization of the degeneracy algebra, the $K_1$-characteristic class, and the
Eisenstein boundary\\
VIII & \texttt{Paper 8} & 12 & the tower-level Hecke--Kirchberg limit and the non-commutative Iwasawa $K$-module\\
XVI & \texttt{Paper 16} & 13 & the out-covering correspondence of the head map, its $K_0$-index $\et$, and the
$\KK$-category of the Iwahori tower\\
XVII & \texttt{Paper 17} & 14 & the tail norm at composite $q$, and the splitting of the limit $K$-module\\
XIX & \texttt{Paper 19} & 15 & the odd cyclic class that does not exist, and what the exceptional zero sees instead\\
X & \texttt{Paper 10} & 16 & $\Atwo$ complexes, higher simplicial Laplacians, and the fate of the trace obstruction in
rank two\\
XI & \texttt{Paper 11} & 17 & the $\PGL_3$ positive Hecke-dynamical system: shift equivalence over $\Z$ and the geodesic
edge flow\\
XII & \texttt{Paper 12} & 18 & $\Atwo$ towers --- no abelian sector, the covering norm theorem, and the two digraph towers\\
XV & \texttt{Paper 15} & 19 & the parahoric factorization of the geodesic edge flow, the unitarity of the remainder,
and the chamber companion\\
XVIII & \texttt{Paper 18} & 20 & the unitary remainder in rank $d$\\
XX & \texttt{Paper 20} & 21 & the invariant subspace in rank $d$, and one flag count that produces it\\
Phys1 & \texttt{Phys 1} & 22 & Effective mass from spanning trees, and the absence of Bloch phases on rank-two hyperbolic lattices\\
Holo1 & \texttt{Holo 1} & 23 & What the boundary algebra of $p$-adic AdS/CFT does not see\\
Net1 & \texttt{Net 1} & 24 & The sign of $\lambda$: an algebraic certificate for how a graph tower is generated\\
Ctrl1 & \texttt{Ctrl 1} & 25 & Power-sum lower bounds on hidden eigenvalues in positive realizations\\
Bio1 & \texttt{Bio 1} & 26 & Assembly ambiguity is not determined by branch statistics: critical groups of de Bruijn graphs\\
Bio2 & \texttt{Bio 2} & 27 & The sandpile group as a coordinate system on genome reconstructions\\
Bio3 & \texttt{Bio 3} & 28 & The sandpile group of a de Bruijn graph splits along its repeats\\
Bio4 & \texttt{Bio 4} & 29 & Multi-copy repeats: where the two-copy theory of assembly ambiguity stops\\
Bio5 & \texttt{Bio 5} & 30 & The bundle-chain decomposition of a bacterial genome: \emph{Escherichia coli} K-12 at five $k$-mer lengths\\
Bio6 & \texttt{Bio 6} & 31 & Double strands without a new lemma: the reverse-complement double cover of a de Bruijn graph\\
Bio7 & \texttt{Bio 7} & 32 & The reverse-complement quotient is a signed graph: what Zaslavsky's theory gives, and what it cannot\\
\bottomrule
\end{longtable}

\mainmatter
\chapter{Introduction: one group, two channels, and what towers do to it}\label{chap:intro}

\section{The object}

Everything in this book starts from a finite $(q+1)$-regular graph $Y$ --- in the
arithmetic cases a quotient $\Gamma\backslash X$ of the Bruhat--Tits tree $X$ of
$\mathrm{PGL}_2$ over a non-archimedean local field with residue field of order $q$ --- and
from three integer matrices attached to it: the adjacency (Hecke) operator $T_p$, the
Laplacian $\Delta_p=(q+1)I-T_p$, and the non-backtracking (Hashimoto) matrix
$A^{\mathrm{nb}}$ on directed edges. The finite abelian group
\[
\Jac(Y)\;=\;\operatorname{Tor}\operatorname{coker}\Delta_p ,
\]
the \emph{sandpile} or \emph{critical} group, has order $\kappa(Y)$, the number of
spanning trees, and its structure --- not just its order --- is read off from the Smith
normal form of $\Delta_p$. On an arithmetic quotient, $\kappa(Y)=\frac1{|V|}\prod_f
\bigl(p+1-a_p(f)\bigr)$ is a product of inverse local $L$-values over the eigenforms
$f$ of level $Y$: Frobenius traces precipitate as torsion. This single identity
(Chapter~\ref{chap:I}, Theorem~1.3.5) is the seed of the programme; the word
``algorithmic'' in the title refers to the fact that the group, and everything built on
it, is computed exactly, by integer elimination, on every instance we discuss.

\section{The two channels}

There are two natural non-commutative algebras over $Y$. The \emph{boundary channel}
is the crossed product $C(\partial X)\rtimes\Gamma$, equivalently the Cuntz--Krieger
algebra of $A^{\mathrm{nb}}$; it carries a Satake-twisted KMS dynamics whose partition
function is exactly the local $L$-factor. The \emph{Laplacian channel} is the
$\Z[u^{\pm1}]$-module $\coker(1-uT_p+qu^2)$ whose Fitting ideal is the Ihara zeta
function and whose specialization at $u=1$ has torsion $\Jac(Y)$.

The first principle of the book is that \emph{these two channels do not see the same
thing}. The boundary algebra is arithmetically blind: its $K$-theory is
$\Z^r\oplus\Z/(r-1)$ with $r$ the rank of $\Gamma$, whatever the Hecke data
(Chapter~\ref{chap:I}, Theorem~1.3.7); and its KMS \emph{states} see only the Perron
branch of the spectrum, so tempered Satake parameters survive only as functionals
(Theorem~1.2.5, the weight-gap theorem 1.2.11). The Laplacian channel sees the
arithmetic but is a module, not an algebra. The obvious repair --- realize the Ihara
companion matrix $C=\bigl(\begin{smallmatrix}T_p&-qI\\ I&0\end{smallmatrix}\bigr)$ by a
non-negative matrix, hence by a Cuntz--Krieger algebra --- is impossible:
$\operatorname{Tr}C^2=-n(q-1)<0$, and shift equivalence preserves the nonzero spectrum,
so no dilation of any size works (Chapter~\ref{chap:II}). What \emph{can} be done is
to realize the group $\Z\oplus\Jac(Y)$ as $K_0$ of a Kirchberg algebra by a
$2n\times2n$ non-negative ``shadow'' matrix (Chapter~\ref{chap:VI}); this shares only
Bowen--Franks data, not spectrum, with the pencil, which is why it evades the
obstruction. In the projective limit the obstruction dissolves into a type
$\mathrm{III}_{1/q}$ factor with no trace, but so does $K_0$ (Chapter~\ref{chap:III}).

\section{Towers}

The second principle is that new content appears only along towers. For a finite
graph the Hecke--Dirac Fredholm module is linear algebra in $M_n(\Z)$; along a tower
$Y_0\leftarrow Y_1\leftarrow\cdots$ the groups $\Jac(Y_k)$ grow, and the growth law is
where the arithmetic lives.

\emph{Abelian towers} (Part~II, Chapters~\ref{chap:IV}, \ref{chap:IX}, \ref{chap:XIV}).
For a $\Z/\ell^k$ voltage tower the inverse limit of the Laplacian cokernels is
$\Lambda^n/D(1+T)\Lambda^n$, $\Lambda=\Z_\ell[[T]]$, with exact control and
characteristic ideal generated by $\Theta(T)=\det D(1+T)$; the Iwasawa invariants are
read off the Newton polygon. $\Theta$ has an exceptional zero of order exactly two at
the trivial character, and its second derivative is $-\kappa(Y_0)\|h_\alpha\|^2$, where
$h_\alpha$ is the harmonic representative of the voltage class: a graph
Mazur--Tate--Teitelbaum formula whose $L$-invariant $-\|h_\alpha\|^2$ is a rational
number. This $L$-invariant is, at once, the effective mass of the bottom Bloch band of
the $\Z$-cover, the diffusion constant of the accumulated voltage under simple random
walk, and the Albanese length of the class; it is \emph{not} a cyclic-cohomology index,
and Chapter~\ref{chap:XIV} proves that no such index exists on the commutative triple.

\emph{The Iwahori tower} (Chapters~\ref{chap:V}, \ref{chap:VII}, \ref{chap:XIII}). The
tower of iterated line digraphs (non-backtracking paths of length $k$) has no deck
group. Its sandpile orders obey $\operatorname{ord}_p|\Jac(Y_k)|=\mu p^k-k+\nu_0$ with
$\lambda=-1$, which no $\Z_p[[T]]$-module can produce; the $-k$ is an accumulated Euler
characteristic of a two-term degeneracy complex. The tower is governed instead by two
degeneracy maps $\tau$ (tail) and $\eta$ (head) whose incidence calculus
$\tau_*\eta^*=A_k$, $\eta^*\tau_*=A_{k+1}$, $\tau_*\tau^*=qI$ replaces the group ring.
The upper Hecke operator factors through the tail norm and therefore kills its kernel,
which makes the ``arc determinant'' exactly $q^{n_k(q-1)}$ and turns the growth law into
a snake-lemma identity. A universal (Cohn) localization of the degeneracy algebra carries
a $K_1$-class whose level evaluations are $\pm n_k\kappa(Y_k)$ and whose boundary is a
non-split extension of $\Jac(Y_k)$ by $\Z/n_k$; the analytic side of a ``main
conjecture'' for this class is the central open problem of the book.

\emph{Operator algebras of the tower} (Part~III). The tail map is an in-covering of
shadow graphs and induces unital embeddings $\mathcal O_k\to\mathcal O_{k+1}$ of the
level Kirchberg algebras, whose inductive limit is a Kirchberg algebra with
$K_1\cong\Z$ and $\operatorname{Tor}K_0$ the Pontryagin dual of the sandpile pro-module.
For $q=p$ prime the kernel of the tail norm on $\Jac(Y_{k+1})$ is exactly the
$p$-torsion, so the $p$-part of the pro-module is torsion-free of infinite rank --- a
third certificate that no commutative Iwasawa mechanism is at work. The head map is an
out-covering and induces no homomorphism, but a proper $C^*$-correspondence whose
$K_0$-index is $\eta_*$; the two Kasparov round trips are the Hecke operators.

\section{Rank two}

Part~IV replaces the tree by $\widetilde A_2$ buildings. The cubic Hecke pencil
$P(u)=I-uA_1+qu^2A_2-q^3u^3$ carries the arithmetic, but the $\Z/3$ type grading kills
the trace obstruction of rank one, and the realization problem becomes an instance of
the (open) weak Generalized Spectral Conjecture over $\Z$. Thick quotients have
$b_1=0$, so there are no abelian towers and no rank-two Iwasawa theory: the theory is
non-commutative from the start. The geodesic edge flow $L_E$ --- the rank-two
Hashimoto matrix --- carries the whole pencil through an explicit integer intertwiner,
and its complement is $\sqrt q$ times an orthogonal matrix, so the ``parahoric'' zeros
of the edge zeta lie on the critical circle for \emph{every} complex, Ramanujan or not;
the Ramanujan property lives entirely in the vertex pencil (Chapter~\ref{chap:XV}).

\section{Applications}

Part~V applies the sandpile/tower machinery outside number theory: the effective-mass
reading of the $L$-invariant becomes a statement about tree lattices in condensed
matter; the boundary-blindness theorem becomes a statement about what the boundary
algebra of $p$-adic AdS/CFT cannot see; the sign of the Iwasawa $\lambda$ certifies
whether a graph tower is generated by a normal subgroup; the trace obstruction gives
power-sum lower bounds on hidden eigenvalues in positive realizations; and, most
substantially, the critical group of a de Bruijn graph is shown to measure genome
assembly ambiguity beyond branch statistics, to split exactly along repeats, and --- for
two-copy repeats --- to be the critical group of a one-vertex map, with a real phage
genome as the running example.

\section{How to read the book}

Chapter~\ref{chap:I} is required for everything. Part~II needs Chapter~\ref{chap:II}
\S5 (the towers) and Chapter~\ref{chap:III} \S2 (the growth law). Part~III needs
Chapters~\ref{chap:V}--\ref{chap:VII} and \ref{chap:VI}. Part~IV needs
Chapter~\ref{chap:II} for the obstruction and Chapter~\ref{chap:VI} for the shadow
construction, and is otherwise self-contained. The applied chapters need only the
definitions of Chapter~\ref{chap:I}, \S3 and, for the biology, nothing beyond the
sandpile group of a digraph. The closing chapter collects every ``solved / open /
impossible'' verdict in one place, with the corrections that later chapters made to
earlier ones.

\part{The bridge over the tree: two channels and one obstruction}
% generated from Paper 1/anccft.tex
\chapter[I. The Langlands--thermodynamic bridge over Bruhat--Tits trees]{Algorithmic non-commutative class field theory, I: the Langlands--thermodynamic bridge over Bruhat--Tits trees}\label{chap:I}
\chaptermark{I}
\begin{chapterabstract}
We open the non-abelian, algorithmic half of the programme: the localized Bost--Connes
formalism is transported from the one-dimensional valuation space to the Bruhat--Tits
tree of $\GL_2$ over a non-archimedean local field, and the arithmetic input is upgraded
from Frobenius classes in an abelian Galois group to Satake parameters of unramified
automorphic representations.

Section~\ref{I:sec:geometry} constructs the state space $\Omega$ (the space of lattice
chains, equivalently the tree boundary $\partial X\cong\PPone(F)$), the truncated Hecke
groupoids $\cG^{(N)}$ and their limit $\cG=\bigcup_N\cG^{(N)}$, and the tail equivalence
relation $\cR$ with its Hecke-degree cocycle; the groupoid is \'etale, amenable and
essentially principal, and its algebra is a Cuntz algebra locally and a Cuntz--Krieger
algebra $\cO_{A^{\mathrm{nb}}_Y}\simeq C(\partial X)\rtimes\Gamma$ globally.

Section~\ref{I:sec:thermo} inserts the Satake parameters into the Radon--Nikodym cocycle.
The partition function of the twisted system is \emph{exactly} the local $L$-factor,
$Z_p(\beta)=L_p(\beta,\pi_p)$, an identity of the Hecke recursion, and the abscissa of
convergence $\beta_c(\pi_p)=\log_q\max_i|\alpha_{p,i}|$ is the largest pole of $L_p$. Two
structural facts then follow, and they cut in opposite directions. Positivity
(Perron--Frobenius) shows that $\KMS$ \emph{states} see only the Perron branch of the
unramified spectrum, so cuspidal Satake data survive only as $\KMS$ \emph{functionals};
and the Kakutani dichotomy for the associated conformal measures has obstruction series
$\Xi_\beta(\pi,\pi';S)=\sum_{p\in S}p^{-\beta}\|\,\cdot\,\|^2$ governed by
Rankin--Selberg, giving a measure-theoretic form of strong multiplicity one. Combining
them yields the \emph{weight gap theorem}: equilibrium requires
$\beta>\beta_c(S)+w/2$ while the obstruction is confined to $\beta\le\beta_c(S)$, where
$w$ is the motivic weight; the two ranges meet if and only if $w=0$. The abelian theory
is the case $w=0$; Paper XXI's universal gap of width $1$ is the Eisenstein case $w=2$;
tempered weight two gives $w=1$ and a residual gap of exactly $\tfrac12$, the Ramanujan
exponent.

Section~\ref{I:sec:ktheory} is the algorithmic bridge, and it is where the higher-weight
arithmetic survives. Pimsner's exact sequence for the Hecke correspondence, the
non-commutative replacement of the equivariant Pimsner--Voiculescu sequence, identifies
$K_*$ with the kernel and cokernel of $1-u\,T_p+qu^2$ over $\Z[u^{\pm1}]$ --- the
\emph{Ihara module}, whose Fitting ideal is the Ihara zeta function by Bass's identity.
Specializing $u=1$ gives the Hecke--Laplacian $\Delta_p=(q+1)I-T_p$ and the structure
theorem
\[
K_0\cong\Z\oplus\Jac(Y),\qquad K_1\cong\Z,\qquad
|\Jac(Y)|=\frac1{|V|}\prod_f\bigl(p+1-a_p(f)\bigr)=\frac1{|V|}\prod_f\#E_f(\F_p),
\]
the torsion being computed exactly, as a group and not merely as an order, by the Smith
normal form of $\Delta_p$: Frobenius traces precipitate as $K$-theoretic torsion. By
contrast the boundary crossed product is arithmetically blind,
$K_0(C(\partial X)\rtimes\Gamma)\cong\Z^r\oplus\Z/(r-1)$ with $r$ the rank of $\Gamma$,
a theorem proved here by collapsing the reversal relation
$[f]+[\bar f]=c_{o(f)}$ to a single class of order $|E_Y|-|V_Y|$. Eleven independent
machine computations confirm both channels. Along Benjamini--Schramm convergent families
the torsion density converges exactly to Lyons' tree entropy,
$\frac1n\log|\Tors K_0|\to q\log q-\frac{q-1}2\log(q^2-1)$, and the Ramanujan bound
$2\log(1+\sqrt q)$, while unconditional, is strict and exceeds it by
$2q^{-1/2}+O(q^{-1})$ as $q\to\infty$ --- a Taylor expansion, not a fit --- because it
replaces a Kesten--McKay integral by its spectral edge. The Laplacian channel is realised as a
$\Z[u^{\pm1}]$-module and as the index of a Hecke--Dirac operator, not as a
$C^*$-algebra: the class $T_p-qI$ is a legitimate $KK$-class but not a correspondence,
and closing that gap is the paper's main open problem.
\end{chapterabstract}

\section*{Introduction, and the status of each statement}

The thirty-nine papers of \emph{Localized Bost--Connes systems} settle the
$\GL_1$ problem: the $\KMS_\beta$ simplex of $\cA_{K,S}$ is $\Prob(G_S/\Xi_\beta^\perp)$,
the obstruction group $\Xi_\beta$ is cut out by a Dirichlet series in the Frobenius
classes, the dichotomy underneath is Kakutani's, and the integral invariants are Smith
normal forms of $\Ns\pp\cdot 1-\Lambda^r\Pi$. Every ingredient of that theory has a
$\GL_2$ counterpart, and the counterparts are not formal analogies: valuations become
lattice chains, the profinite completion becomes the tree boundary, Frobenius classes
become Satake parameters, Dirichlet series become automorphic $L$-functions, and the
Smith normal form of a norm-minus-multiplication matrix becomes the Smith normal form of
a Hecke--Laplacian. This paper writes the dictionary and proves the entries that can be
proved.

It is written against a negative result of the earlier series, and we state at once how
it stands with respect to it. Paper XXI of that series proves that the untwisted $\GL_2$
system has a \emph{universal gap of width one}: the automorphic obstruction can be
non-trivial only for $\beta\le\beta_c(S)$, Gibbs states exist only for
$\beta>1+\beta_c(S)$, and no temperature does both. Nothing below contradicts that.
What changes here is \emph{where} the automorphic parameter is inserted --- into the
dynamics, not into the state --- and the effect of the change is quantitative and exact:
the gap is $w/2$, where $w$ is the motivic weight of the representation
(Theorem~\ref{I:thm:gap}). Paper XXI is the case $w=2$; the abelian theory of the monograph
is the case $w=0$, where the gap closes and the classification succeeds; tempered weight
two is $w=1$ and the residual gap is the Ramanujan exponent $\tfrac12$. The thermodynamic
channel therefore classifies exactly the weight-zero, that is Artin, part of the Langlands
correspondence. That is a sharp statement about what thermodynamics can and cannot do,
and it is the reason Section~\ref{I:sec:ktheory} exists: $K$-theory carries no positivity
constraint and no temperature, and it does see the higher-weight data.

\smallskip
\noindent\textbf{Status.} We label every numbered statement in one of three ways.
\emph{Theorem/Proposition}: proved here or an immediate assembly of quoted results, with
the assembly given. \emph{Verified}: an integral statement established by exact machine
computation on a stated finite list of cases, with a proof sketch but no proof; these are
marked $[\mathrm{V}]$ and the verification table is Table~\ref{I:tab:snf}; in the present
version no numbered statement carries this label, Theorem~\ref{I:thm:blind} having been
promoted to a proof, and Table~\ref{I:tab:snf} now serves as independent confirmation
rather than as evidence. \emph{Conjecture}:
neither. No statement below is asserted at a higher level than it is entitled to. In
particular Theorem~\ref{I:thm:kms} is a theorem, its extension to ramified and to $\GL_n$
data is Conjecture~\ref{I:conj:gln}, and the identification of the full extremal simplex
with the polar divisor of an $L$-function --- the form in which such a correspondence is
usually stated informally --- is \emph{false} as stated, for the reason isolated in
Theorem~\ref{I:thm:perron}.

\section{The geometric state space over Bruhat--Tits trees}\label{I:sec:geometry}

\subsection{The local field and its tree}

Throughout, $F$ is a non-archimedean local field with ring of integers $\cO$, maximal
ideal $\mathfrak{m}=(\varpi)$, uniformiser $\varpi$, normalised valuation
$v\colon F^\times\twoheadrightarrow\Z$, and residue field $k_F=\cO/\mathfrak{m}$ of
cardinality $q$. In the arithmetic applications $F=\Q_p$ and $q=p$. Write
$G=\GL_2(F)$, $K=\GL_2(\cO)$, $Z\cong F^\times$ for the centre.

\begin{definition}[Lattices and the tree]\label{I:def:tree}
An \emph{$\cO$-lattice} in $V=F^2$ is a free $\cO$-submodule $L\subset V$ of rank $2$.
Two lattices are \emph{homothetic}, $L\sim L'$, if $L'=cL$ for some $c\in F^\times$. The
Bruhat--Tits tree $X=X_F$ of $\PGL_2(F)$ has
\[
X^{(0)}=\{\text{lattices}\}/\!\sim,\qquad
[L]\text{ adjacent to }[L']\iff\exists\,L\in[L],\ L'\in[L']\ \text{with}\
\varpi L\subsetneq L'\subsetneq L .
\]
$X$ is a tree, it is $(q+1)$-regular, $G$ acts on it by $g\cdot[L]=[gL]$ with $Z K$ the
stabiliser of the standard vertex $x_0=[\cO^2]$, and the vertex set is
$X^{(0)}\cong G/ZK$.
\end{definition}

\begin{definition}[Boundary]\label{I:def:boundary}
A \emph{ray} from $x_0$ is a sequence $\xi=(x_0,x_1,x_2,\dots)$ of vertices with
$x_i\sim x_{i+1}$ and $x_{i+1}\neq x_{i-1}$ (\emph{non-backtracking}). The
\emph{boundary} $\partial X$ is the set of rays from $x_0$, topologised by the cylinder
sets $C(x_0,\dots,x_n)$. Then $\partial X$ is a compact totally disconnected space
without isolated points --- a Cantor set --- and the map sending a ray to the
corresponding line yields a $G$-equivariant homeomorphism
\[
\partial X\;\xrightarrow{\ \cong\ }\;\PPone(F)=\PPone(\cO)/\!\sim .
\]
The cylinder $C(x_0,\dots,x_n)$ has $q+1$ children at level $1$ and exactly $q$ children
at every later level.
\end{definition}

The two branching numbers $q+1$ and $q$ are the whole geometric content of the section,
and they are the reason the boundary and not the lattice space is the correct state
space; see Remark~\ref{I:rem:offbyone}.

\begin{definition}[Chain space]\label{I:def:chains}
Let
\[
\widetilde\Omega:=\Bigl\{(L_n)_{n\ge0}\ :\ L_0=\cO^2,\ L_{n+1}\subsetneq L_n,\
[L_n:L_{n+1}]=q\ \ \forall n\Bigr\},
\]
the space of \emph{descending lattice chains of step one}, with the inverse-limit
topology. Since the sublattices of index $q$ in a lattice $L$ are in bijection with
$\PPone(k_F)$, which has $q+1$ elements, $\widetilde\Omega$ is homeomorphic to the
boundary of the rooted $(q+1)$-ary tree, hence a Cantor set, and
$[L_0],[L_1],[L_2],\dots$ is a walk in $X$ issuing from $x_0$. The subspace
\[
\Omega:=\{(L_n)\in\widetilde\Omega:\ [L_{n+1}]\neq[L_{n-1}]\ \forall n\ge1\}
\]
of \emph{geodesic} chains is closed and $\Omega\cong\partial X\cong\PPone(F)$ by
Definition~\ref{I:def:boundary}. We call $\widetilde\Omega$ the \emph{Toeplitz} (or full)
state space and $\Omega$ the \emph{boundary} state space.
\end{definition}

\begin{remark}\label{I:rem:offbyone}
$\widetilde\Omega$ is the $\GL_2$ analogue of the space $Y_S$ of the abelian theory: its
level-$n$ set is the set of sublattices of index $q^n$ in $\cO^2$, of cardinality
$\sigma_1(q^n)=1+q+\dots+q^n$, and the associated Dirichlet series is the Solomon zeta
$\zeta_{M_2(\cO)}(\beta)=(1-q^{-\beta})^{-1}(1-q^{1-\beta})^{-1}$ of Paper XXI. The
passage $\widetilde\Omega\rightsquigarrow\Omega$ removes exactly the excess
$d_\pp=q+1>q$ that Paper XXI, Theorem 1.2 identifies as the source of its off-by-one:
on $\Omega$ the branching is $q$ per step and matches the norm.
\end{remark}

\subsection{The Hecke monoid and the truncated Hecke groupoids}

Let $\cH(G,K)=C_c(K\backslash G/K)$ be the spherical Hecke algebra with convolution, and
\[
T:=\bigl[K\begin{psmallmatrix}\varpi&0\\0&1\end{psmallmatrix}K\bigr],\qquad
R^{\pm1}:=\bigl[K\begin{psmallmatrix}\varpi^{\pm1}&0\\0&\varpi^{\pm1}\end{psmallmatrix}K\bigr],
\qquad
T_n:=\bigl[K\begin{psmallmatrix}\varpi^{n}&0\\0&1\end{psmallmatrix}K\bigr]\ (n\ge0),
\]
so that $\cH(G,K)=\C[T,R^{\pm1}]$ and the Satake isomorphism is
$\Sat\colon\cH(G,K)\xrightarrow{\cong}\C[\alpha_1^{\pm1},\alpha_2^{\pm1}]^{S_2}$ with
$\Sat(T)=q^{1/2}(\alpha_1+\alpha_2)$ and $\Sat(R)=\alpha_1\alpha_2$. Acting on functions
on $X^{(0)}$, $T$ is the adjacency (nearest-neighbour) operator and $R$ acts trivially.

\begin{lemma}[Hecke recursion]\label{I:lem:recursion}
On $X$ one has $T_1=T$, $T\,T_1=T_2+(q+1)T_0$ and, for $n\ge2$,
\[
T\,T_n=T_{n+1}+q\,T_{n-1},
\qquad\text{equivalently}\qquad
\sum_{n\ge0}T_n\,u^n=\bigl(1-Tu+qu^2\bigr)^{-1}
\]
as an identity of formal power series with coefficients in $\cH(G,K)$ normalised by
$\Sat(R)=q$. Consequently, for a character $\lambda$ of $\cH(G,K)$ with Satake
parameters $(\alpha_1,\alpha_2)$, $\alpha_1\alpha_2=q$,
\begin{equation}\label{I:eq:heckegen}
\sum_{n\ge0}\lambda(T_n)\,u^n=\frac1{(1-\alpha_1u)(1-\alpha_2u)} .
\end{equation}
\end{lemma}

The pencil $1-Tu+qu^2$ of Lemma~\ref{I:lem:recursion} is the single object that reappears
in every section below: in Section~\ref{I:sec:thermo} as the inverse of a partition
function, in Section~\ref{I:sec:ktheory} as the presentation matrix of a $K$-group.

\begin{definition}[Hecke monoid and degree]\label{I:def:monoid}
Let $S_\varpi:=\bigl(M_2(\cO)\cap G\bigr)/\cO^\times$, a monoid acting on lattices by
$L\mapsto gL$, and let $\deg\colon S_\varpi\to\Z_{\ge0}$, $\deg(g)=v(\det g)$. The
$K$-double cosets of $S_\varpi$ are indexed by elementary divisors, and
$\{\deg=n\}$ is the union of the double cosets of $T_a R^b$ with $a+2b=n$.
\end{definition}

\begin{definition}[Truncated Hecke groupoids]\label{I:def:groupoid}
Let $\Gamma\le\PGL_2(F)$ be a discrete, torsion-free, cocompact subgroup; such $\Gamma$
exist and are free of finite rank, and $Y:=\Gamma\backslash X$ is a finite
$(q+1)$-regular graph. Write $E_Y$ for its set of oriented edges, $\bar e$ for the
reversal of $e$, and
\[
A^{\mathrm{nb}}_{ef}=\begin{cases}1,&t(e)=o(f)\text{ and }f\neq\bar e,\\
0,&\text{otherwise},\end{cases}
\qquad e,f\in E_Y,
\]
for the \emph{non-backtracking} (edge) matrix. Let
$\Sigma_Y:=\{\omega\in E_Y^{\N}:A^{\mathrm{nb}}_{\omega_i\omega_{i+1}}=1\}$ with shift
$\sigma$. For $N\in\N$ set
\begin{equation}\label{I:eq:truncgroupoid}
\cG^{(N)}:=\bigl\{(\omega,i-j,\omega')\in\Sigma_Y\times\Z\times\Sigma_Y\ :\
0\le i,j\le N,\ \sigma^i\omega=\sigma^j\omega'\bigr\},
\end{equation}
and $\cG_\Gamma:=\bigcup_{N\ge0}\cG^{(N)}$, the \emph{Hecke groupoid}, with the
composition $(\omega,k,\omega')(\omega',k',\omega'')=(\omega,k+k',\omega'')$, the
topology generated by the sets
$Z(U,i,j,V)=\{(\omega,i-j,\omega'):\omega\in U,\ \omega'\in V,\ \sigma^i\omega=\sigma^j\omega'\}$
for $U,V$ cylinders, and the \emph{Hecke degree cocycle}
\[
c\colon\cG_\Gamma\longrightarrow\Z,\qquad c(\omega,k,\omega')=k .
\]
The same formula with $\Sigma_Y$ replaced by $\widetilde\Omega$ and $\sigma$ by the
one-sided shift on the full $(q+1)$-ary tree defines the \emph{Toeplitz Hecke groupoid}
$\widetilde\cG$.
\end{definition}

The truncation index $N$ is the number of Hecke moves allowed; $\cG^{(N)}$ is an open
subgroupoid, $\cG^{(N)}\subset\cG^{(N+1)}$, and $C^*(\cG_\Gamma)=\varinjlim_N C^*(\cG^{(N)})$.
This inductive limit is what carries the thermodynamic limit of
Section~\ref{I:sec:thermo}: a measure class on $\Omega$ is a compatible family of
quasi-invariant classes for the finite-level relations $\cR^{(N)}$ below.

\begin{definition}[Tail equivalence]\label{I:def:tail}
The \emph{tail (Hecke) equivalence relation} is the image of $\cG_\Gamma$ under
$(r,s)$,
\[
\cR:=\bigl\{(\omega,\omega')\in\Omega\times\Omega\ :\ \exists\, i,j\ge0,\
\sigma^i\omega=\sigma^j\omega'\bigr\},
\]
with the finite-level subrelations $\cR^{(N)}$ obtained by restricting $i,j\le N$, and
the \emph{degree-zero} subrelation $\cR_0:=\{(\omega,\omega'):\exists\,i,\ \sigma^i\omega=\sigma^i\omega'\}$.
A probability measure $\mu$ on $\Omega$ is \emph{$\beta$-quasi-invariant} for $\cR$ if
the measure class is preserved and the Radon--Nikodym cocycle of $\cG_\Gamma$ is
$e^{-\beta c}$; $\cR_0$ is a hyperfinite (AF) subrelation, and $\cR/\cR_0\cong\Z$ through
$c$.
\end{definition}

\begin{proposition}\label{I:prop:groupoid}
$\cG_\Gamma$ is a second countable, locally compact, Hausdorff, \'etale groupoid with
unit space $\Omega$; it is amenable, minimal and essentially principal, and
\begin{enumerate}
\item[(i)] $\cG_\Gamma$ is canonically isomorphic to the transformation groupoid
$\partial X\rtimes\Gamma$ under $\Omega\cong\partial X$, the Hecke degree $c$
corresponding to the Busemann (horocycle) cocycle of $\Gamma$ acting on $\partial X$;
\item[(ii)] $C^*(\cG_\Gamma)=C^*_r(\cG_\Gamma)$ is the Cuntz--Krieger algebra
$\cO_{A^{\mathrm{nb}}_Y}$, simple, purely infinite, nuclear, and Morita equivalent to
$C(\partial X)\rtimes\Gamma$;
\item[(iii)] in the rooted local model, $C^*(\widetilde\cG)\cong\cO_{q+1}$ and its
boundary (geodesic) reduction is stably isomorphic to $\cO_q$, so that
\[
K_0\bigl(C^*(\widetilde\cG)|_\Omega\bigr)\cong\Z/(q-1),\qquad
K_1\bigl(C^*(\widetilde\cG)|_\Omega\bigr)=0 .
\]
\end{enumerate}
\end{proposition}

\begin{remark}\label{I:rem:unitclass}
Part (iii) is the exact $\GL_2$ analogue of the unit-class torsion $\Z/(\Ns\pp^h-1)$ of
Paper XXXVIII: locally the boundary quotient contributes $\Z/(q-1)$ and nothing else. All
arithmetic beyond the residue cardinality is therefore invisible at a single place, and
must be sought either in the global assembly (Section~\ref{I:sec:thermo}) or in the
quotient graph $Y$ (Section~\ref{I:sec:ktheory}). This is the structural reason the
$\GL_2$ theory cannot be purely local, and it has no counterpart in the $\GL_1$ theory,
where the local boundary already carries the class group.
\end{remark}

\section{The thermodynamic--automorphic correspondence}\label{I:sec:thermo}

\subsection{Satake weights in the Radon--Nikodym cocycle}

Fix an unramified irreducible admissible representation $\pi_p$ of $G$ with Satake
parameters $(\alpha_{p,1},\alpha_{p,2})$ normalised arithmetically,
$\alpha_{p,1}\alpha_{p,2}=q^{w}$ with $w$ the motivic weight ($w=0$ for Artin
parameters, $w=k-1$ for weight-$k$ forms; $w=1$ in the classical weight-two case), and
local $L$-factor
\[
L_p(s,\pi_p)=\prod_{i=1,2}\bigl(1-\alpha_{p,i}\,q^{-s}\bigr)^{-1}.
\]

\begin{definition}[Satake potential and twisted dynamics]\label{I:def:dynamics}
Let $h_\pi\colon X^{(0)}\to\C^\times$ be a Hecke eigenfunction, $T h_\pi=\lambda_\pi h_\pi$
with $\lambda_\pi=\alpha_{p,1}+\alpha_{p,2}$, normalised by $h_\pi(x_0)=1$. Define the
\emph{Satake potential} on oriented edges and the associated \emph{Ruelle--Hecke transfer
operator} on $C(\Omega)$ by
\[
\varphi_\pi(e):=\log\frac{h_\pi(t(e))}{h_\pi(o(e))},\qquad
\bigl(\cL_{\beta,\pi}f\bigr)(\omega):=\sum_{\sigma(\eta)=\omega}
q^{-\beta\,c(\eta)}\,e^{\varphi_\pi(\eta_1)}f(\eta).
\]
The \emph{twisted dynamics} $\sigma^{(\pi)}$ on $C^*(\cG_\Gamma)$ is the one-parameter
group determined on $C_c(\cG_\Gamma)$ by
\[
\sigma^{(\pi)}_t(f)(\gamma)=e^{\,it\,(\beta\text{-independent})\,c_\pi(\gamma)}f(\gamma),
\qquad
c_\pi(\gamma):=c(\gamma)\log q-\varphi_\pi(\gamma),
\]
so that the Radon--Nikodym cocycle of a $\KMS_\beta$ state is
\begin{equation}\label{I:eq:RN}
\frac{d\mu\circ\gamma}{d\mu}(\omega)=e^{-\beta c_\pi(\gamma)}
=q^{-\beta c(\gamma)}\Bigl(\frac{h_\pi(t\gamma)}{h_\pi(o\gamma)}\Bigr)^{\beta}.
\end{equation}
\end{definition}

Equation~\eqref{I:eq:RN} is the precise sense in which the Satake parameter enters the
Radon--Nikodym derivative: it is the Doob $h$-transform of the geometric cocycle by the
Hecke eigenfunction. The untwisted system of Paper XXI is $h_\pi\equiv1$.

\subsection{The partition function is the local \texorpdfstring{$L$}{L}-factor}

\begin{theorem}[Partition function]\label{I:thm:partition}
Let $Z_p(\beta;\pi):=\sum_{n\ge0}\lambda_\pi(T_n)\,q^{-n\beta}$ be the partition function
of the truncated tower $(\cG^{(N)})_N$ with the dynamics of
Definition~\ref{I:def:dynamics}. Then the series converges absolutely if and only if
$\beta>\beta_c(\pi_p)$, where
\[
\beta_c(\pi_p):=\log_q\max_i|\alpha_{p,i}|,
\]
and in that range
\[
\boxed{\;Z_p(\beta;\pi)=\prod_{i=1,2}\bigl(1-\alpha_{p,i}q^{-\beta}\bigr)^{-1}
=L_p(\beta,\pi_p).\;}
\]
Moreover $\beta_c(\pi_p)$ is exactly the largest real part of a pole of
$s\mapsto L_p(s,\pi_p)$, so the critical inverse temperature of the twisted system is the
abscissa of convergence of the local $L$-factor. In particular
$\beta_c=w/2$ when $\pi_p$ is tempered $($Ramanujan$)$, and $\beta_c=w$ when $\pi_p$ is the
Eisenstein point $(\alpha_1,\alpha_2)=(1,q^{w})$.
\end{theorem}

\begin{proof}
Substitute $u=q^{-\beta}$ in \eqref{I:eq:heckegen} of Lemma~\ref{I:lem:recursion}; the radius
of convergence of the right-hand side is $\min_i|\alpha_{p,i}|^{-1}$. The poles of
$L_p(s,\pi_p)$ are at $q^{-s}=\alpha_{p,i}^{-1}$, i.e. at
$\Re s=\log_q|\alpha_{p,i}|$.
\end{proof}

\begin{remark}
For the trivial (Eisenstein) parameter $w=2$, $(\alpha_1,\alpha_2)=(1,q)$, one recovers
$Z_p(\beta)=(1-q^{-\beta})^{-1}(1-q^{1-\beta})^{-1}$, the Solomon zeta of $M_2(\cO)$ and
the local factor of $\zeta(\beta)\zeta(\beta-1)$: Theorem~\ref{I:thm:partition} contains
the partition function of the Connes--Marcolli $\GL_2$-system, and of Paper XXI, as its
$h_\pi\equiv1$ specialisation.
\end{remark}

\subsection{Equilibrium states, and what they can see}

\begin{theorem}[$\KMS$ classification, local]\label{I:thm:kms}
Let $Y=\Gamma\backslash X$ be connected and non-bipartite, and let
$A^{\mathrm{nb}}_{\beta,\pi}$ be the non-backtracking matrix weighted by
$e^{-\beta c_\pi}$ as in \eqref{I:eq:RN}. Then:
\begin{enumerate}
\item[(i)] $\KMS_\beta\bigl(C^*(\cG_\Gamma),\sigma^{(\pi)}\bigr)\neq\emptyset$ if and only
if the spectral radius satisfies $\rho\bigl(A^{\mathrm{nb}}_{\beta,\pi}\bigr)=1$, i.e.
if and only if the Bowen--Ruelle pressure equation $P(-\beta c_\pi)=0$ holds;
\item[(ii)] in that case the $\KMS_\beta$ simplex is a single point, the state
associated with the Perron eigenvector, and the corresponding measure $\mu_{\beta,\pi}$
is the unique $e^{-\beta c_\pi}$-conformal measure for $\cR$;
\item[(iii)] the associated von Neumann factor is of type $\mathrm{III}_{q^{-\beta}}$ if
$c_\pi(\cG_\Gamma)\subseteq\Z\log q$, and of type $\mathrm{III}_1$ if the closed group
generated by $e^{-\beta c_\pi(\cG_\Gamma)}$ is $\R_{>0}$;
\item[(iv)] for the untwisted potential the equation in $(\mathrm{i})$ reads
$q^{1-\beta}=1$, so the transition is at $\beta_c=1$ on the boundary state space
$\Omega$, against $\beta_c=2$ on the Toeplitz state space $\widetilde\Omega$.
\end{enumerate}
\end{theorem}

Part (iv) is the promised removal of Paper XXI's off-by-one: it is bought by the passage
from $\widetilde\Omega$ to $\Omega$, i.e. by the boundary quotient, and by nothing else.
The next theorem says what the price is.

\begin{theorem}[Perron obstruction: states cannot see cuspidal parameters]\label{I:thm:perron}
Let $\sigma^{(\pi)}$ be as in Definition~\ref{I:def:dynamics}. A $\KMS_\beta$ state exists
only if the eigenfunction $h_\pi$ can be chosen strictly positive on $X^{(0)}$. On a
$(q+1)$-regular tree a positive $T$-eigenfunction with eigenvalue $\lambda$ exists if and
only if $|\lambda|\ge2\sqrt q$. Consequently:
\begin{enumerate}
\item[(i)] if $\pi_p$ is tempered $(|\lambda_\pi|<2\sqrt q$, i.e.
$|\alpha_{p,i}|=q^{w/2})$, the twisted system has \emph{no} $\KMS_\beta$ state at any
$\beta$; the Satake data survive only as the $\KMS_\beta$ \emph{functional}
$\phi_{\beta,\pi}$ obtained by analytic continuation in $\lambda$, whose partition
function is $L_p(\beta,\pi_p)$ by Theorem~\ref{I:thm:partition};
\item[(ii)] the extremal $\KMS_\beta$ states of the untwisted system are parameterised by
the Perron branch alone --- the Eisenstein/complementary-series locus
$\lambda\ge2\sqrt q$ of the unramified spectrum --- and their partition functions are the
corresponding $L$-factors $L_p(\beta,1\boxplus|\cdot|^{w})$;
\item[(iii)] therefore the assertion ``the extremal $\KMS_\beta$ simplex is
parameterised by the poles of the automorphic $L$-function'' is true on the Eisenstein
locus and false on the cuspidal locus; the correct general statement replaces states by
functionals, and the simplex by the polar divisor of a meromorphic family.
\end{enumerate}
\end{theorem}

\begin{remark}[Thermodynamic ghost states]\label{I:rem:ghost}
We propose the name \emph{thermodynamic ghost states} for the non-positive continuations
$\phi_{\beta,\pi}$ of Theorem~\ref{I:thm:perron}(i): they are $\sigma^{(\pi)}$-$\KMS$
functionals with the correct modular property and the correct partition function
$L_p(\beta,\pi_p)$, but they are not states, having no positive extension to the algebra
and hence no GNS representation and no thermal equilibrium interpretation. The
terminology is borrowed, with the usual caution, from the negative-norm states of gauge
quantisation: as there, the ghosts are indispensable to the bookkeeping and absent from
the physical spectrum. The physical content of Theorem~\ref{I:thm:perron} is then this. The
metric geometry of the tree fixes a threshold, $2\sqrt q$, which is exactly the tempered
bound; a Satake parameter satisfying the Ramanujan estimate lies \emph{below} it, and a
positive Hecke eigenfunction --- the Gibbs weight of the would-be equilibrium --- then
simply does not exist on $X$. Cuspidality and thermal equilibrium are mutually exclusive
on the Bruhat--Tits tree. The arithmetic of the cuspidal spectrum is therefore not
suppressed but displaced: it survives as analytic data off the state space, and, as
Section~\ref{I:sec:ktheory} shows, as integral data in $K$-theory, where no positivity is
required.
\end{remark}

\begin{remark}
Theorem~\ref{I:thm:perron} is the operator-algebraic shadow of a familiar fact: positivity
of a state is a Perron--Frobenius condition, and Perron--Frobenius selects the trivial
representation. It is the same phenomenon that makes the constant function the only
positive $\ell^2$-eigenfunction on a finite graph, and the same one that in Paper XXVII
forces a $\KMS$ state to be tracial on the coefficient algebra and thereby erases inner
twists. Any programme that hopes to read cuspidal data off an equilibrium \emph{state}
must first defeat this, and we do not know how to.
\end{remark}

\subsection{The Kakutani obstruction and Rankin--Selberg}

We now globalise. Let $S$ be a set of rational primes, $\beta_c(S)$ the convergence
exponent of $\sum_{p\in S}p^{-\beta}$, and for unramified $\pi=\otimes_p\pi_p$ let
\[
\Omega_S:=\prod_{p\in S}\Omega_p,\qquad
\mu^\pi_\beta:=\bigotimes_{p\in S}\mu_{\beta,\pi_p},
\]
the product of the local conformal measures of Theorem~\ref{I:thm:kms}(ii), taken on the
locus where they exist, that is on the Eisenstein and complementary-series range
$|\lambda_p|\ge2\sqrt q$, or of their $h$-transform representatives there.

\begin{theorem}[Kakutani dichotomy for Satake data]\label{I:thm:kakutani}
For any two such $\pi,\pi'$ the measures $\mu^\pi_\beta$ and $\mu^{\pi'}_\beta$ are
either equivalent or mutually singular, and they are equivalent if and only if
\begin{equation}\label{I:eq:obstruction}
\Xi_\beta(\pi,\pi';S):=\sum_{p\in S}p^{-\beta}\,
\bigl|\widehat\lambda_p(\pi)-\widehat\lambda_p(\pi')\bigr|^2<\infty,
\qquad
\widehat\lambda_p:=\frac{\lambda_p}{2\sqrt p}=\cos\theta_p ,
\end{equation}
the equivalence of the Hellinger criterion with \eqref{I:eq:obstruction} holding with
absolute implied constants depending only on the temperedness bound. The series
\eqref{I:eq:obstruction} is the exact $\GL_2$ analogue of the obstruction series
$\sum_\pp\Ns\pp^{-\beta}|1-\chi(\Frob_\pp)|^2$ of the abelian theory, with the Frobenius
character replaced by the normalised Satake trace.
\end{theorem}

\begin{corollary}[Measure-theoretic strong multiplicity one]\label{I:cor:multone}
Let $\pi,\pi'$ be cuspidal automorphic representations of $\GL_2(\A_\Q)$, unramified
outside a finite set, and $S$ the set of all but finitely many primes, so
$\beta_c(S)=1$. Expanding \eqref{I:eq:obstruction} into Rankin--Selberg logarithms,
\[
\Xi_\beta(\pi,\pi';S)=\log L_S(\beta,\pi\times\widetilde\pi)
+\log L_S(\beta,\pi'\times\widetilde{\pi'})
-2\,\Re\log L_S(\beta,\pi\times\widetilde{\pi'})+O(1),
\]
so that, by the pole of $L(s,\pi\times\widetilde\pi)$ at $s=1$ and the regularity and
non-vanishing there of $L(s,\pi\times\widetilde{\pi'})$ for $\pi\not\cong\pi'$
$($Jacquet--Shalika$)$, one has $\Xi_1(\pi,\pi';S)=\infty$ and hence
\[
\pi\not\cong\pi'\ \Longrightarrow\ \mu^\pi_1\perp\mu^{\pi'}_1 .
\]
Distinct automorphic representations occupy mutually singular thermodynamic phases.
\end{corollary}

\begin{remark}
The Rankin--Selberg $L$-function appears here legitimately, and its appearance is not in
conflict with Paper XXI, which shows that the \emph{one-representation} obstruction
$\Xi_\beta(\rho,S)$ is linear in $\Tr\rho(\Frob_\pp)$ and therefore involves only the
standard $L$-function, the $\Sym^2$-term being spurious. The quantity
\eqref{I:eq:obstruction} is a Hellinger distance \emph{between two} representations; it is
genuinely quadratic, and the second moment of Satake traces is by definition
Rankin--Selberg.
\end{remark}

\subsection{The weight gap}

\begin{theorem}[Weight gap]\label{I:thm:gap}
Let $\pi$ have motivic weight $w$ at almost every place, i.e.
$|\alpha_{p,i}|=p^{w/2}$ for tempered $\pi_p$. Then, for every set $S$ of primes:
\[
\text{Gibbs states exist}\iff\beta>\beta_c(S)+\tfrac{w}{2},
\qquad
\Xi_\beta(\pi,\pi';S)=\infty\ \text{possible}\iff\beta\le\beta_c(S).
\]
The two ranges are disjoint with a gap of width exactly $w/2$, and they meet if and only
if $w=0$. Consequently:
\begin{enumerate}
\item[(i)] $w=0$ (Artin parameters, weight-one forms, the abelian case): the gap
closes, and one recovers the classification of the localized Bost--Connes theory,
namely $\KMS_\beta\cong\Prob(G_S/\Xi_\beta^\perp)$;
\item[(ii)] $w=1$ (tempered weight two): the residual gap is $\tfrac12$, the Ramanujan
exponent;
\item[(iii)] $w=2$ (the Eisenstein/untwisted $\GL_2$ system): the gap is $1$, which is
Paper XXI's universal off-by-one.
\end{enumerate}
\end{theorem}

\begin{proof}[Proof sketch]
By Theorem~\ref{I:thm:partition}, the equilibrium range is the domain of absolute
convergence of the Euler product of the local factors $L_p(\beta,\pi_p)$,
$p\in S$, whose abscissa is $\beta_c(S)+\beta_c(\pi_p)$; in the tempered case this equals
$\beta_c(S)+w/2$. For the obstruction range, note that
$|\widehat\lambda_p|\le1$, so the series \eqref{I:eq:obstruction} is bounded above
by $4\sum_{p\in S}p^{-\beta}$ and is therefore finite for every
$\beta>\beta_c(S)$. Disjointness and the width follow
\end{proof}

Theorem~\ref{I:thm:gap} is the main negative--positive result of the section. It says that
the thermodynamic channel is exactly a \emph{weight-zero} instrument: it classifies the
Artin part of the Langlands correspondence and is blind, by an amount equal to half the
motivic weight, to everything above it. This is a sharpening, not a refutation, of Paper
XXI: the twist moves the equilibrium threshold from $\beta_c(S)+1$ down to
$\beta_c(S)+w/2$, and for weight-one forms it closes the gap altogether.

\begin{conjecture}[$\GL_n$ and ramified data]\label{I:conj:gln}
Let $F$ be local, $G_n=\GL_n(F)$, $X_n$ the Bruhat--Tits building of $\PGL_n(F)$, and
$\pi_p$ unramified with Satake parameters $(\alpha_{p,1},\dots,\alpha_{p,n})$. Then the
partition function of the Hecke groupoid of the chamber boundary of $X_n$, twisted as in
Definition~\ref{I:def:dynamics}, equals $L_p(\beta,\pi_p)=\prod_i(1-\alpha_{p,i}q^{-\beta})^{-1}$,
the critical inverse temperature is $\max_i\log_q|\alpha_{p,i}|$, the weight gap theorem
holds verbatim with $w$ the motivic weight, and for ramified $\pi_p$ the same statements
hold with $\Omega$ replaced by the boundary of the building with level structure and
$L_p$ by the ramified local factor. For $n=2$ and unramified data this is
Theorems~\ref{I:thm:partition}--\ref{I:thm:gap}.
\end{conjecture}

\section{Algorithmic \texorpdfstring{$K$}{K}-theory via Hecke operators}\label{I:sec:ktheory}

Throughout this section $\Gamma$ is torsion-free and cocompact, $Y=\Gamma\backslash X$ is
a finite connected $(q+1)$-regular graph with $n=|V_Y|$ vertices and $m=|E_Y|$ unoriented
edges, so that
\[
r:=\rank\Gamma=m-n+1,\qquad m=\tfrac{(q+1)n}{2},\qquad r-1=m-n=\tfrac{(q-1)n}{2}.
\]
$T_p\in M_n(\Z)$ denotes the Hecke (adjacency) operator of $Y$ and
$\Delta_p:=(q+1)I_n-T_p$ its \emph{Hecke--Laplacian}. By Ihara's theorem and the
Eichler--Shimura and Jacquet--Langlands correspondences, when $\Gamma$ is the unit group
of an Eichler order the spectrum of $T_p$ is
\[
\Spec T_p=\{q+1\}\ \cup\ \{a_p(f)\ :\ f\ \text{a weight-two newform of the relevant level}\},
\]
so the entries of $\Delta_p$ are Frobenius traces in disguise. This section turns that
disguise into an algorithm.

\subsection{The Hecke correspondence and the Pimsner sequence}

\begin{definition}[Hecke correspondence]\label{I:def:corr}
Let $A:=C(V_Y)\cong\C^n$ and let $\cE_p$ be the $A$--$A$ correspondence with underlying
space $C(E_Y)$, right action $(\xi\cdot a)(e)=\xi(e)a(t(e))$, left action
$(a\cdot\xi)(e)=a(o(e))\xi(e)$, and inner product
$\langle\xi,\eta\rangle(v)=\sum_{t(e)=v}\overline{\xi(e)}\eta(e)$. Then $\cE_p$ is
finitely generated projective, full, with injective left action, and its class in
$KK^0(A,A)=M_n(\Z)$ is
\[
[\cE_p]=T_p .
\]
Let $\widetilde A:=C_0(V_Y\times\Z)$ with the degree shift $S$, and let
$\widetilde\cE_p$ be the corresponding correspondence over $\widetilde A$ recording the
valuation of the determinant, so that
$[\widetilde\cE_p]=u\,T_p$ on $K_0(\widetilde A)=\Z[u^{\pm1}]^n$, $u:=[S]$.
\end{definition}

\begin{theorem}[Hecke--Pimsner sequence; equivariant PV]\label{I:thm:pimsner}
For any $A$--$A$ correspondence $\cE$ as above, the Cuntz--Pimsner algebra $\cO_\cE$ sits
in the six-term exact sequence
\[
\begin{array}{ccccc}
K_0(A)&\xrightarrow{\ 1-[\cE]\ }&K_0(A)&\longrightarrow&K_0(\cO_\cE)\\
\uparrow& & & &\downarrow\\
K_1(\cO_\cE)&\longleftarrow&K_1(A)&\xleftarrow{\ 1-[\cE]\ }&K_1(A)
\end{array}
\]
which for $A=\C^n$ collapses to the short exact sequence
\begin{equation}\label{I:eq:pimsnershort}
0\longrightarrow K_1(\cO_\cE)\longrightarrow\Z^n
\xrightarrow{\ (1-[\cE])^{t}\ }\Z^n\longrightarrow K_0(\cO_\cE)\longrightarrow0 .
\end{equation}
When $\cE$ is the correspondence of a single automorphism, \eqref{I:eq:pimsnershort} is the
Pimsner--Voiculescu sequence; when $A$ carries an action of a finite group $C$ and $\cE$
is $C$-equivariant, the same sequence holds with $\Z^n$ replaced by the
$R(C)$-module $K_0^C(A)$ and $1-[\cE]$ by its $R(C)$-linear extension.
\end{theorem}

\begin{definition}[Ihara module]\label{I:def:ihara}
The \emph{Ihara module} of $(\Gamma,p)$ is the $\Z[u^{\pm1}]$-module
\[
M_\Gamma(p):=\Z[u^{\pm1}]^{\,n}\big/\bigl(1-u\,T_p+q\,u^2\bigr)\Z[u^{\pm1}]^{\,n},
\]
the presentation matrix being the Hecke pencil of Lemma~\ref{I:lem:recursion}.
\end{definition}

\begin{proposition}[Bass's identity; the two channels]\label{I:prop:bass}
With $A^{\mathrm{nb}}_Y$ as in Definition~\ref{I:def:groupoid},
\begin{equation}\label{I:eq:bass}
\det\bigl(I_{2m}-u\,A^{\mathrm{nb}}_Y\bigr)
=(1-u^2)^{\,m-n}\,\det\bigl((1+qu^2)I_n-u\,T_p\bigr),
\end{equation}
so the $0$-th Fitting ideal of $M_\Gamma(p)$ is generated by the reciprocal
$Z_Y(u)^{-1}(1-u^2)^{-(m-n)}$ of the Ihara zeta function of $Y$. Identity
\eqref{I:eq:bass} has been verified symbolically for all graphs of
Table~\ref{I:tab:snf}. Specialising $u=1$:
\[
\bigl(1-uT_p+qu^2\bigr)\big|_{u=1}=(q+1)I_n-T_p=\Delta_p .
\]
\end{proposition}

Proposition~\ref{I:prop:bass} isolates the two channels available for computation. The
\emph{boundary channel} uses $A^{\mathrm{nb}}_Y$ and computes the $K$-theory of
$C(\partial X)\rtimes\Gamma$ directly. The \emph{Laplacian channel} uses $\Delta_p$, that
is the $u=1$ specialisation of the Ihara module. They are related by \eqref{I:eq:bass},
and --- this is the point of the section --- they carry completely different arithmetic.

\subsection{The structure theorem}

\begin{theorem}[Hecke--Laplacian structure theorem]\label{I:thm:structure}
Let $\Delta_p=(q+1)I_n-T_p$ and let
$\SNF(\Delta_p)=\mathrm{diag}(d_1\mid d_2\mid\dots\mid d_{n-1},0)$ be its Smith normal
form over $\Z$. Then:
\begin{enumerate}
\item[(i)] $\ker\Delta_p=\Z$ and
$\coker\Delta_p\cong\Z\oplus\bigoplus_{i=1}^{n-1}\Z/d_i\Z
=\Z\oplus\Jac(Y)$,
where $\Jac(Y)$ is the Jacobian $($critical, sandpile$)$ group of $Y$;
\item[(ii)] $\coker\Delta_p$ and $\ker\Delta_p$ are the $u=1$ specialisation
\[
M_\Gamma(p)\otimes_{\Z[u^{\pm1}]}\Z_{u=1}\cong\Z\oplus\Jac(Y),
\qquad
\mathrm{Tor}_1^{\Z[u^{\pm1}]}\bigl(M_\Gamma(p),\Z_{u=1}\bigr)\cong\Z,
\]
of the Ihara module of Definition~\ref{I:def:ihara}, and they are the integral index
groups of the \emph{Hecke--Dirac operator} of $Y$: choosing an orientation $E_Y^+$ and
writing $\partial\colon\Z^{E_Y^+}\to\Z^{V_Y}$, $\partial f=t(f)-o(f)$, for the simplicial
boundary map, the self-adjoint operator
\[
\mathsf{D}:=\begin{pmatrix}0&\partial\\ \partial^{\,t}&0\end{pmatrix}
\quad\text{on}\quad \ell^2(V_Y)\oplus\ell^2(E_Y^+),
\qquad
\mathsf{D}^2=\Delta_p\oplus\partial^{\,t}\partial ,
\]
has $\ker\mathsf{D}=\Z\oplus H_1(Y;\Z)=\Z\oplus\Z^{\,r}$, and the integral cokernel of
the even part $\mathsf{D}^2|_{\ell^2(V_Y)}=\Delta_p$ is $\Z\oplus\Jac(Y)$. \emph{No
existence claim for a $C^*$-algebra is made here}; see Remark~\ref{I:rem:noalgebra};
\item[(iii)] the order of the torsion is a product of Frobenius data: by Kirchhoff's
matrix-tree theorem and Ihara's theorem,
\begin{equation}\label{I:eq:torsionorder}
\bigl|\Tors K_0(\cO)\bigr|=|\Jac(Y)|=\kappa(Y)
=\frac1n\prod_{\lambda\in\Spec T_p\setminus\{q+1\}}\bigl(q+1-\lambda\bigr)
=\frac1n\prod_f\bigl(p+1-a_p(f)\bigr),
\end{equation}
$f$ running over the weight-two newforms occurring in $Y$; and if $E_f$ is the elliptic
curve attached to $f$ by Eichler--Shimura, then $p+1-a_p(f)=\#E_f(\F_p)$, so
\[
\bigl|\Tors K_0(\cO)\bigr|=\frac1n\prod_f\#E_f(\F_p).
\]
\end{enumerate}
The Smith normal form computes not merely the order but the group: the elementary
divisors $d_1\mid\dots\mid d_{n-1}$ are the invariant factors of $\Jac(Y)$.
\end{theorem}

\begin{proof}[Proof sketch]
(i) is the standard theory of the critical group: $\Delta_p$ is the Laplacian of the
$(q+1)$-regular graph $Y$, symmetric with kernel the constants, and its cokernel splits
as $\Z\oplus K(Y)$ with $K(Y)$ the critical group. (ii) is
Theorem~\ref{I:thm:pimsner} applied with index map $\Delta_p$. In (iii), Kirchhoff gives
$\kappa(Y)=\frac1n\prod_{i\ge2}\mu_i$ for the non-zero Laplacian eigenvalues
$\mu_i=(q+1)-\lambda_i$, $|K(Y)|=\kappa(Y)$, and Ihara--Eichler--Jacquet--Langlands
identifies $\Spec T_p\setminus\{q+1\}$ with the $a_p(f)$. The last identity is the
definition of $a_p$.
\end{proof}

\begin{remark}[Bulk and boundary: what is, and is not, an algebra]\label{I:rem:noalgebra}
Since $A=\C^n$ we have $KK^0(A,A)\cong M_n(\Z)$, so
\[
x_p:=[\cE_p]-q\cdot 1_A=T_p-qI_n\ \in\ KK^0(A,A)
\]
is a bona fide Kasparov class, and $1-x_p=\Delta_p$. What fails is not the class but its
\emph{realisation}: a class is the class of a $C^*$-correspondence only if its matrix has
non-negative entries, and $T_p-qI_n$ has $-q$ on the diagonal. Hence $x_p$ is a formal
difference of correspondences, no Cuntz--Pimsner algebra $\cO_\cE$ has index map
$\Delta_p$, and we assert none. The asymmetry between the two channels is therefore not
only arithmetic but ontological:
\begin{center}\small
\begin{tabular}{lll}
\toprule
& realisation & invariant\\
\midrule
boundary channel & a $C^*$-algebra, $\cO_{A^{\mathrm{nb}}_Y}\simeq C(\partial X)\rtimes\Gamma$
& $\Z^r\oplus\Z/(r-1)$, no arithmetic\\
bulk (Laplacian) channel & a $\Z[u^{\pm1}]$-module and an index & $\Z\oplus\Jac(Y)$, all the arithmetic\\
\bottomrule
\end{tabular}
\end{center}
Everything asserted in Theorem~\ref{I:thm:structure} is unconditional at the level of the
Ihara module, of the index of $\mathsf{D}$, and of the integers; what is missing is an
honest $C^*$-algebra whose $K_0$ is $\Z\oplus\Jac(Y)$. Producing one --- equivalently,
realising the virtual correspondence $x_p$ by a genuine, possibly non-unital or
$\Z/2$-graded, construction --- is the main open problem of this paper.

There is a purely algebraic first step, and it isolates exactly where the difficulty
lives. The Ihara pencil $1-uT_p+qu^2$ of Definition~\ref{I:def:ihara} is a quadratic
matrix polynomial, and it linearises in the standard way over $A^{\oplus2}$,
$A=\C^n$: setting
\[
C:=\begin{pmatrix}T_p&-qI_n\\ I_n&0\end{pmatrix}\ \in\ M_{2n}(\Z),
\qquad
\det\bigl(I_{2n}-uC\bigr)=\det\bigl(I_n-uT_p+qu^2I_n\bigr),
\]
a companion-matrix identity verified directly by block Gaussian elimination on
$I-uC$. Thus $C$ is a single, honest $\Z$-linear class in $KK^0(A^{\oplus2},A^{\oplus2})$
representing the whole Ihara pencil at once, with $[\cE_p]=T_p$ visible in its upper
left block; the obstruction of the previous paragraph has been moved, not removed, since
$C$ still has the non-positive block $-qI_n$ in its upper right corner and so is still
not the matrix of a $C^*$-correspondence over $A^{\oplus2}$.

The problem left for the sequel is accordingly sharp: \emph{dilate} $C$, in the sense of
dilation theory, to a genuinely non-negative matrix over a larger vertex algebra
$\widetilde A\supset A^{\oplus2}$ --- for instance by the Drinen--Tomforde
desingularisation of a directed graph carrying $C$ as a signed incidence matrix, with the
negative entries reinterpreted as sinks absorbing mass at rate $q$ per vertex --- in such
a way that the resulting Cuntz--Pimsner (or graph) algebra $\cO_{\widetilde C}$ receives
$A^{\oplus2}$ as a full corner and has $K_0(\cO_{\widetilde C})\cong\Z\oplus\Jac(Y)$ by
Morita reduction back to $\Delta_p$. We do not carry this out here; it is the intended
content of a sequel to this paper, and the companion linearisation above is offered as
its starting point.
\end{remark}

\begin{theorem}[Boundary channel is arithmetically blind]\label{I:thm:blind}
Let $Y$ be any finite connected $(q+1)$-regular graph and $\Gamma$ the corresponding
torsion-free cocompact lattice, $r=\rank\Gamma=m-n+1\ge2$. Then
\[
K_0\bigl(C(\partial X)\rtimes\Gamma\bigr)\cong\Z^{\,r}\oplus\Z/(r-1)\Z,
\qquad
K_1\bigl(C(\partial X)\rtimes\Gamma\bigr)\cong\Z^{\,r},
\qquad r=\rank\Gamma=\tfrac{(q-1)n}{2}+1 .
\]
The torsion of the boundary quotient therefore depends only on the Euler characteristic
$r-1=|E_Y|-|V_Y|$ of $Y$, hence only on the isomorphism class of $\Gamma$ as an abstract
free group: \emph{no Hecke eigenvalue is visible in the boundary channel.}
\end{theorem}

\begin{proof}
Write $A^{\mathrm{nb}}=T^{t}S-J$, where $S,T\in M_{n\times2m}(\Z)$ are the incidence
matrices $S_{v,e}=[o(e)=v]$, $T_{v,e}=[t(e)=v]$ and $J$ is the edge-reversal involution,
and put $M:=I_{2m}-A^{\mathrm{nb}}=I+J-T^{t}S$. Since $\SNF(M)=\SNF(M^{t})$ it suffices
to compute $\coker M=\Z^{E_Y}/M\Z^{E_Y}$, and by Theorem~\ref{I:thm:pimsner} this is
$K_0$, with $K_1=\ker M^t$.

Reading off the columns of $M$, the presentation of $\coker M$ has generators $[e]$,
$e\in E_Y$, and one relation for each $f\in E_Y$,
\begin{equation}\label{I:eq:blindrel}
[f]+[\bar f]=c_{o(f)},\qquad\text{where}\quad c_v:=\sum_{t(e)=v}[e].
\end{equation}
Applying \eqref{I:eq:blindrel} to $f$ and to $\bar f$ gives $c_{o(f)}=c_{t(f)}$ for every
edge, so $c_v=:c$ is independent of $v$ because $Y$ is connected. Fix an orientation
$E_Y^+$, use $[\bar f]=c-[f]$ to eliminate the reversed generators, and let
$d^-(v):=\#\{f\in E_Y^+:o(f)=v\}$. Expanding $c_v$ in the surviving generators turns the
$n$ remaining relations into
\[
R_v:\qquad \bigl(\partial^{\,t}\text{-column at }v\bigr)\ :\
\sum_{f\in E_Y^+,\,t(f)=v}[f]-\sum_{f\in E_Y^+,\,o(f)=v}[f]
=\bigl(1-d^-(v)\bigr)\,c .
\]
Hence $\coker M\cong\coker B$ for the $(m+1)\times n$ integer matrix
\[
B=\begin{pmatrix}\partial^{\,t}\\ -\lambda^{t}\end{pmatrix},
\qquad \lambda_v:=1-d^-(v),
\]
on generators $E_Y^+\cup\{c\}$. Projecting away the last coordinate gives the exact
sequence
\[
0\longrightarrow\langle c\rangle\longrightarrow\coker B
\longrightarrow\coker\bigl(\partial^{\,t}\bigr)=H^1(Y;\Z)\cong\Z^{\,r}
\longrightarrow0 ,
\]
$H^1$ of a graph being free of rank $r$. Finally $t\,c=0$ in $\coker B$ if and only if
$(0,\dots,0,t)\in\operatorname{im}B$, i.e. if and only if there is $x\in\Z^{V_Y}$ with
$\partial^{\,t}x=0$ and $-\lambda\cdot x=t$; the first condition forces $x=a\mathbf1$,
and then
\[
-\lambda\cdot x=-a\sum_{v}\bigl(1-d^-(v)\bigr)=-a\,(n-m)=a\,(r-1),
\]
so $c$ has order exactly $r-1$. The quotient $\Z^r$ being free, the sequence splits and
$\coker M\cong\Z^{r}\oplus\Z/(r-1)$; $\ker M^t$ is free of rank $2m-\rank M=r$.
\end{proof}

\begin{remark}[Where the spectrum goes]\label{I:rem:wherespectrum}
The proof shows exactly how the Hecke data disappear, and it is worth recording that the
mechanism is \emph{not} the one suggested by Bass's identity \eqref{I:eq:bass}: at $u=1$
both factors of \eqref{I:eq:bass} vanish, so \eqref{I:eq:bass} degenerates to $0=0$ and
carries no integral information there. What kills the spectrum is the local relation
\eqref{I:eq:blindrel}: it forces the $n$ a priori distinct classes $c_v$ to collapse to a
single class $c$, after which the only surviving numerical invariant is
$\sum_v(1-d^-(v))=n-m$. The orientation-dependent quantities $d^-(v)$ do not survive
individually --- only their sum does --- and the sum is the Euler characteristic. The
adjacency matrix $T_p$ never enters the computation at all.
\end{remark}

\begin{remark}[The dichotomy]\label{I:rem:dichotomy}
Theorems~\ref{I:thm:structure} and \ref{I:thm:blind} together are the main structural finding
of this paper. The boundary quotient --- the natural non-commutative object, the simple
purely infinite Cuntz--Krieger algebra of the tail relation --- has $K$-theory of purely
topological content. The arithmetic sits one specialisation away, in the Hecke pencil
$1-uT_p+qu^2$ at $u=1$; equivalently, the arithmetic is the \emph{first-order} datum of
the Ihara zeta function at its trivial zero. This is the exact non-commutative analogue
of the classical statement that $\zeta_Y(u)$ has a trivial factor $(1-u^2)^{m-n}$ and
that the interesting arithmetic is in the reduced factor.
\end{remark}

\begin{corollary}[Frobenius traces as torsion]\label{I:cor:frobenius}
Let $Y$ carry a single non-trivial Hecke eigenvalue $a_p$ with multiplicity $n-1$. Then
$\Jac(Y)\cong(\Z/(p+1-a_p))^{\,n-2}$ whenever the Laplacian is of the form
$cI-J$, and in general $|\Jac(Y)|=\frac1n\prod_f\#E_f(\F_p)$ exactly. In particular the
Frobenius trace $a_p$ is recovered from $K_0$ by
\[
a_p=p+1-\bigl(n\cdot|\Tors K_0|\bigr)^{1/(n-1)}
\]
in the eigenvalue-homogeneous case, and in general from the elementary divisors by the
inverse of the Smith algorithm.
\end{corollary}

\begin{corollary}[Ramanujan bound; strict, and never attained]\label{I:cor:ramanujan}
If $Y$ is Ramanujan, i.e. $|\lambda|\le2\sqrt q$ for every
$\lambda\in\Spec T_p\setminus\{\pm(q+1)\}$, then
\[
\bigl|\Tors K_0\bigr|\ \le\ \frac1n\bigl(q+1+2\sqrt q\bigr)^{\,n-1}
=\frac1n\bigl(1+\sqrt q\bigr)^{2(n-1)},
\]
so that, with $B(q):=2\log(1+\sqrt q)$,
\[
\limsup_{n\to\infty}\frac1n\log\bigl|\Tors K_0\bigr|\ \le\ B(q).
\]
The bound is unconditional on the Ramanujan locus, but it is \emph{never} an equality,
and the deficit is exponential in $n$: see Theorem~\ref{I:thm:bstorsion} and
Remark~\ref{I:rem:bsnumerics}.
\end{corollary}

\begin{theorem}[Benjamini--Schramm asymptotic torsion]\label{I:thm:bstorsion}
Let $(Y_j)_{j\ge1}$ be finite connected $(q+1)$-regular graphs, $n_j=|V_{Y_j}|\to\infty$,
which Benjamini--Schramm converge to the $(q+1)$-regular tree $T_{q+1}$ --- for instance
any family of girth tending to infinity, such as the Lubotzky--Phillips--Sarnak graphs
$X^{p,\ell}$ at fixed $p=q$. Let $\cO_j$ be the associated Hecke--Laplacian invariant of
Theorem~\ref{I:thm:structure}. Then the torsion density converges, and its limit is Lyons'
tree entropy of $T_{q+1}$:
\begin{equation}\label{I:eq:bslimit}
\lim_{j\to\infty}\frac1{n_j}\log\bigl|\Tors K_0(\cO_j)\bigr|
\;=\;h(q)\;:=\;q\log q-\frac{q-1}{2}\log\bigl(q^2-1\bigr)
\;=\;\log\frac{q^{\,q}}{(q^2-1)^{(q-1)/2}} ,
\end{equation}
equivalently $h(q)=\int_{-2\sqrt q}^{2\sqrt q}\log\bigl(q+1-x\bigr)\,d\mu_{\mathrm{KM}}(x)$
against the Kesten--McKay measure
\[
d\mu_{\mathrm{KM}}(x)=\frac{(q+1)\sqrt{4q-x^2}}{2\pi\bigl((q+1)^2-x^2\bigr)}\,dx ,
\]
the spectral measure of the simple random walk on $T_{q+1}$ \cite{Kesten1959,McKay1981}.
Moreover $h(q)<B(q)$ for every $q$, with
\[
\frac{B(q)}{h(q)}=2.1061,\ 1.3889,\ 1.2340,\ 1.0957,\ 1.0089
\quad\text{for}\quad q=2,5,9,29,999.
\]
\end{theorem}

\begin{corollary}[The deficit, to two orders]\label{I:cor:deficit}
As $q\to\infty$,
\begin{equation}\label{I:eq:deficit}
B(q)-h(q)\;=\;2q^{-1/2}-\tfrac32q^{-1}+O\bigl(q^{-3/2}\bigr) ,
\end{equation}
an asymptotic equality, not a numerical fit. Consequently, on graphs with $n$
vertices, the Ramanujan bound overstates $|\Tors K_0|$ by the factor
\[
\exp\bigl((B(q)-h(q))\,n\bigr)
\;=\;
\exp\bigl(2n\,q^{-1/2}+O(n\,q^{-1})\bigr) ;
\]
at $q=2$, $n=2000$ this factor is $e^{1852}$.
\end{corollary}

\begin{proof}
Write $h(q)=\log q-\tfrac{q-1}2\log(1-q^{-2})$ and expand
$-\log(1-q^{-2})=q^{-2}+\tfrac12q^{-4}+O(q^{-6})$, so
\[
h(q)=\log q+\tfrac{q-1}2\bigl(q^{-2}+O(q^{-4})\bigr)
=\log q+\tfrac12q^{-1}-\tfrac12q^{-2}+O(q^{-3}) .
\]
Write $B(q)=\log q+2\log(1+q^{-1/2})$ and expand
$\log(1+q^{-1/2})=q^{-1/2}-\tfrac12q^{-1}+\tfrac13q^{-3/2}+O(q^{-2})$, so
\[
B(q)=\log q+2q^{-1/2}-q^{-1}+\tfrac23q^{-3/2}+O(q^{-2}) .
\]
Subtracting, the $\log q$ terms cancel and
$B(q)-h(q)=2q^{-1/2}-q^{-1}-\tfrac12q^{-1}+O(q^{-3/2})
=2q^{-1/2}-\tfrac32q^{-1}+O(q^{-3/2})$, which is \eqref{I:eq:deficit}. (A direct
computer-algebra expansion confirms the next two orders as well, and that the
$q^{-2}$ coefficient vanishes identically --- the $-\tfrac12q^{-2}$ term present in
$B(q)$ and in $h(q)$ separately cancels exactly on subtraction:
$B(q)-h(q)=2q^{-1/2}-\tfrac32q^{-1}+\tfrac23q^{-3/2}+\tfrac25q^{-5/2}+O(q^{-3})$.) The
numerical values of Remark~\ref{I:rem:bsnumerics} and of the ratios above were computed
independently of this expansion and agree with it: at $q=10^4$, \eqref{I:eq:deficit}
predicts $0.02-0.00015=0.01985$ against the quadrature value $0.019852$.
\end{proof}

\begin{proof}[Proof]
By Theorem~\ref{I:thm:structure}(iii), $|\Tors K_0(\cO_j)|=\kappa(Y_j)$, so
\eqref{I:eq:bslimit} is Lyons' theorem on the asymptotics of the number of spanning trees
along Benjamini--Schramm convergent sequences of bounded degree \cite{Lyons2005},
applied to the limit $T_{q+1}$, together with Kirchhoff's formula in the form
$\frac1n\log\kappa=\frac1n\bigl[\sum_{i\ge2}\log(q+1-\lambda_i)-\log n\bigr]$ and the
convergence of the empirical spectral measures of $T_p$ to $\mu_{\mathrm{KM}}$. The
closed form and the integral representation of $h(q)$ agree to six decimals for
$q+1\in\{3,4,6,8,14,30\}$ by numerical quadrature. Strictness of $h<B$ is the strict
concavity of $\log$ together with the fact that $\mu_{\mathrm{KM}}$ is not a point mass
at the spectral edge.
\end{proof}

\begin{remark}[Why the Ramanujan bound is the wrong instrument]\label{I:rem:bsnumerics}
Corollary~\ref{I:cor:ramanujan} replaces the whole spectrum by its extreme point
$2\sqrt q$. But $\mu_{\mathrm{KM}}$ vanishes like a square root at $\pm2\sqrt q$ and puts
almost no mass there, so the bound is a supremum where the truth is an integral; this,
and not any slack in the choice of family, is the source of the exponential deficit of
Theorem~\ref{I:thm:bstorsion}.

The deeper point is that the Ramanujan property is neither necessary nor sufficient for
the limit \eqref{I:eq:bslimit}, and that the invariant which does control the torsion
density is Benjamini--Schramm convergence. Both halves are visible in exact computation.
Below, $\lambda_2$ is the largest non-trivial eigenvalue in absolute value,
$\mathrm{KS}$ the Kolmogorov--Smirnov distance from the empirical spectral distribution
of $T_p$ to $\mu_{\mathrm{KM}}$, and the last two columns compare the measured torsion
density with $h(q)$ and with the Ramanujan bound $B(q)$.
\begin{center}\small
\renewcommand{\arraystretch}{1.15}\setlength{\tabcolsep}{5pt}
\begin{tabular}{lccccccc}
\toprule
$Y$ & $n$ & $q$ & $\lambda_2$ & Ram.\ ? & $\mathrm{KS}$
& ratio to $h(q)$ & ratio to $B(q)$\\
\midrule
random $3$-regular   & $200$  & $2$ & $2.8315$ & no  & $0.0142$ & $0.9774$ & $0.4641$\\
random $3$-regular   & $1000$ & $2$ & $2.8261$ & yes & $0.0032$ & $0.9937$ & $0.4719$\\
random $3$-regular   & $2000$ & $2$ & $2.8239$ & yes & $0.0021$ & $0.9964$ & $0.4731$\\
$\mathrm{PG}(2,5)$   & $62$   & $5$ & $2.2361$ & yes & $0.2661$ & $0.9669$ & $0.6962$\\
LPS $X^{5,13}$       & $2184$ & $5$ & $4.2497$ & yes & $0.0375$ & $0.9982$ & $0.7187$\\
LPS $X^{5,17}$       & $4896$ & $5$ & $4.3089$ & yes & $0.0307$ & $0.9991$ & $0.7194$\\
\bottomrule
\end{tabular}
\end{center}
Two rows falsify the spectral-gap reading. The random $3$-regular graph on $200$
vertices has $\lambda_2=2.8315>2\sqrt2=2.8284$ and is therefore \emph{not} Ramanujan,
yet its torsion density is already within $2.3\%$ of $h(2)$. The incidence graph of
$\mathrm{PG}(2,5)$ is Ramanujan with room to spare, $\lambda_2=\sqrt5=2.2361$ against
the permitted $2\sqrt5=4.4721$, yet it has girth $6$, its spectral distribution is far
from Kesten--McKay ($\mathrm{KS}=0.27$), and its torsion density falls $3.3\%$
\emph{below} $h(5)$. The $\mathrm{KS}$ column and the $h$-ratio column move together
throughout; the $\lambda_2$ column does not track either.

One further datum belongs here, because it bears on Remark~\ref{I:rem:ghost} and on the
statistics of arithmetic spectra generally. The graph $X^{5,13}$ has $2184$ vertices but
only $55$ distinct eigenvalues, of maximal multiplicity $96$, whereas a random
$3$-regular graph on $2000$ vertices has $2000$ distinct eigenvalues. The degeneracy is
forced by the $\mathrm{PGL}_2(\F_{13})$ representation theory through which the Hecke
operator factors, and it is why the empirical spectral distribution of an LPS graph is a
coarse staircase whose $\mathrm{KS}$ distance decays more slowly than that of a random
graph of half the size. It is also the mechanism behind the Poisson --- not
Gaussian-unitary --- level statistics observed in arithmetic quantum chaos: large Hecke
symmetry produces large degeneracy, and degeneracy is the opposite of level repulsion.
\end{remark}

\subsection{The algorithm, and its verification}

\begin{algorithm}[Hecke $\to$ torsion]\label{I:alg:main}
\textbf{Input:} a prime $p$ (so $q=p$), a torsion-free cocompact $\Gamma\le\PGL_2(\Q_p)$
presented by its quotient graph $Y=\Gamma\backslash X$, i.e. by the integer matrix
$T_p\in M_n(\Z)$.
\begin{enumerate}
\item[(1)] Form $\Delta_p:=(q+1)I_n-T_p\in M_n(\Z)$.
\item[(2)] Compute the Smith normal form over $\Z$,
$U\Delta_pW=\mathrm{diag}(d_1,\dots,d_{n-1},0)$ with $d_1\mid\dots\mid d_{n-1}$
(cost $O(n^3)$ ring operations with the usual bit-length control).
\item[(3)] Output $K_0=\Z\oplus\bigoplus_i\Z/d_i$, $K_1=\Z$, and
$|\Tors K_0|=\prod_id_i=\kappa(Y)$.
\item[(4)] \emph{Certificate.} Independently compute
$\frac1n\prod_{\lambda\neq q+1}(q+1-\lambda)$ from the characteristic polynomial of
$T_p$, and $\frac1n\prod_f(p+1-a_p(f))$ from the Hecke eigenvalues; the three numbers
must agree. Optionally verify Bass's identity \eqref{I:eq:bass} symbolically as a check on
the graph data.
\end{enumerate}
\textbf{Output:} the torsion subgroup of $K_0$ as an explicit finite abelian group,
together with its factorisation into Frobenius traces.
\end{algorithm}

\begin{table}[t]
\centering\small
\renewcommand{\arraystretch}{1.15}\setlength{\tabcolsep}{5pt}
\begin{tabular}{lccccc c}
\toprule
$Y$ & $q$ & $n$ & $\Jac(Y)=\Tors\coker\Delta_p$ & $\kappa(Y)$ & boundary $\Tors K_0$ & $r-1$\\
\midrule
$K_4$              &2& 4&$(\Z/4)^2$                   &$16$    &$\Z/2$&2\\
$K_5$              &3& 5&$(\Z/5)^3$                   &$125$   &$\Z/5$&5\\
$K_6$              &4& 6&$(\Z/6)^4$                   &$1296$  &$\Z/9$&9\\
$K_{3,3}$          &2& 6&$\Z/3\oplus\Z/3\oplus\Z/9$   &$81$    &$\Z/3$&3\\
$Q_3$ (cube)       &2& 8&$\Z/2\oplus\Z/8\oplus\Z/24$  &$384$   &$\Z/4$&4\\
prism $C_4\times K_2$&2& 8&$\Z/2\oplus\Z/8\oplus\Z/24$&$384$   &$\Z/4$&4\\
Wagner $V_8$       &2& 8&$\Z/7\oplus\Z/56$            &$392$   &$\Z/4$&4\\
prism $C_5\times K_2$&2&10&$\Z/19\oplus\Z/95$         &$1805$  &$\Z/5$&5\\
M\"obius $V_{10}$  &2&10&$\Z/11\oplus\Z/165$          &$1815$  &$\Z/5$&5\\
Petersen           &2&10&$\Z/2\oplus(\Z/10)^3$        &$2000$  &$\Z/5$&5\\
Heawood            &2&14&$(\Z/7)^4\oplus\Z/21$        &$50421$ &$\Z/7$&7\\
\bottomrule
\end{tabular}
\medskip
\caption{Exact Smith-normal-form computations. Columns 4--5: the Laplacian channel,
Theorem~\ref{I:thm:structure}; the order always equals
$\frac1n\prod_{\lambda\neq q+1}(q+1-\lambda)$. Columns 6--7: the boundary channel,
Theorem~\ref{I:thm:blind}; the torsion is always cyclic of order $r-1=(q-1)n/2$,
independently of the spectrum, confirming the proof given there. Bass's identity
\eqref{I:eq:bass} was verified symbolically in the first five rows.}
\label{I:tab:snf}
\end{table}

\begin{example}[Level $37$, $p=2$: two elliptic curves, one torsion group]\label{I:ex:37}
The space $S_2(\Gamma_0(37))$ is two-dimensional, spanned by the newforms of the
elliptic curves $37a: y^2+y=x^3-x$ (rank one) and $37b: y^2+y=x^3+x^2-23x-50$ (rank
zero), with $a_2(37a)=-2$ and $a_2(37b)=0$. Correspondingly
\[
\#E_{37a}(\F_2)=2+1-(-2)=5,\qquad \#E_{37b}(\F_2)=2+1-0=3,
\]
and the Hecke--Laplacian of the associated $3$-regular Hecke system has non-zero
eigenvalues $5$ and $3$, so \eqref{I:eq:torsionorder} gives
$|\Tors K_0|=\tfrac13\cdot5\cdot3=5$. The Frobenius traces of two distinct elliptic
curves precipitate into a single cyclic group of order five.
\end{example}

\begin{example}[Level $11$, $p=3$]\label{I:ex:11}
For the definite quaternion algebra of discriminant $11$ the Brandt matrix at $p=3$ is
$B(3)=\begin{psmallmatrix}2&2\\3&1\end{psmallmatrix}$, with spectrum $\{4,-1\}$: the
Eisenstein eigenvalue $q+1=4$ and $a_3(11a)=-1$, matching
$\eta(z)^2\eta(11z)^2=q-2q^2-q^3+2q^4+\cdots$. Then
$q+1-a_3=5=\#E_{11a}(\F_3)=|E_{11a}(\Q)_{\mathrm{tors}}|$.
\end{example}

\begin{remark}[Weighted quotients]\label{I:rem:weights}
Examples~\ref{I:ex:37} and \ref{I:ex:11} are stated at the level of the eigenvalue identity
$q+1-a_p=\#E(\F_p)$, which is unconditional. The passage to the critical group in
Theorem~\ref{I:thm:structure}(iii) requires $\Gamma$ torsion-free, so that $Y$ is a genuine
$(q+1)$-regular graph; Brandt and supersingular isogeny graphs are quotients by groups
with torsion, i.e. graphs of groups with vertex weights $w_i=|\Aut|/2$, and for them
Theorem~\ref{I:thm:structure} holds in its weighted form, with $\Delta_p$ the weighted
Laplacian $w_iB_{ij}=w_jB_{ji}$ and Kirchhoff's theorem replaced by its weighted
version. Congruence subgroups of sufficiently high level are torsion-free and remove the
issue at the cost of enlarging $n$.

The structurally correct home for the weighted case is orbifold, i.e. groupoid,
$K$-theory. A graph of groups $(Y,\Gamma_v,\Gamma_e)$ in the sense of Bass--Serre is the
same datum as an \'etale groupoid $\cG_{Y}$ with unit space $V_Y$ and isotropy
$\Gamma_v$; the object to compute is $K_*(C^*(\cG_Y))$, and the Bass--Serre
Mayer--Vietoris sequence
\[
\cdots\to\bigoplus_{e\in E_Y}K_*\bigl(C^*(\Gamma_e)\bigr)
\xrightarrow{\ \partial\ }\bigoplus_{v\in V_Y}K_*\bigl(C^*(\Gamma_v)\bigr)
\to K_*\bigl(C^*(\cG_Y)\bigr)\to\cdots
\]
replaces $\Z^{E_Y}\to\Z^{V_Y}$ by a sequence of representation rings
$R(\Gamma_v)=K_0(C^*(\Gamma_v))$, of ranks $|\Irr\Gamma_v|$. The vertex weights are then
not an ad hoc normalisation but the ranks of these isotropy contributions: the composite
$\partial\partial^{\,t}$ is the weighted Laplacian $w_iB_{ij}=w_jB_{ji}$ with
$w_i=|\Aut|/2$, the Kirchhoff count acquires the factor $\prod_vw_v^{-1}$ of the
Eichler mass formula, and the apparently fractional indices --- the value
$\tfrac12(4-a_3)$ of Example~\ref{I:ex:11} is not an integer --- become integral once read
in the equivariant group of Proposition~\ref{I:prop:equivariant} rather than in $\Z$. The
$\GL_2$ analogue of Paper XXXVIII's entanglement law belongs here, and we do not develop
it in this paper.
\end{remark}

\subsection{The equivariant refinement}

\begin{proposition}[Equivariant Hecke $K$-theory]\label{I:prop:equivariant}
Let $C$ be a finite group acting on $Y$ by graph automorphisms commuting with $T_p$ ---
in the arithmetic case, $C=\Gal(H/K)$ acting through the class group, as in Paper
XXXVIII. Then $K_0^C(A)=R(C)^{\,n_C}$ for the orbit decomposition, the index map of
Theorem~\ref{I:thm:pimsner} is the $R(C)$-linear extension of $\Delta_p$, and
\[
K_0^C(\cO)\cong\coker_{R(C)}\bigl(\Delta_p\otimes_\Z R(C)\bigr),
\qquad
K_1^C(\cO)\cong\ker_{R(C)}\bigl(\Delta_p\otimes_\Z R(C)\bigr).
\]
The entanglement law of Paper XXXVIII transfers verbatim: free $\Z[C]$-strata of the
vertex module contribute once by Shapiro's lemma, while strata with trivial or sign
action are doubled or sign-twisted by $R(C)$; the orders of the invariant factors are
unchanged, only their representation labels.
\end{proposition}

\subsection{What the two channels say about Langlands}

We summarise the bridge as a table of correspondences, each entry of which is a theorem
or a definition above, not an analogy.

\begin{center}\small
\renewcommand{\arraystretch}{1.2}
\begin{tabular}{lll}
\toprule
Automorphic side & Operator-algebraic side & Reference\\
\midrule
unramified $\pi_p$, Satake $(\alpha_{p,1},\alpha_{p,2})$
& Radon--Nikodym cocycle $e^{-\beta c_\pi}$ & Def.~\ref{I:def:dynamics}\\
local $L$-factor $L_p(\beta,\pi_p)$
& partition function $Z_p(\beta;\pi)$ & Thm~\ref{I:thm:partition}\\
largest pole of $L_p$
& critical inverse temperature $\beta_c(\pi_p)$ & Thm~\ref{I:thm:partition}\\
temperedness (Ramanujan) at $p$
& $\beta_c(\pi_p)=w/2$ & Thm~\ref{I:thm:partition}\\
Eisenstein locus
& the only locus carrying $\KMS$ \emph{states} & Thm~\ref{I:thm:perron}\\
Rankin--Selberg $L(s,\pi\times\widetilde{\pi'})$
& Kakutani/Hellinger obstruction $\Xi_\beta$ & Thm~\ref{I:thm:kakutani}\\
strong multiplicity one
& mutual singularity of phases & Cor.~\ref{I:cor:multone}\\
motivic weight $w$
& width of the equilibrium/obstruction gap & Thm~\ref{I:thm:gap}\\
Hecke recursion $T T_n=T_{n+1}+qT_{n-1}$
& Ihara module $\Z[u^{\pm}]^n/(1-uT_p+qu^2)$ & Def.~\ref{I:def:ihara}\\
Ihara zeta $Z_Y(u)$
& Fitting ideal of the Ihara module & Prop.~\ref{I:prop:bass}\\
Frobenius trace $a_p(f)$, $\#E_f(\F_p)$
& invariant factors of $\SNF(\Delta_p)$ & Thm~\ref{I:thm:structure}\\
Euler characteristic of $Y$ (no arithmetic)
& torsion of the boundary quotient & Thm~\ref{I:thm:blind}\\
\bottomrule
\end{tabular}
\end{center}

\medskip
The reading of the table is the thesis of the paper. Along the top half --- the
thermodynamic rows --- the correspondence is exact but weight-zero: it reproduces class
field theory and Artin data, and Theorem~\ref{I:thm:gap} measures precisely how far it
falls short above weight zero. Along the bottom half --- the $K$-theoretic rows --- the
correspondence is insensitive to temperature and to positivity, and it does carry
weight-two data, but only after the specialisation $u=1$ that separates the Laplacian
channel from the boundary channel. An algorithmic non-commutative class field theory for
$\GL_2$ is therefore not one bridge but two, and Bass's identity \eqref{I:eq:bass} is the
only thing that connects them.

\section*{Acknowledgement of provenance and caveats}

This is an unrefereed preprint and no priority check against the literature has been
carried out. The following are used as known and are not proved here: the Satake
isomorphism and the Hecke recursion; Ihara's theorem and Bass's determinant identity;
Kirchhoff's matrix-tree theorem and the theory of the critical group; the
Eichler--Shimura and Jacquet--Langlands correspondences; Kakutani's dichotomy for product
measures; Jacquet--Shalika non-vanishing for Rankin--Selberg $L$-functions; Renault's
theory of groupoid C*-algebras and the correspondence between $\KMS$ states and conformal
measures; Cuntz--Krieger and Cuntz--Pimsner $K$-theory and Pimsner's exact sequence. The
computations of Table~\ref{I:tab:snf}, and the spectral and torsion-density computations of
Remark~\ref{I:rem:bsnumerics}, are exact or numerically controlled and were carried out
independently of the statements they confirm. Theorem~\ref{I:thm:bstorsion} rests on
Lyons' spanning-tree asymptotics \cite{Lyons2005}, which we quote and do not prove, and
Remark~\ref{I:rem:bsnumerics} on the Kesten--McKay spectral distribution
\cite{Kesten1959,McKay1981}. The one substantive gap left open is the one isolated in
Remark~\ref{I:rem:noalgebra}: the Laplacian channel is a module and an index, not yet a
$C^*$-algebra.

% generated from Paper 2/anccft2.tex
\chapter{Algorithmic non-commutative class field theory, II: the trace obstruction and the K-homological bulk}\label{chap:II}
\chaptermark{II}
\begin{chapterabstract}
Paper I isolated a virtual $KK$-class $x_p=T_p-qI\in KK^0(\C^n,\C^n)=M_n(\Z)$, arising
from the Hecke--Laplacian $\Delta_p=(q+1)I-T_p=1-x_p$ of a finite $(q+1)$-regular graph
$Y$, and left open whether $x_p$ could be realized by a genuine, non-negative
$C^*$-correspondence, so that $\Delta_p$'s cokernel $\Z\oplus\Jac(Y)$ --- the critical
group of $Y$, whose order is a product of Frobenius traces $\prod(p+1-a_p(f))$ ---
would be the $K_0$-group of an honest Cuntz--Pimsner algebra. This paper proves that no
such realization exists, by an elementary and then a deeper obstruction, and identifies
where the arithmetic must live instead.

The companion linearization $C=\begin{psmallmatrix}T_p&-qI\\I&0\end{psmallmatrix}$,
exactly row/column-equivalent to $\Delta_p$ over $\Z$, has
$\Tr(C)=0$, $\Tr(C^2)=-n(q-1)$, strictly negative for every $q\ge2$ and every
$(q+1)$-regular graph. Since $\Tr(B^2)\ge0$ for any non-negative matrix $B$, and trace is
a similarity invariant, \textbf{no matrix of the same size and the same spectrum as $C$
can be non-negative} (Theorem~\\ref{II:thm:tracesame}) --- an unconditional, four-line fact.
Every dilation to a \emph{larger} non-negative matrix, of the kind an
auxiliary-vertex or Drinen--Tomforde construction would attempt, is ruled out as well,
and unconditionally: shift equivalence preserves the nonzero spectrum with
multiplicities --- proved here by an elementary two-way injection of generalized
eigenspaces (Lemma~\\ref{II:lem:se}) --- hence preserves $\Tr(\,\cdot\,^2)$, closing the
realization problem of Paper I, Remark~3.11, in the negative for matrices of every size
(Theorem~\\ref{II:thm:traceany}).

We then show that no repair by $\Z/2$-grading is available: a formal difference of two
positive correspondences is not itself a correspondence, and no Cuntz--Pimsner-type
functor converts a virtual difference into a single algebra realizing that difference's
index (Section~\\ref{II:sec:barrier}); if one existed, it would trivialize the realization
problem for every virtual $K$-theory class whatsoever, which is itself evidence against
its existence.

The arithmetic is instead formalized where no positivity is required: the Hecke--Dirac
operator $\mathsf D_p$ defines an honest Fredholm module over $C(V_Y)$ with
$\ind(\mathsf D_p)=\Z\oplus\Jac(Y)$ exactly (Theorem~\\ref{II:thm:index}), inside
$KK^0(C(V_Y),C(V_Y))\cong M_n(\Z)$. We record, rather than disguise, that for a single
finite graph this $KK$-group is nothing but a matrix ring and carries no content beyond
ordinary linear algebra; the paper closes by locating the correct venue for a genuinely
non-commutative sequel in the projective limit over a congruence tower of graphs, where
the coefficient algebra becomes an infinite, totally disconnected profinite space and
$KK^0$ ceases to be finite-dimensional. As a first exact instance of the tower
philosophy, we compute four abelian $\Z/\ell^k$ voltage towers over $K_4$ to nine levels:
the torsion obeys the Iwasawa law $\operatorname{ord}_\ell|\Jac(Y_k)|=\mu\ell^k+\lambda
k+\nu$ exactly at every computed level, with $(\mu,\lambda)=(3,1),(0,3),(2,3),(0,1)$,
positive $\mu$ occurring freely, and the Smith normal forms exhibiting $\lambda$ as the
number of persistently growing cyclic factors. The canonical combinatorial model of the
congruence tower itself --- the Iwahori path tower --- obeys instead an exact
pseudo-Iwasawa law with $\lambda=-1<0$, impossible for a $\Z_p[[T]]$-characteristic
ideal: the non-normal tower is not an Iwasawa module, and the conjecture for it is
restated over the Iwahori--Hecke algebra, with its naive commutative form withdrawn.
\end{chapterabstract}

\section*{Status of the results, stated in advance}

As in Paper I, each numbered result is labelled by how it is entitled to be believed.
Theorem~\\ref{II:thm:tracesame} is proved here in full, elementarily, using nothing beyond
the fact that trace is a similarity invariant and that non-negative matrices have
non-negative squared-trace. Theorem~\\ref{II:thm:traceany}, which rules out the auxiliary-vertex and Drinen--Tomforde
schemes of any size, is likewise unconditional: the fact it needs --- that shift
equivalence preserves the nonzero spectrum with multiplicities --- is proved here as the
elementary Lemma~\\ref{II:lem:se} by a two-way injection of generalized eigenspaces, and no
external citation is load-bearing for it. Theorem~\\ref{II:thm:index} is proved in
full and requires no citation beyond standard finite-dimensional linear algebra dressed
in Fredholm-module language. Section~\\ref{II:sec:tower} is exploratory: it states a
research programme, not a theorem, and we say so throughout.

\section{The trace obstruction}\label{II:sec:trace}

\subsection{Setup, recalled from Paper I}

$Y$ is a finite connected $(q+1)$-regular graph on $n$ vertices, $T_p\in M_n(\Z)$ its
adjacency (Hecke) operator, $\Delta_p:=(q+1)I_n-T_p$ its Hecke--Laplacian, and
$\Jac(Y)$ its critical group, with
\[
\coker_\Z(\Delta_p)\cong\Z\oplus\Jac(Y),\qquad
\bigl|\Jac(Y)\bigr|=\kappa(Y)=\frac1n\prod_{f}\bigl(p+1-a_p(f)\bigr)=\frac1n\prod_f\#E_f(\F_p),
\]
the product over the weight-two newforms $f$ whose Hecke eigenvalues occur in
$\Spec T_p$, by Kirchhoff's theorem together with the Eichler--Shimura and
Jacquet--Langlands correspondences. Introduce the companion linearization
\[
C:=\begin{pmatrix}T_p&-qI_n\\I_n&0\end{pmatrix}\ \in\ M_{2n}(\Z),
\]
of the Ihara pencil $1-uT_p+qu^2$, so that $\det(I_{2n}-uC)=\det(I_n-uT_p+qu^2I_n)$.
By Schur elimination against the invertible block $I_n$ in the lower right of
$I_{2n}-C=\begin{psmallmatrix}I_n-T_p&qI_n\\-I_n&I_n\end{psmallmatrix}$, one has the exact
integral equivalence
\[
\begin{pmatrix}I_n&-qI_n\\0&I_n\end{pmatrix}
\bigl(I_{2n}-C\bigr)
\begin{pmatrix}I_n&0\\I_n&I_n\end{pmatrix}
=\begin{pmatrix}\Delta_p&0\\0&I_n\end{pmatrix},
\qquad\text{hence}\qquad
\coker_\Z(I_{2n}-C)\cong\Z\oplus\Jac(Y),
\]
with no denominators at any step, since both flanking matrices are unimodular. So $C$ is
not merely spectrally related to $\Delta_p$; it is honestly row/column-equivalent to it,
and any non-negative realization of $C$ solves the realization problem for $\Delta_p$
with no loss.

\subsection{The elementary, unconditional obstruction}

\begin{lemma}\label{II:lem:tracecomp}
For every $(q+1)$-regular simple graph $Y$ on $n$ vertices, $\Tr(C)=0$ and
\[
\Tr(C^2)=-n(q-1).
\]
\end{lemma}

\begin{proof}
$\Tr(C)=\Tr(T_p)+\Tr(0)=0$ since $Y$ is simple (no loops). In block form,
\[
C^2=\begin{pmatrix}T_p^2-qI_n&-qT_p\\T_p&-qI_n\end{pmatrix},
\]
so $\Tr(C^2)=\Tr(T_p^2)-qn-qn=\Tr(T_p^2)-2qn$. Since $T_p$ is symmetric with entries in
$\{0,1\}$, $\Tr(T_p^2)=\sum_{i,j}(T_p)_{ij}^2=2|E_Y|=(q+1)n$ by regularity. Hence
$\Tr(C^2)=(q+1)n-2qn=-n(q-1)$.
\end{proof}

This has been verified directly on the eleven graphs of Table~\\ref{II:tab:trace}, in every
case with exact agreement; the lemma is general and the table is confirmation, not
evidence.

\begin{theorem}[Unconditional obstruction: no equal-size, equal-spectrum realization]
\label{II:thm:tracesame}
Let $q\ge2$. No matrix $\widetilde C\in M_{2n}(\Z_{\ge0})$ similar to $C$ --- in
particular, no non-negative matrix with the same characteristic polynomial as $C$ ---
exists.
\end{theorem}

\begin{proof}
Trace is invariant under similarity, so $\Tr(\widetilde C)=\Tr(C)=0$ and
$\Tr(\widetilde C^2)=\Tr(C^2)=-n(q-1)<0$ would be forced. But for any matrix
$B\in M_N(\Z_{\ge0})$, $\Tr(B^2)=\sum_{i,j}B_{ij}B_{ji}\ge0$, a sum of products of
non-negative integers. Contradiction.
\end{proof}

\begin{remark}
Theorem~\\ref{II:thm:tracesame} needs nothing beyond Lemma~\\ref{II:lem:tracecomp} and the
one-line non-negativity of $\Tr(B^2)$ for non-negative $B$; it requires no appeal to
symbolic dynamics, and it already rules out any construction that keeps $C$'s size
fixed at $2n$ --- for instance, any attempt to directly reinterpret the entries of $C$
itself, without adding vertices, as a graph adjacency matrix.
\end{remark}

\subsection{The obstruction to dilation of any size}

The constructions actually attempted in the course of this project (auxiliary sink
vertices, Drinen--Tomforde-style desingularization, cascading intermediate nodes) all
enlarge the matrix, replacing $C$ by some $\widetilde C\in M_N(\Z_{\ge0})$, $N>2n$,
related to $C$ not by similarity but by \emph{shift equivalence}: matrices $R\in
M_{2n\times N}(\Z)$, $S\in M_{N\times2n}(\Z)$ and a lag $\ell$ with
\[
CR=R\widetilde C,\qquad SC=\widetilde CS,\qquad RS=C^\ell,\qquad SR=\widetilde C^\ell.
\]
Shift equivalence is the correct and standard notion for comparing presentation
matrices of possibly different sizes whose associated shifts of finite type (equivalently,
Cuntz--Krieger algebras) are to be considered the same dynamical, and $K$-theoretic,
object: the Bowen--Franks group $\coker_\Z(I-\widetilde C)$ is a shift-equivalence
invariant, which is exactly why enlarging the matrix could in principle have been a way
to keep $\coker(I-\widetilde C)\cong\Z\oplus\Jac(Y)$ while repairing non-negativity.

\begin{lemma}[Shift equivalence preserves the nonzero spectrum]\label{II:lem:se}
Let $A\in M_a(\C)$, $B\in M_b(\C)$ be lag-$\ell$ shift equivalent: there exist
$R\in M_{a\times b}$, $S\in M_{b\times a}$ with
\[
AR=RB,\qquad SA=BS,\qquad RS=A^\ell,\qquad SR=B^\ell .
\]
Then for every $\lambda\neq0$ the generalized eigenspaces satisfy
$\dim V_\lambda(A)=\dim V_\lambda(B)$; hence the nonzero spectra of $A$ and $B$ agree as
multisets, and $\Tr(A^k)=\Tr(B^k)$ for \emph{every} $k\ge1$, including $k<\ell$.
\end{lemma}

\begin{proof}
Let $v\in V_\lambda(A)=\ker(A-\lambda)^{k_0}$, $\lambda\neq0$. Then
$(B-\lambda)^{k_0}Sv=S(A-\lambda)^{k_0}v=0$, so $S$ maps $V_\lambda(A)$ into
$V_\lambda(B)$. If $Sv=0$ then $A^\ell v=RSv=0$; but $A$ restricted to $V_\lambda(A)$
has sole eigenvalue $\lambda\neq0$, hence is invertible there, forcing $v=0$. So
$S|_{V_\lambda(A)}$ is injective. Symmetrically, $R$ maps $V_\lambda(B)$ injectively
into $V_\lambda(A)$. Two finite-dimensional spaces injecting into each other have equal
dimension. Algebraic multiplicity equals the dimension of the generalized eigenspace,
and $\Tr(M^k)=\sum_\lambda m_\lambda\lambda^k$ with zero eigenvalues contributing
nothing for $k\ge1$; the trace identity follows.
\end{proof}

\begin{theorem}[Unconditional obstruction: no dilation of any size]\label{II:thm:traceany}
Let $q\ge2$. No matrix $\widetilde C\in M_N(\Z_{\ge0})$, of any size $N$, is shift
equivalent to $C$. Consequently no auxiliary-vertex, desingularization, or other
enlargement scheme produces a non-negative integer matrix presenting the same
$K$-theoretic and dynamical data as $C$.
\end{theorem}

\begin{proof}
By Lemma~\\ref{II:lem:se}, shift equivalence would force
$\Tr(\widetilde C^2)=\Tr(C^2)=-n(q-1)<0$, contradicting
$\Tr(\widetilde C^2)=\sum_{i,j}\widetilde C_{ij}\widetilde C_{ji}\ge0$.
\end{proof}

\begin{remark}[Relation to the spectral conjecture literature]
An earlier draft stated Theorem~\\ref{II:thm:traceany} conditionally on the resolution of
the Boyle--Handelman spectral conjecture by Kim--Ormes--Roush. That was an error of
economy, not of substance: the deep content of that literature is the \emph{converse}
direction --- which trace-positive spectra \emph{are} realizable by non-negative
matrices --- whereas the direction needed here, invariance of the nonzero spectrum under
shift equivalence, is the elementary Lemma~\\ref{II:lem:se}. The theorem is unconditional
and self-contained; the spectral-conjecture literature remains relevant context for the
converse question, and its citations remain unverified by us.
\end{remark}

\begin{remark}[What is, and is not, ruled out]\label{II:rem:scope}
Theorem~\\ref{II:thm:traceany} rules out any \emph{faithful} dilation of $C$ --- one
preserving its nonzero spectrum, hence its actual Ihara/Hecke dynamical content, which is
the only kind of dilation that would let the resulting graph algebra be honestly called a
realization of the Hecke correspondence. It does not, and cannot, rule out the bare
abstract existence of \emph{some} unrelated non-negative matrix $D$ with
$\coker_\Z(I-D)\cong\Z\oplus\Jac(Y)$: by Kirchberg--Phillips classification, a Kirchberg
algebra realizing any prescribed pair $(K_0,K_1)$ with $K_0$ finitely generated and $K_1$
free exists unconditionally, with no relation to $T_p$ required. Such an algebra would
answer no question of interest here; the content of this paper is that the
\emph{natural} one, built from $C$ itself, does not exist.
\end{remark}

\section{The virtual correspondence barrier}\label{II:sec:barrier}

Given Theorems~\\ref{II:thm:tracesame}--\\ref{II:thm:traceany}, it is natural to ask whether the
obstruction can be sidestepped by weakening what counts as a ``correspondence,'' for
instance by allowing a $\Z/2$-graded or super-structure in which $T_p$ is placed in even
degree and $qI_n$ in odd degree, with the hope that a graded Cuntz--Pimsner-type
construction applied to $\cE^+\oplus\cE^-$ (with $[\cE^+]=T_p$, $[\cE^-]=qI_n$, both
honest positive correspondences) might produce a single algebra whose $K_0$-group is the
cokernel of $1-([\cE^+]-[\cE^-])=\Delta_p$.

\begin{proposition}\label{II:prop:nograding}
No such construction exists. More precisely: for any $C^*$-algebra $A$ and any two
genuine $A$--$A$ correspondences $\cE^+,\cE^-$, the Cuntz--Pimsner algebra of the direct
sum $\cE^+\oplus\cE^-$ --- graded or not --- has $K$-theoretic index map
$1-[\cE^+]-[\cE^-]$, not $1-[\cE^+]+[\cE^-]$; a $\Z/2$-grading records parity for the
purpose of assigning classes to $KK^0$ versus $KK^1$, and does not introduce a sign into
the index computation of a single Cuntz--Pimsner algebra built from an honestly positive
module.
\end{proposition}

\begin{proof}[Discussion, in place of a proof]
The Cuntz--Pimsner construction is a functor from the category of genuine Hilbert
$A$--$A$-bimodules (with a fixed left action by compacts, or by adjointable operators, as
appropriate) to $C^*$-algebras, and its $K$-theory is computed, via Pimsner's exact
sequence, from the class $[\cE]$ of the \emph{single} bimodule fed to it, and the root
of the constraint is the axiom $\langle\xi,\xi\rangle\ge0$ of the Hilbert $C^*$-module
inner product: every channel of the completed module carries non-negative norm, so no
summand can act with a negative sign in the completion. There is no step in the
construction at which a summand of the bimodule can be assigned a formal minus sign: the bimodule $\cE^+\oplus\cE^-$ is a perfectly good, positive, correspondence
in its own right, with class $[\cE^+\oplus\cE^-]=[\cE^+]+[\cE^-]=T_p+qI_n$, and its
Cuntz--Pimsner algebra has index map $1-(T_p+qI_n)$, the \emph{opposite} sign
adjustment from the one wanted. Introducing a $\Z/2$-grading changes which Kasparov
group ($KK^0$ or $KK^1$) a resulting class is assigned to; it does not alter this
computation. We are not aware of, and have not been able to construct, any variant
Cuntz--Pimsner-type functor that takes a \emph{formal difference} of two correspondences
as input and returns an algebra realizing that difference as an index. Were such a
functor to exist, it would apply to every virtual class in every $KK^0(A,A)$, converting
every formal difference of finitely generated projective modules into a single algebra
realizing it as an index map, and would in particular apply to the underlying integer
matrix problem addressed by Theorems~\\ref{II:thm:tracesame}--\\ref{II:thm:traceany}, trivializing
it; the fact that those theorems hold, and are elementary, is itself evidence that no
such functor exists.
\end{proof}

\begin{remark}[The obstruction is ontological, not technical]
The negative block $-qI_n$ in $C$ is not an artifact of a particular presentation that a
cleverer bookkeeping device --- grading, desingularization, or otherwise --- could
absorb. It is a faithful record of the non-backtracking correction subtracted in the
Ihara recursion $T\cdot T_1=T_2+(q+1)T_0$: the class $x_p=T_p-qI_n$ is the class of a
correspondence \emph{minus} the class of $q$ copies of the identity correspondence, and
Theorem~\\ref{II:thm:tracesame} shows this difference is not, for any $q\ge2$, again the
class of a correspondence. A virtual class that is provably not effective cannot be made
effective by re-describing it; it can only be relocated to a category that does not
demand effectivity. That category is $KK$-theory, to which we now turn.
\end{remark}

\section{The K-homological bulk, and the congruence tower}\label{II:sec:tower}

\subsection{The Hecke--Dirac operator and its index}

\begin{definition}[Hecke--Dirac operator]\label{II:def:dirac}
Fix an orientation $E_Y^+$ of the edges of $Y$ and let
$\partial\colon\Z^{E_Y^+}\to\Z^{V_Y}$, $\partial(f)=t(f)-o(f)$, be the simplicial
boundary map, so $\partial^t\colon\Z^{V_Y}\to\Z^{E_Y^+}$ is its transpose (coboundary).
Define the self-adjoint operator
\[
\mathsf D_p:=\begin{pmatrix}0&\partial\\\partial^t&0\end{pmatrix}\quad\text{on}\quad
\ell^2(V_Y)\oplus\ell^2(E_Y^+),\qquad
\mathsf D_p^2=\begin{pmatrix}\partial\partial^t&0\\0&\partial^t\partial\end{pmatrix},
\]
i.e.\ $\mathsf D_p=d+d^*$ is the discrete Hodge--de Rham operator of the graph, with
even part the vertex (Hodge $0$-)Laplacian $\partial\partial^t=\Delta_p$ and odd part
the edge (Hodge $1$-)Laplacian $\partial^t\partial$.
\end{definition}

\begin{lemma}\label{II:lem:laplacian}
$\partial\partial^t=\Delta_p=(q+1)I-T_p$ on $\ell^2(V_Y)$.
\end{lemma}

\begin{proof}
$(\partial\partial^t)_{uv}=\sum_{f\in E_Y^+}\partial_{f}(u)\,\partial_f(v)$; the diagonal
term is the degree $q+1$ (each incident edge contributes $\pm1$ squared) and the
off-diagonal term at $u\ne v$ is $-1$ for each edge joining $u,v$ (with a sign
cancellation independent of orientation choice, since $\partial_f(u)\partial_f(v)=-1$
whenever $f$ joins $u,v$ regardless of which endpoint is called $o(f)$), giving exactly
$(q+1)I-T_p$.
\end{proof}

\begin{theorem}[Fredholm module and index]\label{II:thm:index}
$(C(V_Y),\ \ell^2(V_Y)\oplus\ell^2(E_Y^+),\ \mathsf D_p)$ is an even Fredholm module over
$A:=C(V_Y)$, defining a class $[\mathsf D_p]\in KK^0(A,A)\cong M_n(\Z)$, $n=|V_Y|$, and
\[
\ind(\mathsf D_p)\ :=\ \bigl[\ker\mathsf D_p\bigr]-\bigl[\coker\mathsf D_p\bigr]
\ \in\ K_0(A),
\qquad
\ker\mathsf D_p\cong\Z\oplus\Z^{\,\rank H_1(Y)},
\]
\[
\coker\bigl(\mathsf D_p^2|_{\ell^2(V_Y)}\bigr)\cong\Z\oplus\Jac(Y),
\]
so that the even part of the index recovers $\Delta_p$'s cokernel exactly:
$K_0$-theoretically, $[\mathsf D_p^2|_{\ell^2(V_Y)}]$ has cokernel $\Z\oplus\Jac(Y)$, with
\emph{no positivity requirement anywhere in the construction} --- $\mathsf D_p$ is
self-adjoint, not positive, and $\partial,\partial^t$ are integer matrices of arbitrary
sign.
\end{theorem}

\begin{proof}
$\mathsf D_p$ is a finite self-adjoint matrix, hence trivially Fredholm (bounded, and
every finite-rank perturbation of $0$ is compact), giving a class in $KK^0(A,A)$ since
$A=C(V_Y)=\C^n$ acts diagonally on $\ell^2(V_Y)$ and trivially, via a chosen embedding of
$V_Y$-labels, on $\ell^2(E_Y^+)$ compatible with $\partial$'s $A$-bimodule structure.
$\ker\mathsf D_p=\ker\partial^t\oplus\ker\partial$; $\ker\partial^t$ is the space of
harmonic $0$-cochains, $\cong\Z$ (the constants, $Y$ connected), and $\ker\partial\cong
H_1(Y;\Z)\cong\Z^{\rank H_1(Y)}$, the cycle space. The even part follows from
Lemma~\\ref{II:lem:laplacian} and Theorem~3.5(i) of Paper I.
\end{proof}

\begin{remark}[The honest scope of this theorem]\label{II:rem:scope2}
This is a correct and complete restatement, in Fredholm-module language, of results
already proved by direct linear algebra in Paper I. We record it because the language
matters for what comes next: $\mathsf D_p$'s self-adjointness, not positivity, is all
that $KK$-theory demands, and Theorem~\\ref{II:thm:tracesame} shows this is not a matter of
convenience but of necessity --- no positive substitute for $\mathsf D_p^2-qI$-type data
exists. It is not, however, a new theorem about $Y$; every computation in it already
appears in Paper I under other names.
\end{remark}

\subsection{Why nothing new happens for a single graph}

\begin{remark}[Finite-dimensional triviality]\label{II:rem:trivial}
For $A=C(V_Y)=\C^n$, $KK^0(A,A)\cong M_n(\Z)$ as a ring under the Kasparov product, and
the product of two classes represented by integer matrices is exactly their matrix
product. Consequently, any identity one might hope to derive by taking a Kasparov
product of $[\mathsf D_p]$ with the Hecke correspondence class $[\cE_p]=T_p$ is nothing
more than the corresponding matrix identity, already available by direct computation. We
tested the specific proposed identity $[\mathsf D_p]\otimes[\cE_p]\overset{?}{=}
[\mathsf D_{p^2}]+q[\mathsf D_0]$ under both natural readings of $[\mathsf D_p]$ (as
$\Delta_p$ and as $T_p-qI$) against the matrices directly: neither matches. What is true,
and is already Lemma~2.1 of Paper I, is the unadorned recursion
$T_p\cdot T_p=T_2+(q+1)I$, with no reference to $\mathsf D_p$ or to any Kasparov product
required to state or prove it. \textbf{$KK$-theory for a single finite graph adds no
content beyond $M_n(\Z)$}; dressing an identity in Kasparov-product language does not
make it a new theorem when the underlying ring is already $M_n(\Z)$ with ordinary matrix
multiplication.
\end{remark}

\subsection{The congruence tower: where new content could appear}\label{II:subsec:tower}

The finite-dimensionality identified in Remark~\\ref{II:rem:trivial} is specific to working
at a single level $Y$. The genuine non-commutative geometry of the Langlands
correspondence, if it is to be found at all along this route, should appear only in the
limit over a \emph{tower} of graphs.

\begin{conjecture}[Congruence tower and the profinite $KK$-limit]\label{II:conj:tower}
Let $(Y_k)_{k\ge0}$ be the tower of quotient graphs $Y_k=\Gamma_0(p^k)\backslash X$ for
the Bruhat--Tits tree $X$ of $\PGL_2(\Q_p)$, with covering maps $Y_{k+1}\to Y_k$ inducing
transfer maps on vertex algebras $A_k=C(V_{Y_k})$ and hence a directed system
\[
A_0\hookrightarrow A_1\hookrightarrow A_2\hookrightarrow\cdots,\qquad
A_\infty:=\varinjlim_kA_k\ \cong\ C\bigl(\varprojlim_kV_{Y_k}\bigr),
\]
with $\varprojlim_kV_{Y_k}$ a profinite (Cantor) set carrying a residual action of
$\PGL_2(\Z_p)$. Then:
\begin{enumerate}
\item[(i)] the Hecke--Dirac operators $\mathsf D_{p,k}$ are compatible under the transfer
maps up to a controlled error, defining a class
$[\mathsf D_p]_\infty\in KK^0(A_\infty,A_\infty)$;
\item[(ii)] $KK^0(A_\infty,A_\infty)$ is no longer a finite matrix ring, but a genuine
inductive limit incorporating the $p$-adic representation-theoretic structure of
$\PGL_2(\Z_p)$-modules, in which the individual finite invariants
$\Jac(Y_k)=\Tors\coker(\Delta_{p,k})$ assemble into an inverse system with growth of
Iwasawa shape. Here the two kinds of tower must be separated, and the separation is
dictated by our own data (Section~\\ref{II:sec:abtower}), not by preference:
\begin{itemize}
\item[(ii-a)] \emph{Abelian deck towers.} For towers carrying a genuine
$\Z_p$-deck action (the voltage towers of Table~\\ref{II:tab:tower}), the limit
$\varprojlim_k\Jac(Y_k)[p^\infty]$ is conjectured to be a finitely generated torsion
$\Z_p[[T]]$-module with characteristic element the cleared voltage pencil
$g(T)=\det D(1+T)$; the finite-level evidence is exact --- the growth law and the
Newton-polygon identities $\mu_{\mathrm{meas}}=\mu(g)$,
$\lambda_{\mathrm{meas}}=\lambda(g)-1$ of Remarks~\\ref{II:rem:towerhonest} and
\\ref{II:rem:newton} hold to the integer in all four computed towers --- but the
limit-module statement itself remains conjectural at the level of this paper, and its
attribution in the literature is flagged in the caveats.
\item[(ii-b)] \emph{The congruence tower.} The tower $\Gamma_0(p^k)\backslash X$ is
non-normal and carries no $\Z_p$-deck action, and its canonical combinatorial model, the
Iwahori path tower of Remark~\\ref{II:rem:iwahori}, exhibits the exact pseudo-Iwasawa law
$\nu_2=6\cdot2^k-k-4$ with $\lambda=-1<0$, which is impossible for the characteristic
ideal of any finitely generated torsion $\Z_p[[T]]$-module. \emph{The naive
$\Z_p[[T]]$-form of this conjecture, including the tentative identification of the
Ihara variable $u$ with $1+T$ proposed in an earlier draft, is therefore withdrawn as
refuted in shape by its own model.} What we conjecture in its place is that the correct
base object for the congruence tower is the Iwahori--Hecke algebra
$\cH(\PGL_2(\Q_p),\mathcal I)$ of the tower, acting through the degeneracy
correspondences, with the finite invariants governed by a characteristic ideal over
that non-commutative base, and the observed $-k$ term arising as an Euler-characteristic
correction --- the removal of one Eisenstein rank per level --- external to the module
invariants proper. We do not know the correct precise formulation, and we prefer to say
so rather than to state a formula the data has already falsified;
\end{itemize}
\item[(iii)] the trace obstruction of Theorem~\\ref{II:thm:traceany} persists at every finite
level $k$, so the entire tower consists of virtual, non-realizable classes; whether the
\emph{limiting} object $[\mathsf D_p]_\infty$ is subject to an analogous obstruction, or
whether the profinite limit process itself provides the missing room for a genuine
realization, is open.
\end{enumerate}
\end{conjecture}

\begin{remark}
We state Conjecture~\\ref{II:conj:tower} as a research programme, not a theorem, and we have
proved none of its parts; part (ii-a) is the best supported, part (ii-b) is deliberately
left without a formula. For part (iii) one sharp dichotomy can already be posed, and it
is the intended opening theorem of a sequel: if the limit carries a faithful invariant
trace $\tau$ (type $\mathrm{II}$ regime), the obstruction of
Theorem~\\ref{II:thm:traceany} should persist as non-negativity of the $\tau$-weighted
second moment of any would-be correspondence, killing realizations exactly as at finite
level; if the limit is traceless (type $\mathrm{III}$ regime, as the boundary crossed
products of Paper I in fact are), the numerical obstruction has no formulation at all,
and the question becomes whether absence of an obstruction can be converted into an
actual construction. Part (iii) is, in our view, the most important open
question raised by this paper: it asks whether the No-Go theorems of
Section~\\ref{II:sec:trace}, proved level-by-level, survive or dissolve in the projective
limit. We do not know the answer, and we think this --- rather than any further
finite-graph computation --- is the correct next object of study.
\end{remark}

\section{Computations}\label{II:sec:computations}

\begin{table}[ht]
\centering\small
\renewcommand{\arraystretch}{1.15}
\begin{tabular}{lcccc}
\toprule
$Y$ & $n$ & $q$ & $\Tr(C^2)$ (computed) & $-n(q-1)$ (formula)\\
\midrule
$K_4$              & $4$  & $2$ & $-4$  & $-4$\\
$K_5$              & $5$  & $3$ & $-10$ & $-10$\\
$K_6$              & $6$  & $4$ & $-18$ & $-18$\\
$K_{3,3}$          & $6$  & $2$ & $-6$  & $-6$\\
$Q_3$ (cube)       & $8$  & $2$ & $-8$  & $-8$\\
prism $C_4\times K_2$&$8$ & $2$ & $-8$  & $-8$\\
Wagner $V_8$       & $8$  & $2$ & $-8$  & $-8$\\
prism $C_5\times K_2$&$10$& $2$ & $-10$ & $-10$\\
M\"obius $V_{10}$  & $10$ & $2$ & $-10$ & $-10$\\
Petersen           & $10$ & $2$ & $-10$ & $-10$\\
Heawood            & $14$ & $2$ & $-14$ & $-14$\\
\bottomrule
\end{tabular}
\medskip
\caption{$\Tr(C^2)$ computed directly from the $2n\times2n$ matrix $C$, against the
closed form of Lemma~\\ref{II:lem:tracecomp}, for every graph of Paper I's Table~1. Exact
agreement in all eleven cases, confirming a fact already proved in general; the table is
illustrative, not evidential.}
\label{II:tab:trace}
\end{table}

\section{A first computed instance of the tower: abelian
\texorpdfstring{$\Z_\ell$}{Z-l}-covers}\label{II:sec:abtower}

Conjecture~\\ref{II:conj:tower} concerns the congruence tower $\Gamma_0(p^k)\backslash X$,
which is neither Galois nor abelian and whose explicit construction requires quaternionic
lattice data we do not assemble here. But part (ii) of the conjecture makes a
\emph{shape} prediction --- Iwasawa-type growth of the torsion --- and that shape can be
tested rigorously today in the nearest fully constructible situation: abelian
$\Z/\ell^k$-covers of a fixed base graph, built from a voltage assignment
$\alpha\colon E_Y\to\Z$ on the chords of a spanning tree, with derived graphs
$Y_k$ on $n\ell^k$ vertices and compatible covering maps $Y_{k+1}\to Y_k$. Every $Y_k$ is
again $(q+1)$-regular, so the entire theory of Papers I--II applies level by level; in
particular the obstruction density of Theorem~\\ref{II:thm:traceany} is constant along the
tower, $\Tr(C_k^2)/n_k=-(q-1)$ for every $k$, so no na\"ive normalization dilutes the
No-Go theorem in the limit.

Two independent exact methods were used and cross-checked at every level where both run:
the reduced-Laplacian determinant and Smith normal form of the derived graph directly,
and the character factorization
\[
\kappa(Y_k)\;=\;\frac{\kappa(Y_0)}{\ell^{k}}\prod_{j=1}^{\ell^k-1}
\det D\bigl(\zeta_{\ell^k}^{\,j}\bigr),
\qquad
D(t):=(q+1)I_n-A(t),
\]
with $A(t)$ the voltage adjacency pencil, the product computed exactly as a resultant of
integer polynomials. The two methods agreed exactly in all twelve cross-checked levels.

\begin{table}[ht]
\centering\small
\renewcommand{\arraystretch}{1.2}
\begin{tabular}{lccccl}
\toprule
voltages on $K_4$ & $\ell$ & $\mu$ & $\lambda$ & $\nu$ & exact law, verified range\\
\midrule
$(1,0,0)$ & $2$ & $3$ & $1$ & $1$ & $\operatorname{ord}_2\kappa(Y_k)=3\cdot2^k+k+1$,\ $2\le k\le9$\\
$(1,1,1)$ & $2$ & $0$ & $3$ & $4$ & $\operatorname{ord}_2\kappa(Y_k)=3k+4$,\ $0\le k\le9$\\
$(1,2,0)$ & $2$ & $2$ & $3$ & $3$ & $\operatorname{ord}_2\kappa(Y_k)=2^{k+1}+3k+3$,\ $2\le k\le9$\\
$(1,0,0)$ & $3$ & $0$ & $1$ & $0$ & $\operatorname{ord}_3\kappa(Y_k)=k$,\ $0\le k\le6$\\
\bottomrule
\end{tabular}
\medskip
\caption{Exact $\ell$-adic valuations of $|\Jac(Y_k)|=\kappa(Y_k)$ along four
$\Z/\ell^k$ voltage towers over $K_4$ ($q=2$). In each tower the Iwasawa law
$\operatorname{ord}_\ell\kappa(Y_k)=\mu\ell^k+\lambda k+\nu$ holds \emph{exactly} for
every computed level in the stated range; the parameters were fitted from two increments
and the remaining six to eight levels are genuine verification, not fit.}
\label{II:tab:tower}
\end{table}

\begin{remark}[What the table shows, and what it does not]\label{II:rem:towerhonest}
Three points, in decreasing order of confidence. (1) The growth law itself, for abelian
$\ell$-towers of graphs, is a known theorem of Vallières and of McGown--Vallières
\cite{Vallieres,McGownVallieres}, with closely related Iwasawa-theoretic structure in
Gonet \cite{Gonet}; the bibliographic records are now verified and their titles
corroborate the attribution, though the theorems themselves have not been checked
against the papers' texts here. Our computations \emph{confirm} that theory in four
instances; they do not discover it. (2) The Smith normal forms refine the order-growth
law structurally: in the $\mu=0$ tower $(1,1,1)$ the $2$-primary part of $\Jac(Y_k)$ is,
for $1\le k\le3$,
\[
\Jac(Y_k)[2^\infty]\;\cong\;\Z/2^{\,k}\oplus\Z/2^{\,k+2}\oplus\Z/2^{\,k+2},
\]
exactly $\lambda=3$ cyclic factors each growing by one power of $\ell$ per level, while
in the $\mu=3$ tower the \emph{number} of cyclic factors itself grows ($5$ factors at
$k=2$, $9$ at $k=3$): structurally, $\lambda$ counts persistently growing cyclic factors
and $\mu>0$ manifests as exponential proliferation of new ones. We do not know whether
this structural statement in this form appears in the graph-Iwasawa literature.
(3) Positive $\mu$ occurs freely here ($\mu=3$ and $\mu=2$ above), in genuine contrast to
the cyclotomic number-field setting, where $\mu=0$ is the Ferrero--Washington theorem;
whatever the correct formulation of Conjecture~\\ref{II:conj:tower}(ii) for the congruence
tower turns out to be, it cannot import the $\mu=0$ expectation from number fields.
\end{remark}

\begin{remark}[The Iwahori path tower: exact data against the naive conjecture]
\label{II:rem:iwahori}
The canonical combinatorial model of the congruence tower itself is constructible: take
level-$k$ vertices to be the non-backtracking paths of length $k$ in $Y_0$ (level $1$ is
the Hashimoto digraph of $A^{\mathrm{nb}}$, level $k\ge2$ its iterated line digraphs),
so that $|V(Y_k)|=n(q+1)q^{k-1}$ reproduces the index
$[\Gamma(1):\Gamma_0(p^k)]=(p+1)p^{k-1}$ exactly. For $Y_0=K_4$, $q=2$, the sandpile
groups of these Eulerian digraphs were computed exactly for $0\le k\le6$ (up to $384$
vertices). Two exact laws emerge. First, a Levine-type recursion,
$|\mathrm{sandpile}(Y_{k+1})|=2^{\,n_k-1}\,|\mathrm{sandpile}(Y_k)|$, verified at five
consecutive levels. Second, the valuation obeys a pseudo-Iwasawa law exactly for
$2\le k\le6$:
\[
\nu_2\bigl(|\mathrm{sandpile}(Y_k)|\bigr)=6\cdot2^k-k-4,
\qquad\text{i.e.}\qquad(\mu,\lambda,\nu)=(6,\,-1,\,-4).
\]
The invariant $\lambda=-1$ is \emph{negative}, which is impossible for the
characteristic ideal of a finitely generated torsion $\Z_p[[T]]$-module. The canonical
combinatorial model of the congruence tower therefore \emph{refutes the naive
$\Z_p[[T]]$-form of part (ii) of Conjecture~\\ref{II:conj:tower} as stated above}: whatever
governs the congruence tower, it is not an ordinary Iwasawa module --- consistently with
the fact that this tower carries no $\Z_p$-deck action, being non-normal. The $-1$ has a
transparent source, the removal of one Eisenstein (trivial-direction) rank per level, an
Euler-characteristic correction rather than a module invariant; the corrected form of
the conjecture should be stated either over the Iwahori--Hecke algebra of the tower or
as the abelian-deck statement, which the data of Table~\\ref{II:tab:tower} and the Newton
test below confirm exactly. Structurally, the Smith normal forms at $k\le3$ show the
same three persistent growing factors $\Z/2^{\,k}\oplus(\Z/2^{\,k+2})^2$ as the abelian
$(1,1,1)$ tower, plus an exponentially proliferating cloud of bounded factors ($8$
copies of $\Z/2$ at $k=2$; $12$ copies of $\Z/2$ and $8$ of $\Z/4$ at $k=3$): $\lambda$
counts persistent factors, $\mu$ the proliferation rate, and the negative correction
lives outside both.
\end{remark}

\begin{remark}[Newton-polygon confirmation of the characteristic element]
\label{II:rem:newton}
In the four abelian towers of Table~\\ref{II:tab:tower} the characteristic-element mechanism
was tested directly: with $g(T)$ the cleared voltage pencil $\det D(1+T)\in\Z[T]$, the
Newton data give, in all four cases exactly,
\[
\mu_{\mathrm{measured}}=\mu(g),\qquad
\lambda_{\mathrm{measured}}=\lambda(g)-1 ,
\]
and the uniform offset $-1$ is fully accounted: $t\mapsto t^{-1}$ symmetry of the pencil
forces a double zero $T^2\mid g$ at the trivial character, contributing $+2k$ to the
valuation sum, while the normalization $\kappa(Y_k)=\kappa(Y_0)\ell^{-k}\prod_{\zeta\neq1}
\det D(\zeta)$ removes $k$; net $+k$, i.e.\ $\lambda(g)-2+1$. Within the abelian regime
the graph-theoretic $p$-adic $L$-function is thus confirmed to the integer.
\end{remark}

\begin{remark}[Relation to Conjecture~\\ref{II:conj:tower}]
These towers are a rigorous toy, not the conjecture: the congruence tower is non-normal
and its covering combinatorics are those of the full tree, not of an abelian voltage
group. What the computation establishes is only this: the Iwasawa shape asserted in
part (ii) is exactly right in the nearest constructible abelian analogue, down to the
integer, with the characteristic-ideal mechanism (here, the resultant of the voltage
pencil against $t^{\ell^k}-1$) functioning precisely as the graph-theoretic
$p$-adic $L$-function. Part (iii) --- the fate of the trace obstruction in the projective
limit --- is untouched by any of this, and remains the paper's central open question.
\end{remark}

\section*{Acknowledgement of provenance and caveats}

This is an unrefereed preprint. Theorem~\\ref{II:thm:tracesame} and Theorem~\\ref{II:thm:index}
are proved here in full from elementary linear algebra and require no external citation.
Theorem~\\ref{II:thm:traceany} depends on the claim that shift equivalence forces agreement
of the full nonzero spectrum of two integer matrices, which we attribute to the
resolution of the Boyle--Handelman spectral conjecture by Kim, Ormes and Roush (both
names corrected from an earlier draft's ``Boyer'' and ``Ozawa''); this
attribution has not been checked against the primary literature, since the authors had
no network access while writing this paper, and should be verified independently before
further citation. Proposition~\\ref{II:prop:nograding} is an argument, not a theorem in the
usual sense: it asserts the non-existence of a construction, supported by an
impossibility argument (such a functor would trivialize
Theorems~\\ref{II:thm:tracesame}--\\ref{II:thm:traceany}) rather than a proof ruling out every
conceivable construction. Conjecture~\\ref{II:conj:tower} is entirely open; none of its three
parts is established here, and part (iii) is explicitly posed as the paper's main
unresolved question. All computations of Section~\\ref{II:sec:computations} are exact and
were checked independently of the general proof of Lemma~\\ref{II:lem:tracecomp}. The tower
computations of Section~\\ref{II:sec:abtower} are exact integer computations, cross-checked
between two independent methods at twelve levels. Bibliography status: every external
entry is now verified against its primary source or its Crossref record, with the DOI
recorded in place. The three entries that had returned mutually contradictory data in
two earlier out-of-band passes (Vallières, McGown--Vallières, Gonet) are resolved, and
the resolution is chastening: \emph{all three} earlier variants were wrong, in the
Gonet case down to the title and the journal. Residual soft spots are noted inline:
the Williams end page (Crossref records the start page only), and the
Pimsner chapter year (1996 electronic versus 1997 volume). One distinction is
maintained deliberately: what has been verified is the bibliographic \emph{metadata};
the attribution of the abelian growth law to these works rests on their (now
corroboratively titled) records and on recollection of their content, not on a reading
of their theorems performed within this verification pass. The possibility that the
structural refinement of Remark~\\ref{II:rem:towerhonest}(2) is known to their authors
stands as before. Two author-name errata were applied: ``Boyer--Handelman'' corrected to
Boyle--Handelman, and ``Kim--Ozawa--Roush'' corrected to Kim--Ormes--Roush; the latter
was the present authors' own error, and both corrected names should themselves be
confirmed at the primary-source stage.

% generated from Paper 6/anccft6.tex
\chapter{Algorithmic non-commutative class field theory, VI: the shadow algebra --- a functorial Kirchberg realization of the critical group}\label{chap:VI}
\chaptermark{VI}
\begin{chapterabstract}
The central open problem of this series --- posed in [Ch.~\ref{chap:III}, Rem.~1.10] and sharpened
in [Ch.~\ref{chap:IV}, \S6] --- asked for a Kirchberg $C^*$-algebra, functorial in the Hecke data of
a finite $(q+1)$-regular graph $Y$, whose $K_0$-group realizes the Laplacian-channel
invariant $\Z\oplus\Jac(Y)$, with $|\Jac(Y)|=\frac1n\prod_f\#E_f(\F_p)$ a product of
Frobenius point counts. Papers II--III proved two No-Go theorems that appeared to
surround the problem: no non-negative matrix of any size is shift equivalent to the
Ihara companion $C$ (the trace obstruction $\Tr(C^2)=-n(q-1)<0$), and the natural
tail-groupoid algebra is arithmetically blind. This paper solves the single-level
problem, inside the exact gap those theorems leave open: shift equivalence governs
spectra, not cokernels. Define the \emph{shadow graph} $\shadow$ on $2n$ vertices,
with adjacency
\[
B(Y)\;=\;\begin{pmatrix}T_p& I_n\\ qI_n & 2I_n\end{pmatrix}:
\]
the original graph, one arc from each vertex to its shadow, $q$ return arcs, and a
double loop at each shadow. The double loops are unimodular pivots of determinant
$-1$, and a two-line block factorization gives
$U\,(I-B^t)\,V=\Delta_p\oplus(-I_n)$ with explicit integer unimodular $U,V$; hence
the graph algebra $\cO_{B(Y)}$ is a unital Kirchberg algebra with
\[
K_0\bigl(\cO_{B(Y)}\bigr)\cong\Z\oplus\Jac(Y),\qquad
K_1\bigl(\cO_{B(Y)}\bigr)\cong\Z,
\]
functorially in $Y$. Being purely infinite it carries no trace, so the obstruction of
[Ch.~\ref{chap:II}] has no formulation on it; and it does not contradict [Ch.~\ref{chap:II}], because $B(Y)$ shares
with the Ihara pencil only its Bowen--Franks data, not its spectrum --- the exact
reconciliation is the per-eigenvalue identity
$(1-2u)(1-u\lambda)-qu^2\big|_{u=1}=-(q+1-\lambda)$. What remains open, and what [Ch.~\ref{chap:II}]
proves is impossible, are stated with equal precision: tower-level functoriality is
open; a \emph{Hecke-dynamical} realization is impossible, which is the sense in which
the present construction is maximal.
\end{chapterabstract}

\section*{Status of the results}

Theorems~\\ref{VI:thm:ktheory} and \\ref{VI:thm:kirchberg} are proved in full;
the unimodular factorization at the heart of Theorem~\\ref{VI:thm:ktheory} is a two-line
block computation reproduced below and machine-verified, together with the Smith-form
match against [Ch.~\ref{chap:I}, Table~1], on $K_4$, $K_5$, $K_{3,3}$ and Petersen
(Table~\\ref{VI:tab:battery}). The classical inputs --- simplicity and pure infiniteness
criteria for Cuntz--Krieger/graph algebras and their $K$-theory --- are quoted with a
flagged citation. Section~\\ref{VI:sec:scope} states exactly what is solved, what is
open, and what is impossible; one error in the task specification that prompted this
paper is corrected there (Remark~\\ref{VI:rem:k1}).

\section{The shadow graph}\label{VI:sec:shadow}

Throughout, $Y$ is a finite connected $(q+1)$-regular graph on $n$ vertices, $q\ge2$,
with adjacency (Hecke) matrix $T_p$ and Hecke--Laplacian $\Delta_p=(q+1)I-T_p$, whose
cokernel is $\Z\oplus\Jac(Y)$ with $\ker\Delta_p=\Z$ ([Ch.~\ref{chap:I}, Thm.~3.5]).

\begin{definition}[Shadow graph]\label{VI:def:shadow}
The \emph{shadow graph} $\shadow$ is the directed graph on vertex set
$V_Y\sqcup V_Y^{\,\prime}$ (each $v$ paired with its shadow $v'$) with arcs:
\begin{enumerate}
\item[(i)] $v\to w$ for every edge $vw$ of $Y$, in both orientations (the $T_p$
block);
\item[(ii)] one arc $v\to v'$ for every $v$;
\item[(iii)] $q$ parallel arcs $v'\to v$ for every $v$;
\item[(iv)] two loops at every shadow $v'$.
\end{enumerate}
Its adjacency matrix is $B(Y)=\begin{psmallmatrix}T_p&I\\qI&2I\end{psmallmatrix}$,
and the \emph{shadow algebra} is the graph $C^*$-algebra $\cO_{B(Y)}$.
\end{definition}

The design principle is worth stating in one sentence, because it is precisely the
loophole left by the No-Go theorems of [Ch.~\ref{chap:II}]. The obstruction
$\Tr(C^2)=-n(q-1)<0$ forbids any non-negative matrix from carrying the
\emph{spectrum} of the Ihara companion --- the $-qI$ block cannot be dilated away.
But $K$-theory sees only the \emph{cokernel}, which is a unimodular-equivalence
invariant, and unimodular row operations, unlike similarities, may be implemented by
pivots of determinant $-1$: a shadow vertex with a double loop contributes the
$1\times1$ block $1-2=-1$, and Schur elimination against it \emph{adds} $q$ back to
the diagonal --- exactly the operation that is spectrally impossible and
cohomologically free.

\section{The \texorpdfstring{$K$}{K}-theory}\label{VI:sec:ktheory}

\begin{theorem}[Functorial realization]\label{VI:thm:ktheory}
With $U=\begin{psmallmatrix}I&-qI\\0&I\end{psmallmatrix}$ and
$V=\begin{psmallmatrix}I&0\\-I&I\end{psmallmatrix}$, both integer unimodular,
\begin{equation}\label{VI:eq:factor}
U\,\bigl(I_{2n}-B(Y)^t\bigr)\,V\;=\;\Delta_p\oplus(-I_n).
\end{equation}
Consequently
\[
K_0\bigl(\cO_{B(Y)}\bigr)=\coker_\Z\bigl(I-B(Y)^t\bigr)\cong\Z\oplus\Jac(Y),
\qquad
K_1\bigl(\cO_{B(Y)}\bigr)=\ker_\Z\bigl(I-B(Y)^t\bigr)\cong\Z ,
\]
and the assignment $Y\mapsto\cO_{B(Y)}$ is functorial under graph isomorphisms: any
automorphism $\pi$ of $Y$ satisfies $(\pi\oplus\pi)B(Y)(\pi\oplus\pi)^{-1}=B(Y)$,
inducing an action on $\cO_{B(Y)}$ compatible with the induced action on
$\Z\oplus\Jac(Y)$.
\end{theorem}

\begin{proof}
$B^t=\begin{psmallmatrix}T_p&qI\\I&2I\end{psmallmatrix}$, so
$I-B^t=\begin{psmallmatrix}I-T_p&-qI\\-I&-I\end{psmallmatrix}$. Two lines of block
multiplication:
\[
U(I-B^t)=\begin{pmatrix}I&-qI\\0&I\end{pmatrix}
\begin{pmatrix}I-T_p&-qI\\-I&-I\end{pmatrix}
=\begin{pmatrix}(q+1)I-T_p&0\\-I&-I\end{pmatrix},
\]
\[
U(I-B^t)V=\begin{pmatrix}\Delta_p&0\\-I&-I\end{pmatrix}
\begin{pmatrix}I&0\\-I&I\end{pmatrix}
=\begin{pmatrix}\Delta_p&0\\0&-I\end{pmatrix}.
\]
Cokernel and kernel are unimodular-equivalence invariants, and
$\coker(\Delta_p\oplus(-I))=\coker\Delta_p=\Z\oplus\Jac(Y)$,
$\ker=\ker\Delta_p=\Z$ by [Ch.~\ref{chap:I}, Thm.~3.5(i)]. The $K$-theory formula for graph
$C^*$-algebras is the classical
$K_0=\coker(I-B^t)$, $K_1=\ker(I-B^t)$ \cite{CK}. Equivariance is the visible block
symmetry, verified explicitly for a transposition on $K_4$.
\end{proof}

\section{The Kirchberg properties}\label{VI:sec:kirchberg}

\begin{theorem}[Shadow algebra is Kirchberg]\label{VI:thm:kirchberg}
For $Y$ connected and $q\ge1$, the matrix $B(Y)$ is irreducible (the shadow graph is
strongly connected), is not a permutation matrix, and has a vertex --- any shadow
$v'$ --- carrying two distinct loops. Consequently $\cO_{B(Y)}$ is separable,
nuclear, simple, purely infinite and in the UCT class: a unital Kirchberg algebra.
Being purely infinite, it admits no nonzero trace, finite or semifinite; the trace
obstruction of {\rm[Ch.~\ref{chap:II}, Thms.~1.5--1.6]} therefore has no formulation on
$\cO_{B(Y)}$.
\end{theorem}

\begin{proof}
Strong connectivity: within the first layer all vertices are mutually reachable
through the arcs of the connected graph $Y$; layers communicate through
$v\to v'\to v$. The double loop at $v'$ gives two distinct return paths at one
vertex and period $1$. Irreducibility with a non-permutation matrix and the
double-loop witness place $B(Y)$ under the classical simplicity and pure
infiniteness criteria for Cuntz--Krieger algebras, and graph algebras of finite
graphs are separable, nuclear and satisfy the UCT \cite{CK}. Pure infiniteness
excludes traces. Machine confirmation of the witnesses on all four test graphs:
Table~\\ref{VI:tab:battery}.
\end{proof}

\section{Non-contradiction, and the exact scope of the solution}\label{VI:sec:scope}

\begin{proposition}[Why this does not contradict the No-Go theorems]\label{VI:prop:noncontra}
$B(Y)$ is not shift equivalent to the Ihara companion $C$ of {\rm[II]}: already
$\Tr B(Y)=2n\neq0=\Tr C$, and Lemma~{\rm[II,~1.4]} makes traces shift-equivalence
invariants. The shadow algebra shares with the Ihara pencil its Bowen--Franks data
and nothing spectral: per eigenvalue $\lambda$ of $T_p$,
\[
\det\bigl(I-uB(Y)\bigr)=\prod_{\lambda}\Bigl[(1-2u)(1-u\lambda)-qu^{2}\Bigr],
\qquad
(1-2u)(1-u\lambda)-qu^{2}\Big|_{u=1}=-(q+1-\lambda),
\]
so the $u=1$ specialization of the shadow zeta reproduces the Laplacian spectrum up
to sign, while for $u\neq1$ the two polynomial families differ. The construction
lives exactly in the gap the No-Go theorems left open: shift equivalence constrains
spectra; $K$-theory is a cokernel invariant; and Remark~[II, after Thm.~1.6]
anticipated precisely this distinction when it separated existence from canonicity.
\end{proposition}

\begin{proof}
The block determinant is computed with the commuting-block formula (three of the
four blocks are scalars), and the $u=1$ evaluation is
$(1-2)(1-\lambda)-q=\lambda-1-q$. Traces: immediate.
\end{proof}

\begin{remark}[A correction to the task specification]\label{VI:rem:k1}
The specification that prompted this paper asked for $K_1\cong\Z^{\,r}$. That is the
\emph{boundary}-channel invariant ([Ch.~\ref{chap:I}, Thm.~3.7]), not the Laplacian-channel one:
the target fixed by [Ch.~\ref{chap:I}, Thm.~3.5], [Ch.~\ref{chap:III}, Thm.~1.x] and [Ch.~\ref{chap:IV}] is
$K_1=\ker\Delta_p=\Z$, and that is what Theorem~\\ref{VI:thm:ktheory} delivers.
Realizing $\Z^r$ in $K_1$ alongside $\Z\oplus\Jac(Y)$ in $K_0$ would require a
non-square-type presentation and a different target; it is not the series' problem.
\end{remark}

\begin{remark}[What is solved, what is open, what is impossible]\label{VI:rem:scope}
\emph{Solved:} the single-level canonicity problem of [Ch.~\ref{chap:III}, Rem.~1.10] --- a
Kirchberg algebra functorial in the Hecke datum $T_p$, traceless, with
$K_0\cong\Z\oplus\Jac(Y)$, $K_1\cong\Z$, the torsion order
$\frac1n\prod_f\#E_f(\F_p)$ carried on the nose. \emph{Open:} tower-level
functoriality --- whether the tail norms of [Ch.~\ref{chap:V}, Thm.~3.1] lift to
$*$-homomorphisms or correspondences
$\cO_{B(Y_{k+1})}\to\cO_{B(Y_k)}$ inducing the norm maps of [IV--V] on $K_0$; nothing
here constructs such lifts. \emph{Impossible:} a \emph{Hecke-dynamical} realization,
i.e.\ an algebra whose gauge dynamics is the Ihara flow itself --- that is exactly
what [Ch.~\ref{chap:II}, Thm.~1.6] forbids for every non-negative presentation of every size. The
shadow algebra is Hecke-functorial, not Hecke-dynamical, and by [Ch.~\ref{chap:II}] no algebra can
be both: the construction is maximal in a theorem-certified sense.
\end{remark}

\begin{remark}[Normalization]\label{VI:rem:norm}
The construction uses the minimal unimodular pivot: loop weight $2$ is the least $m$
with $\det(1-m)=-1$, and the off-diagonal pair $(1,q)$ is one of the factorizations
$\alpha\beta=q$; any such choice yields the same $K$-theory by the same
factorization. We fix $(1,q,2)$ as the normal form and make no uniqueness claim
beyond it.
\end{remark}

\section{The verification battery}\label{VI:sec:battery}

\begin{table}[ht]
\centering\small
\renewcommand{\arraystretch}{1.2}
\begin{tabular}{lccccc}
\toprule
$Y$ & $n$ & $q$ & (S) factorization \\eqref{VI:eq:factor} & $\Jac(Y)$ from
$\mathrm{SNF}(I-B^t)$ & vs [Ch.~\ref{chap:I}, Table 1]\\
\midrule
$K_4$      & 4  & 2 & exact & $(\Z/4)^2$                 & match\\
$K_5$      & 5  & 3 & exact & $(\Z/5)^3$                 & match\\
$K_{3,3}$  & 6  & 2 & exact & $\Z/3\oplus\Z/3\oplus\Z/9$ & match\\
Petersen   & 10 & 2 & exact & $\Z/2\oplus(\Z/10)^3$      & match\\
\bottomrule
\end{tabular}
\medskip
\caption{Machine verification (\texttt{hecke\_kirchberg.py}): the explicit
factorization \\eqref{VI:eq:factor} holds as an exact matrix identity; the invariant
factors of $I-B(Y)^t$ reproduce [Ch.~\ref{chap:I}, Table~1] exactly; strong connectivity, the
double-loop witness, $\Tr B=2n\neq0$, and the $u{=}1$ reconciliation identity all
verified; automorphism equivariance spot-checked on $K_4$.}
\label{VI:tab:battery}
\end{table}

\section*{Acknowledgement of provenance and caveats}

Unrefereed preprint. Theorems~\\ref{VI:thm:ktheory}--\\ref{VI:thm:kirchberg} and
Proposition~\\ref{VI:prop:noncontra} are proved; the two-line factorization
\\eqref{VI:eq:factor} can be checked by hand and was machine-verified on four graphs.
The classical Cuntz--Krieger inputs (graph-algebra $K$-theory; simplicity and pure
infiniteness criteria) are quoted through \cite{CK}. Remark~\\ref{VI:rem:scope} delimits the result: single-level canonicity
solved, tower functoriality open, Hecke-dynamical realization impossible by [Ch.~\ref{chap:II}].

\input{chapters/III}
\part{Towers: commutative Iwasawa theory and its failure}
% generated from Paper 4/anccft4.tex
\chapter[IV. The limit module of abelian towers and the characteristic identity]{Algorithmic non-commutative class field theory, IV: the limit module of abelian towers and the characteristic identity}\label{chap:IV}
\chaptermark{IV}
\begin{chapterabstract}
For an abelian $\Z/\ell^k$ voltage tower $(Y_k)$ over a finite base graph, Papers
II--III established the exact Iwasawa growth law of the sandpile orders and its
Newton-polygon mechanism, leaving the limit-module statement of
[Ch.~\ref{chap:II}, Conj.~3.6(ii-a)] open. This paper closes it. Writing $D(t)$ for the voltage
Laplacian pencil and $\La=\Zl[[T]]$, we prove:
\[
X_\infty:=\varprojlim_k\bigl(\coker_\Z(\Delta_{Y_k})\otimes\Zl,\ \pi_*\bigr)
\;\cong\;\La^n\big/D(1+T)\,\La^n,
\]
a finitely generated torsion $\La$-module with
$\car(X_\infty)=\bigl(\det D(1+T)\bigr)$, together with an \emph{exact} control
theorem $X_\infty/\omega_kX_\infty\cong\coker_\Z(\Delta_{Y_k})\otimes\Zl$ --- exact,
not merely pseudo-isomorphic, because the tower's presentations are literally the
$\omega_k$-specializations of one $\La$-presentation. The proof rests on three
structural identities --- the derived Laplacian is the block-circulant $D(\sigma)$;
the covering pushforward is the group-ring coefficient projection; pushforward is
surjective on $\ell$-primary torsion with kernel of the exact predicted order --- each
proved and additionally machine-verified at every level of two towers. Combined with
[Ch.~\ref{chap:II}, Rem.~5.4], the measured invariants $(\mu,\lambda)=(\mu(g),\lambda(g)-1)$,
$g=\det D(1+T)$, are now consequences of a theorem rather than observations. The
result is stated with explicit deference to the graph-Iwasawa literature of
Vallières, McGown--Vallières and Gonet, whose verified titles indicate substantial
overlap; what is new here, at minimum, is the exactness of the control isomorphism at
the level of Smith normal forms and the integration into the trilogy's $K$-theoretic
framework. The non-normal congruence tower ([Ch.~\ref{chap:II}, Conj.~3.6(ii-b)]) and the canonicity
problem ([Ch.~\ref{chap:III}, Rem.~1.10]) are untouched and remain the open frontier.
\end{chapterabstract}

\section*{Status of the results}

Every numbered statement in this paper is proved in full. Lemmas~\ref{IV:lem:block},
\ref{IV:lem:norm} and \ref{IV:lem:torsurj} are elementary and were additionally verified by
exact machine computation at every level $k\le3$ of the tower with voltages $(1,1,1)$,
$\ell=2$, and every level $k\le2$ of the tower $(1,0,0)$, $\ell=3$
(Table~\ref{IV:tab:battery}). The single external algebraic input is the structure
theory of finitely generated $\La$-modules, quoted with a flagged citation. The
relation of Theorem~\ref{IV:thm:limit} to the published graph-Iwasawa literature is
discussed, deliberately and prominently, in Section~\ref{IV:sec:lit}.

\section{Setup and the three structural identities}\label{IV:sec:setup}

Fix a finite connected base graph $Y_0$ on $n$ vertices, a prime $\ell$, a spanning
tree, and a voltage assignment $\alpha$ on the chords, giving the derived
$\Z/\ell^k$-covers $Y_k$ on $n\ell^k$ vertices with deck generator $\sigma$ and
covering maps $\pi\colon Y_{k+1}\to Y_k$. Let
\[
D(t)\;=\;(q+1)I_n-A(t)\ \in\ M_n\bigl(\Z[t,t^{-1}]\bigr)
\]
be the voltage Laplacian pencil ($A(t)$ the voltage adjacency), $\Delta_k$ the
ordinary Laplacian of $Y_k$, $G_k=\Z/\ell^k$, and order $V_{Y_k}$ as
$(v,g)\mapsto vN+g$, $N=\ell^k$. Write $S_N$ for the cyclic shift on $\Zl^N$ and
$Q_N\colon\Zl^{\ell N}\to\Zl^{N}$ for reduction of the group coordinate,
$e_g\mapsto e_{g\bmod N}$.

\begin{lemma}[Block-circulant identification]\label{IV:lem:block}
In the ordering above, $\Delta_k=D(S_N)$ exactly: the level-$k$ Laplacian is the
image of the pencil under $t\mapsto S_N$. Equivalently, as a
$\Z[G_k]$-module with the deck action,
$\coker_\Z(\Delta_k)\cong\Z[G_k]^n\big/D(\sigma)\Z[G_k]^n$.
\end{lemma}

\begin{proof}
This is the definition of the derived graph read backwards: a base edge $(u,v)$ of
voltage $a$ contributes the arcs $(u,g)\!-\!(v,g{+}a)$ for all $g$, i.e.\ the block
$E_{uv}\otimes S_N^{\,a}$ plus its transpose, and the degree term is
$(q+1)I_n\otimes I_N$. Summing gives $\Delta_k=D(S_N)$; identifying
$\Zl^N\cong\Zl[G_k]$ with $S_N=$ multiplication by $\sigma$ gives the module form.
Verified as an exact matrix identity at all computed levels.
\end{proof}

\begin{lemma}[Pushforward is coefficient projection]\label{IV:lem:norm}
Let $P_k:=I_n\otimes Q_{\ell^k}$, the matrix of the covering pushforward
$\pi_*\colon\Z^{V_{Y_{k+1}}}\to\Z^{V_{Y_k}}$ (sum over fibres). Then
\[
P_k\,\Delta_{k+1}=\Delta_k\,P_k,
\qquad
P_k\,R_{k+1}=R_k\,P_k,
\]
where $R_k=I_n\otimes S_{\ell^k}$ is the deck action; under
Lemma~\ref{IV:lem:block} the induced map on cokernels is the group-ring coefficient
projection $\Z[G_{k+1}]\to\Z[G_k]$. Moreover $P_k$ is surjective.
\end{lemma}

\begin{proof}
Both identities reduce to the one-line facts $Q\,S_{\ell N}=S_N\,Q$ and
$Q\,S_{\ell N}^{\,a}=S_N^{\,a}\,Q$ applied blockwise to $\Delta=D(S)$; surjectivity
is immediate since every $e_{(v,h)}$ is hit. Verified as exact matrix identities at
all computed levels.
\end{proof}

\begin{lemma}[Norm surjectivity on torsion]\label{IV:lem:torsurj}
The induced norm
$\pi_*\colon\coker_\Z(\Delta_{k+1})\to\coker_\Z(\Delta_k)$ is
surjective; it preserves the degree functional $\deg$ (total mass), hence restricts
to a surjection of torsion subgroups, with kernel of order
\[
\bigl|\ker(\pi_*|_{\Tor_\ell})\bigr|
=\ell^{\,\ord_\ell|\Jac(Y_{k+1})|\,-\,\ord_\ell|\Jac(Y_k)|}.
\]
\end{lemma}

\begin{proof}
Surjectivity of $P_k$ plus the intertwining gives surjectivity on cokernels. The
column sums of $\Delta_k$ vanish, so $\deg([y])=\sum_iy_i$ is well defined and
realizes $\coker_\Z(\Delta_k)\cong\Z\oplus\Tor$ with $\Tor=\ker(\deg)$ (a finitely
generated group surjecting onto $\Z$ with rank-zero kernel has kernel equal to its
torsion). Pushforward preserves total mass, so $\deg\circ\pi_*=\deg$; hence for
torsion $t$ downstairs any preimage $a$ has $\deg(a)=0$, i.e.\ $a$ is torsion:
surjectivity on $\Tor$. The kernel order is then forced by counting. Verified,
including the exact kernel orders $\ell^{3},\ell^{3},\ell^{3}$ and
$\ell^{1},\ell^{1}$ at the computed transitions (Table~\ref{IV:tab:battery}).
\end{proof}

\section{The limit module and the characteristic identity}\label{IV:sec:limit}

\begin{theorem}[Limit module]\label{IV:thm:limit}
Let $g(T):=\det D(1+T)\in\Zl[T]$ and assume $g\neq0$ $($true in every tower of
{\rm[Ch.~\ref{chap:II}, \S5]}$)$. Then, with transition maps the norms $\pi_*$ of
Lemma~\ref{IV:lem:torsurj},
\[
X_\infty:=\varprojlim_k\bigl(\coker_\Z(\Delta_{Y_k})\otimes\Zl\bigr)
\;\cong\;\La^n\big/\,D(1+T)\,\La^n ,
\qquad \La=\Zl[[T]],
\]
under the standard identification $\varprojlim_k\Zl[G_k]\cong\La$,
$\sigma\mapsto1+T$ $($the Laurent entries $t^{\pm a}$ of the pencil landing in
$\La^\times$-multiples since $1+T\in\La^\times)$. The module $X_\infty$ is finitely
generated and torsion over $\La$, and
\[
\car(X_\infty)\;=\;\bigl(g(T)\bigr)\;=\;\bigl(\det D(1+T)\bigr).
\]
\end{theorem}

\begin{proof}
By Lemmas~\ref{IV:lem:block}--\ref{IV:lem:norm} the tower of cokernels is, level by level
and compatibly with the transition maps, the tower of presentations
$\Zl[G_k]^n\xrightarrow{D(\sigma)}\Zl[G_k]^n\to C_k\to0$ with transitions the
coefficient projections. All modules are finitely generated $\Zl$-modules, hence
compact, and the transition maps are continuous surjections on the presenting
copies; inverse limits are exact on such systems (Mittag--Leffler for compact
groups), so $\varprojlim C_k=\coker\bigl(\varprojlim D(\sigma)\bigr)
=\La^n/D(1+T)\La^n$. Since $g=\det D(1+T)\neq0$, the map $D(1+T)$ is injective over
the fraction field, so the cokernel is $\La$-torsion; it is finitely generated by
presentation. For the characteristic ideal, localize at a height-one prime
$P\subset\La$: $\La_P$ is a discrete valuation ring, elementary divisors give
$\operatorname{length}_{\La_P}\bigl((X_\infty)_P\bigr)=v_P\bigl(\det D(1+T)\bigr)$,
and taking the product over height-one primes yields
$\car(X_\infty)=(\det D(1+T))$ by the structure theory of finitely generated
torsion $\La$-modules \cite{Washington}.
\end{proof}

\begin{corollary}[Exact control]\label{IV:cor:control}
For every $k\ge0$, with $\omega_k=(1+T)^{\ell^k}-1$,
\[
X_\infty\big/\omega_k X_\infty\;\cong\;\coker_\Z(\Delta_{Y_k})\otimes\Zl
\;\cong\;\bigl(\Z\oplus\Jac(Y_k)\bigr)\otimes\Zl ,
\]
an isomorphism on the nose --- not merely up to bounded kernels and cokernels ---
because $\La/\omega_k\cong\Zl[G_k]$ carries the level-$k$ presentation exactly.
\end{corollary}

\begin{corollary}[Invariants, now as a theorem]\label{IV:cor:invariants}
The Iwasawa invariants of the tower are $\mu=\mu(g)$ and $\lambda=\lambda(g)-1$,
where $(\mu(g),\lambda(g))$ are read off the Newton polygon of $g$, the $-1$ being
the trivial-zero and normalization accounting of {\rm[Ch.~\ref{chap:II}, Rem.~5.4]}. In particular
the four exact laws of {\rm[Ch.~\ref{chap:II}, Table~2]} are consequences of
Theorem~\ref{IV:thm:limit}, no longer observations.
\end{corollary}

\section{The verification battery}\label{IV:sec:battery}

\begin{table}[ht]
\centering\small
\renewcommand{\arraystretch}{1.2}
\begin{tabular}{lccccc}
\toprule
tower & levels & (B) block & (S) deck & (P)+(C) norm identities & (N) torsion norm\\
\midrule
$(1,1,1)$, $\ell=2$ & $k\le3$ & exact & exact & exact & surj.; $|\ker|=2^3,2^3,2^3$\\
$(1,0,0)$, $\ell=3$ & $k\le2$ & exact & exact & exact & surj.; $|\ker|=3^1,3^1$\\
\bottomrule
\end{tabular}
\medskip
\caption{Machine verification of Lemmas~\ref{IV:lem:block}--\ref{IV:lem:torsurj}:
(B) $\Delta_k=D(S_N)$; (S) $R_k\Delta_k=\Delta_kR_k$; (P) $P_k\Delta_{k+1}
=\Delta_kP_k$; (C) $P_kR_{k+1}=R_kP_k$; (N) surjectivity of $\pi_*$ on
$\ell$-primary torsion with kernel orders exactly
$\ell^{\Delta\ord}$. Every check passed as an exact identity; the lemmas are proved
above independently, so the table is confirmation, not evidence.}
\label{IV:tab:battery}
\end{table}

\section{Relation to the literature, stated prominently}\label{IV:sec:lit}

The verified bibliographic records of Vallières \cite{Vallieres}, McGown--Vallières
\cite{McGownVallieres} and, above all, Gonet --- whose paper is titled precisely
\emph{Iwasawa theory of Jacobians of graphs} \cite{Gonet} --- make it very likely
that Theorem~\ref{IV:thm:limit}, or a statement equivalent to it, is present in that
literature; the growth law certainly is, as acknowledged since [Ch.~\ref{chap:II}, \S5]. Under the
series' policy, we have verified those papers' metadata but have not read their
theorems, and we therefore claim priority for nothing in
Sections~\ref{IV:sec:setup}--\ref{IV:sec:limit} beyond what a comparison with their texts
will leave standing. What we believe is ours at minimum: the \emph{exactness} of
Corollary~\ref{IV:cor:control} in the form stated, its verification at the resolution
of Smith normal forms across [Ch.~\ref{chap:II}, Table~2] and Table~\ref{IV:tab:battery}, and the
integration of the statement into the trilogy's $K$-theoretic framework, where
$\coker_\Z(\Delta_{Y_k})$ is the $K_0$-group of the level-$k$ Laplacian channel
$($[Ch.~\ref{chap:I}, Thm.~3.5]$)$, so that Theorem~\ref{IV:thm:limit} reads: \emph{the
$K$-theoretic torsion of the Laplacian channel along an abelian tower is one
finitely generated torsion Iwasawa module, with characteristic ideal the voltage
pencil determinant}.

\section{What remains open}\label{IV:sec:open}

Nothing here touches the two hard problems. The congruence tower is non-normal, its
pseudo-Iwasawa law has $\lambda=-1$, and the correct base object --- conjecturally a
module theory over the Iwahori--Hecke algebra with an intrinsic Euler-characteristic
correction ([Ch.~\ref{chap:II}, Conj.~3.6(ii-b)], [Ch.~\ref{chap:III}, \S2]) --- still lacks a definition of
characteristic ideal. And the canonicity problem of [Ch.~\ref{chap:III}, Rem.~1.10] --- a Kirchberg
algebra over $C(\Omega)$, functorial in the Hecke data, with
$K_0\cong(\Z\oplus\Jac(Y_k))_k$ --- is exactly as open as before;
Theorem~\ref{IV:thm:limit} sharpens its target, since the prescribed $K$-theory is now
known to be a single $\La$-module, but supplies no construction.

\section*{Acknowledgement of provenance and caveats}

Unrefereed preprint. All numbered statements are proved; the machine battery of
Table~\ref{IV:tab:battery} confirms the three lemmas at seven levels across two towers.
External inputs: the structure theory of finitely generated torsion
$\La$-modules and the identification $\varprojlim\Zl[G_k]\cong\La$, both standard,
cited to \cite{Washington}, whose bibliographic record is now primary-source
verified (only the chapter-level pointer remains a recollection-based reading aid);
the graph-Iwasawa literature
\cite{Vallieres,McGownVallieres,Gonet}, bibliographic data verified in [II--III],
texts unread, priority explicitly ceded in Section~\ref{IV:sec:lit}.

% generated from Paper 9/anccft9.tex
\chapter{Algorithmic non-commutative class field theory, IX: exceptional zeros, the graph $\cL$-invariant, and the operator-algebraic
Mazur--Tate--Teitelbaum formula}\label{chap:IX}
\chaptermark{IX}
\begin{chapterabstract}
Section~3 of [Ch.~\ref{chap:III}] recorded, without proof, that the cleared voltage pencil
$g(T)=T^{2}h(T)$ of an abelian $\Z/\ell^{k}$ tower over a $(q{+}1)$-regular graph has a
double zero at the trivial character, and that the constant term of the Iwasawa growth
law obeys $\nu_0=\ord_\ell\kappa(Y_0)-\mu+\sum_j\varepsilon_j$ in all four towers of
[Ch.~\ref{chap:II}, Table~2]. This paper turns both observations into theorems and identifies the
number they conceal. The double zero is forced by the symmetry $D(t^{-1})=D(t)^{t}$ of
the pencil; its second-order coefficient has the closed form
\[
\tfrac12\Theta''(0)\;=\;-\,\kappa(Y_0)\,\|h_\alpha\|^{2},
\]
where $h_\alpha$ is the harmonic representative of the voltage cochain $\alpha$ --- its
orthogonal projection onto the cycle space, equivalently onto the kernel of the
Hecke--Dirac operator of [Ch.~\ref{chap:II}] --- and $\kappa(Y_0)=|\Jac(Y_0)|=\frac1n\prod_f\#E_f(\F_p)$
is the algebraic $L$-value of [Ch.~\ref{chap:I}]. We therefore define the graph $\cL$-invariant
$\cL(Y_\bullet):=\tfrac12\Theta''(0)/\kappa(Y_0)=-\|h_\alpha\|^{2}$, a rational number
depending only on the class $[\alpha]\in H^{1}(Y_0;\Z)$ and not on $\ell$; for a
single-chord voltage, $-\tfrac12\Theta''(0)$ counts the spanning trees avoiding the
chord. The bookkeeping identity of [Ch.~\ref{chap:III}] becomes a theorem with an exact description of
the finitely supported defects $\varepsilon_j$, and the specification's proposed
alignment is replaced by the correct one: in the exceptional-zero regime $\lambda_h=0$
all defects vanish and $\nu_0=-\ord_\ell\cL(Y_\bullet)$, while in general
$\nu_0+\ord_\ell\cL(Y_\bullet)=(\ord_\ell a_2-\mu)+\sum_j\varepsilon_j$. On the
operator-algebraic side we prove a twisted Bass identity for the character-twisted
Satake dynamics of [Ch.~\ref{chap:I}] and show that at the critical inverse temperature $\beta=1$ the
twisted Bass determinant equals $q^{-n}\Theta(t-1)$: its zero at the trivial character
is simple in the temperature direction with derivative proportional to $\kappa(Y_0)$,
of order two in the character direction with second derivative proportional to
$\cL(Y_\bullet)\kappa(Y_0)$, and the partition function of the critical KMS family has a
double pole with leading coefficient $n(q-q^{-1})/\cL(Y_\bullet)$ --- the
Mazur--Tate--Teitelbaum and Greenberg--Stevens shapes, with the roles of algebraic
$L$-value and $\cL$-invariant made explicit. The abelian shadow limit of [Ch.~\ref{chap:VIII}] carries
$\Theta$ as the characteristic element of its equivariant $K_0$. For the non-normal
Iwahori tower there is no character direction: its zero is simple at every level, and
the constancy of its residue $n_k\kappa(Y_k)q^{-n_k}$ along the tower is exactly the
pseudo-Iwasawa law, with the Euler term $-k$ appearing as the $p$-adic valuation of the
vertex count. No $\cL$-invariant exists for that tower, and we say so.
\end{chapterabstract}

\section*{Status of the results}

Every numbered statement is proved. The proofs use only results of Papers I--VIII,
finite-dimensional linear algebra (a Schur complement replaces analytic perturbation
theory), the Weierstrass preparation theorem for $\Zl[[T]]$ \cite{Wa}, and the character
product formula for abelian covers already used in [Ch.~\ref{chap:II}, \S5] and [Ch.~\ref{chap:IV}]. The twisted Bass
identity is the Artin--Ihara $L$-function formula of Stark--Terras \cite{ST}; we give a
six-line proof so that nothing is quoted. Every identity was machine-verified exactly
on the four abelian towers of [Ch.~\ref{chap:II}, Table~2] --- growth laws recomputed directly from the
reduced Laplacians of the derived graphs up to $128$ vertices --- and on the Iwahori tower
at levels $k\le3$ (Table~\ref{IX:tab:battery}). The Mazur--Tate--Teitelbaum \cite{MTT} and
Greenberg--Stevens \cite{GS} formulas are cited for the shape of the analogy only;
nothing here depends on them. Section~\ref{IX:sec:scope} records what is open and five
corrections to the specification that prompted this paper.

\section{The exceptional zero of the voltage pencil}\label{IX:sec:zero}

Setting as in [Ch.~\ref{chap:II}, \S5] and [Ch.~\ref{chap:IV}]. $Y_0$ is a finite connected $(q{+}1)$-regular graph
with $n$ vertices and $m=\frac{(q+1)n}2$ edges, $r=m-n+1$; fix an orientation $E^+$ of
the edges and a voltage $\alpha\colon E^+\to\Z$, extended to oriented edges by
$\alpha(\bar e)=-\alpha(e)$, so that $\alpha$ is an integer $1$-cochain. The voltage
adjacency and Laplacian pencils are
\[
A(t)_{uv}=\sum_{e\colon u\to v}t^{\alpha(e)},\qquad
D(t)=(q{+}1)I_n-A(t)\ \in M_n\bigl(\Z[t,t^{-1}]\bigr),
\]
and $\Theta(T):=\det D(1+T)\in\Z[T,(1+T)^{-1}]\subset\La:=\Zl[[T]]$ is the graph
$\ell$-adic $L$-function, the characteristic element of the limit module of
[Ch.~\ref{chap:IV}, Thm.~2.1]. Its cleared form $g(T)=(1+T)^{d}\Theta(T)\in\Z[T]$, $d$ minimal, is the
polynomial of [Ch.~\ref{chap:II}, Rem.~5.4] and [Ch.~\ref{chap:III}, \S3]; since $(1+T)^{d}=1+dT+\cdots$ and
$\Theta(0)=\Theta'(0)=0$ (Theorem~\ref{IX:thm:zero}), the coefficients
$a_i$ of $T^{i}$ in $\Theta$ and in $g$ agree for $i\le2$. We write
$\kappa_0:=\kappa(Y_0)=|\Jac(Y_0)|=\frac1n\prod_f\#E_f(\F_p)$, the algebraic $L$-value
of [Ch.~\ref{chap:I}, Thm.~3.5], and $\Delta_0=D(1)$ for the Laplacian of $Y_0$.

Hodge data. Let $\partial\colon\Z^{E^+}\to\Z^{V}$, $\partial e=t(e)-o(e)$, be the
boundary map, so $\Delta_0=\partial\partial^{t}$ [Ch.~\ref{chap:II}, Lem.~3.2]; the cycle space
$\ker\partial=H_1(Y_0;\Z)$ and the cut space $\operatorname{im}\partial^{t}$ are
orthogonal complements in $\Q^{E^+}$. With $\Delta_0^{+}$ the Moore--Penrose inverse
(the inverse of $\Delta_0$ on $\mathbf1^{\perp}$, zero on $\mathbf1$), the orthogonal
projection onto the cycle space is
\[
\Pi:=I-\partial^{t}\Delta_0^{+}\partial ,
\]
which is also the projection onto $\ker\mathsf D_p\cap\ell^{2}(E^+)$ for the
Hecke--Dirac operator $\mathsf D_p$ of [Ch.~\ref{chap:II}, Def.~3.1]. Put $h_\alpha:=\Pi\alpha$, the
harmonic representative of $\alpha$, and $\|h_\alpha\|^{2}=\alpha^{t}\Pi\alpha
=\|\alpha\|^{2}-(\partial\alpha)^{t}\Delta_0^{+}(\partial\alpha)\in\Q_{\ge0}$.

\begin{lemma}[Symmetry and switching]\label{IX:lem:sym}
$D(t^{-1})=D(t)^{t}$, hence $f(t):=\det D(t)$ satisfies $f(t^{-1})=f(t)$. For
$\varphi\in\Z^{V}$, the switched voltage $\alpha+\partial^{t}\varphi$ has pencil
$\operatorname{diag}(t^{\varphi_u})\,D(t)\,\operatorname{diag}(t^{-\varphi_u})$, so $f$
depends only on the class $[\alpha]\in H^{1}(Y_0;\Z)$.
\end{lemma}

\begin{proof}
$A(t^{-1})_{uv}=\sum_{e\colon u\to v}t^{-\alpha(e)}=\sum_{e\colon v\to u}t^{\alpha(e)}
=A(t)_{vu}$. Switching multiplies the entry $t^{\alpha(e)}$, $e\colon u\to v$, by
$t^{\varphi_v-\varphi_u}$.
\end{proof}

\begin{theorem}[Exceptional zero of order two]\label{IX:thm:zero}
\begin{enumerate}
\item[(i)] $\Theta(0)=0$ and $\Theta'(0)=0$.
\item[(ii)] $a_2=\tfrac12\Theta''(0)=-\,\kappa(Y_0)\,\|h_\alpha\|^{2}$.
\item[(iii)] The zero has order exactly two if and only if $h_\alpha\neq0$, i.e.\ if and
only if $[\alpha]\neq0$ in $H^{1}(Y_0;\Q)$.
\item[(iv)] If $\alpha$ is the indicator of a single edge $e$, then
$\|h_\alpha\|^{2}=1-\Reff(e)$ and $-a_2=\kappa(Y_0)\bigl(1-\Reff(e)\bigr)$ is the number
of spanning trees of $Y_0$ not containing $e$.
\end{enumerate}
\end{theorem}

\begin{proof}
Write $s=T$ and expand $D(1+s)=\Delta_0+sD_1+s^{2}D_2+O(s^{3})$ with
$D_1=D'(1)=-A'(1)=:-W$ and $D_2=\frac12D''(1)=-\frac12A''(1)$. By
Lemma~\ref{IX:lem:sym}, $A'(1)^{t}=-A'(1)$, so $W$ is antisymmetric, and
$(W\mathbf1)_u=\sum_{o(e)=u}\alpha(e)-\sum_{t(e)=u}\alpha(e)=-(\partial\alpha)_u$, i.e.\
$W\mathbf1=-\partial\alpha$. Also
$\mathbf1^{t}A''(1)\mathbf1=\sum_{e\in E^+}\bigl[\alpha(\alpha-1)+\alpha(\alpha+1)\bigr]
=2\|\alpha\|^{2}$, so $\mathbf1^{t}D_2\mathbf1=-\|\alpha\|^{2}$.

Decompose $\Q^{n}=\Q\mathbf1\oplus\mathbf1^{\perp}$ orthogonally, with unit vector
$\mathbf1/\sqrt n$. In this decomposition $\Delta_0=0\oplus\Delta'$ with $\Delta'$
invertible and $\det\Delta'=\prod_{\mu\neq0}\mu=n\kappa_0$ [Ch.~\ref{chap:III}, eq.~(2.2)]. Writing
$E:=sD_1+s^{2}D_2+O(s^{3})$ in block form $\begin{psmallmatrix}e_{00}&e_{0*}\\
e_{*0}&E_{**}\end{psmallmatrix}$, the Schur complement gives
\[
f(1+s)=\det(\Delta'+E_{**})\cdot\bigl(e_{00}-e_{0*}(\Delta'+E_{**})^{-1}e_{*0}\bigr).
\]
Here $e_{00}=\frac1n\mathbf1^{t}E\mathbf1=\frac{s}{n}\mathbf1^{t}D_1\mathbf1
+\frac{s^{2}}n\mathbf1^{t}D_2\mathbf1+O(s^{3})=-\frac{s^{2}}n\|\alpha\|^{2}+O(s^{3})$,
because $\mathbf1^{t}W\mathbf1=0$ by antisymmetry; and $e_{0*},e_{*0}=O(s)$ with
$e_{*0}=\frac{s}{\sqrt n}D_1\mathbf1|_{\mathbf1^\perp}+O(s^{2})$,
$e_{0*}=\frac{s}{\sqrt n}(D_1^{t}\mathbf1)^{t}|_{\mathbf1^\perp}+O(s^{2})$; note
$D_1\mathbf1=-W\mathbf1=\partial\alpha\in\mathbf1^{\perp}$ and
$D_1^{t}\mathbf1=W\mathbf1=-\partial\alpha$. Since $(\Delta'+E_{**})^{-1}=\Delta_0^{+}
+O(s)$ on $\mathbf1^{\perp}$,
\[
e_{0*}(\Delta'+E_{**})^{-1}e_{*0}
=\frac{s^{2}}n(-\partial\alpha)^{t}\Delta_0^{+}(\partial\alpha)+O(s^{3}) .
\]
Hence
$f(1+s)=\bigl(n\kappa_0+O(s)\bigr)\cdot\frac{s^{2}}{n}\bigl(-\|\alpha\|^{2}
+(\partial\alpha)^{t}\Delta_0^{+}(\partial\alpha)\bigr)+O(s^{3})
=-\kappa_0\|\Pi\alpha\|^{2}s^{2}+O(s^{3})$, which is (i) and (ii). (iii): $\Pi\alpha=0$
iff $\alpha$ lies in the cut space $\operatorname{im}\partial^{t}\otimes\Q$ iff
$[\alpha]=0$ in $H^{1}(Y_0;\Q)$. (iv): for $\alpha=e$,
$\|\Pi e\|^{2}=1-e^{t}\partial^{t}\Delta_0^{+}\partial e=1-\Reff(e)$, and
$\kappa_0\Reff(e)$ is the number of spanning trees containing $e$ (Kirchhoff). On $K_4$:
$\Reff=\frac12$, $16\cdot\frac12=8$ trees avoid the chord, and $a_2=-8$ in both
$(1,0,0)$ towers (Table~\ref{IX:tab:battery}).
\end{proof}

\begin{definition}[Graph $\cL$-invariant]\label{IX:def:L}
$\displaystyle\cL(Y_\bullet):=\frac{\Theta''(0)}{2\,\kappa(Y_0)}=-\|h_\alpha\|^{2}
=-\|\Pi_{\ker\mathsf D_p}\alpha\|^{2}\ \in\Q_{\le0}$, nonzero iff $[\alpha]\neq0$
rationally. It depends only on $[\alpha]\in H^{1}(Y_0;\Z)$ and not on $\ell$: the same
rational number serves every prime. On $K_4$ the four towers of [Ch.~\ref{chap:II}, Table~2] have
$\cL=-\frac12,\,-1,\,-\frac32,\,-\frac12$ for $(1,0,0)$, $(1,1,1)$, $(1,2,0)$, $(1,0,0)$
respectively.
\end{definition}

\begin{theorem}[The graph Mazur--Tate--Teitelbaum formula]\label{IX:thm:mtt}
$\displaystyle\tfrac12\Theta''(0)=\cL(Y_\bullet)\cdot\kappa(Y_0)
=\cL(Y_\bullet)\cdot\frac1n\prod_f\#E_f(\F_p)$: the leading coefficient of the graph
$\ell$-adic $L$-function at its exceptional zero is the $\cL$-invariant times the
algebraic $L$-value.
\end{theorem}

\begin{proof}
Theorem~\ref{IX:thm:zero}(ii), Definition~\ref{IX:def:L}, and [Ch.~\ref{chap:I}, Thm.~3.5] for the product
formula for $\kappa(Y_0)$.
\end{proof}

\begin{remark}[The analogy, stated as an analogy]\label{IX:rem:analogy}
For an elliptic curve with split multiplicative reduction at $p$, the $p$-adic
$L$-function has an exceptional zero at $s=1$ and \cite{MTT} conjectured, \cite{GS}
proved, $L_p'(E,1)=\cL(E)\,L(E,1)/\Omega_E$ with $\cL(E)=\log_pq_E/\ord_pq_E$ the
$\cL$-invariant of the Tate period. Theorem~\ref{IX:thm:mtt} has that shape: a zero forced
by a symmetry, an algebraic part that is a product of Frobenius point counts, and a
transcendental-looking factor that is in fact a period-type quantity --- here the Hodge
energy of the class $[\alpha]$, i.e.\ its length in the harmonic metric of $Y_0$. The
zero is of order two rather than one because the graph pencil is self-dual under
$t\mapsto t^{-1}$ (Lemma~\ref{IX:lem:sym}), which kills the first derivative; the analogue
in the classical setting would be a functional equation forcing an even order. We
claim the shape and nothing more.
\end{remark}

\section{The $\ell$-adic bookkeeping}\label{IX:sec:book}

Fix a prime $\ell$ with the tower $Y_k$ connected for all $k$ (equivalently
$\alpha\not\equiv0$ in $H^{1}(Y_0;\F_\ell)$, true in the four towers). Write
$h(T):=g(T)/T^{2}\in\Z[T]$, so $h(0)=a_2$, and let $h=\ell^{\mu}P(T)U(T)$ be its
Weierstrass decomposition in $\La$: $P$ distinguished of degree $\lambda_h$, $U$ a unit
\cite{Wa}. Thus $\mu=\min_i\ord_\ell(\text{coefficients of }h)$ and $\lambda_h$ is the
index of the first coefficient attaining it; $[Ch.~\ref{chap:II}, Rem.~5.4]$ and $[Ch.~\ref{chap:IV}, Cor.~2.3]$ give
the Iwasawa invariants of the tower as $\mu$ and $\lambda=\lambda(g)-1=\lambda_h+1$.
For $j\ge1$ put
\[
N_j:=\ord_\ell\prod_{\zeta\ \mathrm{prim.}\ \ell^{j}}h(\zeta-1)
=\ord_\ell\Res_T\bigl(\Phi_{\ell^{j}}(1+T),\,h(T)\bigr),
\qquad
\varepsilon_j:=N_j-\bigl(\mu\,\varphi(\ell^{j})+\lambda_h\bigr),
\]
the exceptional low-level defects of [Ch.~\ref{chap:III}, \S3].

\begin{lemma}[Character product]\label{IX:lem:char}
For every $k\ge0$,
$\kappa(Y_k)=\kappa(Y_0)\,\ell^{-k}\prod_{\zeta^{\ell^{k}}=1,\ \zeta\neq1}\det D(\zeta)$
{\rm[Ch.~\ref{chap:II}, \S5]}, and consequently
\[
\ord_\ell\kappa(Y_k)=\ord_\ell\kappa_0-k+\sum_{j=1}^{k}\bigl(2+N_j\bigr).
\]
\end{lemma}

\begin{proof}
$\det D(\zeta)=\zeta^{-d}g(\zeta-1)$ and $\prod_{\zeta\neq1}\zeta^{-d}$ is a root of
unity; $g(\zeta-1)=(\zeta-1)^{2}h(\zeta-1)$ and
$\prod_{\zeta\ \mathrm{prim.}\ \ell^{j}}(\zeta-1)=\pm\Phi_{\ell^{j}}(1)=\pm\ell$, of
valuation $1$; group the nontrivial roots by exact order $\ell^{j}$, $1\le j\le k$.
\end{proof}

\begin{theorem}[The bookkeeping identity of {\rm[Ch.~\ref{chap:III}, \S3]}]\label{IX:thm:book}
Let $\rho_1,\dots,\rho_{\lambda_h}$ be the roots of $P$ (all of positive valuation) and
let $j_0\ge1$ be least with $1/\varphi(\ell^{j_0})<\min_i\ord_\ell\rho_i$ ($j_0=1$ if
$\lambda_h=0$). Then:
\begin{enumerate}
\item[(i)] $\varepsilon_j=0$ for all $j\ge j_0$; the defects are finitely supported;
\item[(ii)] for all $k\ge j_0-1$,
\[
\ord_\ell\kappa(Y_k)=\mu\,\ell^{k}+(\lambda_h+1)\,k+\nu_0,\qquad
\nu_0=\ord_\ell\kappa(Y_0)-\mu+\sum_{j\ge1}\varepsilon_j ;
\]
\item[(iii)] in particular $\lambda=\lambda_h+1=\lambda(g)-1$, recovering
{\rm[Ch.~\ref{chap:II}, Rem.~5.4]} and {\rm[Ch.~\ref{chap:IV}, Cor.~2.3]}, and the constant term is the identity
recorded in {\rm[Ch.~\ref{chap:III}, \S3]}.
\end{enumerate}
\end{theorem}

\begin{proof}
(i) For $\zeta$ primitive of order $\ell^{j}$, $\ord_\ell(\zeta-1)=1/\varphi(\ell^{j})$;
if this is smaller than every $\ord_\ell\rho_i$ then
$\ord_\ell(\zeta-1-\rho_i)=1/\varphi(\ell^{j})$ for all $i$, so
$\ord_\ell h(\zeta-1)=\mu+\lambda_h/\varphi(\ell^{j})$, and summing over the
$\varphi(\ell^{j})$ primitive roots gives $N_j=\mu\varphi(\ell^{j})+\lambda_h$. (ii) By
Lemma~\ref{IX:lem:char},
$\ord_\ell\kappa(Y_k)=\ord_\ell\kappa_0-k+2k+\sum_{j\le k}(\mu\varphi(\ell^{j})
+\lambda_h+\varepsilon_j)=\ord_\ell\kappa_0+k+\mu(\ell^{k}-1)+\lambda_hk
+\sum_{j\le k}\varepsilon_j$, and for $k\ge j_0-1$ the last sum is the full sum. (iii)
is read off. Machine confirmation (Table~\ref{IX:tab:battery}): the law holds with the
predicted $(\mu,\lambda,\nu_0)=(3,1,1),(0,3,4),(2,3,3),(0,1,0)$ --- the four rows of
[Ch.~\ref{chap:II}, Table~2] --- at every computed level $k\ge2$, and in fact at every $k\ge0$ except
in the $(1,2,0)$ tower, where $(\varepsilon_1,\varepsilon_2)=(-1,+2)$ and the law
starts at $k=2$, exactly as [Ch.~\ref{chap:III}, \S3] reported.
\end{proof}

\begin{corollary}[The exceptional-zero regime]\label{IX:cor:regime}
The following are equivalent: $\lambda_h=0$; $\ord_\ell a_2=\mu$;
$\ord_\ell\cL(Y_\bullet)=\mu-\ord_\ell\kappa(Y_0)$. In this regime $\varepsilon_j=0$ for
every $j\ge1$, the law of Theorem~\ref{IX:thm:book}(ii) holds for every $k\ge0$, and
\[
\nu_0\;=\;-\,\ord_\ell\cL(Y_\bullet):
\]
the constant term of the Iwasawa law is minus the $\ell$-adic valuation of the
$\cL$-invariant. This is the case in both $(1,0,0)$ towers: $\nu_0=1=-\ord_2(-\tfrac12)$
and $\nu_0=0=-\ord_3(-\tfrac12)$.
\end{corollary}

\begin{proof}
$\mu\le\ord_\ell a_2$ always, with equality iff the constant coefficient attains the
minimum iff $\lambda_h=0$; the third condition is the second rewritten. If $\lambda_h=0$
then $P=1$, $\ord_\ell h(\zeta-1)=\mu$ for every $\zeta$, so $N_j=\mu\varphi(\ell^{j})$
and $\varepsilon_j=0$, $j_0=1$; then $\nu_0=\ord_\ell\kappa_0-\mu
=\ord_\ell\kappa_0-\ord_\ell a_2=-\ord_\ell\cL$.
\end{proof}

\begin{theorem}[The $\cL$-defect identity]\label{IX:thm:defect}
In general,
\[
\nu_0+\ord_\ell\cL(Y_\bullet)\;=\;\bigl(\ord_\ell a_2-\mu\bigr)+\sum_{j\ge1}\varepsilon_j
\;=:\;\delta(Y_\bullet,\ell),
\]
with $\ord_\ell a_2-\mu\ge0$ and equality iff $\lambda_h=0$. On the four towers
$\delta=0,4,2,0$.
\end{theorem}

\begin{proof}
Substitute $\ord_\ell\cL=\ord_\ell a_2-\ord_\ell\kappa_0$ into
Theorem~\ref{IX:thm:book}(ii).
\end{proof}

\begin{remark}[A correction to the specification]\label{IX:rem:align}
The specification asked to prove that $\ord_\ell\cL(Y_\bullet)$ ``aligns exactly'' with
$\nu_0-\ord_\ell\kappa(Y_0)+\mu$. By Theorem~\ref{IX:thm:book} the latter quantity is
$\sum_j\varepsilon_j$, which is $0,0,1,0$ on the four towers, whereas
$\ord_\ell\cL=-1,0,-1,0$: there is no such identity. The statements that hold are
Corollary~\ref{IX:cor:regime} and Theorem~\ref{IX:thm:defect}. The specification's other
claim, that $\ord_\ell a_2=\mu$ ``controls the main term in the $\lambda_h=0$ regime'',
is the definition of that regime; its content is the vanishing of all defects and the
formula $\nu_0=-\ord_\ell\cL$.
\end{remark}

\begin{proposition}[The exceptional zero inside the Iwasawa module]\label{IX:prop:module}
Let $X_\infty\cong\La^{n}/D(1+T)\La^{n}$ be the limit module of {\rm[Ch.~\ref{chap:IV}, Thm.~2.1]}, with
$\car(X_\infty)=(\Theta)$. Then $X_\infty/TX_\infty\cong(\Z\oplus\Jac(Y_0))\otimes\Zl$
has $\Zl$-rank one, while $\ord_T\Theta=2$; at the augmentation prime the localization
is
\[
X_\infty\otimes_{\La}\La_{(T)}\;\cong\;\La_{(T)}\big/T^{2}\La_{(T)} ,
\]
cyclic of length two. The exceptional zero is the statement that the Eisenstein
eigenspace of the Iwasawa module is not semisimple.
\end{proposition}

\begin{proof}
The first isomorphism is the exact control theorem [Ch.~\ref{chap:IV}, Cor.~2.2] at $k=0$.
$\La_{(T)}$ is a discrete valuation ring with uniformizer $T$ and residue field $\Ql$,
so $D(1+T)$ has elementary divisors over it; reducing modulo $T$ gives
$D(1)=\Delta_0$ of rank $n-1$ over $\Ql$, so exactly one elementary divisor is divisible
by $T$, and since the $T$-adic order of the determinant is $2$
(Theorem~\ref{IX:thm:zero}) that divisor is $T^{2}$ times a unit.
\end{proof}

\section{The operator-algebraic formulation}\label{IX:sec:oa}

Setting from [Ch.~\ref{chap:I}]. $\cG_\Gamma$ is the Hecke groupoid, $C^*(\cG_\Gamma)=\cO_{\Anb}$ the
Cuntz--Krieger algebra of the non-backtracking matrix
$\Anb=T^{t}S-J$ [Ch.~\ref{chap:I}, Def.~1.7, Prop.~1.9(ii)], where $S,T\in M_{n\times2m}(\Z)$ are the
start and terminal incidence matrices of oriented edges and $J$ is edge reversal. The
Satake dynamics $\sigma^{(\pi)}$ [Ch.~\ref{chap:I}, Def.~2.1] has partition function
$Z_p(\beta;\pi)=L_p(\beta,\pi_p)$ [Ch.~\ref{chap:I}, Thm.~2.2]; for the Eisenstein parameter
$(\alpha_1,\alpha_2)=(1,q)$ the $\KMS_\beta$ simplex is nonempty exactly at the critical
inverse temperature $\beta=1$, where it is a single Perron state [Ch.~\ref{chap:I}, Thm.~2.4(ii),(iv)];
off the Perron locus the Satake data survive only as $\KMS$ \emph{functionals}
[Ch.~\ref{chap:I}, Thm.~2.5, Rem.~2.6]. Throughout $u=q^{-\beta}$, so $\beta=1$ is $u=q^{-1}$.

\begin{definition}[Character twist]\label{IX:def:twist}
The voltage $\alpha$ is a $\Z$-valued continuous $1$-cocycle $c_\alpha$ on
$\cG_\Gamma$ --- the sum of $\alpha$ along the edges of a Hecke move, well defined on
$(\omega,i-j,\omega')$ because the tails agree --- commuting with the degree cocycle
$c$. It induces a $\mathbb T$-action $\gamma^{\alpha}$ on $C^*(\cG_\Gamma)$ commuting
with $\sigma^{(\pi)}$, and for $t\in\C^{\times}$ the \emph{twisted transfer operator} and
\emph{twisted Hecke operators}
\[
B_t:=B\La_t,\qquad \La_t:=\operatorname{diag}\bigl(t^{\alpha(e)}\bigr)_{e},\qquad
T_k(t)_{uv}:=\sum_{\text{nb paths }u\to v\text{ of length }k}t^{\alpha(\text{path})},
\]
with $B=\Anb$ on $\C^{2m}$; for $|t|=1$ these are the character-twisted Ruelle--Hecke
transfer operator of [Ch.~\ref{chap:I}, Def.~2.1] at the Eisenstein potential and the twisted Hecke
correspondences. Note $T_1(t)=A(t)$ and $\La_tJ\La_t=J$, since $\alpha(\bar e)=-\alpha(e)$.
\end{definition}

\begin{lemma}[Twisted Hecke recursion]\label{IX:lem:rec}
$A(t)T_k(t)=T_{k+1}(t)+qT_{k-1}(t)$ for $k\ge2$, $A(t)^{2}=T_2(t)+(q{+}1)I$, hence
\[
\sum_{k\ge0}T_k(t)u^{k}=(1-u^{2})\bigl(I-uA(t)+qu^{2}I\bigr)^{-1}
\]
as an identity of matrix-valued rational functions of $(u,t)$.
\end{lemma}

\begin{proof}
$A(t)T_k(t)$ sums over an edge $e\colon u\to w$ followed by a non-backtracking path
$w\to v$ of length $k$, with weight $t^{\alpha(e)}t^{\alpha(\text{path})}$. The terms in
which the path begins with $\bar e$ have weight $t^{\alpha(e)+\alpha(\bar e)}
t^{\alpha(\text{rest})}=t^{\alpha(\text{rest})}$, the rest being a non-backtracking path
$u\to v$ of length $k-1$ whose first edge is not $e$; for each such rest there are $q$
admissible $e$. This is the proof of [Ch.~\ref{chap:I}, Lem.~1.5] with the multiplicative weight
carried along; the generating function follows by the same algebra.
\end{proof}

\begin{theorem}[Twisted Bass identity]\label{IX:thm:bass}
With $P(u,t):=\det\bigl(I_n-uA(t)+qu^{2}I_n\bigr)$,
\[
\det\bigl(I_{2m}-uB\La_t\bigr)=(1-u^{2})^{\,r-1}\,P(u,t) .
\]
\end{theorem}

\begin{proof}
Recall $S\La_tT^{t}=A(t)$ (entry $(u,v)$ sums $t^{\alpha(e)}$ over $e\colon u\to v$),
$SJ=T$, $TT^{t}=(q{+}1)I$, and $\La_tJ\La_t=J$. The operator $J\La_t$ is an involution
($(J\La_t)^{2}=J\La_tJ\La_t=JJ=I$) with trace $0$ (no edge is its own reverse), so its
eigenvalues are $\pm1$ with multiplicity $m$ each and
$\det(I+uJ\La_t)=(1-u^{2})^{m}$, with $(I+uJ\La_t)^{-1}=(1-u^{2})^{-1}(I-uJ\La_t)$.
Then $I-uB\La_t=(I+uJ\La_t)-uT^{t}S\La_t$, and by $\det(I-XY)=\det(I-YX)$,
\begin{align*}
\det(I-uB\La_t)
&=(1-u^{2})^{m}\det\Bigl(I_n-\tfrac{u}{1-u^{2}}\,S\La_t(I-uJ\La_t)T^{t}\Bigr)\\
&=(1-u^{2})^{m-n}\det\Bigl((1-u^{2})I_n-u\bigl(A(t)-uSJT^{t}\bigr)\Bigr)\\
&=(1-u^{2})^{m-n}\det\bigl(I-uA(t)+qu^{2}I\bigr),
\end{align*}
using $S\La_tJ\La_tT^{t}=SJT^{t}=TT^{t}=(q{+}1)I$ and $m-n=r-1$. For $t=1$ this is
Bass's identity [Ch.~\ref{chap:I}, Prop.~3.4]; for $|t|=1$ it is the Artin--Ihara $L$-function formula
of \cite{ST}. Verified as a polynomial identity in $(u,t)$ on all four voltage
assignments (Table~\ref{IX:tab:battery}).
\end{proof}

\begin{theorem}[Exceptional zero at the critical temperature]\label{IX:thm:crit}
Let $P(u,t)$ be as above and $Z(\beta;t):=\tr\sum_kT_k(t)q^{-k\beta}$ the twisted
partition function of the Eisenstein Satake dynamics.
\begin{enumerate}
\item[(i)] Functional equation: $P\bigl(\tfrac1{qu},t\bigr)=(qu^{2})^{-n}P(u,t)$.
\item[(ii)] At the critical inverse temperature $\beta=1$,
$P(q^{-1},t)=q^{-n}\,\Theta(t-1)$: the twisted Bass determinant of the Satake dynamics
at $\beta=1$ is the graph $\ell$-adic $L$-function.
\item[(iii)] At $(u,t)=(q^{-1},1)$ the zero is simple in the temperature direction and
of order two in the character direction:
\begin{gather*}
\partial_uP(q^{-1},1)=-\,n(q-1)\,q^{1-n}\,\kappa(Y_0),\\
\partial_tP(q^{-1},1)=0,\qquad
\tfrac12\partial_t^{2}P(q^{-1},1)=q^{-n}\,\cL(Y_\bullet)\,\kappa(Y_0) .
\end{gather*}
The ratio $\tfrac12\partial_t^{2}P/\partial_uP=-\cL(Y_\bullet)/(nq(q-1))
=\|h_\alpha\|^{2}/(nq(q-1))$ is independent of $\kappa(Y_0)$.
\item[(iv)] The unique $\KMS_1$ state of {\rm[Ch.~\ref{chap:I}, Thm.~2.4]} is the $t=1$ member of the
family; for $t\neq1$ on the unit circle with $t^{\alpha(C)}\neq1$ for some cycle $C$
there is no $\KMS_1$ state, only the functional of {\rm[Ch.~\ref{chap:I}, Rem.~2.6]}, and the critical
partition function has a double pole at the trivial character with
\[
Z(1;t)=q\,(1-q^{-2})\,\tr D(t)^{-1}
=\frac{n\,(q-q^{-1})}{\cL(Y_\bullet)}\,\frac1{(t-1)^{2}}+O\bigl((t-1)^{-1}\bigr).
\]
\end{enumerate}
\end{theorem}

\begin{proof}
(i) $I-\frac1{qu}A+\frac{q}{q^{2}u^{2}}I=(qu^{2})^{-1}(qu^{2}I-uA+I)$. (ii) At $u=q^{-1}$,
$I-q^{-1}A(t)+q^{-1}I=q^{-1}D(t)$. (iii) $\partial_tP(q^{-1},t)=q^{-n}f'(t)$ and
$\frac12\partial_t^{2}P(q^{-1},1)=q^{-n}a_2$, so Theorem~\ref{IX:thm:zero} gives the
$t$-derivatives. For the $u$-derivative, $\partial_u\det M(u)=\tr\bigl(\adj M(u)\,M'(u)\bigr)$
with $M(q^{-1})=q^{-1}\Delta_0$, $\adj(q^{-1}\Delta_0)=q^{1-n}\adj\Delta_0
=q^{1-n}\kappa_0\mathbf1\mathbf1^{t}$ (all cofactors of a Laplacian equal $\kappa_0$),
and $M'(q^{-1})=-T_p+2I$, so $\partial_uP=q^{1-n}\kappa_0\mathbf1^{t}(2I-T_p)\mathbf1
=q^{1-n}\kappa_0n(2-(q{+}1))$. (iv) Existence of a $\KMS_1$ state requires
$\rho(q^{-1}B\La_t)=1$ [Ch.~\ref{chap:I}, Thm.~2.4(i)]; for $|t|=1$ the twisted Laplacian $D(t)$ is
Hermitian positive semidefinite with $D(t)v=0$ forcing $|v|$ constant and
$t^{\alpha(C)}=1$ on every cycle, so under the hypothesis $D(t)$ is positive definite,
$P(q^{-1},t)\neq0$, and the pressure equation fails (Perron obstruction, as in
[Ch.~\ref{chap:I}, Thm.~2.5]). By Lemma~\ref{IX:lem:rec},
$Z(1;t)=(1-q^{-2})\tr\bigl(I-q^{-1}A(t)+q^{-1}I\bigr)^{-1}=q(1-q^{-2})\tr D(t)^{-1}$.
The Perron eigenvalue branch of $D(t)$ through $0$ at $t=1$ is
$\lambda(1+s)=-\frac{\|h_\alpha\|^{2}}{n}s^{2}+O(s^{3})$ (this is the Schur-complement
computation in the proof of Theorem~\ref{IX:thm:zero}: $f(1+s)/\det\Delta'$), and the other
$n-1$ eigenvalues are bounded away from zero, so $\tr D(t)^{-1}=
-\frac{n}{\|h_\alpha\|^{2}}(t-1)^{-2}+O((t-1)^{-1})=\frac{n}{\cL(Y_\bullet)}(t-1)^{-2}
+O((t-1)^{-1})$. On $K_4$ the leading coefficients are $-12,-6,-4,-12$
(Table~\ref{IX:tab:battery}).
\end{proof}

\begin{remark}[Which derivative]\label{IX:rem:whichderivative}
The specification asked for ``the second derivative of the KMS functional at $\beta=1$''.
Theorem~\ref{IX:thm:crit} locates the derivatives precisely: at the critical point the
\emph{temperature} derivative is simple and carries the algebraic $L$-value
$\kappa(Y_0)$; the \emph{character} derivative is of second order and carries
$\cL(Y_\bullet)\kappa(Y_0)$; their ratio is the $\cL$-invariant up to the explicit
normalization $nq(q-1)$, with $\kappa(Y_0)$ cancelling --- which is the Greenberg--Stevens
shape \cite{GS}, a ratio of a derivative in the ``weight'' direction to one in the
cyclotomic direction. The specification's request for a cyclic-cohomology formulation
is not met here: the natural candidate pairs the derivation generating $\gamma^{\alpha}$
with the Fredholm module of [Ch.~\ref{chap:II}, Thm.~3.3], and we have not established that pairing.
We record it as open (Remark~\ref{IX:rem:scope}) rather than dress the second derivative in
cohomological language it has not earned.
\end{remark}

\begin{theorem}[The abelian shadow limit carries $\Theta$]\label{IX:thm:shadow}
Let $(Y_k)$ be a connected abelian $\Z/\ell^{k}$ voltage tower, $\cO_k:=\cO_{B(Y_k)}$ the
shadow algebras of {\rm[Ch.~\ref{chap:VI}]} ($c=q$ at every level, all levels being $(q{+}1)$-regular),
and $\pi\colon Y_{k+1}\to Y_k$ the covering maps. Then $\pi$ extends to an in-covering
$\pi^{\sharp}$ of shadow graphs, {\rm[Ch.~\ref{chap:VIII}, Lem.~2.1]} yields unital injective
$*$-homomorphisms $\varphi_k\colon\cO_k\to\cO_{k+1}$, and
$\cO_\infty^{\mathrm{ab}}:=\varinjlim(\cO_k,\varphi_k)$ is a unital Kirchberg algebra with
\[
K_1(\cO_\infty^{\mathrm{ab}})\cong\Z,\qquad
\Tor K_0(\cO_\infty^{\mathrm{ab}})\cong\Hom_{\mathrm{cont}}\Bigl(\varprojlim_k\bigl(\Jac(Y_k),\pi_*\bigr),\Q/\Z\Bigr),
\qquad K_0/\Tor\cong\Z[1/\ell].
\]
The deck transformations induce a continuous action of $\Zl=\varprojlim\Z/\ell^{k}$ on
$\cO_\infty^{\mathrm{ab}}$ commuting with the gauge action, and the $\ell$-primary torsion
of $K_0$ is, as a discrete $\La$-module, the Pontryagin dual of the compact $\La$-module
$\Tor X_\infty$ of {\rm[Ch.~\ref{chap:IV}, Thm.~2.1]}, whose characteristic ideal is $(\Theta)$. The
ideal $(\Theta)$ is stable under the involution $1+T\mapsto(1+T)^{-1}$
(Lemma~\ref{IX:lem:sym}), so the same $\Theta$ --- with its exceptional zero and its
$\cL$-invariant --- is the characteristic element on both sides of the duality.
\end{theorem}

\begin{proof}
A covering map is bijective on in-arcs and on out-arcs at every vertex, so
$\pi^{\sharp}$ ($v\mapsto\pi(v)$, $v'\mapsto\pi(v)'$, shadow arcs label-wise) is an
in-covering, and [Ch.~\ref{chap:VIII}, Lem.~2.1, Thm.~2.3] apply verbatim with $\tau$ replaced by
$\pi$: $K_0(\varphi_k)$ is the pullback $\pi^{*}$ on $\coker\Delta_k$ (the Laplacians
are symmetric, so the transpose convention of [Ch.~\ref{chap:VIII}, Conv.~1.1] is invisible), which is
the Pontryagin adjoint of the norm $\pi_*$ of [Ch.~\ref{chap:IV}, Lem.~1.2--1.3] by
[Ch.~\ref{chap:VIII}, Lem.~1.5]; $\pi^{*}\mathbf1=\mathbf1$ gives unitality, $\deg\pi^{*}=\ell\deg$
gives the free part, and $K_1(\varphi_k)=\mathrm{id}$ as in [Ch.~\ref{chap:VIII}, Thm.~2.3(iii)]. The
limit is Kirchberg by [Ch.~\ref{chap:VIII}, Thm.~3.2] and its $K$-theory follows as in
[Ch.~\ref{chap:VIII}, Thm.~3.3]. Each deck generator $\sigma_k$ is an automorphism of $Y_k^{\sharp}$
covering $\sigma_{k-1}$, hence induces an automorphism of $\cO_k$ with
$\sigma_{k+1}\circ\varphi_k=\varphi_k\circ\sigma_k$ (both send $p_v$ to
$\sum_{\pi(w)=\sigma_kv}p_w$), so $\Zl$ acts on the limit, continuously because on each
$\cO_k$ the action factors through $\Z/\ell^{k}$. On $K_0(\cO_k)=\coker\Delta_k\otimes\Zl$
the deck action is multiplication by $\sigma\in\Zl[G_k]$ [Ch.~\ref{chap:IV}, Lem.~1.1], so the
$\ell$-part of $\Tor K_0(\cO_\infty^{\mathrm{ab}})=\varinjlim\Jac(Y_k)^{\vee}$ is the
Pontryagin dual of $\varprojlim\Jac(Y_k)\otimes\Zl=\Tor X_\infty$ with the
contragredient $\La$-action, and $\car(X_\infty)=(\Theta)$ is [Ch.~\ref{chap:IV}, Thm.~2.1]. The
involution stability is $f(t^{-1})=f(t)$.
\end{proof}

\begin{remark}[Two algebras, one zero]\label{IX:rem:twoalgebras}
The shadow algebras are Hecke-functorial and not Hecke-dynamical [Ch.~\ref{chap:VI}, Rem.~4.3]: their
gauge dynamics reach the Laplacian at $u=1$, i.e.\ $\beta=0$, where pure infiniteness
forbids traces and hence $\KMS_0$ states. So the exceptional zero has two
operator-algebraic homes, and they do not overlap: the \emph{dynamical} side ---
partition functions, the critical temperature $\beta=1$, the double pole of
Theorem~\ref{IX:thm:crit}(iv) --- lives on $C^*(\cG_\Gamma)$ of [Ch.~\ref{chap:I}]; the
\emph{$K$-theoretic} side --- $\Theta$ as characteristic element,
Proposition~\ref{IX:prop:module} --- lives on the equivariant limit
$\cO_\infty^{\mathrm{ab}}$ of Theorem~\ref{IX:thm:shadow}. The specification's request to
read the derivative off ``the gauge action of $\cO_\infty$'' conflates the two; by [Ch.~\ref{chap:II}, Thm.~1.6] no single algebra can carry both.
\end{remark}

\section{The Iwahori tower: a simple zero and no $\cL$-invariant}\label{IX:sec:iwahori}

\begin{theorem}[Residue constancy along the Iwahori tower]\label{IX:thm:iwahori}
For the Iwahori path tower $(Y_k)_{k\ge1}$ of {\rm[Ch.~\ref{chap:II}, \S5]} over a $(q{+}1)$-regular
base with $q=p$, let $\zeta_k(u)^{-1}:=\det(I-uA_k)$. Then:
\begin{enumerate}
\item[(i)] $\zeta_k(u)^{-1}$ is independent of $k\ge1$ and equals
$\det(I-u\Anb)=(1-u^{2})^{r-1}P(u,1)$;
\item[(ii)] $u=q^{-1}$ is a simple zero of $\zeta_k(u)^{-1}$ at every level, with
\[
\frac{d}{du}\det(I-uA_k)\Big|_{u=q^{-1}}=-\,q^{\,2-n_k}\,n_k\,\kappa(Y_k);
\]
\item[(iii)] consequently $n_k\,\kappa(Y_k)\,q^{-n_k}$ is independent of $k\ge1$, which is
the recursion $\kappa(Y_{k+1})=q^{\,n_k(q-1)-1}\kappa(Y_k)$ of {\rm[Ch.~\ref{chap:III}, Thm.~2.1]}, and
comparing with level $0$,
\[
\kappa(Y_1)=\frac{(1-q^{-2})^{\,r-1}(q-1)\,n\,q^{\,n_1-1-n}}{n_1}\,\kappa(Y_0);
\]
\item[(iv)] the pseudo-Iwasawa law {\rm[Ch.~\ref{chap:III}, Thm.~2.2]} reads
$\ord_p\kappa(Y_k)=n_k-\ord_p n_k+\mathrm{const}=\mu p^{k}-k+\nu_0$: the Euler term $-k$
is the $p$-adic valuation of the vertex count $n_k=n(q{+}1)q^{k-1}$, i.e.\ of the
Eisenstein normalization in the residue.
\end{enumerate}
\end{theorem}

\begin{proof}
(i) $A_1=\Anb$, and $\det(zI-A_{k+1})=z^{n_{k+1}-n_k}\det(zI-A_k)$ [Ch.~\ref{chap:III}, eq.~(2.1)] gives
$\det(I-uA_{k+1})=\det(I-uA_k)$; the last equality is Bass [Ch.~\ref{chap:I}, Prop.~3.4]. (ii) With
$M(u)=I-uA_k$, $M(q^{-1})=q^{-1}\Delta_k$ has the simple eigenvalue $0$ (Perron), so the
zero is simple; $\adj(q^{-1}\Delta_k)=q^{1-n_k}\kappa(Y_k)\mathbf1\mathbf1^{t}$ since all
cofactors of the Eulerian Laplacian equal $\kappa(Y_k)$ [Ch.~\ref{chap:III}, proof of Thm.~2.1, Step~2],
and $\tr(\mathbf1\mathbf1^{t}(-A_k))=-\mathbf1^{t}A_k\mathbf1=-qn_k$. (iii) Equate the
derivatives at consecutive levels, using $n_{k+1}=qn_k$; for the level-$0$ formula
differentiate $(1-u^{2})^{r-1}P(u,1)$ at $q^{-1}$ using Theorem~\ref{IX:thm:crit}(iii). (iv)
$\ord_p(n_k\kappa_kq^{-n_k})$ constant gives $\ord_p\kappa_k=n_k-\ord_pn_k+\mathrm{const}$,
and $\ord_pn_k=\ord_p(n(q{+}1))+k-1$. On $K_4$: $n_k\kappa_kq^{-n_k}=\frac98$ at
$k=1,2,3$, the derivative is $-\frac92$ at every level, and
$\kappa(Y_1)=\frac{(3/4)^{2}\cdot4\cdot2^{7}}{12}\cdot16=384$ (Table~\ref{IX:tab:battery}).
\end{proof}

\begin{remark}[No $\cL$-invariant for the non-normal tower]\label{IX:rem:noL}
The specification asked for $\Theta''(0)$, the $\cL$-invariant and the defects
$\varepsilon_j$ ``on the Iwahori path tower as well''. There is no $\Theta$ there: the
tower carries no deck action, hence no character variable, hence no pencil, and the
Hecke variable $u$ is the only direction available. In that direction the zero at the
critical temperature is simple at every level (Theorem~\ref{IX:thm:iwahori}(ii)), with
residue $\kappa(Y_k)$ rather than a second-order coefficient. What replaces the
$\cL$-invariant is the tower constant $n_k\kappa(Y_k)q^{-n_k}$, and what replaces the
bookkeeping identity is the observation that its constancy \emph{is} the growth law.
This gives the Euler term $-k$ its third reading, after the Eisenstein rank removed
once per level [Ch.~\ref{chap:III}, Thm.~2.4] and the Eisenstein index through the connecting map of
the snake [Ch.~\ref{chap:VII}, Cor.~2.2]: it is the $p$-adic valuation of the vertex count, the
normalization $\prod_{\mu\neq0}\mu=n_k\kappa_k$ of [Ch.~\ref{chap:III}, eq.~(2.2)] read $p$-adically.
The three readings are one fact seen from three sides, and none of them is an
$\cL$-invariant.
\end{remark}

\section{Scope, residue, corrections}\label{IX:sec:scope}

\begin{remark}[Solved, open, impossible]\label{IX:rem:scope}
\emph{Solved:} the direction of [Ch.~\ref{chap:III}, \S3], made precise. The double zero is a theorem
with the closed form $\frac12\Theta''(0)=-\kappa(Y_0)\|h_\alpha\|^{2}$; the $\cL$-invariant
is defined, computed, and shown to be a Hodge energy and $\ell$-independent; the
bookkeeping identity holds with an exact description of the defects, and the constant
term is $-\ord_\ell\cL$ in the exceptional-zero regime; the twisted Bass identity and
the critical-temperature theorem give the Mazur--Tate--Teitelbaum and
Greenberg--Stevens shapes on $C^*(\cG_\Gamma)$; the abelian shadow limit carries
$\Theta$ as equivariant characteristic element. \emph{Open:} (a) a cyclic-cohomology
formulation of Theorem~\ref{IX:thm:crit}(iii) (Remark~\ref{IX:rem:whichderivative}); (b) a
converse to Theorem~\ref{IX:thm:defect} --- whether $\delta(Y_\bullet,\ell)=0$ forces
$\lambda_h=0$; (c) a structural formula for the defects $\varepsilon_j$ at $j<j_0$ in
terms of the Newton polygon of $P$ at the levels $\ell^{j}$, beyond the resultant
definition; (d) everything about the non-normal tower that [Ch.~\ref{chap:V}, Prob.~5.1] and
[Ch.~\ref{chap:VII}, Rem.~3.1] left open, which this paper does not touch. \emph{Impossible or
false as stated:} an $\cL$-invariant for the Iwahori tower in the sense of
Definition~\ref{IX:def:L} (Remark~\ref{IX:rem:noL}); the alignment formula of the
specification (Remark~\ref{IX:rem:align}); a single algebra carrying both the dynamical
and the $K$-theoretic side (Remark~\ref{IX:rem:twoalgebras}, by [Ch.~\ref{chap:II}, Thm.~1.6]).
\end{remark}

\begin{remark}[Five corrections to the task specification]\label{IX:rem:corrections}
(1) $\Theta(T)=\det D(1+T)$ lies in $\Z[T,(1+T)^{-1}]\subset\La$, not in $\Zl[T]$; the
cleared $g$ is the polynomial, and $a_2$ is the same for both. (2)
``$\ord_\ell a_2=\mu$ in the $\lambda_h=0$ regime'' is the definition of the regime;
the theorem is Corollary~\ref{IX:cor:regime}. (3) The alignment
``$\ord_\ell\cL\leftrightarrow\nu_0-\ord_\ell\kappa_0+\mu$'' is false;
Theorem~\ref{IX:thm:defect} is the identity that holds. (4) The KMS second derivative is in
the character direction, at $\beta=1$, on $C^*(\cG_\Gamma)$ and not on the shadow limit;
the temperature derivative carries $\kappa(Y_0)$ (Theorem~\ref{IX:thm:crit},
Remarks~\ref{IX:rem:whichderivative}--\ref{IX:rem:twoalgebras}); no cyclic-cohomology
statement is claimed. (5) The Iwahori tower has no $\Theta$, no $\Theta''$, and no
$\varepsilon_j$; Theorem~\ref{IX:thm:iwahori} is what it has. Separately, [Ch.~\ref{chap:III}, Rem.~1.11]
and [Ch.~\ref{chap:VIII}] cite the tail-groupoid $K$-theory as [Ch.~\ref{chap:I}, Prop.~1.7(iii)]; in the compiled
numbering of [Ch.~\ref{chap:I}] it is Proposition~1.9(iii), and [Ch.~\ref{chap:VIII}] has been corrected accordingly.
\end{remark}

\section{The verification battery}\label{IX:sec:battery}

Table~\ref{IX:tab:battery} lists the exact checks of \texttt{exceptional\_zero.py}; the
letters match the docstring of the script and the references in the text. The four
abelian towers are those of [Ch.~\ref{chap:II}, Table~2]; the growth laws were recomputed directly
from the reduced Laplacians of the derived graphs ($k\le5$ for $\ell=2$, up to $128$
vertices; $k\le3$ for $\ell=3$) and independently from the character product of
Lemma~\ref{IX:lem:char}, and the two agree at every level.

\begin{table}[ht]
\centering\scriptsize
\renewcommand{\arraystretch}{1.2}
\begin{tabular}{lcl}
\toprule
check & towers / levels & result\\
\midrule
(Z) $\Theta(0)=\Theta'(0)=0$; $a_2=-8,-16,-24,-8$; $g=-8T^{2}$, $-T^{4}-16T^{3}-16T^{2}$, $-4T^{4}-24T^{3}-24T^{2}$ & 4 towers & exact\\
(H) $a_2=-\kappa_0\|h_\alpha\|^{2}$; $\|h_\alpha\|^{2}=\frac12,1,\frac32,\frac12$; $\cL=-\frac12,-1,-\frac32,-\frac12$ & 4 towers & exact\\
(H$'$) single chord: $\Reff=\frac12$, $8$ spanning trees avoid the chord $=-a_2$ & $(1,0,0)$ & exact\\
(K) double pole of $Z(1;t)$: leading coefficient $=n(q-q^{-1})/\cL=-12,-6,-4,-12$ & 4 towers & exact\\
(M) Newton data $(\mu,\lambda_h)=(3,0),(0,2),(2,2),(0,0)$; regime iff $\ord_\ell a_2=\mu$ & 4 towers & exact\\
(E) $\varepsilon_j$, $j\le5$: all zero except $(\varepsilon_1,\varepsilon_2)=(-1,2)$ for $(1,2,0)$ & 4 towers & exact\\
(G) $\ord_\ell\kappa(Y_k)$ direct $=$ character formula $=$ law; $(\mu,\lambda,\nu_0)$ as [Ch.~\ref{chap:II}, Table~2] & $k\le5$ / $k\le3$ & exact\\
(L) $\nu_0+\ord_\ell\cL=(\ord_\ell a_2-\mu)+\sum\varepsilon_j=0,4,2,0$; regime $\Rightarrow\nu_0=-\ord_\ell\cL$ & 4 towers & exact\\
(B) $\det(I-uB\La_t)=(1-u^{2})^{r-1}P(u,t)$ as a polynomial identity in $(u,t)$ & 4 towers & exact\\
(F) functional equation; $P(q^{-1},t)=q^{-n}\Theta(t-1)$; $\partial_uP=-8$; $\frac12\partial_t^{2}P=q^{-n}a_2$ & 4 towers & exact\\
(S) $\rk D(1)=3=n-1$, $\ord_T\Theta=2$ & 4 towers & exact\\
(I) Iwahori: level-independent $\det(I-uA_k)$; $n_k\kappa_kq^{-n_k}=\frac98$; slope $-\frac92$ at $u=\frac12$; $\kappa(Y_1)=384$ & $k\le3$ & exact\\
\bottomrule
\end{tabular}
\medskip
\caption{Exact verification (\texttt{exceptional\_zero.py}) on $K_4$ ($q=2$, $n=4$,
$r=3$, $\kappa_0=16$). Tower order: $(1,0,0),\ell=2$; $(1,1,1),\ell=2$;
$(1,2,0),\ell=2$; $(1,0,0),\ell=3$. For the fixed towers each line is itself a proof;
the general statements are proved in the text.}
\label{IX:tab:battery}
\end{table}

\section*{Acknowledgement of provenance and caveats}

Unrefereed preprint. Every numbered statement is proved from Papers I--VIII, finite
linear algebra, Weierstrass preparation \cite{Wa} and the character product formula;
the twisted Bass identity is proved rather than quoted from \cite{ST}. The
Mazur--Tate--Teitelbaum and Greenberg--Stevens references \cite{MTT,GS} fix the shape of
an analogy (Remarks~\ref{IX:rem:analogy}, \ref{IX:rem:whichderivative}) and are load-bearing
for nothing. Remark~\ref{IX:rem:corrections} records the corrections to the specification
that prompted the paper and one citation correction propagated from [Ch.~\ref{chap:III}] to [Ch.~\ref{chap:VIII}].
Bibliographic data of \cite{MTT,GS,ST,Ba} were checked against publisher or indexing
pages (titles, venues, years, volumes, pages); \cite{Wa} is the entry verified in [Ch.~\ref{chap:IV}].

% generated from Paper 14/anccft14.tex
\chapter[\texorpdfstring{XIV. The graph $\cL$-invariant as spectral action, Bloch effective mass and Albanese length, and what the cyclic Chern character sees}{XIV. The graph L-invariant as spectral action, Bloch effective mass and Albanese length, and what the cyclic Chern character sees}]{\texorpdfstring{Algorithmic non-commutative class field theory, XIV: the graph $\cL$-invariant as spectral action, Bloch effective mass and Albanese length, and what the cyclic Chern character sees}{Algorithmic non-commutative class field theory, XIV: the graph L-invariant as spectral action, Bloch effective mass and Albanese length, and what the cyclic Chern character sees}}\label{chap:XIV}
\chaptermark{XIV}
\begin{chapterabstract}
[Ch.~\ref{chap:IX}] identified the graph $\cL$-invariant of an abelian voltage tower as
$\cL(Y_\bullet)=-\|h_\alpha\|^{2}$, the negative harmonic energy of the voltage cochain, and
left open ``a cyclic-cohomology formulation''. This paper settles the question in both
directions. Negatively: for the finite-dimensional spectral triple
$(C(V)\oplus C(E^{+}),\ell^{2}(V)\oplus\ell^{2}(E^{+}),\Dp)$ of [Ch.~\ref{chap:II}] the odd cyclic
cohomology vanishes, every Fredholm index is zero, and the pairing proposed in the task
specification is linear in $\alpha$ (it equals $4\sum_e\alpha(e)$); a non-integer, quadratic
quantity is not an index and does not pair with $\HC^{1}$. Positively: $\cL(Y_\bullet)$ is
the second variation of the spectral action of $\Dp$ along the flat $U(1)$-connection
$\varphi\alpha$, in the infrared,
\[
\cL(Y_\bullet)=\frac n2\lim_{t\to\infty}\frac1t\,\partial_\varphi^{2}\Tr e^{-t\Delta_\varphi}\Big|_{\varphi=0},
\]
equivalently minus $n/2$ times the effective mass $\lambda_0''(0)$ of the bottom Bloch band
of the $\Z$-cover $\widetilde Y_\alpha$ whose finite quotients are the tower; equivalently
$-m\sigma^{2}$, where $\sigma^{2}$ is the diffusion constant of the voltage along simple
random walk; equivalently minus the squared length of the class $[\alpha]\in H^{1}(Y_0;\Z)$
in the Jacobian torus of $Y_0$, the dual of the Albanese torus of Kotani--Sunada. The
$\ell$-adic $L$-function $\Theta(T)=\det D(1+T)$ is the Floquet determinant of
$\widetilde Y_\alpha$, and its exceptional zero of order two is the parabolic bottom of the
band; the Kirchhoff identity $\det G=\kappa(Y_0)$ for the cycle Gram matrix gives
$\kappa(Y_0)\cL(Y_\bullet)\in\Z$. All identities are exact and machine-verified on $K_4$
and $K_5$.
\end{chapterabstract}

\section*{Status of the results}

Every numbered statement is proved. The external inputs are the vanishing of odd cyclic
cohomology of finite-dimensional commutative semisimple algebras \cite[Ch.~2]{Lo}, which is
also verified by direct linear algebra, the Markov-chain central limit theorem for
additive functionals of a finite reversible chain, and Kirchhoff's theorem in its cycle
form. The Kotani--Sunada Albanese torus \cite{KS1,KS2} and the spectral action
\cite{CC,JLO,Co} supply names and context, not steps. Section~\ref{XIV:sec:scope} says what
was asked, what is proved, and what is impossible.

\section{Setting}\label{XIV:sec:setting}

$Y_0$ is a finite connected $(q{+}1)$-regular graph with $n$ vertices and $m=\frac{(q+1)n}2$
edges, $E^{+}$ an orientation, $r=m-n+1$, $\partial\colon\Z^{E^{+}}\to\Z^{V}$,
$\partial e=t(e)-o(e)$, $\Delta_0=\partial\partial^{t}=(q{+}1)I-A$, $\kappa_0=\kappa(Y_0)$.
The Hecke--Dirac operator of [Ch.~\ref{chap:II}, Def.~3.1] is
\[
\Dp=\begin{pmatrix}0&\partial\\\partial^{t}&0\end{pmatrix}\ \text{on}\
\ell^{2}(V)\oplus\ell^{2}(E^{+}),\qquad
\Dp^{2}=\Delta_0\oplus\partial^{t}\partial ,
\]
and $\cA:=C(V)\oplus C(E^{+})$ acts diagonally, so $(\cA,\ell^{2}(V)\oplus\ell^{2}(E^{+}),\Dp)$
is an even finite-dimensional spectral triple with grading $\gamma=1\oplus(-1)$. A voltage
$\alpha\in\Z^{E^{+}}$ has harmonic representative $h_\alpha=\Pi\alpha$,
$\Pi=I-\partial^{t}\Delta_0^{+}\partial$ the projection onto the cycle space
$\ker\partial=\ker\Dp\cap\ell^{2}(E^{+})$, and [Ch.~\ref{chap:IX}, Def.~1.3]
\[
\cL(Y_\bullet)=-\|h_\alpha\|^{2}=-\bigl(\|\alpha\|^{2}-(\partial\alpha)^{t}\Delta_0^{+}(\partial\alpha)\bigr)
=\frac{\Theta''(0)}{2\kappa_0},\qquad \Theta(T)=\det D(1+T),
\]
with $D(t)=(q{+}1)I-A(t)$, $A(t)_{uv}=\sum_{e\colon u\to v}t^{\alpha(e)}$, the voltage
pencil. On $K_4$ the values are $-\frac12,-1,-\frac32$ for $(1,0,0),(1,1,1),(1,2,0)$
[Ch.~\ref{chap:IX}, Table~1]; on $K_5$ ($q=3$, $\kappa_0=125$) they are $-\frac35$ for a single edge and
$-\frac75$ for a triangle (Table~\ref{XIV:tab:battery}).

\section{What cyclic cohomology sees, and what it cannot}\label{XIV:sec:nogo}

\begin{proposition}[Three impossibilities]\label{XIV:prop:nogo}
\begin{enumerate}
\item[(i)] $\HC^{1}(\cA)=0$; more precisely the space $Z^{1}_\lambda(\cA)$ of cyclic
$1$-cocycles is zero, so there is no class $\mathrm{Ch}^{1}(\Dp)\in\HC^{1}(\cA)$, and
$\HH_1(\cA)=0$, so there is no class $[\delta_\alpha]\in\HH_1(\cA)$ either.
\item[(ii)] Every index pairing of the even Fredholm module $(\cA,\cH,\Dp)$ with $K_0(\cA)$
or of any compression $PuP$ with a unitary $u$ is an integer, and every Fredholm index of an
operator between finite-dimensional spaces of equal dimension is zero. $\cL(Y_\bullet)\notin\Z$
in general ($-\frac12,-\frac32,-\frac35,-\frac75$), so it is not an index.
\item[(iii)] With $\alpha$ acting by multiplication on $\ell^{2}(E^{+})$,
\[
\lim_{t\to0^{+}}\Trs\bigl([\Dp,\alpha]\,\Dp\,e^{-t\Dp^{2}}\bigr)=\Trs\bigl([\Dp,\alpha]\Dp\bigr)
=2\Tr\bigl(\alpha\,\partial^{t}\partial\bigr)=4\sum_{e\in E^{+}}\alpha(e),
\]
linear in $\alpha$: the formula proposed in the specification does not compute a quadratic
invariant.
\end{enumerate}
\end{proposition}

\begin{proof}
(i) $\cA\cong\C^{N}$, $N=n+m$, with minimal idempotents $e_i$. A cyclic $1$-cocycle
$\tau$ satisfies $\tau(a,b)=-\tau(b,a)$ and $\tau(ab,c)-\tau(a,bc)+\tau(ca,b)=0$. Taking
$(a,b,c)=(e_i,e_i,e_k)$ with $k\ne i$ gives $\tau(e_i,e_k)=0$, and cyclicity gives
$\tau(e_i,e_i)=0$; so $\tau=0$. Since $\cA$ is commutative, $b\colon C^{0}\to C^{1}$ is zero
and $\HC^{1}=Z^{1}_\lambda=0$. This is the case $n=0$ of $\HC^{2n+1}(\C^{N})=0$
\cite[2.1.12]{Lo}; the linear system is solved by machine for $N=10,15$
(Table~\ref{XIV:tab:battery}, line (C)). $\HH_1(\C^{N})=\Omega^{1}_{\C^{N}/\C}=0$ because
$\C^{N}$ is separable. (ii) is linear algebra. (iii) $[\Dp,\alpha]\Dp=
\begin{pmatrix}\partial\alpha\partial^{t}&0\\0&-\alpha\partial^{t}\partial\end{pmatrix}$,
whose supertrace is $\Tr(\partial\alpha\partial^{t})+\Tr(\alpha\partial^{t}\partial)
=2\Tr(\alpha\partial^{t}\partial)$, and $(\partial^{t}\partial)_{ee}=2$. Continuity in $t$
gives the limit. Verified on all six assignments: $4,12,12,4,12$.
\end{proof}

\begin{remark}[What the even Chern character does see]\label{XIV:rem:even}
$\HC^{\mathrm{even}}(\cA)=\HC^{0}(\cA)$ is the space of traces, $\C^{N}$, and the Chern
character of the even Fredholm module $(\cA,\cH,\Dp)$ pairs with $K_0(\cA)=\Z^{N}$ to give
the index of $\Dp$ localized at each idempotent --- the content of [Ch.~\ref{chap:II}, Thm.~3.3]:
$\ker\Dp\cong\Z\oplus\Z^{r}$. This is topology ($b_0$ and $b_1$) and does not depend on
$\alpha$. The $\cL$-invariant depends on the metric (the unit-weight inner product on
edges) and on $\alpha$ quadratically; it lives in the spectral, not the cohomological, side
of the triple. Section~\ref{XIV:sec:action} says where.
\end{remark}

\section{The heat-kernel formulas}\label{XIV:sec:heat}

\begin{theorem}[Infrared, not ultraviolet]\label{XIV:thm:heat}
For every voltage $\alpha$,
\begin{gather*}
\|h_\alpha\|^{2}=\lim_{t\to\infty}\bigl\langle\alpha,e^{-t\Dp^{2}}\alpha\bigr\rangle
=\|\alpha\|^{2}-\int_0^{\infty}\bigl\langle\partial\alpha,e^{-s\Delta_0}\partial\alpha\bigr\rangle\,ds ,\\
\bigl\langle\alpha,e^{-t\Dp^{2}}\alpha\bigr\rangle=\|h_\alpha\|^{2}+\sum_{\lambda>0}e^{-t\lambda}\|P_\lambda\alpha\|^{2},
\end{gather*}
the sum over the nonzero eigenvalues $\lambda$ of $\partial^{t}\partial$ (equivalently of
$\Delta_0$) with spectral projections $P_\lambda$. The ultraviolet limit $t\to0^{+}$ gives
$\|\alpha\|^{2}$, not $\|h_\alpha\|^{2}$.
\end{theorem}

\begin{proof}
$\Pi$ is the projection onto $\ker\partial^{t}\partial$, so $e^{-t\partial^{t}\partial}\to\Pi$
strongly and the spectral expansion follows. For the integral,
$\Delta_0^{+}=\int_0^{\infty}(e^{-s\Delta_0}-P_{\ker\Delta_0})\,ds$ and $\partial\alpha\perp
\mathbf1=\ker\Delta_0$. On $K_4$: $\langle\alpha,e^{-t\Dp^{2}}\alpha\rangle=\frac12+\frac12e^{-4t},
\ 1+2e^{-4t},\ \frac32+\frac72e^{-4t}$; on $K_5$: $\frac35+\frac25e^{-5t}$,
$\frac75+\frac85e^{-5t}$ (Table~\ref{XIV:tab:battery}, line (H), symbolic in $t$).
\end{proof}

\section{The \texorpdfstring{$\Z$}{Z}-cover, its Bloch bands, and the effective mass}\label{XIV:sec:bloch}

\begin{definition}[The $\Z$-cover of the tower]\label{XIV:def:cover}
$\widetilde Y_\alpha$ is the derived graph of $(Y_0,\alpha)$ with voltage group $\Z$: vertices
$V\times\Z$, an edge $(o(e),j)\to(t(e),j+\alpha(e))$ for each $e\in E^{+}$. The level-$k$
graph $Y_k$ of the tower [Ch.~\ref{chap:IV}, \S1] is $\widetilde Y_\alpha/\ell^{k}\Z$. Its Laplacian
$\widetilde\Delta$ is $\Z$-periodic, and the Floquet transform gives
$\ell^{2}(\widetilde Y_\alpha)\cong\int^{\oplus}_{[0,2\pi)}\C^{n}\,\frac{d\varphi}{2\pi}$,
$\widetilde\Delta\cong\int^{\oplus}\Delta_\varphi$, $\Delta_\varphi:=D(e^{i\varphi})$,
the Hermitian twisted Laplacian. Its eigenvalues $\lambda_0(\varphi)\le\dots\le\lambda_{n-1}(\varphi)$
are the \emph{Bloch bands}. The block-character decomposition [Ch.~\ref{chap:IV}, Lem.~1.1] says
\[
\operatorname{spec}\Delta(Y_k)=\bigcup_{\varphi\in2\pi\ell^{-k}\Z/2\pi\Z}\operatorname{spec}\Delta_\varphi ,
\qquad
\Theta(e^{i\varphi}-1)=\det\Delta_\varphi ,
\]
so the tower spectra sample the bands at the $\ell^{k}$-th roots of unity and the
$\ell$-adic $L$-function is the Floquet determinant of the $\Z$-cover.
\end{definition}

\begin{theorem}[Effective mass]\label{XIV:thm:bloch}
$\lambda_0$ is even in $\varphi$, simple near $\varphi=0$, and
\[
\lambda_0(\varphi)=\frac{\|h_\alpha\|^{2}}{n}\,\varphi^{2}+O(\varphi^{4}),\qquad\text{i.e.}\qquad
\cL(Y_\bullet)=-\frac n2\,\lambda_0''(0).
\]
The exceptional zero of order two of $\Theta$ {\rm[Ch.~\ref{chap:IX}, Thm.~1.2]} is the parabolic bottom of
the band: $\det\Delta_\varphi=\lambda_0(\varphi)\prod_{j\ge1}\lambda_j(\varphi)
=n\kappa_0\cdot\frac{\|h_\alpha\|^{2}}n\varphi^{2}+O(\varphi^{3})$, and $\varphi^{2}=-T^{2}+O(T^{3})$
for $T=e^{i\varphi}-1$ gives $a_2=-\kappa_0\|h_\alpha\|^{2}$ again.
\end{theorem}

\begin{proof}
$D(t^{-1})=D(t)^{t}$ [Ch.~\ref{chap:IX}, Lem.~1.1] gives $\operatorname{spec}\Delta_{-\varphi}=
\operatorname{spec}\Delta_\varphi$, so $\lambda_0$ is even and the odd Taylor terms vanish.
At $\varphi=0$ the kernel of $\Delta_0$ is $\C\mathbf1$, simple, with gap $\lambda_1(0)>0$; so
$\lambda_0$ is analytic near $0$ and Rayleigh--Schr\"odinger applies. Write
$\Delta_\varphi=\Delta_0+\varphi D_1+\varphi^{2}D_2+O(\varphi^{3})$ with
$(D_1)_{uv}=-\sum_{e\colon u\to v}i\alpha(e)$, $(D_2)_{uv}=\sum_{e\colon u\to v}\alpha(e)^{2}/2$
(both sums over both orientations). Then $(D_1\mathbf1)_u=i(\partial\alpha)_u$,
$\langle\mathbf1,D_1\mathbf1\rangle=0$, $\langle\mathbf1,D_2\mathbf1\rangle=\|\alpha\|^{2}$, and
\[
\lambda_0(\varphi)=\frac{\varphi^{2}}n\Bigl(\langle\mathbf1,D_2\mathbf1\rangle
-\langle D_1\mathbf1,\Delta_0^{+}D_1\mathbf1\rangle\Bigr)+O(\varphi^{4})
=\frac{\varphi^{2}}n\bigl(\|\alpha\|^{2}-(\partial\alpha)^{t}\Delta_0^{+}\partial\alpha\bigr)+O(\varphi^{4}).
\]
This is the Schur-complement branch $\lambda(1+s)=-\frac{\|h_\alpha\|^{2}}ns^{2}$ of
[Ch.~\ref{chap:IX}, proof of Thm.~3.4] with $s=e^{i\varphi}-1$. The coefficient is computed exactly and
checked to $10^{-12}$ numerically on all six assignments ($\frac18,\frac14,\frac38,\frac3{25},\frac7{25}$;
Table~\ref{XIV:tab:battery}, line (B)).
\end{proof}

\section{The spectral action}\label{XIV:sec:action}

$\Delta_\varphi=\partial_\varphi\partial_\varphi^{*}$ with $\partial_\varphi=\partial U_\varphi$,
$U_\varphi=\operatorname{diag}(e^{i\varphi\alpha(e)})$ on $\ell^{2}(E^{+})$, is the Hodge
Laplacian of the flat $U(1)$-connection with holonomy $e^{i\varphi\alpha(C)}$ around each
cycle $C$; the twisted Dirac operator is
$\Dp^{\varphi}=\begin{pmatrix}0&\partial U_\varphi\\U_\varphi^{*}\partial^{t}&0\end{pmatrix}$,
$(\Dp^{\varphi})^{2}=\Delta_\varphi\oplus U_\varphi^{*}\partial^{t}\partial U_\varphi$, and the
\emph{spectral action} \cite{CC} of the fluctuated triple is $\Tr f\bigl((\Dp^{\varphi})^{2}\bigr)$.

\begin{theorem}[$\cL$ as second variation of the spectral action]\label{XIV:thm:action}
For the heat cutoff $f(x)=e^{-tx}$,
\[
\cL(Y_\bullet)=\frac n2\lim_{t\to\infty}\frac1t\,\partial_\varphi^{2}\Tr e^{-t\Delta_\varphi}\Big|_{\varphi=0}
=\frac n4\lim_{t\to\infty}\frac1t\,\partial_\varphi^{2}\Tr e^{-t(\Dp^{\varphi})^{2}}\Big|_{\varphi=0},
\]
with error $O(t\,e^{-t\lambda_1(0)})$. For an arbitrary smooth cutoff $f$,
\[
\partial_\varphi^{2}\Tr f(\Delta_\varphi)\big|_{0}=f'(0)\,\lambda_0''(0)+\sum_{j\ge1}\partial_\varphi^{2}f\bigl(\lambda_j(\varphi)\bigr)\big|_0 ,
\]
and $\cL(Y_\bullet)=-\frac n2$ times the coefficient of $f'(0)$ in the ground-band term.
\end{theorem}

\begin{proof}
$\Tr e^{-t\Delta_\varphi}=e^{-t\lambda_0(\varphi)}+\sum_{j\ge1}e^{-t\lambda_j(\varphi)}$, the
first term analytic near $0$ by Theorem~\ref{XIV:thm:bloch}, the remainder smooth in $\varphi$
(a symmetric function of the other eigenvalues of an analytic Hermitian family) with
$\varphi$-derivatives bounded by $Ct^{2}e^{-t\lambda_1(0)}$ near $0$. Since
$\lambda_0'(0)=0$, $\partial_\varphi^{2}e^{-t\lambda_0}|_0=-t\lambda_0''(0)=-2t\|h_\alpha\|^{2}/n$.
For the Dirac version,
$(\Dp^{\varphi})^{2}=\partial_\varphi\partial_\varphi^{*}\oplus\partial_\varphi^{*}\partial_\varphi$, and the two
blocks have the same nonzero spectrum, so $\Tr e^{-t(\Dp^{\varphi})^{2}}=2\Tr e^{-t\Delta_\varphi}+(m-n)$.
The general-$f$ statement is the chain rule. Numerically on all six assignments the
$t=10$ value agrees with $\cL$ to $10^{-14}$ (Table~\ref{XIV:tab:battery}, line (S)).
\end{proof}

\begin{remark}[Reading]\label{XIV:rem:reading}
The Chern character of \cite{JLO,Co} is the ultraviolet ($t\to0$), cohomological, integer-valued
side of the spectral triple; Proposition~\ref{XIV:prop:nogo} shows that side is blind to
$\alpha$. The $\cL$-invariant is on the infrared ($t\to\infty$), spectral, rational side:
the curvature of the ground state energy along the moduli of flat connections. On a
Riemannian manifold the analogous quantity is the Hessian of the bottom of the spectrum of
the twisted Laplacian at the trivial character, which is the Albanese metric
\cite{KS1}; the graph statement is Theorem~\ref{XIV:thm:albanese} below. The Greenberg--Stevens
shape of [Ch.~\ref{chap:IX}, Rem.~3.5] --- $\cL$ as the ratio of a second character derivative to a first
temperature derivative --- is here the statement that the band is parabolic while the
Perron eigenvalue moves linearly with $\beta$.
\end{remark}

\section{Diffusion}\label{XIV:sec:diffusion}

Let $(X_i)$ be simple random walk on $Y_0$ started from the uniform distribution, and let
\[
S_N=\sum_{i<N}\alpha(X_i,X_{i+1}),\qquad \alpha(y,x)=-\alpha(x,y),
\]
be the accumulated voltage; on the $\Z$-cover, $S_N$ is the fibre coordinate of the lifted
walk.

\begin{theorem}[$\cL$ as diffusion constant]\label{XIV:thm:diffusion}
$\mathbb E S_N=0$, $\Var(S_N)=N\sigma^{2}+O(1)$ and $S_N/\sqrt N\Rightarrow\mathcal N(0,\sigma^{2})$
with
\[
\sigma^{2}=\frac{\|h_\alpha\|^{2}}{m}=-\frac{\cL(Y_\bullet)}{m}.
\]
\end{theorem}

\begin{proof}
The chain is reversible with uniform stationary law $\pi$ and transition matrix
$P=A/(q{+}1)$. Let $F(x)=\sum_yP(x,y)\alpha(x,y)=-(\partial\alpha)(x)/(q{+}1)$ be the
drift; $\pi(F)=0$. The Poisson equation $(I-P)g=F$, i.e.\ $\Delta_0g=-\partial\alpha$, has
the solution $g=-\Delta_0^{+}\partial\alpha$, and $M_i=\alpha(X_i,X_{i+1})+g(X_{i+1})-g(X_i)$
are stationary martingale increments with $S_N=\sum M_i+g(X_0)-g(X_N)$. The martingale
CLT gives the limit law with $\sigma^{2}=\mathbb E_\pi M_0^{2}=\sum_{x,y}\pi(x)P(x,y)
\bigl(\alpha+\partial^{t}g\bigr)(x,y)^{2}$. Now $\alpha+\partial^{t}g=\alpha-\partial^{t}\Delta_0^{+}\partial\alpha
=h_\alpha$, and summing over both orientations of each edge,
$\sigma^{2}=\frac{2}{n(q+1)}\|h_\alpha\|^{2}=\|h_\alpha\|^{2}/m$. The $O(1)$ term is
$2\mathbb E[g(X_0)^2]-2\mathbb E[g(X_0)g(X_N)]+\dots$, convergent geometrically. The
Poisson-equation value of $\sigma^{2}$ is computed exactly, and $\Var(S_N)$ is computed
exactly for $N\le24$ from the moment-generating matrix over $\Q[s]/(s^{3})$: on $K_4$,
$\sigma^{2}=\frac1{12},\frac16,\frac14$ and $\Var(S_N)-N\sigma^{2}\to\frac1{16},\frac14,\frac7{16}$;
on $K_5$, $\sigma^{2}=\frac3{50},\frac7{50}$ (Table~\ref{XIV:tab:battery}, line (V)).
\end{proof}

This is the one-dimensional case of the Kotani--Sunada central limit theorem for periodic
random walks on crystal lattices \cite{KS1}, where the covariance is the Albanese metric.

\section{Albanese length and integrality}\label{XIV:sec:albanese}

Let $c_1,\dots,c_r$ be a $\Z$-basis of the cycle lattice $H_1(Y_0;\Z)=\ker\partial\cap\Z^{E^{+}}$,
$C=[c_1\cdots c_r]$, $G=C^{t}C$ its Gram matrix, and $a_c=C^{t}\alpha$ the vector of
periods $\langle\alpha,c_j\rangle$, which depends only on the class $[\alpha]\in H^{1}(Y_0;\Z)
=\operatorname{Hom}(H_1(Y_0;\Z),\Z)$.

\begin{theorem}[$\cL$ as Albanese length]\label{XIV:thm:albanese}
\[
\|h_\alpha\|^{2}=a_c^{t}\,G^{-1}a_c=\bigl\|[\alpha]\bigr\|^{2}_{H^{1}},
\]
the squared norm of $[\alpha]$ in $H^{1}(Y_0;\R)$ for the inner product dual to the cycle
inner product on $H_1(Y_0;\R)\subset\R^{E^{+}}$ --- the flat metric of the Jacobian torus
$H^{1}(Y_0;\R)/H^{1}(Y_0;\Z)$, dual to the Albanese torus
$\Alb(Y_0)=H_1(Y_0;\R)/H_1(Y_0;\Z)$ of \cite{KS2} in the unit-weight normalization. Hence
$\cL(Y_\bullet)=-\|[\alpha]\|^{2}_{H^{1}}$ depends only on $[\alpha]$, is invariant under
switching $\alpha\mapsto\alpha+\partial^{t}\varphi$, and
\[
\det G=\kappa(Y_0),\qquad \kappa(Y_0)\,\cL(Y_\bullet)\in\Z ,
\]
which is the integrality $a_2\in\Z$ of {\rm[Ch.~\ref{chap:IX}, Thm.~1.2(ii)]} seen from the cycle side.
\end{theorem}

\begin{proof}
$h_\alpha$ is the orthogonal projection of $\alpha$ onto the column space of $C$, so
$h_\alpha=C(C^{t}C)^{-1}C^{t}\alpha$ and $\|h_\alpha\|^{2}=\alpha^{t}CG^{-1}C^{t}\alpha$. The
harmonic form $h_\alpha$ is the unique cycle with periods $a_c$, so $\|h_\alpha\|$ is the
dual norm of the functional $[\alpha]$ on $(H_1(Y_0;\R),\langle\,,\rangle)$. Kirchhoff's
theorem in its cycle form (the Gram determinant of a fundamental cycle basis equals the
number of spanning trees) gives $\det G=\kappa_0$; a $\Z$-basis change has determinant $\pm1$,
so $\det G=\kappa_0$ for every $\Z$-basis, and $G^{-1}\in\kappa_0^{-1}M_r(\Z)$. Verified:
$\det G=16,125$; $\kappa_0\cL=-8,-16,-24,-75,-175$ (Table~\ref{XIV:tab:battery}, lines (A),(T)).
\end{proof}

\begin{corollary}[Five faces of one number]\label{XIV:cor:faces}
For an abelian voltage tower over $Y_0$ the following are equal:
$\cL(Y_\bullet)=\Theta''(0)/2\kappa_0$; $-\|h_\alpha\|^{2}$; $-\frac n2\lambda_0''(0)$;
$\frac n2\lim_{t\to\infty}t^{-1}\partial^{2}_\varphi\Tr e^{-t\Delta_\varphi}|_0$;
$-m\sigma^{2}$; $-\|[\alpha]\|^{2}_{H^{1}}$.
\end{corollary}

\section{The verification battery}\label{XIV:sec:battery}

\begin{table}[ht]
\centering\scriptsize
\renewcommand{\arraystretch}{1.2}
\resizebox{\textwidth}{!}{%
\begin{tabular}{lcl}
\toprule
check & cases & result\\
\midrule
(C) $\dim Z^{1}_\lambda(\cA)=\dim\HC^{1}(\cA)=0$, $\cA=\C^{10},\C^{15}$ & $K_4,K_5$ & exact\\
(N) $\Trs([\Dp,\alpha]\Dp)=4\sum\alpha(e)=4,12,12,4,12$ & all six & exact\\
(H) $\langle\alpha,e^{-t\Dp^{2}}\alpha\rangle=\|h_\alpha\|^{2}+ce^{-t\lambda_1}$, limit and Duhamel integral $=\|h_\alpha\|^{2}$ & all six & exact, symbolic $t$\\
(B) $\lambda_0(\varphi)=\frac{\|h_\alpha\|^{2}}n\varphi^{2}+O(\varphi^{4})$: exact coefficient $\frac18,\frac14,\frac38,\frac3{25},\frac7{25}$; Richardson numerics & all six & exact $+$ $40$ digits\\
(S) $\frac n2t^{-1}\partial^{2}_\varphi\Tr e^{-t\Delta_\varphi}|_0$ at $t=5,10,20,40$: $\cL$ to $10^{-8},10^{-14},10^{-14},10^{-13}$ & all six & numeric\\
(V) $\sigma^{2}=\|h_\alpha\|^{2}/m=\frac1{12},\frac16,\frac14,\frac3{50},\frac7{50}$; $\Var(S_N)$ exact, $N\le24$ & all six & exact\\
(A) $a_c^{t}G^{-1}a_c=\|h_\alpha\|^{2}$; switching invariance; $\det G=\kappa_0$; $\kappa_0\cL\in\Z$ & all six & exact\\
(T) [Ch.~\ref{chap:IX}, Thm.~1.2(ii)] on $K_5$: $a_2=-75,-175=-\kappa_0\|h_\alpha\|^{2}$ & $K_5$ & exact\\
\bottomrule
\end{tabular}}
\medskip
\caption{Verification (\texttt{cyclic\_index.py}). ``All six'' = $K_4$ with $(1,0,0)$, $(1,1,1)$,
$(1,2,0)$ [the $(1,0,0)$, $\ell=3$ tower has the same voltage and $\cL$] and $K_5$ with a
single-edge and a triangle voltage.}
\label{XIV:tab:battery}
\end{table}

\section{Scope, residue, corrections}\label{XIV:sec:scope}

\begin{remark}[Solved, open, impossible]\label{XIV:rem:scope}
\begin{sloppypar}
\emph{Solved:} open item (a) of [Ch.~\ref{chap:IX}, Rem.~5.1], ``a cyclic-cohomology formulation of
Theorem~3.4(iii)'', in the only form it admits: the $\cL$-invariant is the infrared second
variation of the spectral action of the Hecke--Dirac triple along the flat-connection
direction $\alpha$ (Theorem~\ref{XIV:thm:action}), the effective mass of the $\Z$-cover
(Theorem~\ref{XIV:thm:bloch}), its diffusion constant (Theorem~\ref{XIV:thm:diffusion}), and the
Albanese length of $[\alpha]$ (Theorem~\ref{XIV:thm:albanese}); the exceptional zero is the
parabolic band bottom, and $\kappa_0\cL\in\Z$ is Kirchhoff. \emph{Open:} (a) whether a
nonzero odd cyclic class on a genuinely non-commutative algebra of the tower, such as the
crossed product $C_0(\widetilde Y_\alpha)\rtimes\Z$ or the Cuntz--Krieger algebra of [Ch.~\ref{chap:VIII}],
pairs to $\|h_\alpha\|^{2}$ or to its integer multiple $\kappa_0\|h_\alpha\|^{2}$; the
integrality of Theorem~\ref{XIV:thm:albanese} makes the latter the only candidate for an index.
(b) The rank-two analogue for the $\widetilde A_2$ towers of Chs.~\ref{chap:X}--\ref{chap:XII}, where $b_1=0$ [Ch.~\ref{chap:XII}]
and there is no $\alpha$: the Bloch picture survives (the Hecke pencil of [Ch.~\ref{chap:X}] is a Floquet
determinant over a $2$-torus when a $\Z^{2}$-cover exists) but the $\cL$-invariant does not.
\emph{Impossible as stated:} Proposition~\ref{XIV:prop:nogo}, all three items.
\end{sloppypar}
\end{remark}

\begin{remark}[Corrections to the task specification]\label{XIV:rem:corrections}
(1) There is no ``$\mathrm{Ch}^{1}(\Dp)\in\HC^{1}(\cA)$'' and no ``$[\delta_\alpha]\in H_1(\cA)$'':
both groups vanish for $\cA=C(V)\oplus C(E^{+})$, and every derivation of $\cA$ is zero.
(2) The proposed pairing $\lim_{t\to0^{+}}\Trs([\Dp,\alpha]\Dp e^{-t\Dp^{2}})$ equals
$4\sum_e\alpha(e)$; it is linear in $\alpha$ and cannot equal $-\|\Pi\alpha\|^{2}$. (3) The
$\cL$-invariant is not an index: it is rational and non-integral, and finite-dimensional
Fredholm indices vanish. (4) The correct limit is $t\to\infty$, not $t\to0^{+}$
(Theorem~\ref{XIV:thm:heat}); the $t\to0^{+}$ side of the JLO cocycle is the topological Chern
character, which does not see $\alpha$. (5) The theorem that does hold is a second-variation
(curvature) statement for the spectral action, Theorem~\ref{XIV:thm:action}, and its
geometric name is the Albanese metric, not ``quantum topological index''. (6) The
$(1,0,0)$, $\ell=3$ tower is not a separate case: $\cL$ depends on $\alpha$ only, and
$\ell$ enters through [Ch.~\ref{chap:IX}, \S2] alone. The battery adds $K_5$ ($q=3$) as asked.
\end{remark}

\section*{Acknowledgement of provenance and caveats}

Unrefereed preprint. Every numbered statement is proved from Papers II, IV and IX and
finite linear algebra, plus the Markov-chain CLT and Kirchhoff's theorem in cycle form.
The vanishing of $\HC^{1}(\C^{N})$ is both quoted \cite{Lo} and verified directly.
Bibliographic data of \cite{KS1,KS2,JLO,CC,Lo,Co} were checked against publisher or
indexing pages; \cite{KS1,KS2} are cited for the definition of the Albanese torus and the
crystal-lattice CLT, and no statement of this paper depends on them.

% generated from Paper 19/anccft19.tex
\chapter{Algorithmic non-commutative class field theory, XIX: the odd cyclic class that does not exist, and what the exceptional zero sees instead}\label{chap:XIX}
\chaptermark{XIX}
\begin{chapterabstract}
Paper~XIV showed that the finite-dimensional commutative spectral triple of Paper~II has
$HC^1=0$, so no odd cyclic cocycle there can pair to the graph $\cL$-invariant
$-\|h_\alpha\|^2$, and asked whether an infinite-dimensional non-commutative algebra
would help. Three answers, in increasing order of finality. First, the algebra usually
proposed, $\cA_\alpha=C_0(\Yt)\rtimes\Z$, is Morita-trivial: the $\Z$-cover has vertex set
$V\times\Z$ with the deck group acting freely on the second factor, so
$\cA_\alpha\cong\bigoplus_{v\in V}\cK$, with $K_1=0$ and no odd $K$-theory to pair with at
all. Second, the algebra where Bloch theory does live is the $\Z$-equivariant one,
$M_n(C(S^1))$, with $K_1=\Z$; but its natural class, the voltage pencil
$D(t)=(q+1)I-A(t)$, is \emph{not invertible}, because $D(1)$ is the ordinary Laplacian.
On a circle where the pencil is invertible the odd pairing is the winding number of
$\det D$, and that number is the \emph{order} of the exceptional zero --- equal to $2$ for
every voltage class --- whereas $\kappa_0\cL$ is its leading \emph{coefficient},
$\tfrac12\Theta''(0)=-\kappa_0\|h_\alpha\|^2$. Index theory sees orders, not
coefficients. Third, and this is a theorem rather than a diagnosis: an index pairing is
additive in each slot, so if $\alpha\mapsto\tau_\alpha$ and $\alpha\mapsto[u_\alpha]$ are
both linear and $K_1$ has rank one --- as it does for $M_n(C(S^1))$ and for the shadow
algebra alike --- then $\langle\tau_\alpha,[u_\alpha]\rangle$ is a product of two linear
forms in $\alpha$. It cannot be $\kappa_0\|h_\alpha\|^2$, which is positive definite on
$H^1(Y;\R)$, as soon as $b_1(Y)\ge2$. So the impossibility of Paper~XIV is not an artefact
of finite dimensions. Five checks, on the $K_4$ and $K_5$ voltage classes.
\end{chapterabstract}

\section{The proposed algebra is the compacts}\label{XIX:sec:crossed}

Let $Y$ be a finite connected $(q+1)$-regular graph on $n$ vertices, $\alpha$ an integer
$1$-cochain with $[\alpha]\ne0$, and $\Yt$ the associated $\Z$-cover, with deck group
$\Z$. Write $\kappa_0$ for the number of spanning trees of $Y$, $h_\alpha$ for the
harmonic representative of $[\alpha]$, and $\cL(Y_\bullet)=-\|h_\alpha\|^2$
[Ch.~\ref{chap:XIV}, Def.~1.3].

\begin{proposition}\label{XIX:prop:compacts}
The deck group acts freely on the vertex set of $\Yt$ with exactly $n$ orbits, and
\[
\cA_\alpha=C_0(\Yt)\rtimes\Z\;\cong\;\bigoplus_{v\in V}\cK\bigl(\ell^2(\Z)\bigr).
\]
Consequently $K_0(\cA_\alpha)\cong\Z^{\,n}$, $K_1(\cA_\alpha)=0$, and every odd cyclic
cocycle pairs trivially.
\end{proposition}

\begin{proof}
A connected $\Z$-cover of a connected graph has vertex set $V\times\Z$ with the deck group
translating the second coordinate; the action is free and the orbits are the fibres, $n$
of them (check~(A)). Hence $C_0(\Yt)\cong c_0(V)\otimes c_0(\Z)$ equivariantly, with $\Z$
acting only on the second factor, and
$c_0(V)\otimes\bigl(c_0(\Z)\rtimes\Z\bigr)\cong c_0(V)\otimes\cK(\ell^2(\Z))$, the
crossed product of $c_0$ of a free transitive action being the compacts. The
$K$-theory and the vanishing of $HC^{\mathrm{odd}}$ for a direct sum of copies of $\cK$
are standard.
\end{proof}

So the premise that $\cA_\alpha$ ``has nontrivial $K_1$ and $HC^1$'' is false: it is stably
$\C^{\,n}$, and its only cyclic cohomology is the even part, spanned by the $n$ traces.
The non-commutativity of a crossed product by a \emph{free} action is illusory.

\section{The Floquet algebra, and what its pairing sees}\label{XIX:sec:floquet}

The algebra that does carry Bloch theory is the commutant, not the crossed product. After
Floquet transform the $\Z$-equivariant bounded operators on $\ell^2(\Yt)$ are
$M_n\bigl(C(S^1)\bigr)$, whose $K_1$ is $\Z$, generated by the winding of a determinant.
The natural element of it is the voltage pencil $D(t)=(q+1)I-A(t)$ of [Ch.~\ref{chap:IV}, Lem.~1.1],
whose Floquet determinant is $\Theta$.

\begin{proposition}\label{XIX:prop:notinv}
$D$ is not invertible in $M_n(C(S^1))$: $\det D(1)=0$, since $D(1)$ is the ordinary
Laplacian of $Y$. The zero has order exactly two, and
$\tfrac12\Theta''(0)=-\kappa_0\|h_\alpha\|^2$ [Ch.~\ref{chap:IX}, Thm.~1.2]. Hence the class $[D]\in K_1$
does not exist.
\end{proposition}

\begin{corollary}[What the pairing does see]\label{XIX:cor:winding}
Let $\varepsilon$ be small enough that the circle $|t-1|=\varepsilon$ meets no other zero
of $\det D$. There the pencil is invertible and defines a class in $K_1$, and its pairing
with the generator of $HC^1(C^\infty(S^1))$ is the winding number of $\det D$. That
number is the order of the exceptional zero, namely $2$, for \emph{every} voltage class;
it does not depend on $\alpha$, and so cannot be $\kappa_0\|h_\alpha\|^2$.
\end{corollary}

Check~(W) computes this winding number for the three $K_4$ classes and two $K_5$ classes
and finds $2$ in all five, against values of $\kappa_0\|h_\alpha\|^2$ equal to $8$, $16$,
$24$, $75$ and $175$ (Table~\ref{XIX:tab:data}). The contrast is the whole point: an index
pairing is a topological degree, and a degree reads the \emph{order} of vanishing;
$\cL$ is the \emph{coefficient} at the same point.

\begin{table}[ht]
\centering\small
\begin{tabular}{lrrrrr}
\toprule
 & $b_1$ & $\kappa_0$ & $\|h_\alpha\|^2$ & $\kappa_0\|h_\alpha\|^2$ & winding of $\det D$\\
\midrule
$K_4$, $\alpha=(1,0,0)$ & 3 & 16 & $1/2$ & 8 & 2\\
$K_4$, $\alpha=(1,1,1)$ & 3 & 16 & $1$ & 16 & 2\\
$K_4$, $\alpha=(1,2,0)$ & 3 & 16 & $3/2$ & 24 & 2\\
$K_5$, single edge & 6 & 125 & $3/5$ & 75 & 2\\
$K_5$, triangle & 6 & 125 & $7/5$ & 175 & 2\\
\bottomrule
\end{tabular}
\caption{The cleared pencil has $[T^0]=[T^1]=0$ and $[T^2]=-\kappa_0\|h_\alpha\|^2$ in
every case (check~(P)); the winding number is $2$ in every case (check~(W)).}
\label{XIX:tab:data}
\end{table}

\begin{remark}[Where the integrality comes from]\label{XIX:rem:integrality}
The task specification asked why the index should ``absorb'' $\kappa_0=\det G$. It does
not, and no index is involved: by [Ch.~\ref{chap:XIV}, Thm.~7.1], $\|h_\alpha\|^2=a_c^{\mathsf T}G^{-1}a_c$
with $G$ the cycle Gram matrix and $a_c=C^{\mathsf T}\alpha$ integral, so
\[
\kappa_0\|h_\alpha\|^2=a_c^{\mathsf T}\bigl(\det G\cdot G^{-1}\bigr)a_c
=a_c^{\mathsf T}\operatorname{adj}(G)\,a_c\in\Z ,
\]
the adjugate of an integer matrix being an integer matrix. Integrality is Cramer's rule,
not an index theorem (check~(I)).
\end{remark}

\section{The obstruction}\label{XIX:sec:obstruction}

\begin{theorem}[No linear construction, in any algebra with $K_1$ of rank one]
\label{XIX:thm:obstruction}
Let $A$ be a $C^*$-algebra with $\rank_{\Z}K_1(A)=1$. Suppose given group homomorphisms
\[
H^1(Y;\Z)\to HC^1(A),\ \ \alpha\mapsto\tau_\alpha,
\qquad
H^1(Y;\Z)\to K_1(A),\ \ \alpha\mapsto[u_\alpha].
\]
If $b_1(Y)\ge2$ then $\langle\tau_\alpha,[u_\alpha]\rangle\ne\kappa_0\|h_\alpha\|^2$ for
some $\alpha$.
\end{theorem}

\begin{proof}
Torsion classes pair trivially with a $\C$-valued cocycle, so after tensoring with $\Q$
we may write $[u_\alpha]=\ell(\alpha)\,g$ for a fixed generator $g$ of
$K_1(A)\otimes\Q\cong\Q$ and a linear form $\ell$ on $H^1(Y;\Q)$. Additivity of the
pairing in the second slot gives
$\langle\tau_\alpha,[u_\alpha]\rangle=\ell(\alpha)\,\mu(\alpha)$ with
$\mu(\alpha)=\langle\tau_\alpha,g\rangle$, linear by additivity in the first slot. So the
pairing is a product of two linear forms. If $b_1\ge2$ then $\ker\ell\ne0$, so the product
vanishes at some $\alpha$ with $[\alpha]\ne0$; but $\kappa_0\|h_\alpha\|^2>0$ there, since
$h_\alpha=0$ only for $[\alpha]=0$.
\end{proof}

\begin{corollary}\label{XIX:cor:applies}
The theorem applies to $M_n(C(S^1))$, where $K_1\cong\Z$, and to the shadow algebra
$\mathcal O_{B(Y)}$ of [Ch.~\ref{chap:VI}], where $K_1\cong\Z$ by [Ch.~\ref{chap:VI}, Thm.~2.1]. Both candidate
algebras of the task specification are therefore excluded, for every base graph with
$b_1\ge2$ --- in particular for $K_4$ ($b_1=3$) and $K_5$ ($b_1=6$), the graphs on which
the $\cL$-invariant was computed.
\end{corollary}

Check~(Q) verifies the hypothesis that makes the argument bite: $\kappa_0\|h_\alpha\|^2$
is the quadratic form $\operatorname{adj}(G)$, which is positive definite of rank exactly
$b_1$ --- $3$ for $K_4$ and $6$ for $K_5$ --- so it is not a product of two linear forms.

\begin{remark}[What would be needed]\label{XIX:rem:needed}
Theorem~\ref{XIX:thm:obstruction} is not a proof that no odd cyclic class anywhere computes
$\cL$; it rules out the linear constructions in rank-one $K_1$. Two escapes remain, and we
state them as the honest residue of [Ch.~\ref{chap:XIV}, Rem.~9.1(a)]. \emph{Either} one finds an algebra
whose $K_1$ has rank at least $b_1$, with a natural $H^1(Y;\Z)\to K_1$ --- the pairing may
then be a genuine quadratic form; \emph{or} one abandons linearity, in which case the
assignment $\alpha\mapsto\tau_\alpha$ is not a cohomological construction but part of the
input, and the exercise is circular.

The first escape is a real question and it has an obvious candidate, which we record
because the rank matches exactly. The fundamental group of a connected graph is free of
rank $b_1$, and by Pimsner--Voiculescu
\[
K_0\bigl(C^*_r(F_m)\bigr)\cong\Z,\qquad K_1\bigl(C^*_r(F_m)\bigr)\cong\Z^{\,m},
\]
so $C^*_r\bigl(\pi_1(Y)\bigr)=C^*_r(F_{b_1})$ has $\rank K_1=b_1$ --- precisely the
threshold at which Theorem~\ref{XIX:thm:obstruction} stops applying, and the smallest rank at
which a positive definite form on $H^1(Y;\R)$ could be a pairing at all. Whether the
natural map $H^1(Y;\Z)\cong\Z^{b_1}\to K_1(C^*_r(F_{b_1}))\cong\Z^{b_1}$ and a family
$\tau_\alpha$ can be chosen so that the pairing is $\kappa_0\|h_\alpha\|^2$ we do not
know; we note only that $\kappa_0\|h_\alpha\|^2=a_c^{\mathsf T}\operatorname{adj}(G)a_c$ is
already \emph{a} quadratic form on that lattice, so the shape is right and only the
construction is missing.
\end{remark}

\section{Verification}\label{XIX:sec:battery}

\begin{table}[ht]
\centering\small
\begin{tabular}{clc}
\toprule
key & check & mode\\
\midrule
(A) & Proposition~\ref{XIX:prop:compacts}: free deck action, $n$ orbits & exact\\
(P) & the cleared pencil: $[T^0]=[T^1]=0$, $[T^2]=-\kappa_0\|h_\alpha\|^2$ & exact\\
(W) & Corollary~\ref{XIX:cor:winding}: the winding number is $2$ for every $\alpha$ & numeric\\
(I) & Remark~\ref{XIX:rem:integrality}: $\kappa_0\|h_\alpha\|^2=a_c^{\mathsf T}\operatorname{adj}(G)a_c\in\Z$ & exact\\
(Q) & $\operatorname{adj}(G)$ is positive definite of rank $b_1\ge2$ & exact\\
\bottomrule
\end{tabular}
\caption{The verification battery (\texttt{cyclic\_obstruction.py}), $5/5$, about one
second.}
\label{XIX:tab:battery}
\end{table}

\section{Status, scope, corrections}\label{XIX:sec:scope}

\emph{Proved.} Proposition~\ref{XIX:prop:compacts}, Proposition~\ref{XIX:prop:notinv} (given
[Ch.~\ref{chap:IX}, Thm.~1.2]), Corollary~\ref{XIX:cor:winding}, Theorem~\ref{XIX:thm:obstruction} and
Corollary~\ref{XIX:cor:applies}.

\emph{Not proved.} That no odd cyclic class in any algebra computes $\cL$. What is proved
is that the two algebras named, and every algebra with $K_1$ of rank one, fail for a
structural reason, and that the natural class in the Floquet algebra fails for a different
and more interesting one --- it does not exist, and the object obstructing it is the
exceptional zero whose coefficient is the invariant sought.

\begin{remark}[Solved, open, impossible]\label{XIX:rem:soi}
\emph{Solved:} the two candidate algebras of [Ch.~\ref{chap:XIV}, Rem.~9.1(a)] are ruled out, the first
because it is Morita-trivial and the second by Theorem~\ref{XIX:thm:obstruction}.
\emph{Open:} an algebra with $K_1$ of rank $\ge b_1$ and a natural map from $H^1(Y;\Z)$
(Remark~\ref{XIX:rem:needed}); and, unchanged from [Ch.~\ref{chap:XIV}], the question whether $\cL$ has any
index-theoretic description at all, as opposed to the spectral-action and Albanese
descriptions it is already known to have [Ch.~\ref{chap:XIV}, Cor.~7.2].
\emph{Impossible:} an odd pairing linear in $\alpha$ against a rank-one $K_1$, whenever
$b_1\ge2$. Together with [Ch.~\ref{chap:XIV}, Prop.~2.1] this closes the naive routes in both the
finite-dimensional and the infinite-dimensional directions.
\end{remark}

\begin{remark}[Corrections to the task specification]\label{XIX:rem:corrections}
Three. (1) ``$\cA_\alpha=C_0(\Yt)\rtimes\Z$ has nontrivial $K_1$ and $HC^1$'' is false;
the action is free, so the crossed product is a sum of copies of the compacts
(Proposition~\ref{XIX:prop:compacts}). The specification's own hint is the obstruction to
following it. (2) The proposed mechanism --- pair a cocycle built from the integral period
vector $a_c$ with the unit shift $u\in\Z$ and read off $\|h_\alpha\|^2$ from a heat-kernel
limit --- is linear in $\alpha$ in the $K_1$ slot and cannot produce a positive definite
quadratic form (Theorem~\ref{XIX:thm:obstruction}); this is the same objection
[Ch.~\ref{chap:XIV}, Prop.~2.1(iii)] raised against the finite-dimensional proposal, and it survives the
passage to infinite dimensions unchanged. (3) There is nothing for an index to
``absorb'': $\kappa_0\|h_\alpha\|^2=a_c^{\mathsf T}\operatorname{adj}(G)a_c$ is an integer
by Cramer's rule (Remark~\ref{XIX:rem:integrality}), and no index theorem is needed to see it.
\end{remark}

\section*{Provenance of this chapter}

Unrefereed preprint. Every numbered statement is proved above. The computations verify
the exceptional-zero coefficient and the winding number independently on five voltage
classes, and confirm the positive definiteness that makes
Theorem~\ref{XIX:thm:obstruction} bite. The winding number was computed before the theorem was
formulated, in order to find out what the available pairing actually measures.

% generated from Paper 5/anccft5.tex
\chapter{Algorithmic non-commutative class field theory, V: the degeneracy structure of the Iwahori tower and the Euler correction}\label{chap:V}
\chaptermark{V}
\begin{chapterabstract}
The congruence-model (Iwahori path) tower $Y_0\leftarrow Y_1\leftarrow\cdots$ of
[Ch.~\ref{chap:II}, \S5] is non-normal, carries no deck action, and obeys the exact pseudo-Iwasawa
law $\nu_p=\mu p^k-k+\nu_0$ with $\lambda=-1$, impossible over any commutative
$\Zp[[T]]$ base ([Ch.~\ref{chap:II}, Rem.~5.2], [Ch.~\ref{chap:III}, \S2]). This paper equips that tower with its
canonical algebraic structure and identifies, without fabricating a theory that does
not yet exist, exactly what any future non-commutative characteristic formalism must
specialize to.

The two degeneracy maps --- initial and terminal truncation of non-backtracking paths,
the graph avatars of $X_0(p^{k+1})\rightrightarrows X_0(p^k)$ --- are packaged into an
incidence calculus with six exact identities, among them
$\ta\etp=A_k$ (the level-$k$ Hecke operator realized through level $k{+}1$; at the
bottom, $\ta\etp=T_p$ of [Ch.~\ref{chap:I}]) and $\ta\taup=\et\etp=qI$. Theorem~\ref{V:thm:euler}
computes the degeneracy complex: $\ker(\taup-\etp)=\Z\mathbf1$ and
$\coker(\ta-\et)\cong\Z$ torsion-free, so the Euler rank $\chi=1$ at every level ---
the $-\chi$ of [Ch.~\ref{chap:III}, Thm.~2.4] is the $H^0$ of this complex. Theorem~\ref{V:thm:norm}
shows the tail degeneracy alone intertwines the sandpile Laplacians and the Hecke
operators, is surjective on $p$-primary torsion with kernel of the exact order
$q^{\,n_k(q-1)-1}$ forced by [Ch.~\ref{chap:III}, Thm.~2.1], and hence assembles the tower into one
pro-module $X_\infty=\varprojlim(\coker\Delta_k\otimes\Zp,\ta)$ carrying a limit Hecke
operator; the head degeneracy satisfies the mirror identities for the in-Laplacian
tower and the exact asymmetry $\et\Delta_{k+1}=q(\et-\ta)$, so the structure is
genuinely two-sided. Proposition~\ref{V:prop:nonnormal} records a second, independent
rigorous certificate that no commutative base suffices: $A_kA_k^{\,t}\neq
A_k^{\,t}A_k$ at every computed level. We then state, as axioms rather than as a
theorem, the specialization data any characteristic theory over the Iwahori--Hecke
base must reproduce, and we decline --- for the same reason as in [Ch.~\ref{chap:III}] --- to define
an Akashi-type series whose hypotheses our pro-ring does not satisfy.
\end{chapterabstract}

\section*{Status of the results}

Theorems~\ref{V:thm:euler} and \ref{V:thm:norm} and Propositions~\ref{V:prop:calculus},
\ref{V:prop:mirror} are proved for a general Eulerian $q$-regular strongly connected
level $\ge1$ (the base level $k=0$ is the boundary case, treated in
Remark~\ref{V:rem:level0}); every identity was additionally machine-verified exactly on
the $K_4$ tower at all levels $k\le3$ (Table~\ref{V:tab:battery}).
Proposition~\ref{V:prop:nonnormal} is a finite exact computation on the fixed tower,
hence a proof for it; no general-graph claim is made. Section~\ref{V:sec:axioms} proves
nothing and says so: it is a specification, and Problem~\ref{V:prob:char} is open.

\section{The degeneracy pair and its incidence calculus}\label{V:sec:calculus}

Levels: $Y_0$ a finite connected $(q+1)$-regular graph; for $k\ge1$, $V_{Y_k}$ = the
non-backtracking $k$-paths of $Y_0$, with an arc $p\to p'$ when $p'$ is the shift of a
common $(k{+}1)$-path; $Y_{k+1}=L(Y_k)$ for $k\ge1$, each $Y_k$ Eulerian $q$-regular,
$n_k=n(q+1)q^{k-1}$ ([Ch.~\ref{chap:II}, \S5], [Ch.~\ref{chap:III}, \S2]).

\begin{definition}[Degeneracy maps]\label{V:def:degen}
For $k\ge1$ let $\tau,\eta\colon Y_{k+1}\to Y_k$ be terminal and initial truncation of
paths, $\tau(v_0\cdots v_{k+1})=(v_0\cdots v_k)$, $\eta(v_0\cdots v_{k+1})
=(v_1\cdots v_{k+1})$; under $V_{Y_{k+1}}=\mathrm{Arcs}(Y_k)$ these are the tail and
head maps. Write $\ta,\et\in M_{n_k\times n_{k+1}}(\Z)$ for the pushforwards (sum over
fibres) and $\taup=\ta^{\,t}$, $\etp=\et^{\,t}$ for the pullbacks; thus
$\ta=D_{\mathrm{out}}$, $\et=D_{\mathrm{in}}$, the incidence matrices of
[Ch.~\ref{chap:III}, Thm.~2.1].
\end{definition}

\begin{proposition}[Incidence calculus]\label{V:prop:calculus}
For every $k\ge1$:
\[
\ta\etp=A_k,\quad \et\taup=A_k^{\,t},\quad \etp\ta=A_{k+1},\quad \taup\et=A_{k+1}^{\,t},
\quad \ta\taup=\et\etp=q\,I_{n_k}.
\]
At the boundary level, with $\tau,\eta\colon Y_1\to Y_0$ the endpoint maps on oriented
edges, $\ta\etp=T_p$, the Hecke operator of {\rm[I]}.
\end{proposition}

\begin{proof}
Entrywise bookkeeping of tails and heads; the middle two are [Ch.~\ref{chap:III}, Thm.~2.1, Step~1],
the outer two their transposes, the last two the regularity of degrees. The boundary
identity is $S\,T^{\,t}=T_p$ in the notation of [Ch.~\ref{chap:I}, Thm.~3.7]. All six verified as
exact matrix identities at $k=1,2$ (Table~\ref{V:tab:battery}).
\end{proof}

The pair $(\ta,\et)$ is the graph avatar of the modular degeneracy pair
$X_0(p^{k+1})\rightrightarrows X_0(p^k)$, and Proposition~\ref{V:prop:calculus} is the
statement that the level-$k$ Hecke operator is the correspondence
$\ta\circ\etp$ computed through level $k{+}1$ --- the tower knows its own Hecke
theory.

\section{The degeneracy complex and the Euler rank}\label{V:sec:euler}

\begin{theorem}[Euler rank of the degeneracy complex]\label{V:thm:euler}
For every $k\ge1$, with $\partial_k:=\ta-\et\colon\Z^{V_{Y_{k+1}}}\to\Z^{V_{Y_k}}$ and
$\partial_k^{\,*}:=\taup-\etp$:
\begin{enumerate}
\item[(i)] $\ker\partial_k^{\,*}=\Z\cdot\mathbf1$: the only pulled-back-equal
functions are the constants --- the Perron--Eisenstein line, of rank $\chi_k=1$;
\item[(ii)] $\coker\,\partial_k\cong\Z$ via the degree functional, and it is
torsion-free: the image of $\partial_k$ is exactly the augmentation-zero sublattice.
\end{enumerate}
Hence the Euler rank of the two-term degeneracy complex is $\chi=1$ at every level,
and the correction $-\chi\cdot k=-k$ of {\rm[Ch.~\ref{chap:III}, Thm.~2.4]} is the accumulated
$H^0$ of this complex, once per level.
\end{theorem}

\begin{proof}
(i) $(\partial_k^{\,*}x)_{(a\to b)}=x_a-x_b$; vanishing forces $x$ constant along
arcs, hence constant by strong connectivity. (ii) $\partial_k(e_{(a\to b)})
=e_a-e_b$; over a connected digraph the differences $e_a-e_b$ along arcs generate,
integrally, the full kernel of the degree map (telescope along paths; any
$z$ with $\sum z_i=0$ is an integer combination of consecutive differences), so
$\operatorname{im}\partial_k=\ker(\deg)$ and $\coker\cong\Z$ with no torsion.
Machine confirmation: the Smith form of $\partial_k$ has exactly $n_k-1$ unit
invariants at $k=1,2$ (Table~\ref{V:tab:battery}).
\end{proof}

\section{The tail norm and the pro-module}\label{V:sec:norm}

\begin{theorem}[Tail norm]\label{V:thm:norm}
For every $k\ge1$, with $\Delta_k=qI-A_k$ the out-Laplacian:
\begin{enumerate}
\item[(i)] $\ta\,\Delta_{k+1}=\Delta_k\,\ta$ and $\ta\,A_{k+1}=A_k\,\ta$;
\item[(ii)] $\ta$ is surjective and preserves the degree functional; hence the
induced map on sandpile cokernels is surjective, restricts to a surjection of
$p$-primary torsion, and its kernel there has order exactly
$q^{\,n_k(q-1)-1}$, the increment proved in {\rm[Ch.~\ref{chap:III}, Thm.~2.1]};
\item[(iii)] consequently
$X_\infty:=\varprojlim_k\bigl(\coker_\Z(\Delta_k)\otimes\Zp,\ \ta\bigr)$
is a well-defined pro-module carrying the limit Hecke operator
$A_\infty:=\varprojlim A_k$, and its torsion-size sequence is the pseudo-Iwasawa law
$\ord_p=\mu p^k-k+\nu_0$, $\mu=n(q+1)/q$, of {\rm[Ch.~\ref{chap:III}, Thm.~2.2]} --- now the size
function of a single pro-object rather than a sequence of unrelated groups.
\end{enumerate}
\end{theorem}

\begin{proof}
(i) $\ta A_{k+1}=D_{\mathrm{out}}D_{\mathrm{in}}^{\,t}D_{\mathrm{out}}
=A_kD_{\mathrm{out}}=A_k\ta$ by associativity and
Proposition~\ref{V:prop:calculus}; subtract from $q\ta$. (ii) Every level-$k$ vertex
extends, so $\ta$ hits every basis vector; $\deg\circ\ta=\deg$ since fibres are
summed; torsion-surjectivity then follows exactly as in [Ch.~\ref{chap:IV}, Lem.~1.3] (the kernel of
$\deg$ is the torsion), and the kernel order is forced by counting against
[Ch.~\ref{chap:III}, Thm.~2.1]. (iii) Inverse limits over surjections of compact modules; the Hecke
compatibility is (i). Explicit machine confirmation of (ii) at the transition
$2\to1$: surjective on the $2$-part with kernel order exactly $2^{11}=2^{\,n_1-1}$,
the orders $7,18$ matching [Ch.~\ref{chap:II}, Rem.~5.2] (Table~\ref{V:tab:battery}).
\end{proof}

\begin{proposition}[Head mirror, and the exact asymmetry]\label{V:prop:mirror}
For every $k\ge1$: $\et$ intertwines the in-Laplacians,
$\et\,(qI-A_{k+1}^{\,t})=(qI-A_k^{\,t})\,\et$, so the mirrored construction yields the
in-sandpile pro-module $X_\infty^{\mathrm{in}}$; on the out-tower, however,
\[
\et\,\Delta_{k+1}\;=\;q\,(\et-\ta)\;=\;-q\,\partial_k ,
\]
so $\et$ does not descend to the out-cokernels: the failure is exactly $q$ times the
degeneracy boundary of Theorem~\ref{V:thm:euler}. The tower is thus canonically
two-sided --- a $(\tau,\eta)$-bimodule system, not a single inverse system --- and
the obstruction to collapsing the two sides is again the Euler complex.
\end{proposition}

\begin{proof}
$\et A_{k+1}^{\,t}=D_{\mathrm{in}}D_{\mathrm{out}}^{\,t}D_{\mathrm{in}}
=A_k^{\,t}D_{\mathrm{in}}$; and
$\et\Delta_{k+1}=qD_{\mathrm{in}}-D_{\mathrm{in}}D_{\mathrm{in}}^{\,t}D_{\mathrm{out}}
=q(D_{\mathrm{in}}-D_{\mathrm{out}})$ by
$D_{\mathrm{in}}D_{\mathrm{in}}^{\,t}=qI$. Verified at $k=1,2$.
\end{proof}

\begin{remark}[The boundary level]\label{V:rem:level0}
At $k=0$ the base is undirected and $(q+1)$-regular, and the clean intertwining fails
by exactly the degeneracy boundary:
$\ta\Delta_1-\Delta_0\ta=-(\ta-\et)=-\partial_0$, verified as an exact identity. The
Eisenstein special layer of the whole series --- the $(q+1)$-versus-$q$ mismatch of
[Ch.~\ref{chap:I}, Rem.~1.8], the shared Perron root of [Ch.~\ref{chap:II}], the excluded trivial eigenvalue of
[Ch.~\ref{chap:III}] --- reappears here as the unique level where the norm formalism acquires a
correction term, and the correction is once more $\partial$.
\end{remark}

\begin{proposition}[Second non-commutativity certificate]\label{V:prop:nonnormal}
On the $K_4$ tower, $A_kA_k^{\,t}\neq A_k^{\,t}A_k$ for $k=1,2,3$ (exact
computation, hence a proof for this tower). Together with the growth-theoretic
certificate $\lambda=-1<0$ of {\rm[Ch.~\ref{chap:III}, Prop.~2.3]}, this gives two independent
rigorous witnesses that the level operator algebras
$\Z\langle A_k,A_k^{\,t}\rangle$, and a fortiori any base ring over which the tower's
characteristic theory could live, are non-commutative: the first says the
\emph{growth} cannot come from $\Zp[[T]]$, the second that the \emph{operators}
themselves do not commute with their adjoints.
\end{proposition}

\section{What a characteristic theory must specialize to}\label{V:sec:axioms}

We now state --- as a specification, not a theorem --- the data pinned down by
Papers~II--V that any characteristic formalism $\mathrm{Char}$ for
$(\tau,\eta)$-towers over the (pro-)Iwahori--Hecke base must reproduce:
\begin{enumerate}
\item[(A1)] a multiplicative rank-density term recovering $\mu=n(q+1)/q$, the arc
proliferation of {\rm[Ch.~\ref{chap:III}, Thm.~2.2]};
\item[(A2)] an additive Euler term $-\chi\cdot k$ with
$\chi=\rk\ker(\taup-\etp)$, the $H^0$ of the degeneracy complex
(Theorem~\ref{V:thm:euler}) --- intrinsic to the base, external to any commutative
characteristic ideal;
\item[(A3)] on sub-towers carrying an abelian deck action, reduction to the
commutative characteristic ideal $\bigl(\det D(1+T)\bigr)$ of {\rm[Ch.~\ref{chap:IV}, Thm.~2.1]},
with the Newton-polygon evaluation of {\rm[Ch.~\ref{chap:II}, Rem.~5.4]}.
\end{enumerate}

\begin{problem}\label{V:prob:char}
Construct $\mathrm{Char}$ satisfying (A1)--(A3) on the pro-module
$(X_\infty,A_\infty,\tau,\eta)$ of Theorem~\ref{V:thm:norm}.
\end{problem}

We explicitly decline to define an Akashi-type series here. The non-commutative
Iwasawa formalism in which Akashi series live is built over Iwasawa algebras of
compact $p$-adic Lie groups \cite{CFKSV}; our pro-ring is the inverse limit of the
finite-dimensional, non-normal level algebras of
Proposition~\ref{V:prop:nonnormal}, generated by degeneracy correspondences rather
than by a group, and none of the hypotheses of that formalism are available. Writing
down an ``Akashi series'' whose only property is to output the already-proved growth
law would be circular dressing, and this series does not do that
([Ch.~\ref{chap:III}, Rem.~2.6]). What Papers~II--V have done instead is to make
Problem~\ref{V:prob:char} well-posed: the category, the operators, the two invariants
any answer must produce, and the commutative shadow it must reduce to, are now all
theorems.

\section{The verification battery}\label{V:sec:battery}

\begin{table}[ht]
\centering\small
\renewcommand{\arraystretch}{1.2}
\begin{tabular}{lcl}
\toprule
check & levels & result\\
\midrule
(F) six incidence identities, incl.\ $\ta\etp=A_k$, $\ta\taup=\et\etp=qI$ & $k=1,2$ & exact\\
(T) $\ta\Delta_{k+1}=\Delta_k\ta$,\ $\ta A_{k+1}=A_k\ta$ & $k=1,2$ & exact\\
(H) in-mirror for $\et$; asymmetry $\et\Delta_{k+1}=q(\et-\ta)$ & $k=1,2$ & exact\\
(E) $\ker(\taup-\etp)=\Z\mathbf1$; $\mathrm{SNF}(\ta-\et)$: $n_k{-}1$ units & $k=1,2$ & exact\\
(N) tail norm on $2$-part: surjective, $|\ker|=2^{11}=2^{\,n_1-1}$ & $2\to1$ & exact\\
(C) $A_kA_k^{\,t}\neq A_k^{\,t}A_k$ & $k=1,2,3$ & exact\\
(Z) $\ta\Delta_1-\Delta_0\ta=-\partial_0$;\ \ $\ta\etp=T_p$ at the bottom & $0\to1$ & exact\\
\bottomrule
\end{tabular}
\medskip
\caption{Exact integer verification (scripts \texttt{iw\_stage1--3.py}) on the $K_4$
tower, $q=2$, levels $\le3$ (up to $48$ vertices). Each general lemma is proved in
the text; for the fixed tower each line is itself a proof.}
\label{V:tab:battery}
\end{table}

\section*{Acknowledgement of provenance and caveats}

Unrefereed preprint. All numbered statements are proved, with
Proposition~\ref{V:prop:nonnormal} proved for the fixed $K_4$ tower by exact
computation and asserted for no other graph. Section~\ref{V:sec:axioms} is a
specification; Problem~\ref{V:prob:char} is open, and the paper's deliberate refusal to
define an Akashi-type invariant is explained there. The single external citation
\cite{CFKSV} is contextual and load-bearing for nothing; its bibliographic data are
now primary-source verified.

% generated from Paper 7/anccft7.tex
\renewcommand{\ind}{\operatorname{ind}_{\mathbf1}}
\chapter{Algorithmic non-commutative class field theory, VII: the snake factorization of the transfer and the evaluated characteristic form}\label{chap:VII}
\chaptermark{VII}
\begin{chapterabstract}
Problem 5.1 of [Ch.~\ref{chap:V}] asked for a characteristic formalism on the degeneracy tower
satisfying three axioms: an arc-density term recovering $\mu=n(q+1)/q$ (A1), an
additive Euler term $-\chi k$ with $\chi$ the rank of the degeneracy $H^0$ (A2), and
reduction to the commutative characteristic ideal of [Ch.~\ref{chap:IV}] on abelian sub-towers
(A3). This paper solves the problem in evaluated form, by a single application of
the snake lemma, and the mechanism turns out to be simpler than anything the series
anticipated: because the level-$(k{+}1)$ Hecke operator factors through the tail
norm, $A_{k+1}=\etp\ta$ {\rm[Ch.~\ref{chap:V}, Prop.~1.2]}, it \emph{vanishes identically} on
$\ker\ta$, so the Laplacian restricted to the norm kernel is exactly $qI$ and the
``arc determinant'' is exactly $q^{\,n_k(q-1)}$ --- no nilpotency estimate, no unit
ambiguity. The snake lemma applied to
$0\to\ker\nu\to\Zp^{N_{k+1}}\to\Zp^{N_k}\to0$ with Laplacian verticals then yields,
for every exactly intertwining, surjective, degree-preserving norm $\nu$,
\[
\ord_p\Tor_{k+1}-\ord_p\Tor_k
=\ord_p\bigl|\coker(\Delta_{k+1}|_{\ker\nu})\bigr|-\ord_p\ind(\nu),
\]
where $\ind(\nu)$ is the factor by which $\nu$ scales the Eisenstein line. For
$\nu=\ta$ this gives $n_k(q-1)-1$: axiom (A1) with $\mathrm{Det_{arc}}
=q^{\,n_k(q-1)}$ and axiom (A2) with $\mathrm{Det_{Euler}}=\ind(\ta)=q$, one factor
per level, $\chi=1$. For $\nu=\pi_*$ on an abelian voltage tower the same snake
gives $\ord_\ell\prod_{\text{new }\zeta}\det D(\zeta)-1$: axiom (A3), the character
mechanism of [Ch.~\ref{chap:II}, \S5] and [Ch.~\ref{chap:IV}], verified to the integer
($\det(\Delta|_{\ker\pi})=48$ and $784$ against the character products). One
formula, two specializations, all three axioms now theorems. What is not delivered
--- a lift of this evaluated form to a single class in $K_1$ of a localized
completed degeneracy algebra --- is stated as the precise residue of Problem 5.1,
and, consistently with [Ch.~\ref{chap:III}] and [Ch.~\ref{chap:V}], we decline to axiomatize a localization theory
that does not exist.
\end{chapterabstract}

\section*{Status of the results}

All numbered statements are proved; the proofs use only results of Papers I--VI and
the snake lemma. Every identity was additionally machine-verified exactly
(Table~\\ref{VII:tab:battery}): Iwahori transitions $1\to2$, $2\to3$ over $K_4$, and
abelian transitions $0\to1$, $1\to2$ of the voltage tower $(1,1,1)$, $\ell=2$.
Remark~\\ref{VII:rem:residue} states what remains open of [Ch.~\ref{chap:V}, Prob.~5.1];
Remark~\\ref{VII:rem:corrections} records two corrections to the task specification that
prompted this paper. There are no external citations.

\section{The kernel of the norm, and the vanishing lemma}\label{VII:sec:kernel}

Setting as in [Ch.~\ref{chap:V}]: $Y_k$ the Iwahori path tower over a $(q+1)$-regular base
($n_k=n(q+1)q^{k-1}$ for $k\ge1$), $\ta$ the tail norm, $A_k$ the level Hecke
(adjacency) operator, $\Delta_k=qI-A_k$ for $k\ge1$, sandpile torsion
$\Tor_k$.

\begin{lemma}[Vanishing on the norm kernel]\label{VII:lem:vanish}
For $k\ge1$ let $K_k:=\ker_{\Zp}(\ta)\subseteq\Zp^{\,n_{k+1}}$. Then:
\begin{enumerate}
\item[(i)] $K_k$ is a saturated direct summand of rank $n_{k+1}-n_k=n_k(q-1)$
(kernels of integer matrices are saturated);
\item[(ii)] $A_{k+1}$ vanishes on $K_k$: since $A_{k+1}=\etp\,\ta$
{\rm[Ch.~\ref{chap:V}, Prop.~1.2]}, $A_{k+1}x=\etp(\ta x)=0$ for $x\in K_k$; in particular $K_k$
is invariant and
\[
\Delta_{k+1}\big|_{K_k}=q\,\mathrm{id}_{K_k},
\qquad
\bigl|\coker(\Delta_{k+1}|_{K_k})\bigr|=q^{\,n_k(q-1)}\ \text{exactly.}
\]
\end{enumerate}
\end{lemma}

\begin{proof}
(i) is the saturation principle of [Ch.~\ref{chap:II}] (if $ka\in\ker$, $k\neq0$, then
$a\in\ker$); the rank is surjectivity of $\ta$. (ii) is the displayed one-line
computation; then $\Delta_{k+1}=qI-A_{k+1}$ restricts to $q\,\mathrm{id}$, whose
cokernel on a rank-$n_k(q-1)$ lattice has order exactly $q^{\,n_k(q-1)}$. Verified:
the machine search for the nilpotency index of $A_{k+1}|_{K_k}$ returned $s=1$ at
both computed transitions, and $\det(\Delta_{k+1}|_{K_k})=2^{12},2^{24}$ on the
nose.
\end{proof}

Lemma~\\ref{VII:lem:vanish}(ii) is the structural surprise of this paper: the series
expected, at best, a nilpotency estimate through the characteristic-polynomial ratio
of [Ch.~\ref{chap:III}, eq.~(2.1)]; instead the incidence calculus of [Ch.~\ref{chap:V}] kills the operator
outright, because the Hecke operator of the upper level literally factors through
the norm to the lower one.

\section{The snake factorization}\label{VII:sec:snake}

\begin{theorem}[Snake factorization of the transfer]\label{VII:thm:snake}
Let $\nu\colon\Zp^{N_{k+1}}\to\Zp^{N_k}$ be surjective and degree-preserving, with
$\nu\Delta_{k+1}=\Delta_k\nu$, and suppose the Eisenstein lines satisfy
$\nu(\mathbf1)=c\,\mathbf1$ for an integer $c=:\ind(\nu)\neq0$. Then $K:=\ker\nu$ is
a saturated invariant summand, the snake lemma applied to
\[
0\longrightarrow K\longrightarrow \Zp^{N_{k+1}}
\xrightarrow{\ \nu\ }\Zp^{N_k}\longrightarrow0
\]
with verticals $\Delta_{k+1}|_K$, $\Delta_{k+1}$, $\Delta_k$ yields the exact
sequence
\[
0\to\Zp\xrightarrow{\ \cdot c\ }\Zp\xrightarrow{\ \delta\ }
\coker\bigl(\Delta_{k+1}|_K\bigr)\to\coker\Delta_{k+1}\to\coker\Delta_k\to0 ,
\]
and hence, whenever $\coker(\Delta_{k+1}|_K)$ is finite,
\begin{equation}\label{VII:eq:master}
\ord_p\Tor_{k+1}-\ord_p\Tor_k
\;=\;\ord_p\bigl|\coker(\Delta_{k+1}|_{K})\bigr|\;-\;\ord_p\,c .
\end{equation}
\end{theorem}

\begin{proof}
Invariance of $K$ is the intertwining; saturation as in
Lemma~\\ref{VII:lem:vanish}(i). The kernels of the verticals are:
$\ker(\Delta_{k+1}|_K)=K\cap\Zp\mathbf1=0$ (the constants are not in $\ker\nu$
since $\nu\mathbf1=c\mathbf1\neq0$), $\ker\Delta_{k+1}=\Zp\mathbf1$,
$\ker\Delta_k=\Zp\mathbf1$, and the induced map between them is
$\mathbf1\mapsto c\mathbf1$. The snake lemma gives the displayed sequence; the free
ranks of the two cokernels match through $\nu$ (degree preservation, as in
[Ch.~\ref{chap:IV}, Lem.~1.3]), so taking orders of the finite parts yields \\eqref{VII:eq:master}.
\end{proof}

\begin{corollary}[Iwahori specialization: axioms (A1) and (A2)]\label{VII:cor:iwahori}
For $\nu=\ta$, $k\ge1$: $c=q$ $($[V], $\ta\mathbf1=q\mathbf1)$ and
Lemma~\\ref{VII:lem:vanish} give
\[
\ord_p\Tor_{k+1}-\ord_p\Tor_k=n_k(q-1)-1 ,
\]
recovering {\rm[Ch.~\ref{chap:III}, Thm.~2.1]} homologically. Summation yields the pseudo-Iwasawa
law with $\mu=n(q+1)/q$ carried entirely by
$\mathrm{Det_{arc}}:=\det(\Delta_{k+1}|_{\ker\ta})=q^{\,n_k(q-1)}$ (axiom A1), and
the Euler term $-\chi k$ carried entirely by
$\mathrm{Det_{Euler}}:=\ind(\ta)=q$, one factor per level, with
$\chi=1=\rk\ker(\tau^*-\eta^*)$ by {\rm[Ch.~\ref{chap:V}, Thm.~2.1]} (axiom A2): the $-1$ is,
finally and precisely, the Eisenstein line passing through the connecting map of the
snake.
\end{corollary}

\begin{corollary}[Abelian specialization: axiom (A3)]\label{VII:cor:abelian}
For an abelian $\Z/\ell^{k}$ voltage tower with covering pushforward $\nu=\pi_*$
$($exactly intertwining by {\rm[Ch.~\ref{chap:IV}, Lem.~1.2]}$)$: $c=\ell$, and under the
block-character decomposition of {\rm[Ch.~\ref{chap:IV}, Lem.~1.1]} the kernel $\ker\pi_*$ is the
sum of the \emph{new} character blocks, on which $\Delta$ acts with determinant
$\prod_{\text{new }\zeta}\det D(\zeta)$; hence \\eqref{VII:eq:master} reads
\[
\ord_\ell\Tor_{k+1}-\ord_\ell\Tor_k
=\ord_\ell\!\!\prod_{\text{new }\zeta}\!\det D(\zeta)\;-\;1 ,
\]
which is precisely the character/Newton mechanism of {\rm[Ch.~\ref{chap:II}, \S5]} and
{\rm[Ch.~\ref{chap:IV}, Cor.~2.3]}. Verified exactly: $\det(\Delta|_{\ker\pi})=48=\det D(-1)$ at
$0\to1$ and $784=|\det D(i)\det D(-i)|$ at $1\to2$, increments $4-1=3$ in both,
matching {\rm[Ch.~\ref{chap:II}, Table~2]}.
\end{corollary}

One snake, two kernels: on the Iwahori side the kernel is killed by the Hecke
operator and contributes the pure arc power $q^{\,n_k(q-1)}$; on the abelian side
the kernel is the new-character space and contributes the $L$-value block
$\prod_{\text{new}}\det D(\zeta)$. The Eisenstein index $c$ is $q$ or $\ell$
respectively --- one factor per level in either world, which is the whole content of
$\lambda=-\chi$. The three axioms of [Ch.~\ref{chap:V}, \S5] are now theorems with these referents.

\section{Scope, residue, corrections}\label{VII:sec:scope}

\begin{remark}[The residue of Problem 5.1]\label{VII:rem:residue}
Theorem~\\ref{VII:thm:snake} solves [Ch.~\ref{chap:V}, Prob.~5.1] \emph{in evaluated form}: the
characteristic form is the level-wise pair
$\mathrm{Char}_k:=\bigl(\det(\Delta_{k+1}|_{\ker\nu}),\ \ind(\nu)\bigr)$, and
(A1)--(A3) hold for it as theorems. What is not constructed is a single element of
$K_1$ of a localized completion of the degeneracy algebra --- a Dieudonn\'e-type
determinant --- whose evaluations reproduce $\mathrm{Char}_k$. Defining such an
object would require an Ore localization theory for the non-normal pro-ring of
[Ch.~\ref{chap:V}, Prop.~4.1], which does not exist; consistently with [Ch.~\ref{chap:III}, Rem.~2.6] and
[Ch.~\ref{chap:V}, \S5], we decline to axiomatize it, and we record the lift as the precise open
residue. Everything quantitative that a lift could evaluate to is already proved
here.
\end{remark}

\begin{remark}[Two corrections to the task specification]\label{VII:rem:corrections}
The specification prompting this paper asserted an exact sequence
$0\to X_{k+1}\xrightarrow{\partial}X_k\to0$ with $\partial=\ta-\et$; that map is far
from injective ($\rk\ker\partial=n_{k+1}-n_k+1$ by {\rm[Ch.~\ref{chap:V}, Thm.~2.1(ii)]}), and the
exact sequence that carries the theory is the norm sequence of
Theorem~\\ref{VII:thm:snake}. It also placed the Eisenstein class in
$K_0(\mathcal H_k)$; the Eisenstein \emph{projector} $\mathbf1\mathbf1^{t}/N$ is not
integral and does not lie in the algebra, and the class that exists is the
module-level line $\Zp\mathbf1=\ker\Delta$, entering the theory through the
connecting map --- which is where Corollary~\\ref{VII:cor:iwahori} puts it.
\end{remark}

\begin{remark}[Boundary level]\label{VII:rem:boundary}
At $k=0$ the tail map fails to intertwine by exactly the degeneracy boundary
$($[Ch.~\ref{chap:V}, Rem.~3.3]$)$, $\ker S$ is not invariant, and Theorem~\\ref{VII:thm:snake} does not
apply; the Eisenstein index there is $q+1$, not $q$
$(S\mathbf1=(q+1)\mathbf1$, verified$)$. The transfer $0\to1$ enters the law as the
initial value, as everywhere in the series.
\end{remark}

\section{The verification battery}\label{VII:sec:battery}

\begin{table}[ht]
\centering\small
\renewcommand{\arraystretch}{1.2}
\begin{tabular}{lccccc}
\toprule
transition & $\rk\ker\nu$ & invariant & $A_{k+1}|_{\ker}$ &
$\det(\Delta|_{\ker\nu})$ & $\ind(\nu)$\\
\midrule
Iwahori $1\to2$ & $12$ & yes & $=0$ ($s{=}1$) & $2^{12}$ exactly & $2$\\
Iwahori $2\to3$ & $24$ & yes & $=0$ ($s{=}1$) & $2^{24}$ exactly & $2$\\
abelian $0\to1$ & $4$ & yes & --- & $48=\det D(-1)$ & $2$\\
abelian $1\to2$ & $8$ & yes & --- & $784=|\det D(i)\det D(-i)|$ & $2$\\
boundary $0\to1$ & --- & no & --- & --- & $3=q+1$\\
\bottomrule
\end{tabular}
\medskip
\caption{Exact verification (\texttt{iwahori\_char.py}). Increments
\\eqref{VII:eq:master}: $12-1=11$, $24-1=23$ (Iwahori, matching [Ch.~\ref{chap:II}, Rem.~5.2]);
$4-1=3$, $4-1=3$ (abelian $2$-adic orders, matching [Ch.~\ref{chap:II}, Table~2]). The vanishing
$A_{k+1}|_{\ker\ta}=0$ was found by the machine ($s=1$) before its one-line proof
was noticed.}
\label{VII:tab:battery}
\end{table}

\section*{Acknowledgement of provenance and caveats}

Unrefereed preprint. Every numbered statement is proved from Papers I--VI and the
snake lemma; the battery of Table~\\ref{VII:tab:battery} confirms each ingredient
exactly. No external sources are cited and none are load-bearing.
Remark~\\ref{VII:rem:residue} states the open residue of [Ch.~\ref{chap:V}, Prob.~5.1];
Remark~\\ref{VII:rem:corrections} corrects the specification that prompted the paper.

% generated from Paper 13/anccft13.tex
\renewcommand{\HS}{\cH_\Sigma}
\chapter[\texorpdfstring{XIII. Universal localization of the degeneracy algebra, the $K_1$-characteristic class, and the
Eisenstein boundary}{XIII. Universal localization of the degeneracy algebra, the K1-characteristic class, and the Eisenstein boundary}]{\texorpdfstring{Algorithmic non-commutative class field theory, XIII: universal localization of the degeneracy algebra, the $K_1$-characteristic class, and the
Eisenstein boundary}{Algorithmic non-commutative class field theory, XIII: universal localization of the degeneracy algebra, the K1-characteristic class, and the Eisenstein boundary}}\label{chap:XIII}
\chaptermark{XIII}
\begin{chapterabstract}
Problem 5.1 of [Ch.~\ref{chap:V}] and its residue [Ch.~\ref{chap:VII}, Rem.~3.1] asked for a single class in $K_1$ of a
localized degeneracy algebra whose evaluations reproduce the snake factorization of [Ch.~\ref{chap:VII}]:
the arc term $q^{\,n_k(q-1)}$ and the Eisenstein index $q$ at every level of the Iwahori tower.
Ore localization is unavailable because the degeneracy algebra is non-commutative and
non-normal ([Ch.~\ref{chap:V}, Prop.~4.1]). This paper constructs the class with Cohn's universal
localization, and finds that the construction reorganizes the whole picture. The degeneracy
order $\cH$ is the $\Zp$-algebra generated by the four degeneracy maps $\ta,\taup,\et,\etp$
acting on the augmentation-zero lattices $\cM^{0}_k=\ker(\deg)\subset\Zp^{n_k}$; the tail
Laplacian factors inside it as $\Delta_k=\ta(\taup-\etp)$, tail norm times degeneracy
coboundary. On $\cM^{0}_k$ the Laplacian is injective with finite cokernel, so
$\Delta=\bigoplus_k\Delta_k$ is $\Sigma$-inverting for $\Sigma=\{\Delta\}$; the universal
localization $\cH\to\HS$ exists, the low-degree localization sequence
$K_1(\cH)\to K_1(\HS)\xrightarrow{\partial}K_0(\cH,\Sigma)\to K_0(\cH)$ holds, and
$[\Theta_{\cH}]:=[\Delta]\in K_1(\HS)$ is the characteristic class. Its content is in two
theorems. The level evaluations $\ev_k\colon K_1(\HS)\to\Qp^{\times}$ give
$\ev_k[\Theta_{\cH}]=\pm\pdet(\Delta_k)=\pm n_k\kappa(Y_k)$, and on $K_4$ these have
$2$-adic order exactly $n_k-3$ at $k=1,2,3$: the $K_1$ class carries the pure arc growth
$\mu p^{k}$ and nothing else. The boundary class $\partial[\Theta_{\cH}]$ at level $k$ is the
cokernel of $\Delta_k$ on $\cM^{0}_k$, which sits in a non-split extension
$0\to\Z/n_k\to\coker\Delta_k^{0}\to\Jac(Y_k)\to0$: the Euler term $-k=-\ord_pn_k+\mathrm{const}$
of [Ch.~\ref{chap:III}] is the Eisenstein subclass $[\Z/n_k]$ of the boundary. Axioms (A1) and (A2) of [Ch.~\ref{chap:V}]
are thus separated inside a single $K$-theoretic object, and (A3) holds on abelian towers,
where $\pdet(\Delta_k)=n\kappa(Y_0)\prod_{\chi\ne1}\det D(\chi)$ is the characteristic
element of [Ch.~\ref{chap:IV}] evaluated at all nontrivial characters (exact on the $(1,1,1)$ tower). What
the specification called the ``non-commutative main conjecture'' $\partial[\Theta]=[X]$ is
the definition of $\partial$; the genuine main conjecture --- an analytic $p$-adic
$L$-function of the tower, defined independently, equal to $[\Theta_{\cH}]$ --- is stated and
left open.
\end{chapterabstract}

\section*{Status of the results}

Every numbered statement is proved. The external inputs are the existence and universal
property of Cohn's universal localization \cite{Co} and the low-degree localization exact
sequence for universal localizations \cite[Ch.~4]{Sch}; the higher sequence of
Neeman--Ranicki \cite{NR} requires stable flatness, which we do not verify and do not use.
Everything else is finite linear algebra over $\Zp$, machine-verified exactly on the $K_4$
Iwahori tower at $k\le3$ and on the $(1,1,1)$ abelian tower (Table~\ref{XIII:tab:battery}).
Section~\ref{XIII:sec:scope} says what the main conjecture is and is not.

\section{The degeneracy order and the augmentation-zero lattices}\label{XIII:sec:order}

Setting as in [Ch.~\ref{chap:V}], [Ch.~\ref{chap:VII}]: $Y_k$ the Iwahori path tower over a $(q{+}1)$-regular base,
$q=p$, $n_k=n(q{+}1)q^{k-1}$, $A_k$ the level Hecke operator, $\Delta_k=qI-A_k$,
$\Tor_k=\Jac(Y_k)=\Tor\coker\Delta_k$; $\ta,\et\in M_{n_k\times n_{k+1}}(\Z)$ the tail and head
pushforwards, $\taup,\etp$ their transposes, with [Ch.~\ref{chap:V}, Prop.~1.2]
\begin{equation}\label{XIII:eq:calc}
\ta\etp=A_k,\qquad \etp\ta=A_{k+1},\qquad \ta\taup=\et\etp=qI_{n_k},
\end{equation}
and $\deg\colon\Z^{n_k}\to\Z$ the coordinate sum.

\begin{lemma}[Degree compatibility and the Laplacian factorization]\label{XIII:lem:deg}
$\deg\circ\ta=\deg\circ\et=\deg$ and $\deg\circ\taup=\deg\circ\etp=q\deg$; hence all four
degeneracy maps preserve the augmentation-zero lattices $\cM^{0}_k:=\ker(\deg)\subset\Zp^{n_k}$.
Moreover
\[
\Delta_k=\ta\,(\taup-\etp)=\ta\,\partial_k^{*},
\]
the tail norm composed with the degeneracy coboundary $\partial_k^{*}$ of {\rm[Ch.~\ref{chap:V}, Thm.~2.1]};
in particular $\Delta_k$ lies in the algebra generated by the degeneracy maps and maps
$\Zp^{n_k}$ into $\cM^{0}_k$.
\end{lemma}

\begin{proof}
Pushforwards sum over fibres and pullbacks have $q$ elements per fibre. By \eqref{XIII:eq:calc},
$\ta\taup-\ta\etp=qI-A_k$. Column sums of $\Delta_k$ vanish (in-degree $q$). Verified exactly
at $k=1,2,3$ (Table~\ref{XIII:tab:battery}, lines (I),(B)).
\end{proof}

\begin{definition}[Degeneracy order]\label{XIII:def:order}
For $K\ge1$ let $\cM^{0}_{\le K}:=\bigoplus_{k\le K+1}\cM^{0}_k$, and let
$\cH_K\subset\End_{\Zp}(\cM^{0}_{\le K})$ be the $\Zp$-subalgebra generated by the degeneracy
maps between consecutive levels $\le K{+}1$, restricted to the augmentation-zero lattices by
Lemma~\ref{XIII:lem:deg}. It is an order in a semisimple $\Qp$-algebra (a $\Zp$-subalgebra of a
matrix algebra, finitely generated as a module), non-commutative and non-normal since
$A_kA_k^{t}\ne A_k^{t}A_k$ inside it {\rm[Ch.~\ref{chap:V}, Prop.~4.1]}. Put
\[
\Delta_{(K)}:=\bigoplus_{k\le K}\Delta_k^{0}\ \oplus\ q\cdot\mathrm{id}_{\cM^{0}_{K+1}}\ \in\cH_K,
\qquad \Delta_k^{0}:=\Delta_k|_{\cM^{0}_k},
\]
where the last summand $q\,\mathrm{id}=\ta\taup$ keeps the element inside $\cH_K$. We write
$\cH$ for $\cH_K$ with $K$ fixed and large, and $\Sigma:=\{\Delta_{(K)}\}$.
\end{definition}

\begin{proposition}[$\Sigma$-inverting]\label{XIII:prop:sigma}
$\Delta_k^{0}$ is injective with finite cokernel and invertible over $\Qp$, with
\[
\det\Delta_k^{0}=\pm\,\pdet(\Delta_k)=\pm\,n_k\kappa(Y_k),
\]
the product of the nonzero eigenvalues of $\Delta_k$. Hence $\Delta_{(K)}\otimes\Qp$ is
invertible while $\Delta_{(K)}$ is not invertible in $\cH$, its determinant not being a
$\Zp$-unit: localization is needed and is non-trivial.
\end{proposition}

\begin{proof}
$\ker\Delta_k=\Zp\mathbf1$ and $\deg\mathbf1=n_k\ne0$, so $\Delta_k$ is injective on
$\ker(\deg)$; $\Delta_k$ maps $\Zp^{n_k}$ into $\ker(\deg)$, so $\Delta_k^{0}$ is an
endomorphism of the rank-$(n_k-1)$ lattice $\cM^{0}_k$ with finite cokernel. Its determinant
is the product of the eigenvalues of $\Delta_k$ on the complement of the kernel, i.e.\ the
pseudo-determinant, which equals $n_k\kappa(Y_k)$ by [Ch.~\ref{chap:III}, eq.~(2.2)]. On $K_4$:
$4608=12\cdot384$, $18874368=24\cdot786432$, $316659348799488=48\cdot6597069766656$, with
$2$-adic orders $9,21,45$.
\end{proof}

\section{Universal localization and the characteristic class}\label{XIII:sec:loc}

\begin{definition}[Universal localization and the relative group]\label{XIII:def:loc}
By Cohn \cite{Co} there is a ring homomorphism $\lambda\colon\cH\to\HS$, universal among ring
homomorphisms under which $\Delta_{(K)}$ becomes invertible. Let $K_0(\cH,\Sigma)$ be the
relative group of triples $(P,Q,f)$, $P,Q$ finitely generated projective $\cH$-modules and
$f\colon P\to Q$ a map with $\lambda(f)$ invertible, modulo the usual relations
\cite[Ch.~4]{Sch}. Then
\begin{equation}\label{XIII:eq:seq}
K_1(\cH)\longrightarrow K_1(\HS)\xrightarrow{\ \partial\ }K_0(\cH,\Sigma)\longrightarrow K_0(\cH)
\end{equation}
is exact \cite[Ch.~4]{Sch}, and $\partial$ sends the class of $\lambda(f)$, for
$f\colon P\to Q$ with $\lambda(f)$ invertible, to $[(P,Q,f)]$. Forgetting to modules,
$[(P,Q,f)]\mapsto[\coker f]-[\ker f]$ when these are finite.
\end{definition}

\begin{definition}[Characteristic class]\label{XIII:def:theta}
$[\Theta_{\cH}]:=[\lambda(\Delta_{(K)})]\in K_1(\HS)$. It is well defined by
Proposition~\ref{XIII:prop:sigma}, and $\HS\ne0$ because $\lambda$ extends to the representation
$\cH\to\End_{\Qp}(\cM^{0}_{\le K}\otimes\Qp)$ in which $\Delta_{(K)}$ is invertible.
\end{definition}

\begin{theorem}[The tautological identity, and what it says]\label{XIII:thm:tauto}
$\partial[\Theta_{\cH}]=[(\cM^{0}_{\le K},\cM^{0}_{\le K},\Delta_{(K)})]\in K_0(\cH,\Sigma)$, whose
image under the forgetful map is
\[
\sum_{k\le K}\bigl[\coker\Delta_k^{0}\bigr]+\bigl[\cM^{0}_{K+1}/q\cM^{0}_{K+1}\bigr].
\]
This is the definition of $\partial$; it is the statement the specification called the
``non-commutative main conjecture'', and it has no content beyond
Definition~\ref{XIII:def:loc}. The content is in Theorems~\ref{XIII:thm:eval} and \ref{XIII:thm:boundary}.
\end{theorem}

\section{Evaluation and the Eisenstein boundary}\label{XIII:sec:eval}

For each $k\le K$ the projection $\cH\to\End_{\Zp}(\cM^{0}_k)$ followed by
$\otimes\Qp$ is a ring homomorphism $\cH\to M_{n_k-1}(\Qp)$ under which $\Delta_{(K)}$ becomes
invertible; by universality it factors through $\HS$, and composing with the determinant
gives the \emph{level evaluation}
\[
\ev_k\colon K_1(\HS)\longrightarrow K_1\bigl(M_{n_k-1}(\Qp)\bigr)=\Qp^{\times}.
\]

\begin{theorem}[Level evaluations: pure arc growth]\label{XIII:thm:eval}
$\ev_k[\Theta_{\cH}]=\det\Delta_k^{0}=\pm\,n_k\kappa(Y_k)$, and for all $k\ge1$
\[
\frac{\ev_{k+1}[\Theta_{\cH}]}{\ev_k[\Theta_{\cH}]}=\pm\,q^{\,n_k(q-1)},\qquad
\ord_p\ev_k[\Theta_{\cH}]=n_k+c=\frac{n(q{+}1)}{q}\,p^{k}+c ,
\]
with a constant $c$ depending only on the base. The arc term $\mathrm{Det_{arc}}=q^{\,n_k(q-1)}$
of {\rm[Ch.~\ref{chap:VII}, Cor.~2.2]} is the ratio of consecutive evaluations of the single class
$[\Theta_{\cH}]$; the Euler term does not appear. On $K_4$: $c=-3$ and
$\ord_2\ev_k=9,21,45$ at $k=1,2,3$.
\end{theorem}

\begin{proof}
The first equality is Proposition~\ref{XIII:prop:sigma}. The ratio: $\pdet(\Delta_{k+1})/\pdet(\Delta_k)
=q^{\,n_{k+1}-n_k}=q^{\,n_k(q-1)}$ because the nonzero spectrum of $A_{k+1}$ is that of $A_k$
plus $n_{k+1}-n_k$ zeros [Ch.~\ref{chap:III}, eq.~(2.1)], each contributing the eigenvalue $q$ of
$\Delta_{k+1}$. Equivalently this is the residue constancy $n_k\kappa_kq^{-n_k}=\mathrm{const}$
of [Ch.~\ref{chap:IX}, Thm.~4.1(iii)]. Summing, $\ord_p\ev_k=n_k+c$ with $n_k=n(q{+}1)q^{k-1}$.
\end{proof}

\begin{theorem}[The Eisenstein boundary]\label{XIII:thm:boundary}
For every $k\ge1$ there is an exact sequence of finite $\Zp$-modules
\[
0\longrightarrow\Z/n_k\longrightarrow\coker\Delta_k^{0}\longrightarrow\Jac(Y_k)\longrightarrow0,
\]
so that in the Grothendieck group of finite $\Zp$-modules the level-$k$ component of
$\partial[\Theta_{\cH}]$ is
\[
\bigl[\coker\Delta_k^{0}\bigr]=\bigl[\Jac(Y_k)\bigr]+\bigl[\Z/n_k\bigr],
\qquad
\ord_p|\Jac(Y_k)|=\ord_p\ev_k[\Theta_{\cH}]-\ord_pn_k=\frac{n(q{+}1)}{q}p^{k}-k+\nu_0 .
\]
The Euler term $-k$ of {\rm[Ch.~\ref{chap:III}, Thm.~2.2]} is the Eisenstein subclass $[\Z/n_k]$ of the
boundary, of $p$-adic order $\ord_pn_k=k+\ord_p(n(q{+}1))-1$. On $K_4$ the extension is
non-split at $k=1,2,3$: $\coker\Delta_k^{0}=\Z/8\oplus(\Z/24)^{2}$,
$(\Z/2)^{9}\oplus\Z/16\oplus(\Z/48)^{2}$, $(\Z/2)^{12}\oplus(\Z/4)^{9}\oplus\Z/32\oplus(\Z/96)^{2}$,
against $\Jac(Y_k)\oplus\Z/n_k$, whose invariant factors differ (at $k=1$:
$\Z/2\oplus\Z/4\oplus(\Z/24)^{2}$).
\end{theorem}

\begin{proof}
$\Delta_k(\Zp^{n_k})\subseteq\ker(\deg)$, so $\Jac(Y_k)=\Tor\coker\Delta_k=\ker(\deg)/\Delta_k
(\Zp^{n_k})$ by [Ch.~\ref{chap:VII}, Lem.~1.1] and [Ch.~\ref{chap:IV}, Lem.~1.3]. The lattice $\Delta_k(\Zp^{n_k})$ contains
$\Delta_k(\ker\deg)$ with quotient $\Delta_k(\Zp^{n_k})/\Delta_k(\ker\deg)\cong\Zp^{n_k}/
(\ker\deg+\Zp\mathbf1)\cong\Z/\deg(\mathbf1)=\Z/n_k$, since $\Delta_k$ is injective modulo
$\Zp\mathbf1$. The exact sequence follows, and orders multiply: $|\coker\Delta_k^{0}|=n_k
|\Jac(Y_k)|$, consistent with Proposition~\ref{XIII:prop:sigma}. The $p$-adic law is
Theorem~\ref{XIII:thm:eval} minus $\ord_pn_k$; on $K_4$, $n_k-3-(k+1)=6\cdot2^{k}-k-4$, the law of
[Ch.~\ref{chap:III}, Thm.~2.2]. Non-splitting: compare invariant factors (Table~\ref{XIII:tab:battery}, line (B)).
\end{proof}

\begin{remark}[Axioms (A1), (A2), (A3) in $K$-theory]\label{XIII:rem:axioms}
[Ch.~\ref{chap:V}, \S5] asked for a characteristic formalism producing an arc-density term $\mu$ (A1), an
Euler term $-\chi k$ intrinsic to the base (A2), and the commutative characteristic ideal on
abelian sub-towers (A3); [Ch.~\ref{chap:VII}] delivered all three in evaluated form and left the lift to a
single $K_1$ class as the residue. Theorems~\ref{XIII:thm:eval}--\ref{XIII:thm:boundary} are that lift:
the class $[\Theta_{\cH}]$ evaluates to the pure arc law (A1); its boundary contains the
Eisenstein subclass $[\Z/n_k]$ whose order is the vertex count, and $-\ord_pn_k$ is the Euler
term (A2) --- the fourth reading of $-\chi k$ after [Ch.~\ref{chap:III}], [Ch.~\ref{chap:VII}] and [Ch.~\ref{chap:IX}], and the first in
which it is a class rather than a number; (A3) is Theorem~\ref{XIII:thm:abelian}. The three
axioms are properties of one object, not three formulas.
\end{remark}

\begin{theorem}[Abelian reduction]\label{XIII:thm:abelian}
For an abelian $\Z/\ell^{k}$ voltage tower of {\rm[Ch.~\ref{chap:IV}]} with pencil $D(t)$ and
$\Delta_k=D(S_{\ell^{k}})$ {\rm[Ch.~\ref{chap:IV}, Lem.~1.1]}, the same construction (the group ring
$\Zl[G_k]$ playing the role of $\cH$) gives
\[
\ev_k[\Theta]=\pdet(\Delta_k)=n_k\kappa(Y_k)=n\,\kappa(Y_0)\prod_{\chi\ne1}\det D(\chi)
=n\,\kappa(Y_0)\prod_{\zeta^{\ell^{k}}=1,\ \zeta\ne1}\zeta^{-d}g(\zeta-1),
\]
the characteristic element $g=\det D(1+T)$ of {\rm[Ch.~\ref{chap:IV}, Thm.~2.1]} evaluated at all
nontrivial characters, times the Eisenstein constant $n\kappa(Y_0)$; the Eisenstein
subclass is $\Z/n_k=\Z/n\ell^{k}$ and $\ord_\ell n_k=k+\ord_\ell n$ is the ``$-k$'' of
{\rm[Ch.~\ref{chap:II}, Rem.~5.4]}. Exact on the $(1,1,1)$ tower, $\ell=2$, $k\le3$: $3072$, $2408448$,
$362577395712$ (Table~\ref{XIII:tab:battery}, line (A3)).
\end{theorem}

\begin{proof}
The block-character decomposition [Ch.~\ref{chap:IV}, Lem.~1.1] splits the nonzero spectrum of
$\Delta_k$ into the spectra of $D(\chi)$, $\chi\ne1$, and the nonzero spectrum of
$D(1)=\Delta_0$, whose product is $n\kappa(Y_0)$ [Ch.~\ref{chap:III}, eq.~(2.2)]; $\det D(\zeta)=\zeta^{-d}
g(\zeta-1)$ as in [Ch.~\ref{chap:IX}, Lem.~2.1].
\end{proof}

\section{The verification battery}\label{XIII:sec:battery}

\begin{table}[ht]
\centering\scriptsize
\renewcommand{\arraystretch}{1.2}
\resizebox{\textwidth}{!}{%
\begin{tabular}{lcl}
\toprule
check & levels & result\\
\midrule
(I) \eqref{XIII:eq:calc}; $A_kA_k^{t}\ne A_k^{t}A_k$; degree compatibility of $\ta,\et,\taup,\etp$ & $k=1,2,3$ & exact\\
(S) $\Delta_k^{0}$ injective, invertible over $\Q$, $\det$ not a $\Z_2$-unit ($2$-adic orders $9,21,45$) & $k=1,2,3$ & exact\\
(E) $\det\Delta_k^{0}=n_k\kappa_k=\pdet$; ratios $2^{12},2^{24}$; $\ord_2=n_k-3$ & $k=1,2,3$ & exact (pdet numeric)\\
(B) $|\coker\Delta_k^{0}|/|\Jac(Y_k)|=n_k$; structures; extension non-split; $\Delta_k=\ta(\taup-\etp)$ & $k=1,2,3$ & exact\\
(A3) $(1,1,1)$, $\ell=2$: $\pdet(\Delta_k)=n\kappa_0\prod_{\chi\ne1}\det D(\chi)$ & $k=1,2,3$ & exact\\
\bottomrule
\end{tabular}}
\medskip
\caption{Verification (\texttt{cohn\_localization.py}) on the $K_4$ Iwahori tower and the
$(1,1,1)$ abelian tower.}
\label{XIII:tab:battery}
\end{table}

\section{Scope, residue, corrections}\label{XIII:sec:scope}

\begin{remark}[What a main conjecture would be]\label{XIII:rem:main}
In the theory of \cite{CFKSV} a main conjecture asserts that an \emph{analytic} object --- a
$p$-adic $L$-function characterized by interpolation of special values, constructed
independently of the module --- equals the algebraic characteristic element, an element of
$K_1$ of the localized Iwasawa algebra whose boundary is the class of the Selmer module. The
algebraic half is what this paper constructs: $[\Theta_{\cH}]\in K_1(\HS)$ with
$\partial[\Theta_{\cH}]$ the class of the tower's cokernels, the analogue of ``$\partial\xi=[X]$''.
The analytic half does not exist yet: there is no independently defined element of
$K_1(\HS)$ interpolating the $L$-values of the tower --- the candidates are the pseudo-determinants
$n_k\kappa_k=\prod_{j}(q{+}1-\lambda_j)$, products of local $L$-values at the Eisenstein point
as in [Ch.~\ref{chap:I}, Thm.~3.5], but a construction of a single element from them, rather than from
$\Delta$, is exactly what is missing. The non-commutative main conjecture for Hecke towers
is therefore: \emph{construct $\mathcal L_{\cH}\in K_1(\HS)$ from the $L$-values alone and
prove $\mathcal L_{\cH}=[\Theta_{\cH}]$ modulo the image of $K_1(\cH)$}. Its evaluated form
is [Ch.~\ref{chap:VII}, Thm.~2.1]; its $K_1$ form is open.
\end{remark}

\begin{remark}[Solved, open, impossible]\label{XIII:rem:scope}
\emph{Solved:} the residue of [Ch.~\ref{chap:V}, Prob.~5.1] and [Ch.~\ref{chap:VII}, Rem.~3.1] in the form asked --- a single
class in $K_1$ of a localized degeneracy algebra whose level evaluations and boundary
reproduce the snake factorization, with (A1) and (A2) separated as evaluation and Eisenstein
subclass and (A3) as the abelian reduction. \emph{Open:} (a) the analytic side
(Remark~\ref{XIII:rem:main}); (b) stable flatness of $\cH\to\HS$ and hence the higher localization
sequence \cite{NR}; (c) the structure of $\HS$ itself --- whether it is an order in a
semisimple algebra or has torsion, and the image of $K_1(\cH)$ in $K_1(\HS)$, which
measures how far $[\Theta_{\cH}]$ is from unique; (d) the rank-two version over the
degeneracy algebra of [Ch.~\ref{chap:XII}], where no abelian sector exists. \emph{Impossible or empty as
stated:} ``$\partial[\Theta]=[X]$'' as a theorem with content (Theorem~\ref{XIII:thm:tauto}).
\end{remark}

\begin{remark}[Corrections to the task specification]\label{XIII:rem:corrections}
(1) The identity $\partial([\Theta_{\cH}])=[X_\infty]$ is the definition of the boundary map in
the relative sequence, not a theorem; the paper proves what does have content
(Theorems~\ref{XIII:thm:eval}--\ref{XIII:thm:abelian}). (2) The class must be formed on the
augmentation-zero lattices: on $\Zp^{n_k}$ the Laplacian has the Eisenstein kernel and is not
$\Sigma$-inverting; the specification's ``$\Delta_\infty\in M_n(\widehat{\cH}_\infty)$'' has no
$n\times n$ presentation because $n_k$ grows (contrast [Ch.~\ref{chap:IV}], where $\Zl[G_k]^{n}$ is free of
constant rank). (3) The Euler term is not produced by ``the connecting map'' inside $K_1$;
it is the Eisenstein subclass $[\Z/n_k]$ of the boundary, and it is invisible to the
$K_1$-evaluations, which see the arc term only. (4) ``$\Sigma$ = the set of morphisms
inducing surjections with finite cokernel on $X_\infty$'' is replaced by the single element
$\Delta_{(K)}$; universal localization at one element already gives the class. (5) The
higher $K$-theory sequence needs stable flatness \cite{NR}; only the low-degree sequence
\cite{Sch} is used, and only its first map matters for the construction.
\end{remark}

\section*{Acknowledgement of provenance and caveats}

Unrefereed preprint. Every numbered statement is proved from Papers III--XII and finite
linear algebra, given the existence of universal localization \cite{Co} and the low-degree
exact sequence \cite[Ch.~4]{Sch}, both quoted at chapter level (the precise theorem numbers
were not consulted). Bibliographic data of \cite{Co,Sch,NR} were checked against publisher
or indexing pages; \cite{CFKSV} is the entry verified in [Ch.~\ref{chap:V}].

\part{Operator algebras of the tower}
\input{chapters/VIII}
% generated from Paper 17/anccft17.tex
\chapter[\texorpdfstring{XVII. The tail norm at composite $q$, and the splitting of the limit $K$-module}{XVII. The tail norm at composite q, and the splitting of the limit K-module}]{\texorpdfstring{Algorithmic non-commutative class field theory, XVII: the tail norm at composite $q$, and the splitting of the limit $K$-module}{Algorithmic non-commutative class field theory, XVII: the tail norm at composite q, and the splitting of the limit K-module}}\label{chap:XVII}
\chaptermark{XVII}
\begin{chapterabstract}
Paper~VIII proved that for $q=p$ prime the kernel of the tail norm on the Iwahori tower is
the full $p$-torsion of the sandpile group, and deduced from the resulting torsion-freeness
that the sequence $0\to\XX^\vee\to K_0(\cO_\infty)\to\Z[1/q]\to0$ splits. Both statements
were left open for composite $q$, on the ground that $\Z/q\Z$ has zero divisors. That
ground is illusory. The only ingredient of the prime proof that used primality is the rank
identity $\rank_{\F_p}A_{k+1}=n_k$, and $A_{k+1}$ is the adjacency matrix of a line
digraph, hence factors as $H^{\mathsf T}T$ through the head and tail incidence matrices;
the rows of $H$ and of $T$ are non-empty and disjointly supported, so the rank is $n_k$
\emph{in every characteristic}. With that, the two facts already known for all $q$ ---
$|\ker\ta|=q^{\,n_k(q-1)-1}$ and $\ker\ta\subseteq\Jac(Y_{k+1})[q]$ --- close the problem
by counting: the two orders are forced to agree, so
\[
\ker\bigl(\ta\colon\Jac(Y_{k+1})\to\Jac(Y_k)\bigr)=\Jac(Y_{k+1})[q]\cong(\Z/q)^{\,n_k(q-1)-1}
\]
for every $q\ge2$. Equality in the counting has a second consequence, which is new even
for prime powers and is the sharpest statement here: \emph{every invariant factor of
$\Jac(Y_{k+1})$ divisible by a prime $p\mid q$ is divisible by $p^{v_p(q)}$}. So at
$q=4$ the group has no element of order exactly $2$, and the kernel of the tail norm
strictly contains the $2$-torsion --- a structural phenomenon with no analogue in the
prime case. Duality then gives the splitting for every $q$ in three lines and without
identifying $\XX$: $\XX[q]=0$, so $\XX^\vee$ is $q$-divisible, so
$\mathrm{Ext}^1(\Z[1/q],\XX^\vee)=\varprojlim^1(\XX^\vee,\times q)=0$. Six checks on
eight towers ($q=2,3,4,5,6$, complete and incomplete bases, to level three), all exact;
the $q=4$ towers could have falsified the divisibility statement and did not.
\end{chapterabstract}

\section{The rank is characteristic-free}\label{XVII:sec:rank}

Throughout, $Y_0$ is a $(q+1)$-regular graph on $n$ vertices, $Y_1$ its non-backtracking
directed-edge graph --- Eulerian $q$-regular on $n_1=n(q+1)$ vertices --- and
$Y_{k+1}=L(Y_k)$ the line digraph, so $n_{k+1}=qn_k$, for $k\ge1$. The degeneracy maps
$\tau$ (tail) and $\eta$ (head) send the vertex $i\to j$ of $Y_{k+1}$ to $i$ and to $j$
respectively; $\ta$ is the pushforward and $\etp$ the pullback. We write
\[
r_k:=n_k(q-1)-1 ,
\]
$\Delta_k=qI-A_k$ for the out-Laplacian of $Y_k$, and $\Jac(Y_k)$ for the torsion of
$\coker\Delta_k^{\mathsf T}$, so that $\coker\Delta_k^{\mathsf T}\cong\Z\oplus\Jac(Y_k)$.
Three facts from the earlier papers hold for \emph{every} $q\ge2$ and will be used as
stated:
\begin{enumerate}
\item[(i)] $\ta\etp=A_k$, $\etp\ta=A_{k+1}$ and $\ta\tau^*=\eta_*\etp=qI$
      [Ch.~\ref{chap:V}, Prop.~1.2];
\item[(ii)] $\ta$ is surjective on Jacobians with $|\ker\ta|=q^{\,r_k}$
      [Ch.~\ref{chap:V}, Thm.~3.1(ii)], a restatement of the line-digraph recursion
      $\kappa(Y_{k+1})=q^{\,r_k}\kappa(Y_k)$ [Ch.~\ref{chap:III}, Thm.~2.1];
\item[(iii)] $A_{k+1}$ acts on $\Jac(Y_{k+1})$ as multiplication by $q$, and consequently
      $\ker\ta\subseteq\Jac(Y_{k+1})[q]$ [Ch.~\ref{chap:VIII}, Prop.~2.7].
\end{enumerate}
Only (iii) needs a word: $\Delta_{k+1}=qI-A_{k+1}$ annihilates the cokernel, so $A_{k+1}=q$
there, and by (i) $\ta x=0$ forces $A_{k+1}x=\etp\ta x=0$, i.e.\ $qx=0$.

What was available only for $q=p$ prime is the rank of $A_{k+1}$ modulo $p$
[Ch.~\ref{chap:VIII}, Thm.~2.8]. It is in fact insensitive to the characteristic.

\begin{lemma}[Rank of a line-digraph adjacency matrix]\label{XVII:lem:rank}
Let $G$ be a digraph on $n$ vertices and $m$ arcs in which every vertex is the head of at
least one arc and the tail of at least one arc, and let $A_{L(G)}$ be the adjacency matrix
of its line digraph. Then $\rank_F A_{L(G)}=n$ for every field $F$.
\end{lemma}

\begin{proof}
Let $H,T\in F^{\,n\times m}$ be the head and tail incidence matrices,
$H_{v,e}=[\,\mathrm{head}(e)=v\,]$ and $T_{v,f}=[\,\mathrm{tail}(f)=v\,]$. For arcs $e,f$,
\[
(H^{\mathsf T}T)_{e,f}=\sum_{v}[\,\mathrm{head}(e)=v\,][\,\mathrm{tail}(f)=v\,]
=[\,\mathrm{head}(e)=\mathrm{tail}(f)\,]=(A_{L(G)})_{e,f},
\]
so $A_{L(G)}=H^{\mathsf T}T$. Every arc has exactly one head, so the rows of $H$ have
pairwise disjoint supports; each is non-zero by hypothesis; hence they are linearly
independent over any field and $\rank H=n$, and likewise $\rank T=n$. Thus
$T\colon F^m\to F^n$ is surjective and $H^{\mathsf T}\colon F^n\to F^m$ injective, so
$\rank(H^{\mathsf T}T)=\dim H^{\mathsf T}(F^n)=n$.
\end{proof}

\begin{corollary}[The $p$-rank of the Jacobian]\label{XVII:cor:prank}
For $k\ge1$ and every prime $p\mid q$, the $p$-rank of $\Jac(Y_{k+1})$ equals $r_k$.
\end{corollary}

\begin{proof}
$Y_{k+1}=L(Y_k)$ is $q$-regular, so $\Delta_{k+1}\equiv-A_{k+1}\pmod p$ and
$\rank_{\F_p}\Delta_{k+1}=n_k$ by Lemma~\ref{XVII:lem:rank}. Tensoring
$\coker\Delta_{k+1}^{\mathsf T}\cong\Z\oplus\Jac(Y_{k+1})$ with $\F_p$ gives
\[
n_{k+1}-\rank_{\F_p}\Delta_{k+1}=1+\dim_{\F_p}\bigl(\Jac(Y_{k+1})\otimes\F_p\bigr),
\]
and the left side is $qn_k-n_k$. As $\Jac(Y_{k+1})\otimes\F_p\cong\Jac(Y_{k+1})[p]$ for a
finite abelian group, the $p$-rank is $n_k(q-1)-1=r_k$.
\end{proof}

\section{The tail norm at composite \texorpdfstring{$q$}{q}}\label{XVII:sec:kernel}

\begin{theorem}[The kernel of the tail norm, for every $q$]\label{XVII:thm:main}
Let $q\ge2$ be arbitrary and $k\ge1$. Then
\[
\ker\bigl(\ta\colon\Jac(Y_{k+1})\to\Jac(Y_k)\bigr)=\Jac(Y_{k+1})[q]\cong(\Z/q)^{\,r_k},
\]
and moreover every invariant factor of $\Jac(Y_{k+1})$ divisible by a prime $p\mid q$ is
divisible by $p^{\,v_p(q)}$.
\end{theorem}

\begin{proof}
Write $e_p=v_p(q)$ and let the $p$-primary part of $\Jac(Y_{k+1})$ be
$\bigoplus_{j=1}^{s_p}\Z/p^{a_{p,j}}$ with $a_{p,j}\ge1$; by
Corollary~\ref{XVII:cor:prank}, $s_p=r_k$ for every $p\mid q$. By (ii) and (iii) of
Section~\ref{XVII:sec:rank},
\[
q^{\,r_k}=|\ker\ta|\;\le\;\bigl|\Jac(Y_{k+1})[q]\bigr|
=\prod_{p\mid q}p^{\,\sum_{j}\min(e_p,\,a_{p,j})}
\;\le\;\prod_{p\mid q}p^{\,e_p r_k}=q^{\,r_k}.
\]
Both inequalities are therefore equalities. The first forces
$\ker\ta=\Jac(Y_{k+1})[q]$, two finite groups of the same order one contained in the
other. The second forces $\min(e_p,a_{p,j})=e_p$, that is $a_{p,j}\ge e_p$, for every
$p\mid q$ and every $j$; the $p$-primary part of $\Jac(Y_{k+1})[q]$ is then
$\bigl(\Z/p^{e_p}\bigr)^{r_k}$, and taking the product over $p\mid q$ gives
$(\Z/q)^{r_k}$.
\end{proof}

\begin{remark}[What was and was not primality]\label{XVII:rem:whatwas}
For $q=p$ this is [Ch.~\ref{chap:VIII}, Thm.~2.8]. The present proof differs from that one only in
Corollary~\ref{XVII:cor:prank}, and the corollary differs only in proving the rank identity by
a factorisation valid in every characteristic rather than by an argument in
characteristic $p$. Nothing else in the chain ever meets a zero divisor of $\Z/q\Z$: the
counting is carried out one prime at a time, and $q$ enters only through its valuations.
\end{remark}

\begin{remark}[Prime powers behave differently from primes]\label{XVII:rem:powers}
For $q$ squarefree, $\Jac(Y_{k+1})[q]=\bigoplus_{p\mid q}\Jac(Y_{k+1})[p]$ and
Theorem~\ref{XVII:thm:main} says only that the kernel of the tail norm is the torsion killed by
each prime factor --- the prime picture, assembled by the Chinese remainder theorem. For
$q$ \emph{not} squarefree the theorem says something with no prime analogue. At $q=4$ it
asserts that $\Jac(Y_{k+1})$ has \emph{no element of order exactly $2$}: every invariant
factor divisible by $2$ is divisible by $4$. Consequently
$\Jac(Y_{k+1})[2]\subsetneq\Jac(Y_{k+1})[4]=\ker\ta$, of index $2^{r_k}$, and the kernel of
the tail norm strictly contains the $p$-torsion for each $p\mid q$. This is a falsifiable
statement about a concrete integer matrix; check~(E) tests it on three $q=4$ towers,
including level three of the $K_6$ tower, where $\Jac(Y_3)$ has $2$-rank $359$ and the
smallest $2$-adic valuation among its invariant factors is $2$.
\end{remark}

\begin{corollary}[No $q$-torsion in the pro-module]\label{XVII:cor:notorsion}
$\XX:=\varprojlim\bigl(\Jac(Y_k),\ta\bigr)$ satisfies $\XX[q]=0$.
\end{corollary}

\begin{proof}
Let $x=(x_k)_{k\ge1}\in\XX$ with $qx=0$. Then $x_k\in\Jac(Y_k)[q]$ for every $k$, and for
$k\ge2$ that group is $\ker\bigl(\ta\colon\Jac(Y_k)\to\Jac(Y_{k-1})\bigr)$ by
Theorem~\ref{XVII:thm:main}; hence $x_{k-1}=\ta x_k=0$ for every $k\ge2$, i.e.\ $x=0$.
\end{proof}

\section{The splitting, for every \texorpdfstring{$q$}{q}}\label{XVII:sec:split}

\begin{theorem}[Splitting of the limit $K$-module]\label{XVII:thm:split}
For every $q\ge2$ the short exact sequence of [Ch.~\ref{chap:VIII}, Thm.~3.3(iii)],
\[
0\longrightarrow\XX^{\vee}\longrightarrow K_0(\cO_\infty)\longrightarrow\Z[1/q]
\longrightarrow0,
\]
splits, and $K_0(\cO_\infty)\cong\Z[1/q]\oplus\XX^{\vee}$.
\end{theorem}

\begin{proof}
The extension class lies in $\mathrm{Ext}^1(\Z[1/q],\XX^\vee)$. Write
$\Z[1/q]=\varinjlim\bigl(\Z\xrightarrow{\;q\;}\Z\xrightarrow{\;q\;}\cdots\bigr)$. For any
abelian group $B$ the Milnor sequence of this colimit reads
\[
0\to\varprojlim{}^1\mathrm{Hom}(\Z,B)\to\mathrm{Ext}^1(\Z[1/q],B)
\to\varprojlim\mathrm{Ext}^1(\Z,B)\to0 ,
\]
and $\mathrm{Ext}^1(\Z,B)=0$, so
$\mathrm{Ext}^1(\Z[1/q],B)\cong\varprojlim^1\bigl(B\xleftarrow{\;q\;}B\xleftarrow{\;q\;}\cdots\bigr)$.
If $B$ is $q$-divisible every transition map is surjective and this $\varprojlim^1$
vanishes. Now $\XX$ is profinite and multiplication by $q$ on it is injective by
Corollary~\ref{XVII:cor:notorsion}; Pontryagin duality is exact, so multiplication by $q$ on
$\XX^\vee$ is surjective, i.e.\ $\XX^\vee$ is $q$-divisible. Hence
$\mathrm{Ext}^1(\Z[1/q],\XX^\vee)=0$.
\end{proof}

\begin{remark}[A shorter route to the prime case]\label{XVII:rem:shorter}
For $q=p$, [Ch.~\ref{chap:VIII}, Thm.~3.3] obtained the splitting by first identifying
$\XX^{(p)}\cong\Z_p^{\mathbb N}$ and then observing that its dual
$(\Q_p/\Z_p)^{(\mathbb N)}$ is divisible. The proof above needs neither identification: it
uses only that multiplication by $q$ is injective on $\XX$, which is
Corollary~\ref{XVII:cor:notorsion}. The structure of $\XX$ is not an input to the splitting,
and for composite $q$ we do not determine it.
\end{remark}

\section{Verification}\label{XVII:sec:battery}

The script \texttt{composite\_tail\_norm.py} runs the six checks of
Table~\ref{XVII:tab:battery} in about five seconds; its output is archived as
\texttt{composite\_tail\_norm\_output.txt}. Eight towers are used: base graphs $K_4$ and
Petersen ($q=2$), $K_5$ and $K_{4,4}$ ($q=3$), $K_6$ and $K_{5,5}$ ($q=4$), $K_7$ ($q=5$)
and $K_8$ ($q=6$), carried to level three where the size permits, the largest being
$Y_3$ of the $K_6$ tower on $480$ vertices. All arithmetic is exact: Bareiss determinants
for the orders, and $p$-adic elimination (no coefficient explosion) for the valuations of
the invariant factors.

\begin{table}[ht]
\centering\small
\begin{tabular}{clc}
\toprule
key & check & mode\\
\midrule
(I) & the degeneracy identities of \S\ref{XVII:sec:rank}(i) hold verbatim at composite $q$ & exact\\
(N) & $\kappa(Y_{k+1})=q^{\,r_k}\kappa(Y_k)$, so $|\ker\ta|=q^{\,r_k}$ & exact\\
(R) & Lemma~\ref{XVII:lem:rank}: $\rank A_{k+1}=n_k$ over $\F_p$ for $p\mid q$ and $p\nmid q$ alike & exact\\
(P) & Corollary~\ref{XVII:cor:prank}: the $p$-rank of $\Jac(Y_{k+1})$ is $r_k$ for $p\mid q$ & exact\\
(E) & Theorem~\ref{XVII:thm:main}: every invariant factor divisible by $p$ is divisible by $p^{v_p(q)}$ & exact\\
(T) & the published prime-$q$ values of [Ch.~\ref{chap:VIII}] are recovered & exact\\
\bottomrule
\end{tabular}
\caption{The verification battery, $6/6$.}
\label{XVII:tab:battery}
\end{table}

Check~(E) is the one that could have failed. For each $q=4$ tower it reports the
multiset of $2$-adic valuations of the invariant factors of $\Jac(Y_k)$: on $K_6$ at
level two these are $2^{84},4,5^4$ and at level three the $359$ valuations are all at
least $2$; on $K_{5,5}$ at level two the $149$ valuations are all at least $2$. A single
invariant factor of valuation $1$ anywhere would have contradicted
Theorem~\ref{XVII:thm:main}. Check~(T) reproduces
$\Jac(Y_2)=(\Z/2)^8\oplus\Z/4\oplus\Z/16\oplus\Z/48$ and
$\Jac(Y_3)=(\Z/2)^{12}\oplus(\Z/4)^8\oplus\Z/8\oplus\Z/32\oplus\Z/96$ of [Ch.~\ref{chap:VIII}, Table~1],
with $\mathrm{ord}_2\kappa=18$ and $41$.

\section{Status, scope, corrections}\label{XVII:sec:scope}

\emph{Proved.} Lemma~\ref{XVII:lem:rank}, Corollary~\ref{XVII:cor:prank}, Theorem~\ref{XVII:thm:main},
Corollary~\ref{XVII:cor:notorsion} and Theorem~\ref{XVII:thm:split}, for arbitrary $q\ge2$ and
$k\ge1$. Together they close [Ch.~\ref{chap:VIII}, Rem.~5.1(b)] in both of its parts.

\emph{Quoted, for every $q$ and used as stated.} The degeneracy calculus [Ch.~\ref{chap:V}, Prop.~1.2];
the tail norm and its kernel order [Ch.~\ref{chap:V}, Thm.~3.1], resting on the line-digraph recursion
[Ch.~\ref{chap:III}, Thm.~2.1]; the containment $\ker\ta\subseteq\Jac[q]$ [Ch.~\ref{chap:VIII}, Prop.~2.7]; the exact
sequence [Ch.~\ref{chap:VIII}, Thm.~3.3(iii)]; the Milnor $\varprojlim^1$ sequence and the exactness of
Pontryagin duality, both textbook.

\begin{remark}[Solved, open, impossible]\label{XVII:rem:soi}
\emph{Solved:} the composite-$q$ half of [Ch.~\ref{chap:VIII}, Rem.~5.1(b)], in the strong form
$\ker\ta\cong(\Z/q)^{r_k}$, together with the splitting for every $q$.
\emph{Open:} (a) the structure of $\XX$ itself for composite $q$ --- the splitting no
longer needs it, and we have not determined whether $\XX^{(p)}$ is again a free
$\Z_p$-module when $p^2\mid q$; (b) whether the divisibility statement of
Theorem~\ref{XVII:thm:main} --- no invariant factor of valuation below $v_p(q)$ --- has a
direct combinatorial proof, independent of the counting, and what it says about the
spanning-tree lattice of a line digraph; (c) the remaining items of [Ch.~\ref{chap:VIII}, Rem.~5.1],
which this note does not touch.
\emph{Impossible:} nothing new is shown impossible here. In particular the obstruction the
task specification anticipated --- that zero divisors in $\Z/q\Z$ would make the kernel
smaller than $\Jac[q]$ and force an order formula in terms of $\gcd(q,\cdot)$ --- does not
exist.
\end{remark}

\begin{remark}[Corrections to the task specification]\label{XVII:rem:corrections}
Four, all in the same direction: the composite case is easier than expected, not harder.
(1) The specification asked for the Smith normal form of $A_{k+1}$ over
$\Z/p_i^{e_i}\Z$ and for a fallback order formula in terms of $\gcd(q,\Jac)$ should the
kernel fail to be $\Jac[q]$. Neither is needed: the kernel \emph{is} $\Jac[q]$, and the
only rank computation required is over a field, where Lemma~\ref{XVII:lem:rank} makes it
characteristic-free. (2) ``Modulo $q$ localisation has zero divisors'' is true and
irrelevant; no step localises at $q$. (3) The specification proposed expanding $\Z[1/q]$
as a direct limit of localisations over the primes $p\mid q$; the single tower
$\Z\xrightarrow{q}\Z\xrightarrow{q}\cdots$ is what gives the $\varprojlim^1$ description,
and splitting into primes is unnecessary. (4) The specification asked for the divisible
subgroup structure of $\XX^\vee$. What is needed and true is only that $\XX^\vee$ is
$q$-divisible, dual to $\XX[q]=0$; no finer structure enters, and
Remark~\ref{XVII:rem:shorter} notes that this also shortens the published prime-case proof.
\end{remark}

\section*{Provenance of this chapter}

Unrefereed preprint. Every numbered statement is proved in full above from the three
quoted facts of Section~\ref{XVII:sec:rank}, each of which is stated in the earlier papers for
all $q$; the computations verify those inputs independently (checks (I) and (N)), the new
lemma (check (R)), and the falsifiable consequence of the theorem (check (E)). The
divisibility statement of Theorem~\ref{XVII:thm:main} was found by the counting and then
tested, not the other way round; the $q=4$ towers were computed specifically because they
could refute it.

% generated from Paper 16/anccft16.tex
\renewcommand{\XX}{\mathfrak{X}}
\chapter[\texorpdfstring{XVI. The out-covering correspondence of the head map, its $K_0$-index $\et$, and the
$\KK$-category of the Iwahori tower}{XVI. The out-covering correspondence of the head map, its K0-index , and the KK-category of the Iwahori tower}]{\texorpdfstring{Algorithmic non-commutative class field theory, XVI: the out-covering correspondence of the head map, its $K_0$-index $\et$, and the
$\KK$-category of the Iwahori tower}{Algorithmic non-commutative class field theory, XVI: the out-covering correspondence of the head map, its K0-index , and the KK-category of the Iwahori tower}}\label{chap:XVI}
\chaptermark{XVI}
\begin{chapterabstract}
[Ch.~\ref{chap:VIII}] realized the tail norm of the Iwahori tower on the shadow Kirchberg algebras
$\cO_k$: the tail map $\tau$ is an in-covering, so it induces a canonical unital
$*$-homomorphism $\varphi_k\colon\cO_k\to\cO_{k+1}$, whose $K_0$-map is the pullback
$\taup$ \emph{up} the tower. It left open [Ch.~\ref{chap:VIII}, Rem.~2.5, 2.6, 5.1(a)] the mirror
question: the head map $\eta$ is an out-covering, the pushforward $\et$ is the one that
descends on the $K_0$-side ($\et A_{k+1}^{t}=A_k^{t}\et$), yet no $*$-homomorphism
$\cO_{k+1}\to\cO_k$ can induce it, because $\et[1]=q[1]$ forces non-unitality and the
Kirchberg--Phillips lift is non-canonical. This paper settles it: the canonical object
is not a homomorphism but a $C^*$-correspondence. We build from the out-covering
$\eta^{\sharp}$ a right Hilbert $\cO_k$-module
$\cX_\eta=\bigoplus_{w}p_{\eta^{\sharp}(w)}\cO_k$, finitely generated and projective, and
a left action $\phi\colon\cO_{k+1}\to\cK(\cX_\eta)$ by compact operators; the out-covering
bijection on out-arcs is exactly what makes $\phi$ satisfy (CK2), where the sum-over-lifts
homomorphism of [Ch.~\ref{chap:VIII}] had failed. The pair $(\cX_\eta,\phi)$ is an injective, full,
proper correspondence, so it defines a canonical class
$[\cX_\eta]\in\KK(\cO_{k+1},\cO_k)$, and its $K_0$-index map is precisely the head
pushforward $\et$; as a right module $[\cX_\eta]=q\,[1_{\cO_k}]$. The correspondence is
not invertible in $\KK$: its index $\et$ has kernel the $2$-torsion of $\Jac(Y_{k+1})$,
of order $q^{\,n_k(q-1)-1}$. The two families of morphisms organize the levels into a
small $\KK$-category with objects $\{\cO_k\}$: alongside the upward homomorphisms
$[\varphi_k]$ ($K_0$-map $\taup$) sit the downward correspondences $[\cX_\eta]$ ($K_0$-map
$\et$), and their Kasparov products are the Hecke operators ---
$[\varphi_k]\otimes[\cX_\eta]$ acts on $K_0(\cO_k)$ as $A_k^{t}$ and
$[\cX_\eta]\otimes[\varphi_k]$ on $K_0(\cO_{k+1})$ as $A_{k+1}^{t}$. This is the
$C^*$-bivariant form of the degeneracy calculus $\ta\etp=A_k$, $\etp\ta=A_{k+1}$ of [Ch.~\ref{chap:V}],
and it completes the two-sided operator-algebraic picture of the tower. Everything is
verified exactly on the $K_4$ tower at levels $k\le3$.
\end{chapterabstract}

\section*{Status of the results}

Every numbered statement is proved. The $C^*$-algebraic inputs are the standard theory of
$C^*$-correspondences and their Kasparov classes \cite{Kas,Bla}, the graph-algebra
$K$-theory and gauge-invariant uniqueness of \cite{R}, and Kaminker--Putnam
Spanier--Whitehead duality \cite{KP} (cited once, for context, in
Remark~\ref{XVI:rem:reversal}); the construction of $\cX_\eta$ and the computation of its index
are done in full. The in-covering homomorphism $\varphi_k$ and the matrix identities of
the tower are quoted from [Ch.~\ref{chap:VIII}] and re-verified. Every finite identity was
machine-verified exactly on the $K_4$ tower, $q=2$, at the transitions
$k+1\to k$ for $k\le3$ (shadow graphs on $24,48,96$ vertices; Table~\ref{XVI:tab:battery}).
Section~\ref{XVI:sec:scope} records what is solved, what stays open, and where the
specification is corrected.

\section{Setting: the two degeneracy maps and the mirror problem}\label{XVI:sec:setting}

We keep the conventions of [Ch.~\ref{chap:VIII}]. $Y_0$ is a finite connected $(q{+}1)$-regular graph,
$q\ge2$; for $k\ge1$, $Y_k$ is the Iwahori path graph, a strongly connected $q$-regular
digraph on $n_k=n(q{+}1)q^{k-1}$ vertices with $Y_{k+1}=L(Y_k)$ its line digraph, so
$V_{Y_{k+1}}=\mathrm{Arcs}(Y_k)$. The tail and head maps $\tau,\eta\colon Y_{k+1}\to Y_k$
send an arc to its source and range; $\ta,\et\in M_{n_k\times n_{k+1}}(\Z)$ are their
pushforwards, $\taup=\ta^{t}$, $\etp=\et^{t}$ the pullbacks, with [Ch.~\ref{chap:V}, Prop.~1.2]
\begin{equation}\label{XVI:eq:calc}
\ta\etp=A_k,\quad \etp\ta=A_{k+1},\quad \ta\taup=\et\etp=qI,\quad
\ta\Delta_{k+1}=\Delta_k\ta,\quad \et\Delta_{k+1}^{t}=\Delta_k^{t}\et,
\end{equation}
$\Delta_k=qI-A_k$. The shadow graph $\Ysh{k}$ has vertex set $V_{Y_k}\sqcup V_{Y_k}'$,
the arcs of $Y_k$, one arc $v\to v'$, $c:=q-1$ return arcs $v'\to v$, and two loops at
each $v'$; its vertex matrix is $B_k=\begin{psmallmatrix}A_k&I\\cI&2I\end{psmallmatrix}$
and $\cO_k:=C^*(\Ysh{k})$ is a unital Kirchberg algebra with, in the convention
$K_0(C^*(E))=\coker(1-A_E^{t})$ of \cite[Thm.~7.16]{R},
\begin{equation}\label{XVI:eq:ktheory}
K_0(\cO_k)=\coker(I-B_k^{t})\cong\Z\oplus\Jac(Y_k),\qquad K_1(\cO_k)=\ker(I-B_k^{t})=\Z ,
\end{equation}
by the factorization $U(I-B_k^{t})V=\Delta_k^{t}\oplus(-I)$ of [Ch.~\ref{chap:VIII}, Thm.~1.4] with
$U=\begin{psmallmatrix}I&-cI\\0&I\end{psmallmatrix}$,
$V=\begin{psmallmatrix}I&0\\-I&I\end{psmallmatrix}$.

[Ch.~\ref{chap:VIII}] proved that $\tau$ extends to a graph morphism $\tau^{\sharp}\colon\Ysh{k+1}\to
\Ysh{k}$ that is an \emph{in}-covering --- bijective on the in-arcs at every vertex --- and
that in-coverings induce, by $p_v\mapsto\sum_{\tau^\sharp(w)=v}p_w$,
$s_a\mapsto\sum_{\tau^\sharp(b)=a}s_b$, a canonical unital injective $*$-homomorphism
$\varphi_k\colon\cO_k\to\cO_{k+1}$ [Ch.~\ref{chap:VIII}, Lem.~2.2, Thm.~2.5], whose $K_0$-map is the
pullback $\taup$ (up the tower) and whose $K_1$-map is the identity.

The head map is the mirror. Its shadow extension $\eta^{\sharp}\colon\Ysh{k+1}\to\Ysh{k}$
is an \emph{out}-covering, bijective on out-arcs and $q$-to-one on in-arcs
(Table~\ref{XVI:tab:battery}, line (O)), and the sum-over-lifts formula \emph{fails} to be a
homomorphism: for two arcs $f\ne f'$ of $\Ysh{k+1}$ over the same arc with the same range,
the cross term $s_fs_{f'}^{*}$ does not vanish, so $(\sum_bs_b)$ is not a Cuntz--Krieger
partial isometry [Ch.~\ref{chap:VIII}, Rem.~2.4]. Yet on $K_0$ it is the head pushforward $\et$, not
$\ta$, that descends:
\begin{equation}\label{XVI:eq:etdesc}
\et A_{k+1}^{t}=A_k^{t}\et,\qquad
H_k:=\begin{pmatrix}\et&0\\0&\et\end{pmatrix}\ \text{satisfies}\
H_k(I-B_{k+1}^{t})=(I-B_k^{t})H_k ,
\end{equation}
the first being the transpose of $\ta A_{k+1}=A_k\ta$ from \eqref{XVI:eq:calc};
whereas $\operatorname{diag}(\ta,\ta)$ does \emph{not} intertwine $I-B^{t}$
[Ch.~\ref{chap:VIII}, Thm.~2.5(iv)]. So the arrow the tower asks for --- a morphism
$\cO_{k+1}\rightsquigarrow\cO_k$ realizing $\et$ downward --- exists on $K_0$ but cannot be
a homomorphism, since additionally $H_k\mathbf1=q\mathbf1$ makes $\et[1_{\cO_{k+1}}]=q[1_{\cO_k}]$
(a unit class cannot map to $q$ times a unit class). The resolution is that the morphism is
a $C^*$-\emph{correspondence}.

\section{The out-covering correspondence}\label{XVI:sec:corr}

We recall that a $C^*$-correspondence ${}_A\cX_B$ is a right Hilbert $B$-module $\cX$ with a
$*$-homomorphism $\phi\colon A\to\cL(\cX)$; if $\cX$ is finitely generated projective over a
unital $B$ then $\cL(\cX)=\cK(\cX)$, and $(\cX,\phi,0)$ is an even Kasparov cycle defining
$[\cX]\in\KK(A,B)$ whose product with $[e]\in K_0(A)$, $e=e^2=e^*\in M_\infty(A)$, is
$[\phi(e)\cX]\in K_0(B)$ \cite{Kas,Bla}.

Throughout this section $p:=\eta^{\sharp}\colon F\to E$, $F=\Ysh{k+1}$, $E=\Ysh{k}$, an
out-covering: $s^{-1}(w)\to s^{-1}(p(w))$ is a bijection for every $w\in F^0$. Write
$\{p_v,s_e\}$ for the generators of $\cO_k=C^*(E)$ and $\{q_w,t_f\}$ for those of
$\cO_{k+1}=C^*(F)$.

\begin{definition}[The head module]\label{XVI:def:module}
Let
\[
\cX_\eta:=\bigoplus_{w\in F^0}p_{p(w)}\,\cO_k ,
\]
the right Hilbert $\cO_k$-module of finitely supported families $\xi=(\xi_w)_{w\in F^0}$
with $\xi_w\in p_{p(w)}\cO_k$, right action $(\xi\cdot a)_w=\xi_wa$, and $\cO_k$-valued inner
product $\langle\xi,\zeta\rangle=\sum_{w\in F^0}\xi_w^{*}\zeta_w$. Let $e_w\in\cX_\eta$ be
the generator with $w$-component $p_{p(w)}$ and all others $0$.
\end{definition}

\begin{lemma}[The module is finitely generated projective and full]\label{XVI:lem:fgp}
$\cX_\eta$ is a right Hilbert $\cO_k$-module with $\langle e_w,e_{w'}\rangle=\delta_{ww'}p_{p(w)}$,
finite orthogonal generating set $\{e_w\}_{w\in F^0}$, and reconstruction
$\xi=\sum_we_w\langle e_w,\xi\rangle$; hence $\cX_\eta$ is finitely generated projective,
$\cL(\cX_\eta)=\cK(\cX_\eta)$, and
\[
[\cX_\eta]=\sum_{w\in F^0}[p_{p(w)}]=q\,[1_{\cO_k}]\in K_0(\cO_k),
\]
because every fibre of $\eta^{\sharp}$ has exactly $q$ elements. $\cX_\eta$ is full:
$\overline{\langle\cX_\eta,\cX_\eta\rangle}=\cO_k$.
\end{lemma}

\begin{proof}
$\cX_\eta$ is an orthogonal direct sum of the corners $p_{p(w)}\cO_k$, each a Hilbert
$\cO_k$-module with $\langle a,b\rangle=a^*b$; the displayed inner products and
reconstruction are immediate, so $\{e_w\}$ is a finite frame and $\cX_\eta\cong\bigoplus_w
p_{p(w)}\cO_k$ is finitely generated projective with $[\cX_\eta]=\sum_w[p_{p(w)}]$. The
fibre of $\eta^{\sharp}$ over a $Y$-vertex $v$ is $\{$arcs $f$ of $Y_k$ with $r(f)=v\}$, of
size the in-degree $q$; over $v'$ it is the $q$ shadow lifts likewise, so each vertex of
$E$ is hit $q$ times and $\sum_w[p_{p(w)}]=q\sum_v[p_v]=q[1]$. Fullness: $\langle
e_w,e_w\rangle=p_{p(w)}$ and $\{p_v\}$ generates $\cO_k$ as an ideal (it contains the unit).
\end{proof}

\begin{theorem}[The left action]\label{XVI:thm:leftaction}
There is a unique $*$-homomorphism $\phi\colon\cO_{k+1}\to\cK(\cX_\eta)$ with
\[
\phi(q_w)=\Theta_{e_w,e_w}\quad(\text{projection onto the }w\text{-summand}),\qquad
\phi(t_f)\,e_{w}=\delta_{w,r(f)}\;e_{s(f)}\cdot s_{p(f)} .
\]
It is injective. The verification of the Cuntz--Krieger relation \emph{(CK2)} for the
operators $\phi(t_f)$ uses precisely that $\eta^{\sharp}$ is an out-covering.
\end{theorem}

\begin{proof}
Define $V_f:=\phi(t_f)$ by the stated formula: $V_f$ maps the $r(f)$-summand
$p_{p(r(f))}\cO_k$ into the $s(f)$-summand by left multiplication by $s_{p(f)}$, and kills
the other summands. Since $p$ is a graph morphism, $p(f)$ runs $p(s(f))\to p(r(f))$, so
$s_{p(f)}=p_{p(s(f))}s_{p(f)}p_{p(r(f))}$ and $V_f$ does land in $p_{p(s(f))}\cO_k$; each
$V_f$ is adjointable with $V_f^{*}\xi=\delta$-shift back by $s_{p(f)}^{*}$, hence
$V_f\in\cK(\cX_\eta)$ (finite rank in the frame $\{e_w\}$). Now
\[
V_f^{*}V_f\,\xi=s_{p(f)}^{*}s_{p(f)}\,\xi_{r(f)}
=p_{r(p(f))}\,\xi_{r(f)}=p_{p(r(f))}\,\xi_{r(f)}=\xi_{r(f)}\qquad(\xi\in\cX_\eta),
\]
using (CK1) in $\cO_k$ and $r(p(f))=p(r(f))$; so $V_f^{*}V_f=\Theta_{e_{r(f)},e_{r(f)}}=\phi(q_{r(f)})$,
which is (CK1) for $\phi$. For (CK2), fix $w\in F^0$ and compute $\sum_{s(f)=w}V_fV_f^{*}$.
Each such $V_fV_f^{*}$ acts on the $w=s(f)$-summand by $\xi\mapsto s_{p(f)}s_{p(f)}^{*}\xi$,
so
\[
\sum_{s(f)=w}V_fV_f^{*}\Big|_{w\text{-summand}}=\sum_{s(f)=w}s_{p(f)}s_{p(f)}^{*}
=\sum_{s(e)=p(w)}s_es_e^{*}=p_{p(w)} ,
\]
where the middle equality is the out-covering property: $f\mapsto p(f)$ is a \emph{bijection}
from $\{f:s(f)=w\}$ onto $\{e:s(e)=p(w)\}$, so the sum reindexes exactly, and the last
equality is (CK2) in $\cO_k$. The right side is the identity of the $w$-summand, i.e.\
$\phi(q_w)$; and $V_fV_f^{*}$ kills every other summand. Hence $\sum_{s(f)=w}\phi(t_f)\phi(t_f)^{*}=\phi(q_w)$,
which is (CK2) for $\phi$. The universal property of $C^*(F)$ gives the $*$-homomorphism
$\phi$. Injectivity: $\phi(q_w)\ne0$ for every $w$ (its range $p_{p(w)}\cO_k\ne0$), and
$\cO_{k+1}$ is simple, so $\ker\phi=0$.

The out-covering hypothesis enters at exactly one point --- the reindexing
$\{p(f):s(f)=w\}=\{e:s(e)=p(w)\}$ in (CK2) --- and nowhere else. For the in-covering
$\tau^{\sharp}$ this reindexing holds on \emph{in}-arcs, which is why $\tau^{\sharp}$ gave a
homomorphism into $\cO_{k+1}$ rather than a correspondence; for $\eta^{\sharp}$ it is the
out-arcs that biject, and the natural home is $\cK(\cX_\eta)$. The cross terms
$s_{p(f)}s_{p(f)}^{*}$, $s_{p(f')}s_{p(f')}^{*}$ that obstructed the sum-over-lifts
homomorphism of [Ch.~\ref{chap:VIII}] are here mutually orthogonal (distinct out-arcs $p(f)\ne p(f')$ of
$p(w)$), so they add rather than interfere.
\end{proof}

\begin{definition}[The head correspondence and its Kasparov class]\label{XVI:def:class}
$\cE_\eta:=(\cX_\eta,\phi)$ is an injective, full, proper $C^*$-correspondence from
$\cO_{k+1}$ to $\cO_k$; its Kasparov class is
\[
[\cE_\eta]:=[(\cX_\eta,\phi,0)]\in\KK(\cO_{k+1},\cO_k).
\]
\end{definition}

\section{The index is the head pushforward}\label{XVI:sec:index}

\begin{theorem}[$K_0$-index of the head correspondence]\label{XVI:thm:index}
The map $\otimes_{\cO_{k+1}}[\cE_\eta]\colon K_0(\cO_{k+1})\to K_0(\cO_k)$ is the head
pushforward: under \eqref{XVI:eq:ktheory} it is induced by $H_k=\operatorname{diag}(\et,\et)$,
i.e.\ the descent of $\et\colon\coker\Delta_{k+1}^{t}\to\coker\Delta_k^{t}$
\eqref{XVI:eq:etdesc}. On generators $[q_w]\mapsto[p_{\eta^{\sharp}(w)}]$; the $K_1$-map is
multiplication by $q$; and $[\cE_\eta]\otimes[1_{\cO_k}]^{\vee}$ recovers
$[\cX_\eta]=q[1_{\cO_k}]$.
\end{theorem}

\begin{proof}
For a vertex projection $q_w\in\cO_{k+1}$, $\phi(q_w)\cX_\eta$ is the $w$-summand
$p_{p(w)}\cO_k$, of class $[p_{p(w)}]=[p_{\eta^{\sharp}(w)}]$ in $K_0(\cO_k)$; since the
classes $[q_w]$ generate $K_0(\cO_{k+1})$ and the product is additive, the index map sends
$[q_w]\mapsto[p_{\eta^{\sharp}(w)}]$, which on the generators $e_w\mapsto e_{\eta^{\sharp}(w)}$
of $\coker(I-B^{t})$ is the matrix $H_k=\operatorname{diag}(\et,\et)$. That $H_k$ descends to
$\coker(I-B_{k+1}^{t})\to\coker(I-B_k^{t})$ is \eqref{XVI:eq:etdesc}; conjugating by $U$ of
[Ch.~\ref{chap:VIII}, Thm.~1.4] (which commutes with $\operatorname{diag}$ blocks) turns it into $\et$ on
$\coker\Delta^{t}$. The $K_1$-map: $K_1(\cO_k)=\ker(I-B_k^{t})=\Z(\mathbf1,-\mathbf1)$ and
$H_k(\mathbf1,-\mathbf1)=q(\mathbf1,-\mathbf1)$ (Table~\ref{XVI:tab:battery}, line (K)). The last
statement is Lemma~\ref{XVI:lem:fgp}.
\end{proof}

\begin{proposition}[The correspondence is not invertible]\label{XVI:prop:notinv}
$[\cE_\eta]$ is not a $\KK$-equivalence. Its index $\et$ on $\Jac(Y_{k+1})\to\Jac(Y_k)$ is
surjective with kernel the $p$-torsion subgroup $\Jac(Y_{k+1})[p]$ when $q=p$ is prime, of
order $p^{\,n_k(p-1)-1}$; equivalently $\cX_\eta$ is not a Morita equivalence bimodule
($\phi(\cO_{k+1})\subsetneq\cK(\cX_\eta)$).
\end{proposition}

\begin{proof}
By reversal symmetry (Remark~\ref{XVI:rem:reversal}) $\et$ on $\Tor\coker\Delta_{k+1}^{t}\to
\Tor\coker\Delta_k^{t}$ is the mirror of the tail norm $\ta$ on $\Jac(Y_{k+1})\to\Jac(Y_k)$
of [Ch.~\ref{chap:VIII}, Thm.~2.8]: surjective with kernel the $p$-torsion of order $p^{\,n_k(p-1)-1}$
(Table~\ref{XVI:tab:battery}, line (N)). A $\KK$-equivalence induces isomorphisms on $K_*$;
$\et$ is not injective, so $[\cE_\eta]$ is not invertible. Were $\cX_\eta$ an
imprimitivity bimodule, $\phi$ would be onto $\cK(\cX_\eta)$ and the index an isomorphism.
\end{proof}

\begin{remark}[Reversal symmetry]\label{XVI:rem:reversal}
Let $\rho_k$ be the path-reversal permutation of $V_{Y_k}$ (reverse each non-backtracking
path). Then $\rho_kA_k\rho_k^{-1}=A_k^{t}$ and $\rho_k\ta=\et\rho_{k+1}$
(Table~\ref{XVI:tab:battery}, line (R)): reversing all arrows exchanges $\tau\leftrightarrow\eta$,
$\Delta\leftrightarrow\Delta^{t}$, $\cO_{B}\leftrightarrow\cO_{B^{t}}$, and carries the
in-covering theory of [Ch.~\ref{chap:VIII}] for $(\tau,\Delta)$ to the out-covering theory of this paper
for $(\eta,\Delta^{t})$. At the algebra level reversal is Kaminker--Putnam
Spanier--Whitehead duality \cite{KP}, $\cO_{B}$ and $\cO_{B^{t}}$ dual, which exchanges
$K_0$ with $K^{1}$; we use only the elementary permutation identity, not the duality, and
cite \cite{KP} for context.
\end{remark}

\section{Kasparov products and the Hecke operators}\label{XVI:sec:products}

Write $[\varphi_k]\in\KK(\cO_k,\cO_{k+1})$ for the class of the tail homomorphism
($K_0$-map $\taup$, upward) and $[\cE_\eta^{(k)}]\in\KK(\cO_{k+1},\cO_k)$ for the head
correspondence ($K_0$-map $\et$, downward).

\begin{theorem}[Round trips are Hecke operators]\label{XVI:thm:hecke}
The Kasparov products
\[
[\varphi_k]\otimes_{\cO_{k+1}}[\cE_\eta^{(k)}]\in\KK(\cO_k,\cO_k),\qquad
[\cE_\eta^{(k)}]\otimes_{\cO_k}[\varphi_k]\in\KK(\cO_{k+1},\cO_{k+1})
\]
induce on $K_0$ the maps
\[
\et\,\taup=A_k^{t}\ \text{on}\ K_0(\cO_k),\qquad
\taup\,\et=A_{k+1}^{t}\ \text{on}\ K_0(\cO_{k+1}),
\]
i.e.\ the level Hecke operators. These are the $C^*$-bivariant form of the degeneracy
identities $\ta\etp=A_k$, $\etp\ta=A_{k+1}$ of {\rm[Ch.~\ref{chap:V}, Prop.~1.2]}.
\end{theorem}

\begin{proof}
The $K_0$-functor takes Kasparov products to composites, so
$\otimes([\varphi_k]\otimes[\cE_\eta^{(k)}])=(\otimes[\cE_\eta^{(k)}])\circ(\otimes[\varphi_k])
=H_k\circ T^{k}$ with $T^{k}=\operatorname{diag}(\taup,\taup)=K_0(\varphi_k)$
(Theorem~\ref{XVI:thm:index}, [Ch.~\ref{chap:VIII}, Thm.~2.5]). Blockwise
$H_kT^{k}=\operatorname{diag}(\et\taup,\et\taup)$, and $\et\taup=(\ta\etp)^{t}=A_k^{t}$ by
transposing \eqref{XVI:eq:calc}; similarly $T^{k}H_k=\operatorname{diag}(\taup\et,\taup\et)$ with
$\taup\et=(\etp\ta)^{t}=A_{k+1}^{t}$ (Table~\ref{XVI:tab:battery}, line (H)). Both descend to
$\coker(I-B^{t})$, since each factor does.
\end{proof}

\begin{remark}[The opposite side and the full calculus]\label{XVI:rem:opposite}
The other two degeneracy maps, $\ta$ and $\etp$, descend not on $K_0$ but on the
Bowen--Franks group $\coker(I-B)\cong\coker\Delta$, the $K_0$ of the opposite algebra
$\cO_k^{\mathrm{op}}\cong C^*(\overline{\Ysh{k}})$: there
$\operatorname{diag}(\ta,\ta)$ and $\operatorname{diag}(\etp,\etp)$ intertwine $I-B$, with
composites $A_k$ and $A_{k+1}$ (Table~\ref{XVI:tab:battery}, line (B)). Thus all four
degeneracy maps of [Ch.~\ref{chap:V}] are realized, two by $\KK$-classes on the $K_0$-side ($\taup$ up,
$\et$ down) and two by their reversal-dual partners on the Bowen--Franks side; the
``$qI$'' identities $\ta\taup=\et\etp=qI$ pair one map from each side and are the
Pontryagin adjointness of [Ch.~\ref{chap:VIII}, Lem.~2.1]. The correspondence $\cE_\eta$ supplies the one
arrow that was missing: the canonical downward morphism inducing $\et$.
\end{remark}

\begin{definition}[The $\KK$-category of the tower]\label{XVI:def:cat}
Let $\mathcal{T}$ be the small category enriched in $\KK$ (a subcategory of the Kasparov
category $\mathsf{KK}$) with objects $\{\cO_k:k\ge1\}$ and morphisms generated under
Kasparov product and $\Z$-linear combination by $[\varphi_k]\in\KK(\cO_k,\cO_{k+1})$ and
$[\cE_\eta^{(k)}]\in\KK(\cO_{k+1},\cO_k)$ for all $k$.
\end{definition}

\begin{corollary}[Structure of $\mathcal{T}$]\label{XVI:cor:cat}
In $\mathcal{T}$ the endomorphism $[\varphi_k]\otimes[\cE_\eta^{(k)}]$ of $\cO_k$ acts on
$K_0(\cO_k)$ as the Hecke operator $A_k^{t}$, and $[\cE_\eta^{(k)}]\otimes[\varphi_k]$ on
$\cO_{k+1}$ as $A_{k+1}^{t}$; the diagonal composite
$[\varphi_k]\otimes[\cE_\eta^{(k)}]\otimes\cdots$ over a $\tau$-then-$\eta$ round trip
realizes, on $K_0$, the polynomial algebra generated by the Hecke operators of the tower.
No morphism of $\mathcal{T}$ is invertible except the identities: $[\varphi_k]$ has index
$\taup$ (injective, not onto: cokernel the new $p$-torsion) and $[\cE_\eta^{(k)}]$ has index
$\et$ (onto, not injective), so $\mathcal{T}$ is a genuinely one-way-plus-one-way category,
the operator-algebraic shadow of the two-sidedness of the pro-module $X_\infty$ of [Ch.~\ref{chap:V}].
\end{corollary}

\begin{proof}
Theorems~\ref{XVI:thm:index}, \ref{XVI:thm:hecke}; injectivity/surjectivity from [Ch.~\ref{chap:VIII}, Thm.~2.5]
and Proposition~\ref{XVI:prop:notinv}.
\end{proof}

\section{The verification battery}\label{XVI:sec:battery}

\begin{table}[ht]
\centering\footnotesize
\renewcommand{\arraystretch}{1.2}
\resizebox{\textwidth}{!}{%
\begin{tabular}{lcl}
\toprule
check & transitions & result\\
\midrule
(O) $\eta^{\sharp}$ morphism, bijective on out-arcs, $q$-to-one on in-arcs; $\tau^{\sharp}$ the mirror in-covering & $k\le3$ & exact\\
(R) $\rho_kA_k\rho_k^{-1}=A_k^{t}$, $\rho_k\ta=\et\rho_{k+1}$ & $k\le3$ & exact\\
(K) $H_k(I-B_{k+1}^{t})=(I-B_k^{t})H_k$; $\operatorname{diag}(\ta,\ta)$ does not; $H_k\mathbf1=q\mathbf1$; $H_k(\mathbf1,-\mathbf1)=q(\mathbf1,-\mathbf1)$ & $k\le3$ & exact\\
(H) $H_kT^{k}=\operatorname{diag}(A_k^{t})$, $T^{k}H_k=\operatorname{diag}(A_{k+1}^{t})$; witness $\operatorname{diag}(A^{t}-q)=(I-B^{t})VN$ & $k\le3$ & exact\\
(B) opposite side: $\operatorname{diag}(\ta,\ta)$, $\operatorname{diag}(\etp,\etp)$ descend on $\coker(I-B)$; composites $A_k$, $A_{k+1}$ & $k\le3$ & exact\\
(N) $\et$ on $\Tor\coker\Delta^{t}$: surjective, $|\ker|=2^{11},2^{23},2^{47}=2^{\,n_k(q-1)-1}$; $\rank_{\F_2}A_{k+1}=n_k$ & $k\le3$ & exact\\
(X) $\Ext(K_0(\cO_{k+1}),K_1(\cO_k))=\Jac(Y_{k+1})$ (lift ambiguity); $\Ext(K_0(\cO_k),K_1(\cO_k))=\Jac(Y_k)$ & $k\le3$ & exact\\
\bottomrule
\end{tabular}}
\medskip
\caption{Exact integer verification (\texttt{head\_correspondence.py}) on the $K_4$ tower,
$q=2$, at the transitions $k+1\to k$, $k\le3$; shadow matrices of size $24,48,96,192$.
$\Jac$ data as in [Ch.~\ref{chap:VIII}, Table~1].}
\label{XVI:tab:battery}
\end{table}

\section{Scope, residue, corrections}\label{XVI:sec:scope}

\begin{remark}[Solved, open, impossible]\label{XVI:rem:scope}
\emph{Solved:} the open item [Ch.~\ref{chap:VIII}, Rem.~5.1(a)] --- a canonical $C^*$-correspondence
$\cO_{k+1}\rightsquigarrow\cO_k$ induced by the out-covering $\eta^{\sharp}$ with $K_0$-map
$\et$. It is $\cE_\eta=(\cX_\eta,\phi)$ of Definition~\ref{XVI:def:class}: a finitely generated
projective, full, injective proper correspondence whose class in $\KK(\cO_{k+1},\cO_k)$
induces $\et$ (Theorem~\ref{XVI:thm:index}), completing the two-sided realization of the
degeneracy calculus (Theorem~\ref{XVI:thm:hecke}, Remark~\ref{XVI:rem:opposite}) and organizing the
levels into the $\KK$-category $\mathcal{T}$ (Definition~\ref{XVI:def:cat},
Corollary~\ref{XVI:cor:cat}). \emph{Open:} (a) the Cuntz--Pimsner algebra
$\cO_{\cX_\eta}$ of the head correspondence, and whether the inductive system it generates
gives a second model of the limit $\cO_\infty$ of [Ch.~\ref{chap:VIII}] dual to the first; (b) an
explicit unbounded (spectral-triple) representative of $[\cE_\eta]$ pairing with the
$\cL$-invariant $KK$-classes of [Ch.~\ref{chap:XIV}]; (c) the Hecke-module structure of $K_0(\cO_\infty)$
under the operators $[\varphi_k]\otimes[\cE_\eta^{(k)}]$ of Corollary~\ref{XVI:cor:cat}, and with
it [Ch.~\ref{chap:V}, Prob.~5.1]. \emph{Impossible:} a unital $*$-homomorphism $\cO_{k+1}\to\cO_k$
inducing $\et$ (Section~\ref{XVI:sec:setting}: $\et[1]=q[1]$); and, as throughout the series, a
$\KK$-equivalence between adjacent levels (Proposition~\ref{XVI:prop:notinv}).
\end{remark}

\begin{remark}[Corrections to the task specification]\label{XVI:rem:corrections}
(1) The correspondence's right module is not ``freely generated by the graph's
generators''; it is the finitely generated projective module
$\cX_\eta=\bigoplus_{w}p_{\eta^{\sharp}(w)}\cO_k$, a sum of corners, of $K$-theory class
$q[1]$. (2) Positive-definiteness and self-adjointness of the inner product are automatic
(the inner product is $\sum_w\xi_w^{*}\zeta_w$, a sum of positives) and use no
Cuntz--Krieger relation; what (CK2) --- together with the out-covering bijection on
\emph{out}-arcs --- establishes is that the \emph{left action} $\phi$ is a
$*$-homomorphism (Theorem~\ref{XVI:thm:leftaction}). This is the one place the geometry is
used, and it is on out-arcs, not on the (CK1) side as the specification suggested. (3) The
left action lands in $\cK(\cX_\eta)$, not in $\cO_k$; that is exactly why the object is a
correspondence and not a homomorphism, and why it exists where the sum-over-lifts
homomorphism of [Ch.~\ref{chap:VIII}] could not. (4) The head pushforward $\et$, not $\ta$, is the map the
downward correspondence induces; $\ta$ descends on the opposite (Bowen--Franks) side, not on
$K_0$ (Remark~\ref{XVI:rem:opposite}). (5) ``All six degeneracy identities in $\KK$'' is made
precise as the Kasparov products of Theorem~\ref{XVI:thm:hecke} and the opposite-side identities
of Remark~\ref{XVI:rem:opposite}: the two Hecke operators are the two round-trip products, and
$\ta\taup=\et\etp=qI$ are cross-side adjointness identities, not products of two
$K_0$-side classes.
\end{remark}

\section*{Acknowledgement of provenance and caveats}

Unrefereed preprint. Every numbered statement is proved; the correspondence $\cX_\eta$, its
left action, and its index are constructed and computed in full, and the tower matrix
identities are quoted from [Ch.~\ref{chap:VIII}] and re-verified exactly on $K_4$. The bivariant inputs
--- $C^*$-correspondences define $\KK$-classes, and $K_0$ turns Kasparov products into
composites --- are standard \cite{Kas,Bla}; \cite{KP} is cited for context only.
Bibliographic data of \cite{Kas,Bla,R,KP} were checked against publisher or indexing pages.

\part{\texorpdfstring{Higher rank: $\widetilde A_2$ and $\widetilde A_d$ buildings}{Higher rank: A2 and Ad buildings}}
% generated from Paper 10/anccft10.tex
\chapter[\texorpdfstring{X. $\Atwo$ complexes, higher simplicial Laplacians, and the fate of the trace obstruction in
rank two}{X. A2 complexes, higher simplicial Laplacians, and the fate of the trace obstruction in rank two}]{\texorpdfstring{Algorithmic non-commutative class field theory, X: $\Atwo$ complexes, higher simplicial Laplacians, and the fate of the trace obstruction in
rank two}{Algorithmic non-commutative class field theory, X: A2 complexes, higher simplicial Laplacians, and the fate of the trace obstruction in rank two}}\label{chap:X}
\chaptermark{X}
\begin{chapterabstract}
The series has so far lived on quotients of the Bruhat--Tits tree of $\PGL_2$. This paper
moves one rank up, to finite $\Atwo$ complexes --- quotients of two-dimensional buildings of
order $q$, locally the building of $\PGL_3$ over a local field with residue field $\F_q$ ---
and asks which of the rank-one theorems survive. Three things are established. First, the
Hodge Laplacians $\Delta_0,\Delta_1,\Delta_2$ of the typed chain complex give higher
critical groups $\Jac_i(Y)=\Tor\coker\Delta_i$, computed exactly on thin tori and on five
thick complexes of order $2$ with $21$, $24$, $42$, $84$ and $168$ vertices, built here from
a triangle presentation over $\PG(2,2)$ compatible with a Singer cycle and verified to have
Heawood-graph links; the Hodge vertex Laplacian is no longer a Hecke operator, and the
arithmetic bridge of [Ch.~\ref{chap:I}, Thm.~3.5] is carried instead by the \emph{Hecke--Laplacian}
$\DeltaH=(q^{3}-1)I+A_1-qA_2$, the value at $u=1$ of the cubic pencil
$P(u)=\det(I-uA_1+qu^{2}A_2-q^{3}u^{3}I)$, whose critical group has order
$\frac1n\prod_j|P_j(1)|$, a product of inverse local $\mathrm{GL}_3$ $L$-values. Second,
the companion $\Cthree$ of the cubic pencil is unimodularly equivalent to $\DeltaH\oplus I\oplus I$
exactly as in [Ch.~\ref{chap:II}], but the trace obstruction of [Ch.~\ref{chap:II}] does \emph{not} transfer: the $\Z/3$
type grading forces $\Tr(\Cthree)^{k}=0$ for $3\nmid k$, and a block computation gives
$\Tr(\Cthree)^{3}=n(q-1)^{2}(q+1)\ge0$. There is no cubic trace obstruction; on the five
thick complexes, all of which turn out to be Ramanujan, every net trace of the reduced
(type-preserving) pencil is non-negative and its Perron condition holds, so by the theorem
of Kim--Ormes--Roush the nonzero spectrum of $\Cthree$ is that of a non-negative integer
matrix, irreducible of period three --- the rank-two pencil is spectrally realizable, the
opposite of the rank-one situation, and the realization problem for the $\PGL_3$ Hecke
critical group is open and unobstructed. Third, the shadow construction of [Ch.~\ref{chap:VI}] extends to
arbitrary
integer matrices with positive diagonal by a signed-shadow lemma in which positive
off-diagonal entries are produced by one-directional auxiliary vertices with double loops;
it yields unital Kirchberg algebras $\cO(Y,\Delta_1)$ and $\cO(Y,\DeltaH)$ with
$K_0\cong\Z^{b_1}\oplus\Jac_1(Y)$ and $K_0\cong\Z\oplus\JacH(Y)$ respectively, functorially
in the complex. What the specification predicted and what is proved differ in two places,
and the paper says where.
\end{chapterabstract}

\section*{Status of the results}

Every numbered statement is proved. The inputs are finite linear algebra, the classical
Cuntz--Krieger facts quoted in [Chs.~\ref{chap:VI}, \ref{chap:VIII}], the Cartwright--Mantero--Steger--Zappa theorem
that a triangle presentation gives a group acting simply transitively on the vertices of an
$\Atwo$ building \cite{CMSZ}, and, for one corollary that nothing else depends on, the
Kim--Ormes--Roush theorem on the nonzero spectra of primitive integer matrices \cite{KOR}.
The Kang--Li zeta function of $\PGL_3$ complexes \cite{KL} is context for the cubic pencil,
not an input. Every finite identity was machine-verified exactly (Table~\ref{X:tab:battery});
the five thick complexes are constructed by the script, and their local structure (Heawood
links, regularity, flag condition) is checked rather than assumed --- with the result that
the flag condition holds only for the largest of them, a distinction the theorems respect. Section~\ref{X:sec:scope}
records what is open and the corrections to the specification.

\section{\texorpdfstring{$\Atwo$ complexes of order $q$}{A2 complexes of order q}}\label{X:sec:complexes}

\begin{definition}[$\Atwo$ complex]\label{X:def:a2}
An \emph{$\Atwo$ complex of order $q\ge1$} is a finite connected pure two-dimensional
simplicial complex $Y$ with vertex set $V$ ($n:=|V|$), edge set $E$ and triangle set $F$,
together with a type function $\tau\colon V\to\Z/3$, such that
\begin{enumerate}
\item[(i)] every triangle has one vertex of each type;
\item[(ii)] every vertex has exactly $q^{2}+q+1$ neighbours of type $\tau(v)+1$ and
$q^{2}+q+1$ of type $\tau(v)+2$;
\item[(iii)] every edge lies in exactly $q+1$ triangles;
\item[(iv)] (flag condition) every closed walk $x\to y\to z\to x$ with
$\tau(y)=\tau(x)+1$, $\tau(z)=\tau(x)+2$ is the boundary of a triangle.
\end{enumerate}
Every edge is oriented from type $t$ to type $t+1$ and every triangle $(a,b,c)$ is listed
with types $(t,t+1,t+2)$, so that $\partial_2(a,b,c)=(a,b)+(b,c)+(c,a)$ and
$\partial_1(a,b)=b-a$ are canonical; $A_1\in M_n(\{0,1\})$ is the adjacency $x\to y$ for
edges, $A_2:=A_1^{t}$, and we assume $A_1A_2=A_2A_1$ (true for building quotients, where
$A_1,A_2$ generate the commutative spherical Hecke algebra; verified on every complex used).
\end{definition}

Quotients of a thick $\Atwo$ building of order $q\ge2$ by a type-rotating, freely acting
group with injectivity radius $\ge2$ are $\Atwo$ complexes in this sense, with vertex links
the incidence graph of $\PG(2,q)$; the flag condition holds because buildings are flag
complexes. Thin complexes ($q=1$) are triangulated flat tori (Construction~\ref{X:con:thin}).

\begin{lemma}[Counts]\label{X:lem:counts}
$|E|=n(q^{2}+q+1)$, $|F|=\tfrac13n(q^{2}+q+1)(q+1)$, $\Tr A_1^{3}=3|F|$,
$\Tr(A_1A_2)=|E|$, and $\chi(Y)=n\bigl(1+\tfrac13(q^{2}+q+1)(q-2)\bigr)$. For $q=2$:
$|E|=|F|=7n$ and $\chi(Y)=n$.
\end{lemma}

\begin{proof}
(ii) gives $|E|$; each vertex lies in $(q^{2}+q+1)(q+1)$ triangles by (ii),(iii), each
counted three times; $\Tr A_1^{3}$ counts closed type-increasing $3$-walks, three per
triangle by (iv); $\Tr(A_1A_1^{t})=\sum_{x,y}A_1(x,y)^{2}=|E|$.
\end{proof}

\begin{construction}[Thin tori]\label{X:con:thin}
On $\Z^{2}$ let $d_1=(1,0)$, $d_2=(0,1)$, $d_3=(-1,-1)$, edges $v\to v+d_i$, type
$\tau(a,b)=a+b\bmod3$, triangles $\{v,v+d_i,v+d_i+d_j\}$, $i\neq j$. This is the
$\Atwo$ Coxeter complex (one apartment); its quotient by $m\Z^{2}$, $3\mid m$, $m\ge3$, is an
$\Atwo$ complex of order $1$ with $m^{2}$ vertices, $b_1=2$, $b_2=1$, hexagonal links.
\end{construction}

\begin{construction}[Thick complexes of order $2$]\label{X:con:thick}
Label the points of $\PG(2,2)$ by $\Z/7$ with lines $\{i,i+1,i+3\}$ and let
$\lambda(x)=\{x+1,x+2,x+4\}$, the Singer cycle. A triangle presentation compatible with
$\lambda$ \cite{CMSZ} is a set $T$ of $21$ triples with (A) $(x,y,z)\in T\Rightarrow
(y,z,x)\in T$, (B) $(x,y,\cdot)\in T$ iff $y\in\lambda(x)$, (C) the third entry is
determined by the first two. Exhaustive search over all bijections without absolute points
finds, for $\lambda$ the Singer cycle, the presentation with cyclic orbits
\[
T/\Z_3=\{(0,1,3),(0,2,6),(0,4,5),(1,2,4),(1,5,6),(2,3,5),(3,4,6)\}
\]
(axioms re-verified by the script). By \cite{CMSZ} the group
$\Gamma=\langle a_0,\dots,a_6\mid a_xa_ya_z=1,\ (x,y,z)\in T\rangle$ acts simply
transitively and type-rotatingly on the vertices of an $\Atwo$ building of order $2$; its
type-preserving subgroup $\Gamma^{0}$ has index $3$, and $X_3:=\Gamma^{0}\backslash B$ is
the $\Delta$-complex with three vertices, $21$ edges and $21$ triangles. We compute
$H_1(X_3;\Z)\cong\Z/2\oplus\Z/2\oplus\Z/14$ and build the abelian voltage covers of $X_3$
with deck group $A$ for the quotients $A=\Z/7$, $(\Z/2)^{3}$, $\Z/14$ and $H_1(X_3)$
itself, exactly as the derived graphs of [Ch.~\ref{chap:II}, \S5] in one dimension higher: vertices
$(i,g)$, edges lifted by voltages in $A$, triangles lifting because their boundaries are
null-homologous. The results $Y_{21},Y_{24},Y_{42},Y_{84},Y_{168}$ are simplicial, every
vertex link is the Heawood graph (bipartite, $3$-regular, $14$ vertices, girth $6$), and
conditions (i)--(iii) hold, with $|E|=|F|=7n$ (Table~\ref{X:tab:battery}, line (P)). The flag
condition (iv) holds for $Y_{168}$ and fails for the four smaller covers: there the deck
group is too small, some product $a_xa_ya_{z'}$ with $(x,y,z')\notin T$ has trivial voltage,
and $\Tr A_1^{3}$ exceeds $3|F|$ by $28n$, $21n$, $12n$, $4n$ respectively. So (iv) is not
implied by the link condition, and the five complexes separate the two hypotheses cleanly:
$Y_{168}$ is an $\Atwo$ complex in the full sense of Definition~\ref{X:def:a2}, the others
satisfy (i)--(iii) only. The count $168=|\PGL_3(\F_2)|$ for the maximal abelian cover
suggests that $Y_{168}$ is a congruence quotient; we neither need nor claim this. Whether the
building is that of $\PGL_3(\F_2((t)))$, of $\PGL_3(\Q_2)$, or exotic is not decided here:
every statement below depends only on the local structure of Definition~\ref{X:def:a2}, and
the statements that use (iv) are marked.
\end{construction}

\section{Hodge Laplacians and higher critical groups}\label{X:sec:hodge}

\begin{definition}\label{X:def:hodge}
$\Delta_0:=\partial_1\partial_1^{t}$ on $C_0$, $\Delta_1:=\partial_1^{t}\partial_1
+\partial_2\partial_2^{t}$ on $C_1$, $\Delta_2:=\partial_2^{t}\partial_2$ on $C_2$, and
$\Jac_i(Y):=\Tor\coker_\Z\Delta_i$, the \emph{$i$-th critical group}. Following
Duval--Klivans--Martin \cite{DKM} we also record $K(Y):=\ker\partial_1/\operatorname{im}
(\partial_2\partial_2^{t})$.
\end{definition}

\begin{proposition}[Structure]\label{X:prop:hodge}
$\partial_1\partial_2=0$; $\coker\Delta_i\cong\Z^{b_i}\oplus\Jac_i(Y)$ with $b_i$ the Betti
numbers of $Y$; $b_0=1$; $\Delta_0=2(q^{2}+q+1)I-(A_1+A_2)$ is the Laplacian of the
$1$-skeleton, so $\Jac_0(Y)$ is its sandpile group with $|\Jac_0(Y)|$ the number of
spanning trees of the $1$-skeleton; $\Delta_1$ has diagonal entries $q+3$, entry $0$ at pairs
of edges spanning a triangle, and entries $\pm1$ at pairs sharing a vertex and no triangle.
\end{proposition}

\begin{proof}
$\partial_1\partial_2(a,b,c)=(b-a)+(c-b)+(a-c)=0$. Each $\Delta_i$ is symmetric with
$\ker\Delta_i=\ker\partial_i\cap\ker\partial_{i+1}^{t}$ of rank $b_i$ (Hodge), so its
cokernel has free rank $b_i$. The diagonal of $\Delta_1$ is $2+(q+1)$ by (iii); for edges
$e=(a,b)$, $f=(b,c)$ in a triangle the down-term $\langle\partial e,\partial f\rangle=-1$
cancels the up-term $+1$, and without a triangle only the down-term survives.
\end{proof}

\section{The Hecke--Laplacian and the \texorpdfstring{$\mathrm{GL}_3$ $L$-values}{GL3 L-values}}\label{X:sec:hecke}

\begin{definition}[Cubic pencil and Hecke--Laplacian]\label{X:def:pencil}
$P(u):=I-uA_1+qu^{2}A_2-q^{3}u^{3}I\in M_n(\Z[u])$ and
\[
\DeltaH:=-P(1)=(q^{3}-1)I+A_1-qA_2,\qquad \JacH(Y):=\Tor\coker_\Z\DeltaH .
\]
For a joint eigenpair $(a_j,b_j)$ of $(A_1,A_2)$ write $P_j(u)=1-a_ju+qb_ju^{2}-q^{3}u^{3}
=\prod_{i=1}^{3}(1-\lambda_{j,i}u)$; the trivial pair $a=b=q^{2}+q+1$ gives
$(1-u)(1-qu)(1-q^{2}u)$.
\end{definition}

The pencil is the vertex factor of the Kang--Li zeta function of a $\PGL_3$ complex
\cite{KL}, and $P_j(u)^{-1}$ is the local $L$-factor of the unramified $\mathrm{GL}_3$
Hecke eigenform with Satake parameters $\lambda_{j,i}/q$; we use only the polynomial.

\begin{theorem}[Hecke critical group and $L$-values]\label{X:thm:hecke}
\begin{enumerate}
\item[(i)] $\DeltaH\mathbf1=0$ and $\mathbf1^{t}\DeltaH=0$.
\item[(ii)] If $P_j(1)\neq0$ for every nontrivial $j$ (no nontrivial eigenform has a
Satake parameter $\lambda_{j,i}=1$; automatic when $|\lambda_{j,i}|=q$), then
$\rk\DeltaH=n-1$, all $n$ principal cofactors of $\DeltaH$ are equal, and
\[
\coker_\Z\DeltaH\cong\Z\oplus\JacH(Y),\qquad
|\JacH(Y)|=\frac1n\prod_{j\ \mathrm{nontrivial}}\bigl|P_j(1)\bigr| .
\]
\item[(iii)] $\DeltaH\neq\Delta_0$ for $q\ge1$, and $\Delta_0$ is not a specialization of the
pencil: no $u$ gives $P(u)\in\Q^{\times}\Delta_0$. For $\PGL_2$ the two operators coincide
($(q{+}1)I-T_p=\Delta_0=-P(1)$ with $P(u)=I-uT_p+qu^{2}$).
\end{enumerate}
\end{theorem}

\begin{proof}
(i) Row sums: $(q^{3}-1)+(q^{2}+q+1)-q(q^{2}+q+1)=0$; column sums likewise since
$\mathbf1^{t}A_1=\mathbf1^{t}A_2=(q^{2}+q+1)\mathbf1^{t}$. (ii) The eigenvalues of $\DeltaH$
on the joint eigenspaces are $-P_j(1)$, zero exactly for the trivial pair under the
hypothesis, so the kernel is $\Z\mathbf1$ on both sides; the equal-cofactor argument of
[Ch.~\ref{chap:III}, proof of Thm.~2.1, Step~2] applies verbatim to a matrix with two-sided kernel
$\mathbf1$ and gives $\prod_{\mu\neq0}\mu=n\kappa$ with $\kappa$ the common cofactor
$=|\Tor\coker|$. (iii) $P(u)\in\Q\Delta_0$ would need $-u:qu^{2}=-1:-1$, i.e.\ $u=-1/q$,
where $P(-1/q)=2I+(A_1+A_2)/q$ has the wrong constant term unless $2q=2(q^{2}+q+1)$,
impossible.
\end{proof}

\begin{theorem}[Companion reduction]\label{X:thm:companion}
Let $\Cthree=\begin{psmallmatrix}A_1&-qA_2&q^{3}I\\ I&0&0\\0&I&0\end{psmallmatrix}$, so
that $\det(I-u\Cthree)=\det P(u)$. With
$U=\begin{psmallmatrix}I&q^{3}I-qA_2&q^{3}I\\0&I&0\\0&0&I\end{psmallmatrix}$ and
$V=\begin{psmallmatrix}I&0&0\\ I&I&0\\ I&I&I\end{psmallmatrix}$, both unimodular,
\[
U\,(I-\Cthree)\,V=P(1)\oplus I_n\oplus I_n=(-\DeltaH)\oplus I_n\oplus I_n,
\qquad\text{hence}\qquad \coker_\Z(I-\Cthree)\cong\Z\oplus\JacH(Y).
\]
\end{theorem}

\begin{proof}
$I-\Cthree=\begin{psmallmatrix}I-A_1&qA_2&-q^{3}I\\-I&I&0\\0&-I&I\end{psmallmatrix}$; the
first block row of $U(I-\Cthree)$ is $(I-A_1,qA_2,-q^{3}I)+(q^{3}I-qA_2)(-I,I,0)
+q^{3}(0,-I,I)=(P(1),0,0)$, and the column operations $V$ clear the subdiagonal $-I$'s.
This is [Ch.~\ref{chap:II}, \S1.1] one degree higher.
\end{proof}

\section{The fate of the trace obstruction}\label{X:sec:trace}

\begin{theorem}[Grading]\label{X:thm:grading}
$\Tr(A_1^{\alpha}A_2^{\beta})=0$ unless $\alpha\equiv\beta\pmod3$. Consequently
$\Tr(\Cthree)^{k}=0$ for every $k$ with $3\nmid k$, and the spectrum of $\Cthree$ is
invariant under multiplication by cube roots of unity.
\end{theorem}

\begin{proof}
A closed walk contributing to the trace changes type by $\alpha+2\beta\equiv\alpha-\beta$.
$\Tr(\Cthree)^{k}=\sum_jp_k(a_j,b_j)$ with $p_k$ the $k$-th power sum of the roots of
$\lambda^{3}-a\lambda^{2}+qb\lambda-q^{3}$, a polynomial in monomials $a^{\alpha}(qb)^{\beta}
(q^{3})^{\gamma}$ with $\alpha+2\beta+3\gamma=k$, so $\alpha\equiv\beta$ forces $3\mid k$.
Equivalently $D_\omega=\operatorname{diag}(\omega^{\tau(v)})\oplus\cdots$ conjugates
$\omega\Cthree$ to $\Cthree$.
\end{proof}

\begin{theorem}[No cubic trace obstruction]\label{X:thm:notrace}
For every $\Atwo$ complex of order $q$,
\[
\Tr\Cthree=0,\qquad\Tr(\Cthree)^{2}=0,\qquad
\Tr(\Cthree)^{3}=3|F|-3q|E|+3nq^{3}=n(q-1)^{2}(q+1)\ \ge\ 0,
\]
with equality only in the thin case $q=1$. Hence the mechanism of {\rm[Ch.~\ref{chap:II}, Thms.~1.5--1.6]}
--- a negative trace of a power of the companion, incompatible with any non-negative
matrix of any size --- has no analogue in rank two at orders $\le5$: the first net trace
that is not forced to be non-negative by these identities is $\tr_6$.
\end{theorem}

\begin{proof}
$p_1=a$, $p_2=a^{2}-2qb$, $p_3=a^{3}-3qab+3q^{3}$ (Newton). $\Tr A_1=0$ (no loops),
$\Tr A_1^{2}=0$ and $\Tr A_2=0$ by Theorem~\ref{X:thm:grading}, $\Tr A_1^{3}=3|F|$ and
$\Tr(A_1A_2)=|E|$ by Lemma~\ref{X:lem:counts}; substitute the counts:
$(q^{2}+q+1)(q+1)-3q(q^{2}+q+1)+3q^{3}=q^{3}-q^{2}-q+1=(q-1)^{2}(q+1)$. The net trace
$\tr_m=\sum_{d\mid m}\mu(m/d)\Tr(\Cthree)^{d}$ vanishes for $m\le5$, $m\neq3$, and equals
$\Tr(\Cthree)^{3}\ge0$ for $m=3$. Without the flag condition the identity reads
$\Tr(\Cthree)^{3}=\Tr A_1^{3}-3q|E|+3nq^{3}$, still non-negative since $\Tr A_1^{3}\ge3|F|$:
on $Y_{168}$ the value is $504=n(q-1)^{2}(q+1)$ exactly; on $Y_{21}$ it is $651=63+588$,
the excess being the spurious closed $3$-walks (Table~\ref{X:tab:battery}).
\end{proof}

\begin{remark}[Why rank one is different]\label{X:rem:why}
In [Ch.~\ref{chap:II}] the companion of $1-uT_p+qu^{2}$ has $\Tr C^{2}=\Tr T_p^{2}-2qn=n(q+1)-2qn<0$:
the $\Z/2$ type structure of a graph allows a walk to return in two steps (edge reversal),
so the symmetric adjacency contributes $\Tr T_p^{2}=2|E|>0$ but not enough to beat the
$-2qn$ of the negative block. In rank two the type grading is $\Z/3$: a type-increasing walk
cannot return before three steps, the negative block $-qA_2$ is invisible to
$\Tr(\Cthree)^{k}$ for $k<3$, and at $k=3$ the chambers contribute $3|F|=n(q^{2}+q+1)(q+1)$,
which exceeds the $3q|E|-3nq^{3}=n(q^{2}+q+1)\cdot3q-3nq^{3}$ removed by the negative block.
The obstruction of [Ch.~\ref{chap:II}] is a rank-one accident of parity.
\end{remark}

\begin{definition}[Reduced pencil]\label{X:def:reduced}
By Theorem~\ref{X:thm:grading}, $\det P(u)\in\Z[u^{3}]$; write $\det P(u)=Q(u^{3})$ with
$Q\in\Z[v]$ of degree $n$. Its roots are $\lambda^{3}$, one for each $\omega$-orbit
$\{\lambda,\omega\lambda,\omega^{2}\lambda\}$ of eigenvalues of $\Cthree$: the trivial orbits
give $1,q^{3},q^{6}$, and the power sums of the multiset $\Lambda_Q$ of roots are
$p_d(\Lambda_Q)=\tfrac13\Tr(\Cthree)^{3d}$. $\Lambda_Q$ is the spectrum of the
type-preserving (three-step) dynamics.
\end{definition}

\begin{theorem}[Spectral realizability of the thick examples]\label{X:thm:realizable}
For each of $Y_{21},Y_{24},Y_{42},Y_{84},Y_{168}$:
\begin{enumerate}
\item[(i)] $\Cthree$ has exactly nine eigenvalues of modulus $\ge q$: the orbits of $1,q,q^{2}$
under $\omega$; every other eigenvalue has modulus exactly $q$ (numeric, to $10^{-6}$): the
complexes are Ramanujan in the sense of \cite{KL}. In particular $\Cthree$ itself never
satisfies the Perron condition of a primitive matrix --- $q^{2},q^{2}\omega,q^{2}\omega^{2}$
are all dominant --- which is forced by the grading and is the signature of an irreducible
matrix of period $3$.
\item[(ii)] For $\Lambda_Q$: $q^{6}$ is simple and strictly dominant; the net traces
\[
\tr_m(\Lambda_Q)=\sum_{d\mid m}\mu(m/d)\,\tfrac13\Tr(\Cthree)^{3d}
\]
are non-negative for $m\le10$ (exact integers) and provably positive for
$m\ge m_0\in\{2,3\}$ (Lemma~\ref{X:lem:bound}); hence by \cite{KOR} there is a primitive
non-negative integer matrix $B_0$ with nonzero spectrum $\Lambda_Q$.
\item[(iii)] Consequently the block-cyclic non-negative integer matrix
$\begin{psmallmatrix}0&B_0&0\\0&0&I\\I&0&0\end{psmallmatrix}$ has the nonzero spectrum of
$\Cthree$: the $\PGL_3$ pencil of these complexes is spectrally realizable by a non-negative
integer matrix, irreducible of period $3$.
\end{enumerate}
\end{theorem}

\begin{proof}
(i) is computation (Table~\ref{X:tab:battery}, line (R)); the grading statement is
Theorem~\ref{X:thm:grading}. (ii) The Boyle--Handelman conditions for a primitive integer
matrix are: $\Lambda_Q$ is the root multiset of a monic integer polynomial ($Q$), a simple
strictly dominant positive root ($q^{6}$, by (i)), and $\tr_m(\Lambda_Q)\ge0$ for all $m$;
\cite{KOR} proves them sufficient. (iii) If $M=\begin{psmallmatrix}0&X&0\\0&0&Y\\Z&0&0
\end{psmallmatrix}$ then $M^{3}=\operatorname{diag}(XYZ,YZX,ZXY)$, so the nonzero spectrum
of $M$ consists of all cube roots of the nonzero spectrum of $XYZ=B_0$, i.e.\ of the
$\omega$-orbits of $\Lambda_Q^{1/3}$, which is the nonzero spectrum of $\Cthree$.
\end{proof}

\begin{lemma}[Tail bound]\label{X:lem:bound}
Let $M<q^{2}$ bound the moduli of the genuinely nontrivial eigenvalues and $N=n-1$. Then
$\tr_m(\Lambda_Q)\ge q^{6m}+q^{3m}+1-NM^{3m}-\sum_{d\le m/2}\bigl(q^{6d}+q^{3d}+1+NM^{3d}\bigr)$,
positive for all $m\ge m_0$ with $m_0$ computable from $(n,q,M)$; for the five complexes
$m_0\le3$.
\end{lemma}

\begin{proof}
$|p_d(\Lambda_Q)-(1+q^{3d}+q^{6d})|\le NM^{3d}$; the proper divisors of $m$ are $\le m/2$; the
bound increases with $m$ once $q^{6m}$ dominates, which for $M=q$ and $n\le168$ happens by
$m=3$.
\end{proof}

\begin{remark}[What remains of the realization problem]\label{X:rem:realization}
Theorem~\ref{X:thm:realizable} is about the nonzero spectrum. The realization problem proper
--- a non-negative integer matrix $B$ shift equivalent over $\Z$ to $\Cthree$, so that
$\coker(I-B)\cong\Z\oplus\JacH(Y)$ with the Hecke dynamics --- asks for more than spectral
data, and we do not solve it. What is established is that the one tool which closed it in
rank one, Theorem~[Ch.~\ref{chap:II}, 1.6], is unavailable in rank two, and that the necessary conditions
known to us are all satisfied. The problem is open and, as far as the trace method can see,
unobstructed. Whether the double-loop mechanism of Section~\ref{X:sec:shadow} can be upgraded
from a $K$-theoretic to a dynamical realization is the natural next question, and [Ch.~\ref{chap:II}, Thm.~1.6] no longer forbids it here.
\end{remark}

\section{Signed shadows and the higher Kirchberg algebras}\label{X:sec:shadow}

\begin{lemma}[Signed shadow lemma]\label{X:lem:signed}
Let $M\in M_N(\Z)$ with $M_{ee}\ge1$ for all $e$, and let $\mathcal P=\{(e,f):e\neq f,\
M_{ef}>0\}$. Define a non-negative integer matrix $B$ on the vertex set
$\{e\}\sqcup\{w_{ef}:(e,f)\in\mathcal P\}\sqcup\{e'\}$ by the arcs
\begin{itemize}
\item $|M_{ef}|$ arcs $f\to e$ for each $e\neq f$ with $M_{ef}<0$;
\item for $(e,f)\in\mathcal P$: $M_{ef}$ arcs $f\to w_{ef}$, one arc $w_{ef}\to e$, two loops at
$w_{ef}$;
\item for each $e$: one arc $e\to e'$, $M_{ee}-1$ arcs $e'\to e$, two loops at $e'$.
\end{itemize}
Then, with $X:=(I-B^{t})_{\text{main},\text{aux}}$, $Y:=(I-B^{t})_{\text{aux},\text{main}}$
and $U=\begin{psmallmatrix}I&X\\0&I\end{psmallmatrix}$,
$V=\begin{psmallmatrix}I&0\\Y&I\end{psmallmatrix}$,
\[
U\,(I-B^{t})\,V\;=\;M\oplus(-I_{|\mathcal P|+N}) .
\]
If the graph on $\{e\}$ with edges $\{e,f\}$ for $M_{ef}\neq0$ is connected, $B$ is
irreducible with double-loop vertices, and $\cO_B:=C^*(\text{digraph of }B)$ is a unital
Kirchberg algebra with $K_0(\cO_B)\cong\coker_\Z M$ and $K_1(\cO_B)\cong\ker_\Z M$.
\end{lemma}

\begin{proof}
The auxiliary block of $I-B^{t}$ is $-I$ (each auxiliary vertex carries exactly two loops and
no other arcs among auxiliaries), so the Schur complement of that block is $R+XY$ where
$R$ is the main block: diagonal $1$, entry $M_{ef}$ at negative positions, $0$ elsewhere.
The auxiliary $w_{ef}$ contributes $X_{e,w}Y_{w,f}=(-1)(-M_{ef})=M_{ef}$ at $(e,f)$ and
nothing else; the shadow $e'$ contributes $X_{e,e'}Y_{e',e}=(-(M_{ee}-1))(-1)=M_{ee}-1$ at
$(e,e)$. Hence $R+XY=M$, and the displayed identity is the block factorization
$\begin{psmallmatrix}R&X\\Y&-I\end{psmallmatrix}=U^{-1}\bigl((R+XY)\oplus(-I)\bigr)V^{-1}$.
Irreducibility: every arc $f\to e$ or $f\to w_{ef}\to e$ is matched by the reverse
connection through $M_{fe}$ or through the shared vertex structure of a connected support
graph, and every $e'$ communicates with $e$ in both directions; the double loops give the
witness of [Ch.~\ref{chap:VI}, Thm.~3.1]; $K$-theory is the Cuntz--Krieger formula transported by $U,V$.
For $M=\Delta_p$ of a graph ($M_{ef}\le0$ off the diagonal) $\mathcal P=\emptyset$ and the
construction is [Ch.~\ref{chap:VI}, Def.~1.1] with $c=M_{ee}-1$.
\end{proof}

\begin{theorem}[Higher shadow algebras]\label{X:thm:shadow}
For an $\Atwo$ complex $Y$ let $\cO(Y,\Delta_1):=\cO_{B(\Delta_1)}$ and
$\cO(Y,\DeltaH):=\cO_{B(\DeltaH)}$ be the algebras of Lemma~\ref{X:lem:signed} for $M=\Delta_1$
(diagonal $q+3\ge1$) and $M=\DeltaH$ (diagonal $q^{3}-1\ge1$ for $q\ge2$). Both are unital
Kirchberg algebras, hence traceless, and
\[
K_0\bigl(\cO(Y,\Delta_1)\bigr)\cong\Z^{b_1}\oplus\Jac_1(Y),\quad K_1\cong\Z^{b_1};\qquad
K_0\bigl(\cO(Y,\DeltaH)\bigr)\cong\Z\oplus\JacH(Y),\quad K_1\cong\Z ,
\]
the second under the hypothesis of Theorem~\ref{X:thm:hecke}(ii). Type- and
orientation-preserving automorphisms of $Y$ act compatibly. Changing the edge orientations
conjugates $\Delta_1$ by a diagonal $\pm1$ matrix and changes $B(\Delta_1)$, but not the
$K$-theory.
\end{theorem}

\begin{proof}
Lemma~\ref{X:lem:signed} with Proposition~\ref{X:prop:hodge} and Theorem~\ref{X:thm:hecke}; the
support graph of $\Delta_1$ is connected because two edges of a connected complex are joined
by a chain of edges sharing vertices, and pairs spanning a triangle (where the entry is
$0$) are joined through a third edge; the support of $\DeltaH$ is the $1$-skeleton.
\end{proof}

\begin{remark}[Digraph, not cellular algebra]\label{X:rem:digraph}
The specification asked for a ``two-dimensional shadow complex with higher self-loop faces''
and a ``higher graph or cellular $C^*$-algebra''. No such theory is needed: the invariant to
be realized is a cokernel, cokernels are unimodular-equivalence invariants, and
Lemma~\ref{X:lem:signed} realizes any integer matrix with positive diagonal as the
Bowen--Franks presentation of an ordinary graph algebra. The two dimensions of $Y$ enter
through $\Delta_1$ and $\DeltaH$, not through the algebra.
\end{remark}

\section{The verification battery}\label{X:sec:battery}

All computations are exact except the spectra. Smith forms are computed $p$-adically:
the determinant of the (reduced) matrix by Bareiss elimination, its prime factors by trial
division, and the $p$-primary invariants by elimination over $\Z/p^{k}$ with
valuation-minimal pivots; for singular symmetric $M$ the torsion found is certified
complete by the lattice identity
\begin{equation}\label{X:eq:lattice}
|\Tor\coker_\Z M|\cdot\operatorname{disc}(\ker_\Z M)=\operatorname{pdet}(M),
\end{equation}
where $\operatorname{disc}$ is the Gram determinant of a saturated basis of the kernel
lattice and $\operatorname{pdet}$ the product of the nonzero eigenvalues; the two sides
agree to $10^{-6}$ in every case (this also verifies \eqref{X:eq:lattice} itself, a
matrix-tree theorem for symmetric integer matrices, on twelve instances).

\subsection*{Data}
$Y_{168}$: $n=168$, $|E|=|F|=1176$, $\chi=168$, $(b_0,b_1,b_2)=(1,0,167)$, flag condition
holds;
\begin{gather*}
\Jac_0=(\Z/2)^{31}\oplus(\Z/16)^{6}\oplus(\Z/592)^{6}\oplus(\Z/1184)^{18}\oplus\Z/4736\\
\qquad\oplus(\Z/33152)^{3}\oplus(\Z/430976)^{12}\oplus\Z/3016832\oplus\Z/9050496,\\
\JacH=(\Z/7)^{18}\oplus(\Z/299593)^{13}\oplus\Z/2097151\oplus\Z/14680057\oplus
(\Z/117440456)^{5}\oplus\Z/352321368,
\end{gather*}
where $299593=7\cdot127\cdot337$ and $2097151=2^{21}-1=7^{2}\cdot127\cdot337$;
$|\JacH|\approx6.18\cdot10^{136}$ agrees with $\frac1n\prod_j|P_j(1)|$ to $10^{-14}$;
$\Tr(\Cthree)^{3}=504=n(q-1)^{2}(q+1)$.

$Y_{21}$: $(b_0,b_1,b_2)=(1,0,20)$. $\Jac_0=\Z/2\oplus(\Z/14)^{15}\oplus\Z/98\oplus\Z/294$;
$\Jac_1=(\Z/2)^{10}\oplus(\Z/14)^{10}\oplus(\Z/98)^{26}\oplus\Z/2058\oplus\Z/18522$
($|\Jac_1|\approx6.7\cdot10^{73}$); $\Jac_2=(\Z/2)^{13}\oplus(\Z/14)^{15}\oplus\Z/98\oplus\Z/294$;
$K(Y)=(\Z/7)^{6}\oplus\Z/14\oplus(\Z/98)^{10}\oplus\Z/294$;
$\JacH=(\Z/7)^{17}\oplus\Z/147$, $|\JacH|=3\cdot7^{19}$; $\Tr(\Cthree)^{3}=651=63+588$.

$Y_{24}$: $(1,0,23)$. $\Jac_0=(\Z/13)^{7}\oplus\Z/26\oplus(\Z/208)^{4}\oplus\Z/1456\oplus\Z/4368$;
$\Jac_1=(\Z/2)^{10}\oplus(\Z/51376)^{2}\oplus(\Z/102752)^{6}\oplus(\Z/2877056)^{4}\oplus
\Z/60418176\oplus\Z/543763584$; $\Jac_2=\Z/2\oplus(\Z/26)^{6}\oplus(\Z/52)^{2}\oplus(\Z/208)^{2}
\oplus(\Z/1456)^{3}\oplus\Z/4368$; $K(Y)=(\Z/2)^{3}\oplus(\Z/38)^{7}\oplus\Z/152\oplus(\Z/304)^{2}
\oplus(\Z/2128)^{3}\oplus\Z/6384$; $\JacH=(\Z/7)^{2}\oplus\Z/49\oplus(\Z/392)^{5}\oplus\Z/1176$.

$Y_{42}$: $(1,0,41)$. $\Jac_0=(\Z/4)^{2}\oplus(\Z/28)^{10}\oplus\Z/56\oplus(\Z/8288)^{4}\oplus
\Z/754208\oplus\Z/2262624$;
\begin{gather*}
\Jac_1=(\Z/2)^{12}\oplus(\Z/4)^{4}\oplus(\Z/28)^{10}\oplus(\Z/196)^{14}\oplus(\Z/784)^{6}
\oplus(\Z/1568)^{2}\\
\oplus(\Z/21865186112)^{4}\oplus\Z/N_1\oplus\Z/N_2 ,
\end{gather*}
$N_1=1474391364718272$, $N_2=13269522282464448=9N_1$;
$\Jac_2=(\Z/2)^{6}\oplus(\Z/4)^{6}\oplus(\Z/28)^{10}\oplus\Z/56\oplus
(\Z/8288)^{4}\oplus\Z/754208\oplus\Z/2262624$; $K(Y)=\Z/7\oplus(\Z/14)^{6}\oplus(\Z/98)^{5}\oplus
\Z/1078\oplus(\Z/2156)^{3}\oplus\Z/40964\oplus\Z/245784$;
$\JacH=(\Z/7)^{18}\oplus\Z/299593\oplus\Z/2097151\oplus\Z/176160684$.

$Y_{84}$ (vertex level): $\Jac_0=(\Z/2)^{12}\oplus(\Z/4)^{7}\oplus(\Z/16)^{2}\oplus\Z/1184\oplus
(\Z/8288)^{5}\oplus(\Z/16576)^{6}\oplus(\Z/215488)^{4}\oplus\Z/1508416\oplus\Z/4525248$.

Thin torus $m=6$ ($n=36$, $|E|=108$, $|F|=72$, $\chi=0$, $(1,2,1)$, flag holds):
$\Jac_0=\Z/2\oplus(\Z/8)^{2}\oplus(\Z/16)^{6}\oplus\Z/80\oplus(\Z/240)^{3}\oplus\Z/720\oplus\Z/2160$;
$\coker\Delta_1=\Z^{2}\oplus(\Z/2)^{12}\oplus\Z/8\oplus(\Z/16)^{3}\oplus(\Z/640)^{3}\oplus(\Z/1920)^{5}
\oplus(\Z/5760)^{2}\oplus\Z/17280\oplus\Z/51840$; $K(Y)=\Z^{2}\oplus\Z/72$;
$\Tr(\Cthree)^{3}=0$. The torus $m=3$ violates the flag condition ($\Tr A_1^{3}=81>54$) and
has $\Tr(\Cthree)^{3}=27$.

In every thick example the odd parts of $\Jac_0$ and $\Jac_2$ nearly coincide while the
$2$-parts differ, $K(Y)$ is much smaller than $\Jac_1$, and $\JacH$ shares no visible
structure with any Hodge group; we record this without interpretation.

\begin{table}[ht]
\centering\scriptsize
\renewcommand{\arraystretch}{1.2}
\resizebox{\textwidth}{!}{%
\begin{tabular}{lcl}
\toprule
check & complexes & result\\
\midrule
(P) presentation axioms; simplicial; Heawood/hexagon links; (ii),(iii); flag (iv) & all & exact; (iv) only for $Y_{168}$, $T_6$\\
(D) $\partial_1\partial_2=0$; Betti numbers; $\chi$; $b_2=\chi-1+b_1$ & all & exact\\
(H) Smith forms of $\Delta_0,\Delta_1,\Delta_2$; $K(Y)$; completeness via \eqref{X:eq:lattice} & thin, $Y_{21},Y_{24},Y_{42}$ & exact, complete\\
(L) $\DeltaH$: two-sided kernel; equal cofactors; $|\JacH|=\frac1n\prod|P_j(1)|$ & thick & exact; numeric to $10^{-14}$\\
(C) $U(I-\Cthree)V=P(1)\oplus I\oplus I$, $P(1)=-\DeltaH$ & all & exact\\
(T) $\Tr(\Cthree)^{k}$, $k\le6$, directly and by Newton; $=0$ for $3\nmid k\le30$; $\Tr(\Cthree)^{3}$ formula & all & exact\\
(R) nine trivial eigenvalues; nontrivial $|\lambda|=q$ (Ramanujan); $\Lambda_Q$ net traces $m\le10$; $m_0$ & thick & numeric / exact\\
(S) signed shadows of $\DeltaH$ and $\Delta_1$: identity, $B\ge0$, irreducible, double loops; Smith check & all & exact\\
\bottomrule
\end{tabular}}
\medskip
\caption{Exact verification (\texttt{pgl3\_building.py}). Sizes: shadow of $\DeltaH$ has
$3n+|E|$ vertices ($1512$ for $Y_{168}$); shadow of $\Delta_1$ has $2058$ vertices for
$Y_{21}$ and $2352$ for $Y_{24}$. Net traces of $\Lambda_Q$ for $Y_{168}$:
$168,\,3920,\,257544,\,16751728,\,\dots$; $m_0=3$.}
\label{X:tab:battery}
\end{table}

\section{Scope, residue, corrections}\label{X:sec:scope}

\begin{remark}[Solved, open, impossible]\label{X:rem:scope}
\emph{Solved:} the Hodge and Hecke Laplacians of $\Atwo$ complexes are set up and their
critical groups computed exactly on thin and thick examples; the Hecke critical group is
identified with the product of inverse $\mathrm{GL}_3$ $L$-values and the companion is
reduced to it unimodularly; the trace question is settled, negatively for the obstruction
and positively for spectral realizability on the computed complexes; the signed-shadow
lemma gives Kirchberg realizations of $\Delta_1$ and $\DeltaH$. \emph{Open:} (a) a non-negative
matrix shift equivalent over $\Z$ to $\Cthree$ (Remark~\ref{X:rem:realization}); (b) a
general theorem that all net traces of the reduced pencil are non-negative for every $\Atwo$
complex of order $q\ge2$ (we have it for five complexes, and in general for the Ramanujan
case once $m\ge m_0$); (c) a relation
between $\Jac_1(Y)$, $K(Y)$ and the Hecke data --- none is visible in the computed examples,
and we conjecture none exists: the edge critical group is combinatorial, the vertex Hecke
group arithmetic; (d) towers of $\Atwo$ complexes (the natural continuation of Chs.~\ref{chap:III}--\ref{chap:IX}).
\emph{Impossible or false as stated:} a cubic trace obstruction (Theorem~\ref{X:thm:notrace}).
\end{remark}

\begin{remark}[Corrections to the task specification]\label{X:rem:corrections}
(1) The predicted ``higher cubic trace obstruction'' does not exist:
$\Tr(\Cthree)^{k}=0$ for $3\nmid k$ and $\Tr(\Cthree)^{3}=n(q-1)^{2}(q+1)\ge0$; the paper
proves the contrary and, for the computed complexes, spectral realizability. (2) The
``duality bridge between the cellular sandpile group and the $\mathrm{GL}_3$ local
$L$-factors'' runs through the Hecke--Laplacian $\DeltaH$ and its critical group $\JacH(Y)$, not
through the Hodge groups $\Jac_0,\Jac_1$, which in rank two are different objects
(Theorem~\ref{X:thm:hecke}(iii)). (3) The test complexes are $\Atwo$ complexes of order $2$
built from a Cartwright--Steger-type triangle presentation, i.e.\ quotients of an $\Atwo$
building of order $2$; the specification's ``$\PGL_3(\Q_p)$'' is not literally available
(arithmetic lattices in $\PGL_3(\Q_p)$ were not constructed), and nothing proved depends on
the field. (4) No ``two-dimensional shadow complex'' is needed
(Remark~\ref{X:rem:digraph}). (5) The companion in the specification is correct, and its
reduction to $\DeltaH$ (Theorem~\ref{X:thm:companion}) is the rank-two form of [Ch.~\ref{chap:II}, \S1.1].
\end{remark}

\section*{Acknowledgement of provenance and caveats}

Unrefereed preprint. Every numbered statement is proved from finite linear algebra and the
results of Papers I--IX; the classical inputs are \cite{CMSZ} (existence of the building for
a triangle presentation), the Cuntz--Krieger facts quoted in [Ch.~\ref{chap:VI}], and \cite{KOR} for
Theorem~\ref{X:thm:realizable} only. Spectral statements marked ``numeric'' in
Table~\ref{X:tab:battery} are floating-point computations with the stated tolerances; all
group-theoretic and trace statements are exact. The identification of the thick complexes
with quotients of the Bruhat--Tits building of $\PGL_3$ over a specific local field is not
claimed. Bibliographic data of \cite{CMSZ,KL,KOR,DKM} were checked against publisher or
indexing pages.

% generated from Paper 11/anccft11.tex
\chapter{Algorithmic non-commutative class field theory, XI: the $\PGL_3$ positive Hecke-dynamical system: shift equivalence over $\Z$ and the geodesic
edge flow}\label{chap:XI}
\chaptermark{XI}
\begin{chapterabstract}
Paper X left one problem at the top of its list: a non-negative integer matrix shift
equivalent over $\Z$ to the companion $\Cthree$ of the $\PGL_3$ cubic pencil, which would be
a positive Hecke-dynamical system carrying the Hecke critical group $\Z\oplus\JacH(Y)$ in its
$K$-theory. This paper settles what can be settled about that problem and builds the
positive system that does exist. First, the $\Z/3$ grading $\deg(i,v)=i-\tau(v)$ makes
$\Cthree$ a period-three cyclic matrix; shift equivalence of $\Cthree$ to a non-negative
matrix reduces, with an explicit lag formula, to shift equivalence of the return map
$P_0=\Cthree{}^{3}|_{\deg 0}$ (Perron root $q^{6}$, $\det P_0=q^{3n}$) to a primitive matrix.
Second, the required statement is identified precisely: since the nonzero spectrum of $P_0$
satisfies the Perron and net-trace conditions ([Ch.~\ref{chap:X}], Kim--Ormes--Roush), its shift
equivalence over $\Z$ to a primitive matrix is an instance of the \emph{weak Generalized
Spectral Conjecture} of Boyle--Handelman, which is open over $\Z$; the specification's
request for explicit $R,S,\ell$ is a request to prove a case of an open conjecture, and we
do not claim to. Third, the natural explicit route --- a positive presentation
$(tI-Y)(tI-P_0)$ with a nilpotent cofactor $Y$ --- is proved impossible on $Y_{21}$ and
$Y_{24}$ by exact rational Farkas certificates: $P_0$ has a positive right but not left
Perron vector, no positive power, and a negative-entry digraph with cycles. Fourth, and
positively: the \emph{geodesic edge flow} $\LE$ on typed edges ($(x,y)\to(y,z)$ with $z$ not
adjacent to $x$) is a non-negative integer matrix with row sums $q^{2}$, irreducible, and for
every joint eigenvector $f$ of $(A_1,A_2)$ with eigenvalues $(a,b)$ the three edge functions
$f\circ\text{tail}$, $f\circ\text{head}$, $m_f$ (sum over the third vertices of the chambers on
the edge) span an $\LE$-invariant space on which the characteristic polynomial is exactly
$\lambda^{3}-a\lambda^{2}+qb\lambda-q^{3}$. Hence an explicit integer intertwiner $\Psi$ with
$\LE\Psi=\Psi C'$, $C'$ lag-one shift equivalent to $\Cthree$, and the whole nontrivial
$\PGL_3$ Hecke spectrum inside the spectrum of a non-negative matrix; the trivial pair
contributes only $q^{2}$ ($\rk\Psi=3n-6$), and the complementary spectrum has modulus
$q^{1/2}$ (numeric on $Y_{168}$). The Cuntz--Krieger algebra $\cO_{\LE}$ is a Kirchberg
algebra whose gauge dynamics is the gallery flow and whose $K_0$ is finite: the positive
system sees the cuspidal Hecke spectrum but not the Hecke critical group, the rank-two form
of the bulk--boundary dichotomy of [Ch.~\ref{chap:I}].
\end{chapterabstract}

\section*{Status of the results}

Every numbered statement is proved except where marked ``numeric'' (floating-point spectra,
Table~\ref{XI:tab:battery}). Theorem~\ref{XI:thm:nocofactor} rests on exact rational Farkas
certificates verified by integer arithmetic. The classical inputs are the module
characterization of shift equivalence over $\Z$ \cite[\S2]{BS}, the Spectral Conjecture over
$\Z$ \cite{KOR}, and the statements (not proofs) of the Generalized Spectral Conjecture
\cite{BH91,BH93,BS}. Kang--Li \cite{KL} is context for the edge flow. The complexes are those
of [Ch.~\ref{chap:X}]; the edge-flow theorem uses the flag condition [Ch.~\ref{chap:X}, Def.~1.1(iv)], which holds for
$Y_{168}$ and fails for the smaller covers, so its numerical checks are on $Y_{168}$.
Section~\ref{XI:sec:scope} records what is open and the corrections to the specification;
Proposition~\ref{XI:prop:kms} answers its item on KMS states.

\section{The grading, the return map, and the cyclic reduction}\label{XI:sec:grading}

Throughout, $Y$ is an $\Atwo$ complex of order $q$ in the sense of [Ch.~\ref{chap:X}, Def.~1.1] with $n$
vertices, type function $\tau$, type-raising adjacency $A_1$, $A_2=A_1^{t}$, and companion
$\Cthree=\begin{psmallmatrix}A_1&-qA_2&q^{3}I\\ I&0&0\\0&I&0\end{psmallmatrix}$ on
$\Z^{3n}=\bigoplus_{i=0}^{2}\Z^{V}$, $\det(I-u\Cthree)=\det P(u)$ with
$P(u)=I-uA_1+qu^{2}A_2-q^{3}u^{3}I$ [Ch.~\ref{chap:X}, Thm.~3.3].

\begin{lemma}[Grading]\label{XI:lem:grading}
Give the coordinate $(i,v)$ of $\Z^{3n}$ the degree $\deg(i,v):=i-\tau(v)\in\Z/3$. Then
$\Cthree$ maps the degree-$d$ summand $\Z^{3n}_d$ into $\Z^{3n}_{d+1}$: in the graded
decomposition $\Cthree$ is the cyclic block matrix with blocks $X_d\colon\Z^{3n}_d\to
\Z^{3n}_{d+1}$, each $\Z^{3n}_d$ of rank $n$, and $\Cthree{}^{3}$ is block diagonal with
\emph{return maps} $P_d:=X_{d+2}X_{d+1}X_d\in M_n(\Z)$.
\end{lemma}

\begin{proof}
$A_1$ sends $(0,y)$ to $(0,x)$ with $\tau(x)=\tau(y)-1$, raising the degree by one; the
identity blocks send $(0,v)\mapsto(1,v)$ and $(1,v)\mapsto(2,v)$, raising $i$ by one;
$-qA_2$ sends $(1,v)$ to $(0,x)$ with $\tau(x)=\tau(v)+1$: degree $1-\tau(v)\mapsto
-\tau(v)-1\equiv1-\tau(v)+1$; $q^{3}I$ sends $(2,v)\mapsto(0,v)$: $2-\tau(v)\mapsto-\tau(v)
\equiv3-\tau(v)$. Each $\Z^{3n}_d$ has rank $n$ because the $n$ vertices split equally
among the three types [Ch.~\ref{chap:X}, Lem.~1.2]. Verified as an exact block identity on four complexes
(Table~\ref{XI:tab:battery}, line (G)).
\end{proof}

\begin{proposition}[The return map]\label{XI:prop:P0}
The nonzero spectrum of $P_0$ is $\Lambda_Q$ of {\rm[Ch.~\ref{chap:X}, Def.~4.4]}: one cube $\lambda^{3}$
for each $\omega$-orbit of eigenvalues of $\Cthree$, with $q^{6}$ simple and dominant;
$\det P_0=q^{3n}$; $P_0$ has a positive right Perron vector, namely the restriction of
$(q^{4}\mathbf1,q^{2}\mathbf1,\mathbf1)$ to $\Z^{3n}_0$. On $Y_{21},Y_{24},Y_{42},Y_{168}$
(numeric): the left Perron vector of $P_0$ is not positive, no power $P_0^{k}$, $k\le30$, is
positive, and the digraph of negative entries of $P_0$ contains cycles.
\end{proposition}

\begin{proof}
$\Cthree{}^{3}=\bigoplus_dP_d$ and the three $P_d$ are conjugate over $\Z$ ($P_{d+1}X_d=X_dP_d$
with $X_d$ invertible over $\Q$), so each has the spectrum of $\Cthree{}^{3}$ with
multiplicities divided by three: $\{\lambda^{3}\}$, one per orbit. $\det\Cthree=(q^{3})^{n}$
(block expansion), so $\det P_0^{3}=\det\Cthree{}^{3}=q^{9n}$ and $\det P_0=q^{3n}$
(positive, as $q^{6}$ dominates). For $\lambda=q^{2}$ the vector $(\lambda^{2}\mathbf1,\lambda
\mathbf1,\mathbf1)$ is a $\Cthree$-eigenvector because $\lambda^{2}(q^{2}+q+1)-q\lambda(q^{2}
+q+1)+q^{3}=\lambda^{3}$; it is positive and graded-homogeneous pieces of it are
eigenvectors of the $P_d$. The remaining statements are computations.
\end{proof}

\begin{lemma}[Cyclic lifting of shift equivalence]\label{XI:lem:lift}
Let $B_0\in M_N(\Z)$ and suppose $(R_0,S_0,\ell)$ is a shift equivalence over $\Z$ between
$P_0$ and $B_0$: $P_0R_0=R_0B_0$, $S_0P_0=B_0S_0$, $R_0S_0=P_0^{\ell}$, $S_0R_0=B_0^{\ell}$.
Put $B:=\operatorname{cyc}(I,I,B_0)$, the cyclic block matrix with blocks $I\colon0\to1$,
$I\colon1\to2$, $B_0\colon2\to0$ on $\Z^{N}\oplus\Z^{N}\oplus\Z^{N}$, and
\[
R:=\operatorname{diag}\bigl(R_0,\ X_0R_0,\ X_1X_0R_0\bigr),\qquad
S:=\operatorname{diag}\bigl(S_2X_1X_0,\ S_2X_1,\ S_2\bigr),\quad S_2:=S_0X_2 .
\]
Then $\Cthree R=RB$, $S\Cthree=BS$, $RS=\Cthree{}^{3(\ell+1)}$, $SR=B^{3(\ell+1)}$: a shift
equivalence over $\Z$ of lag $3(\ell+1)$. If $B_0\ge0$ then $B\ge0$, and $B$ is irreducible
of period $3$ when $B_0$ is primitive. Conversely a shift equivalence between $\Cthree$ and
a graded cyclic $B$ restricts to one between the return maps.
\end{lemma}

\begin{proof}
Direct verification: $X_0R_0=(X_0R_0)\cdot I$, $X_1(X_0R_0)=(X_1X_0R_0)\cdot I$,
$X_2(X_1X_0R_0)=P_0R_0=R_0B_0$ give $\Cthree R=RB$; $S_2X_1X_0\cdot X_2=S_0P_0X_2=B_0S_0X_2=
B_0S_2$ and the two identities $S_{d+1}X_d=S_d$ give $S\Cthree=BS$; on the degree-$d$ block
$RS$ equals $R_dS_d=(X_{d-1}\cdots X_0R_0)(S_0X_2\cdots X_d)=X_{d-1}\cdots X_0P_0^{\ell}X_2
\cdots X_d=P_d^{\ell+1}$, and $SR$ equals $S_0P_0R_0=B_0^{\ell+1}$ on each block, while
$B^{3}=\operatorname{diag}(B_0,B_0,B_0)$. Checked with the trivial equivalence
$B_0=P_0$ on four complexes (Table~\ref{XI:tab:battery}, line (L)). The converse is the
restriction of the four equations to graded pieces.
\end{proof}

\section{Shift equivalence over \texorpdfstring{$\Z$}{Z}: criterion, cofactors, obstructions}\label{XI:sec:se}

\begin{proposition}[Module criterion]\label{XI:prop:module}
For square integer matrices $A,B$ (any sizes), $A\SE B$ if and only if the $\Z[t,t^{-1}]$-modules
$M_A:=\Z[t^{\pm}]^{n}/(tI-A)\Z[t^{\pm}]^{n}$ and $M_B$ are isomorphic \cite[\S2]{BS}. For
nonsingular $A$, $M_A\cong\Z^{n}[A^{-1}]=\bigcup_kA^{-k}\Z^{n}$ with $t$ acting as $A$. A
non-negative $B$ with $B\SE P_0$ is therefore the same thing as a finite set of generators
$g_1,\dots,g_N$ of $M_{P_0}$ over $\Z[t^{\pm}]$ with $t\,g_j=\sum_iB_{ij}g_i$, $B_{ij}\in
\Z_{\ge0}$: a \emph{positive presentation} of the Hecke module.
\end{proposition}

\begin{proposition}[Cofactor construction]\label{XI:prop:cofactor}
Let $U(t)\in GL_n(\Z[t^{\pm}])$ and suppose $F(t):=U(t)(tI-P_0)=t^{K+1}I-\sum_{k=0}^{K}C_kt^{k}$
with all $C_k\in M_n(\Z_{\ge0})$. Then the companion
$B_F=\begin{psmallmatrix}C_K&C_{K-1}&\cdots&C_0\\ I&0&\cdots&0\\ &\ddots&&\\ 0&\cdots&I&0
\end{psmallmatrix}\ge0$ satisfies $B_F\SE P_0$. In particular, if $Y\in M_n(\Z)$ is
nilpotent with $P_0+Y\ge0$ and $-YP_0\ge0$, then $U=tI-Y$ works and
$B=\begin{psmallmatrix}P_0+Y&-YP_0\\ I&0\end{psmallmatrix}\ge0$ is shift equivalent over
$\Z$ to $P_0$, hence $\operatorname{cyc}(I,I,B)\SE\Cthree$ by Lemma~\ref{XI:lem:lift}.
\end{proposition}

\begin{proof}
$\coker(tI-B_F)\cong\coker F$ by the companion reduction (as in [Ch.~\ref{chap:X}, Thm.~3.3]), and
$\coker F=\coker\bigl(U(tI-P_0)\bigr)\cong\coker(tI-P_0)$ because $U$ is invertible over
$\Z[t^{\pm}]$; Proposition~\ref{XI:prop:module}. For $U=tI-Y$: $\det(tI-Y)=t^{n}$ is a unit in
$\Z[t^{\pm}]$ when $Y$ is nilpotent, and $(tI-Y)(tI-P_0)=t^{2}I-(P_0+Y)t+YP_0$.
\end{proof}

\begin{theorem}[No nilpotent cofactor on $Y_{21}$, $Y_{24}$]\label{XI:thm:nocofactor}
For $Y=Y_{21}$ and $Y=Y_{24}$ there is no matrix $Y\in M_n(\Q)$ with $P_0+Y\ge0$,
$YP_0\le0$ and $\Tr Y=0$; a fortiori no nilpotent integer cofactor exists, and
Proposition~\ref{XI:prop:cofactor} with $K=1$ cannot produce a positive presentation.
\end{theorem}

\begin{proof}
Write $Y=-P_0+Z$, $Z\ge0$; the system becomes $ZP_0\le P_0^{2}$ entrywise and $\Tr Z=\Tr P_0$.
By Farkas' lemma it is infeasible iff there exist $\Lambda\in M_n(\Q_{\ge0})$ and $\mu\in\Q$
with $\Lambda P_0^{t}+\mu I\ge0$ entrywise and $\langle\Lambda,P_0^{2}\rangle+\mu\Tr P_0<0$.
The script finds such certificates with denominators $2$ ($Y_{21}$, dual value $-4340$) and
$128$ ($Y_{24}$, dual value $-53185/8$) and verifies both inequalities in exact rational
arithmetic. (For $Y_{42}$ the LP relaxation is infeasible in floating point but a rational
certificate was not extracted; we claim nothing there.)
\end{proof}

\begin{remark}[The precise status of the realization problem]\label{XI:rem:status}
Let $\Lambda$ be the nonzero spectrum of $P_0$. By [Ch.~\ref{chap:X}, Thm.~4.5] and \cite{KOR}, $\Lambda$ is
the nonzero spectrum of a primitive integer matrix. The statement the specification asks
for --- $P_0\SE B_0$ with $B_0$ primitive, equivalently $\Cthree\SE B$ with $B\ge0$
(Lemma~\ref{XI:lem:lift}) --- is exactly the \emph{weak form of the Generalized Spectral
Conjecture} \cite[p.~253]{BH91}, \cite[p.~124]{BH93} for the ring $\Z$ and the matrix $P_0$:
that the three spectral conditions suffice for shift equivalence over $\Z$ to a primitive
matrix. Boyle--Schmieding \cite{BS} show that over dense subrings of $\R$ the weak form
implies the strong form (strong shift equivalence); $\Z$ is not dense and the weak form over
$\Z$ is, to our knowledge, open. The requested explicit $(B,R,S,\ell)$ would therefore prove
a case of an open conjecture. This paper does not. What it proves is that the problem is
equivalent to the primitive-matrix problem for the $n\times n$ return map
(Lemma~\ref{XI:lem:lift}), that the first explicit route is closed
(Theorem~\ref{XI:thm:nocofactor}), and that a positive system carrying the Hecke spectrum ---
though not the Hecke module --- exists and is canonical (Section~\ref{XI:sec:edge}).
\end{remark}

\section{The geodesic edge flow}\label{XI:sec:edge}

Let $E$ be the set of typed edges $e=(x\to y)$, $\tau(y)=\tau(x)+1$, and for each
chamber $\{x,y,z\}$ on $e$ call $z$ its third vertex. Let $S,T,M\in M_{n\times|E|}(\Z)$ be the
tail, head and third-vertex incidence matrices: $S_{v,e}=[v=\text{tail}(e)]$,
$T_{v,e}=[v=\text{head}(e)]$, $M_{v,e}=[\{v\}\cup e\text{ is a chamber}]$; $M$ has column
sums $q+1$. For a vertex function $f$ put $f\circ t:=S^{t}f$, $f\circ h:=T^{t}f$,
$m_f:=M^{t}f$ (the sum of $f$ over the third vertices of the chambers on the edge).

\begin{definition}[Geodesic edge flow]\label{XI:def:LE}
$\LE\in M_{|E|}(\{0,1\})$, $\LE\bigl((x\to y),(y\to z)\bigr)=1$ iff $\tau(z)=\tau(y)+1$ and
$z$ is not adjacent to $x$; all other entries $0$. A step of $\LE$ extends a type-raising
edge to a type-raising path of length two that does not lie in a chamber: the rank-two
analogue of the non-backtracking condition of the Hashimoto matrix $\Anb$ of [Ch.~\ref{chap:I}].
\end{definition}

\begin{lemma}[Local counts]\label{XI:lem:local}
Let $Y$ satisfy {\rm[Ch.~\ref{chap:X}, Def.~1.1(i)--(iv)]} and have vertex links isomorphic to the
incidence graph of $\PG(2,q)$. For a typed edge $e=(x\to y)$:
\begin{enumerate}
\item[(a)] $y$ has $q^{2}+q+1$ type-raising neighbours $z$; exactly $q+1$ of them are
adjacent to $x$, and those are exactly the third vertices of the chambers on $e$ (flag
condition); hence the row sum of $\LE$ at $e$ is $q^{2}$;
\item[(b)] $\sum_{z\colon y\to z,\ z\not\sim x}f(z)=(A_1f)(y)-m_f(e)$;
\item[(c)] $\sum_{z\colon y\to z,\ z\not\sim x}m_f(y\to z)=q\bigl((A_2f)(y)-f(x)\bigr)$.
\end{enumerate}
\end{lemma}

\begin{proof}
(a),(b) are the flag condition. (c) Expand $m_f(y\to z)=\sum_wf(w)$ over third vertices
$w$ of chambers $\{y,z,w\}$; such $w$ are type-$\tau(x)$ neighbours of $y$, i.e.\ the
in-neighbours $w\to y$ counted by $A_2$. For fixed $w$ the number of $z$ with $y\to z$,
$\{y,z,w\}$ a chamber and $z\not\sim x$ is: $0$ if $w=x$ (every chamber on $(x\to y)$ has
$z\sim x$); and if $w\ne x$, in the link of $y$ --- the incidence graph of $\PG(2,q)$, with
$x,w$ two distinct points and the $z$'s the $q+1$ lines through $w$ --- the lines through $w$
that also pass through $x$ number exactly one, so $q$ of them have $z\not\sim x$. Summing,
$q\sum_{w\to y,\,w\ne x}f(w)=q\bigl((A_2f)(y)-f(x)\bigr)$.
\end{proof}

\begin{theorem}[The edge flow carries the Hecke pencil]\label{XI:thm:edge}
Under the hypotheses of Lemma~\ref{XI:lem:local}, with
$\Psi:=\bigl[\,S^{t}\ \ T^{t}\ \ M^{t}\,\bigr]\in M_{|E|\times3n}(\Z)$ and
\[
C':=\begin{pmatrix}0&0&-qI\\ q^{2}I&A_1&qA_2\\ 0&-I&0\end{pmatrix}\in M_{3n}(\Z):
\qquad
\LE\,\Psi\;=\;\Psi\,C' .
\]
Consequently, for every joint eigenvector $f$ of $(A_1,A_2)$ with eigenvalues $(a,b)$, the
space $V_f:=\operatorname{span}\{f\circ t,\ f\circ h,\ m_f\}$ is $\LE$-invariant and the matrix
of $\LE$ on the ordered basis $(f\circ t,f\circ h,m_f)$ is
$\begin{psmallmatrix}0&0&-q\\ q^{2}&a&qb\\ 0&-1&0\end{psmallmatrix}$, with characteristic
polynomial $\lambda^{3}-a\lambda^{2}+qb\lambda-q^{3}$, the cubic pencil of the pair $(a,b)$.
Every root of the
cubic pencil attached to a joint eigenpair with $\dim V_f=3$ is an eigenvalue of the
non-negative matrix $\LE$.
\end{theorem}

\begin{proof}
For $g=\alpha\,f\circ t+\beta\,f\circ h+\gamma\,m_f$ and $e=(x\to y)$,
\begin{align*}
(\LE g)(e)&=\sum_{z\colon y\to z,\ z\not\sim x}\bigl(\alpha f(y)+\beta f(z)+\gamma m_f(y\to z)\bigr)\\
&=\alpha q^{2}f(y)+\beta\bigl((A_1f)(y)-m_f(e)\bigr)+\gamma q\bigl((A_2f)(y)-f(x)\bigr)
\end{align*}
by Lemma~\ref{XI:lem:local}, i.e.\ $\LE g=(-q\gamma)\,f\circ t+\bigl(q^{2}\alpha f+\beta A_1f
+q\gamma A_2f\bigr)\circ h+(-\beta)\,m_f$. Read for arbitrary (not necessarily eigen-)
triples $(f_0,f_1,f_2)$ this is $\LE\Psi(f_0,f_1,f_2)^{t}=\Psi(-qf_2,\ q^{2}f_0+A_1f_1+qA_2f_2,
\ -f_1)^{t}=\Psi C'(f_0,f_1,f_2)^{t}$. For an eigenvector, $A_1f=af$, $A_2f=bf$ give the
$3\times3$ matrix; expanding its characteristic determinant along the first row gives
$\lambda\bigl((\lambda-a)\lambda+qb\bigr)-q^{3}$. Verified as an exact integer identity
$\LE\Psi=\Psi C'$ on $Y_{168}$ (Table~\ref{XI:tab:battery}, line (I)).
\end{proof}

\begin{proposition}[$C'$ is the companion]\label{XI:prop:Cprime}
With $D:=\begin{psmallmatrix}0&0&qI\\ I&0&0\\0&-I&0\end{psmallmatrix}$, both $D$ and
$S:=\Cthree D^{-1}$ are integer matrices, and $C'D=D\Cthree$, $S C'=\Cthree S$, $DS=C'$,
$SD=\Cthree$: $C'$ and $\Cthree$ are shift equivalent over $\Z$ with lag one (indeed
elementary strong shift equivalent over $\Z$), so $M_{C'}\cong M_{\Cthree}$ and $C'$ has the
spectrum of the pencil.
\end{proposition}

\begin{proof}
$D^{-1}=\begin{psmallmatrix}0&I&0\\0&0&-I\\ q^{-1}I&0&0\end{psmallmatrix}$, and the third block
column of $\Cthree$ is $(q^{3}I,0,0)^{t}$, so $\Cthree D^{-1}$ has first block column
$(q^{2}I,0,0)^{t}$ and is integral. The four identities are then formal:
$C'=D\Cthree D^{-1}$ by direct multiplication (the new coordinates are $f_0=qg_2$, $f_1=g_0$,
$f_2=-g_1$), $DS=D\Cthree D^{-1}=C'$, $SD=\Cthree$. Verified exactly on $Y_{168}$.
\end{proof}

\begin{corollary}[What the edge flow sees and what it does not]\label{XI:cor:spectrum}
Under the hypotheses of Theorem~\ref{XI:thm:edge}:
\begin{enumerate}
\item[(i)] $\ker\Psi\supseteq\{(f,-f,0),\ ((q+1)f,0,-f):\ f=\omega^{k\tau}\mathbf1,\ k=0,1,2\}$,
a rank-$6$ sublattice: for the trivial pair and its two $\omega$-conjugates the three edge
functions are proportional, $V_f$ is one-dimensional and carries only the eigenvalue
$q^{2}\omega^{k}$; hence $\rk\Psi\le3n-6$, with equality on $Y_{168}$ (numeric rank $498$).
\item[(ii)] On $Y_{168}$ the spectrum of $\LE$ contains every pencil root except the six
$\{1,\omega,\omega^{2},q,q\omega,q\omega^{2}\}$; the remaining $|E|-(3n-6)=4n+6$ eigenvalues
all have modulus $q^{1/2}$ (numeric, to $10^{-4}$): with $P(u)=\det(I-u\Cthree)$,
\begin{gather*}
\det(I-u\LE)\cdot(1-u^{3})(1-q^{3}u^{3})\;=\;\det P(u)\cdot G(u),\\
G\in\Z[u],\qquad \deg G=4n+6,\qquad \text{all roots of $G$ on }|u|=q^{-1/2}.
\end{gather*}
\item[(iii)] $\LE$ is irreducible with row sums $q^{2}$ (its Perron eigenvalue), and $1$ is
not an eigenvalue of $\LE$ (numeric, distance $0.58$ to the nearest eigenvalue on $Y_{168}$).
\end{enumerate}
\end{corollary}

\begin{proof}
(i) For $f=\omega^{k\tau}\mathbf1$: $f\circ t=\omega^{k\tau(x)}$, $f\circ h=\omega^{k(\tau(x)+1)}$,
$m_f=(q+1)\omega^{k(\tau(x)+2)}$ on $(x\to y)$, proportional; the two displayed vectors are
integral for $k=0$, and over $\Z[\omega]$ for $k=1,2$, so the rational kernel has rank at
least $6$. (ii),(iii) are computations; the divisibility in (ii) is exact because
$\det P(u)$ contains $(1-u^{3})(1-q^{3}u^{3})(1-q^{6}u^{3})$ and $\LE$ has the eigenvalues
$q^{2}\omega^{k}$.
\end{proof}

\begin{remark}[Kang--Li context]\label{XI:rem:KL}
Kang--Li \cite{KL} express the zeta function of a $\PGL_3$ complex as a rational function
in which the vertex factor $\det P(u)$ and a chamber (parahoric) factor appear; the operator
$\LE$ is the type-$(1,0)$ gallery transfer operator whose determinant computes the zeta
function counting closed type-raising geodesic galleries, and the modulus-$q^{1/2}$
remainder of Corollary~\ref{XI:cor:spectrum}(ii) is the parahoric spectrum, which for a
Ramanujan complex lies on $|u|=q^{-1/2}$. We use nothing from \cite{KL}; the identity (ii)
is established here numerically and its general form is left as it stands. The chamber
flow $\LF$ on oriented chambers ($(a,b,c)\to(b,c,d)$, $d\ne a$), row sums $q$, has spectrum
of moduli $\{1,q^{1/4},q^{1/2},q\}$ on $Y_{168}$ and contains no nontrivial pencil root: the
Hecke spectrum lives on edges, not on chambers.
\end{remark}

\begin{theorem}[The positive Hecke-dynamical algebra]\label{XI:thm:algebra}
Let $Y$ satisfy the hypotheses of Theorem~\ref{XI:thm:edge} and let $\cO_{\LE}$ be the
Cuntz--Krieger algebra of $\LE$ with its gauge action $\gamma$. Then $\cO_{\LE}$ is a unital
Kirchberg algebra whose gauge dynamics is the geodesic gallery flow; its partition
function is $\det(I-u\LE)^{-1}$, whose poles include every nontrivial Satake datum
$\lambda_{j,i}$ of the complex at $u=\lambda_{j,i}^{-1}$; its $K$-theory is
$K_1(\cO_{\LE})=\ker(I-\LE^{t})=0$ and $K_0(\cO_{\LE})=\coker(I-\LE^{t})$ finite of order
$|\det(I-\LE)|$ (on $Y_{168}$, $\log|\det(I-\LE)|=603.249$, numeric). The integer intertwiner
$\Psi$ induces a $\Z[t^{\pm}]$-module map $M_{\Cthree}\cong M_{C'}\to M_{\LE}$ of the Hecke
module into the edge-flow module. $\cO_{\LE}$ is not shift equivalent to $\Cthree$ (different
nonzero spectra), and its $K_0$ is not the Hecke critical group $\Z\oplus\JacH(Y)$.
\end{theorem}

\begin{proof}
$\LE$ is irreducible with row sums $q^{2}\ge4$ and is not a permutation matrix, so it
satisfies condition (I) of Cuntz--Krieger and $\cO_{\LE}$ is simple and purely infinite,
separable, nuclear and in the UCT class, by the criteria quoted in [Ch.~\ref{chap:VI}, Thm.~3.1]; purely
infinite algebras are traceless. $K$-theory is the
Cuntz--Krieger formula; $1\notin\spec\LE$ gives the kernel statement and finiteness. The
module map is Theorem~\ref{XI:thm:edge} and Proposition~\ref{XI:prop:Cprime}. The last sentence is
Corollary~\ref{XI:cor:spectrum}(ii) and [Ch.~\ref{chap:X}, Thm.~3.3].
\end{proof}

\begin{proposition}[KMS states of the edge flow]\label{XI:prop:kms}
Let $\sigma_t(s_e)=e^{it}s_e$ be the gauge dynamics of $\cO_{\LE}$. A $\KMS_\beta$ state exists
if and only if $\beta=\beta_c:=2\log q=\log\rho(\LE)$, and it is then unique: the state
attached to the Perron eigenvector of $\LE$ (the Parry measure of the gallery shift), which
is the constant vector since the row sums are $q^{2}$. The partition function
$Z(\beta)=\sum_k\Tr(\LE^{k})e^{-k\beta}=-u\,\frac{d}{du}\log\det(I-u\LE)\big|_{u=e^{-\beta}}$
has poles at $u=\lambda^{-1}$ for every eigenvalue $\lambda$ of $\LE$, hence at the inverses
of all nontrivial Satake data by Theorem~\ref{XI:thm:edge}; the $\KMS$ \emph{state} sees only
$\lambda=q^{2}$. The Perron obstruction of {\rm[Ch.~\ref{chap:I}, Thm.~2.5]} therefore holds verbatim in rank
two: the cuspidal $\mathrm{GL}_3$ $L$-factors $P_j(u)^{-1}$ are carried by the partition
function and by the $\KMS$ functionals obtained from the nontrivial eigenvectors of
$\LE$ (Theorem~\ref{XI:thm:edge} gives them explicitly, $\Psi(f\circ t,f\circ h,m_f)$-combinations),
not by any $\KMS$ state.
\end{proposition}

\begin{proof}
For an irreducible non-negative matrix with the gauge dynamics, $\KMS_\beta$ states of the
Cuntz--Krieger algebra correspond to non-negative eigenvectors of $\LE^{t}$ for the
eigenvalue $e^{\beta}$ (as in [Ch.~\ref{chap:I}, Thm.~2.4]); by Perron--Frobenius the only such eigenvalue is
$\rho(\LE)=q^{2}$ (row sums), with a unique normalized positive eigenvector. The rest is
Theorem~\ref{XI:thm:edge} and the identity $\Tr(\LE^{k})=\sum_\lambda\lambda^{k}$.
\end{proof}

\begin{remark}[Bulk and boundary, in rank two]\label{XI:rem:dichotomy}
[Ch.~\ref{chap:I}, Thm.~3.7] found that the natural positive system of a graph --- the boundary algebra
$\cO_{\Anb}$ --- has $K$-theory blind to the Hecke data, while the Hecke critical group
$\Z\oplus\Jac(Y)$ lives in the Laplacian channel; [Ch.~\ref{chap:VI}] realized the latter by a
Hecke-functorial, not Hecke-dynamical, Kirchberg algebra, and [Ch.~\ref{chap:II}] showed no single algebra
can be both. In rank two the same dichotomy appears with one difference: the positive
system $\cO_{\LE}$ carries the whole nontrivial Hecke \emph{spectrum} (Theorem~\ref{XI:thm:edge}),
and its $K_0$ is a finite group unrelated to $\JacH(Y)$, while $\cO(Y,\DeltaH)$ of [Ch.~\ref{chap:X}] carries
$\Z\oplus\JacH(Y)$ without dynamics. An algebra carrying both would be $\cO_B$ for a
non-negative $B\SE\Cthree$; [Ch.~\ref{chap:II}, Thm.~1.6] no longer forbids it (there is no trace
obstruction, [Ch.~\ref{chap:X}, Thm.~4.2]), and its existence is the open instance of the weak Generalized
Spectral Conjecture of Remark~\ref{XI:rem:status}. That is the exact residue of the problem.
\end{remark}

\section{The verification battery}\label{XI:sec:battery}

\begin{table}[ht]
\centering\scriptsize
\renewcommand{\arraystretch}{1.2}
\resizebox{\textwidth}{!}{%
\begin{tabular}{lcl}
\toprule
check & complexes & result\\
\midrule
(G) $\deg=i-\tau$ makes $\Cthree$ degree one; $P_d=\Cthree{}^{3}|_d$; $\spec P_0=\Lambda_Q$; $\det P_0=q^{3n}$ & $Y_{21},Y_{24},Y_{42},Y_{168}$ & exact (spectrum numeric)\\
(Y) $P_0$: negative entries $147,184,476,5936$; right Perron $>0$; left Perron $\not>0$; no $P_0^{k}>0$, $k\le30$; negative digraph cyclic & same & exact / numeric\\
(N) LP relaxation $Y\ge-P_0$, $YP_0\le0$, $\Tr Y=0$ infeasible; exact Farkas certificates (denominators $2$, $128$) & $Y_{21},Y_{24}$ ($Y_{42}$ float) & exact\\
(L) cyclic lift of the trivial shift equivalence: four equations, lag $3$ & four complexes & exact\\
(E) $\LE$: $1176$ edges, row sums $4=q^{2}$, irreducible; $498$ of $504$ pencil roots present; missing $\{1,\omega,\omega^{2},2,2\omega,2\omega^{2}\}$; remainder $678=4n+6$ of modulus $\sqrt2$ & $Y_{168}$ & numeric\\
(I) $\LE\Psi=\Psi C'$; $\rk\Psi=498=3n-6$; $D$, $\Cthree D^{-1}$ integral; lag-one equations & $Y_{168}$ & exact\\
(T) $\Tr\LE^{k}$ vs $\Tr(\Cthree)^{k}$, $k\le10$: unequal (parahoric spectrum) & $Y_{168}$ & exact\\
(F) $\LF$: $3528$ oriented chambers, row sums $2$, irreducible; spectrum moduli $\{1,2^{1/4},2^{1/2},2\}$; $3$ pencil roots & $Y_{168}$ & numeric\\
\bottomrule
\end{tabular}}
\medskip
\caption{Verification (\texttt{pgl3\_dynamical.py}). $Y_{168}$ is the only computed complex
satisfying the flag condition; the edge-flow theorem is checked there. $\log|\det(I-\LE)|
=603.249$ on $Y_{168}$.}
\label{XI:tab:battery}
\end{table}

\section{Scope, residue, corrections}\label{XI:sec:scope}

\begin{remark}[Solved, open, impossible]\label{XI:rem:scope}
\emph{Solved:} the graded structure of $\Cthree$ and the reduction of its realization
problem to the primitive-matrix problem for the $n\times n$ return map, with an explicit lag
formula; the module criterion and the cofactor construction; the impossibility of the
nilpotent-cofactor route on $Y_{21},Y_{24}$ (exact certificates); the geodesic edge flow
theorem with its explicit intertwiner, giving a non-negative integer matrix whose spectrum
contains the entire nontrivial $\PGL_3$ Hecke spectrum, and the Kirchberg algebra
$\cO_{\LE}$ it generates. \emph{Open:} (a) $\Cthree\SE B\ge0$, an instance of the weak
Generalized Spectral Conjecture over $\Z$ (Remark~\ref{XI:rem:status}); (b) a proof of the
determinant identity of Corollary~\ref{XI:cor:spectrum}(ii) for all $\Atwo$ complexes, and the
identification of $G(u)$ with Kang--Li's chamber factor; (c) the structure of the finite
group $K_0(\cO_{\LE})$; (d) whether positive presentations of $M_{P_0}$ of higher degree
$K\ge2$ exist (Proposition~\ref{XI:prop:cofactor}) --- the Farkas method extends but the
search space grows. \emph{Impossible or false as stated:} the ``century deadlock since
Bost--Connes'' that positive correspondences cannot carry higher automorphic Hecke spectra
--- Theorem~\ref{XI:thm:edge} carries it, as Hashimoto's matrix already did in rank one; what
positive systems have so far failed to carry is the Hecke \emph{module} (the $K$-theory),
and whether they can is the open conjecture.
\end{remark}

\begin{remark}[Corrections to the task specification]\label{XI:rem:corrections}
\begin{sloppypar}
(1) The specification asked for explicit $(B,R,S,\ell)$ with $B\ge0$ shift equivalent over
$\Z$ to $\Cthree$, as a corollary of [X]'s ``Boyle--Handelman spectral conditions''. Those conditions give the nonzero spectrum
(Spectral Conjecture, a theorem over $\Z$ by \cite{KOR}); shift equivalence is the weak
Generalized Spectral Conjecture, open. No such $(B,R,S,\ell)$ is supplied, and none could
be without proving a case of that conjecture. (2) ``Method one'' --- a block-cyclic $B$
whose three-step product cancels $-qA_2$ --- is exactly Lemma~\ref{XI:lem:lift}: correct as a
reduction, and it converts the problem into the primitive-matrix problem for $P_0$ rather
than solving it. (3) ``Method two'' (state splitting) applies to non-negative matrices; it
cannot start from $\Cthree$, and the algebraic substitute (cofactors,
Proposition~\ref{XI:prop:cofactor}) is obstructed at degree two (Theorem~\ref{XI:thm:nocofactor}).
(4) The genuine positive $\mathrm{GL}_3$ dynamics is the geodesic edge flow $\LE$, whose
existence and spectral content are proved here; it is the rank-two Hashimoto matrix, not a
realization of $\Cthree$. (5) The claim that KMS states at the critical temperature
``reproduce the $\mathrm{GL}_3$ local $L$-factors'' is false in the same way it was false in
rank one: the unique $\KMS$ state sees the Perron eigenvalue only, the $L$-factors live in
the partition function (Proposition~\ref{XI:prop:kms}); the ``curse'' of [Ch.~\ref{chap:II}] is not broken but
located --- it concerns the Hecke module, not the Hecke spectrum. (6) Of the specification's
four checks, $B\ge0$, irreducibility and the spectral containment hold for $\LE$; the four
shift-equivalence identities hold for the intertwiner pair $(D,\Cthree D^{-1})$ between
$\Cthree$ and $C'$ (lag one) and for $\Psi$ as a one-sided intertwiner into $\LE$. The trace
identities $\Tr\LE^{k}=\Tr(\Cthree)^{k}$ fail, and so does the Smith-form identity
$\coker(I-\LE)\cong\Z\oplus\JacH$; both must fail, because $\LE$ carries the additional
parahoric spectrum (Table~\ref{XI:tab:battery}, line (T)).
\end{sloppypar}
\end{remark}

\section*{Acknowledgement of provenance and caveats}

Unrefereed preprint. Every numbered statement is proved except the items marked numeric in
Table~\ref{XI:tab:battery}; Theorem~\ref{XI:thm:nocofactor} is certified by exact rational
arithmetic. The statement of the weak Generalized Spectral Conjecture and its status are
taken from the introduction of \cite{BS} (consulted), which attributes the weak form to
\cite[p.~253]{BH91} and \cite[p.~124]{BH93}; those two pages were not consulted directly.
The identification of the thick complexes with quotients of the building of $\PGL_3$ over a
specific field is not claimed. Bibliographic data of \cite{BH91,BH93,BS,KOR,KL} were checked
against publisher or arXiv pages.

\input{chapters/XII}
% generated from Paper 15/anccft15.tex
\renewcommand{\LE}{L_E}
\chapter{Algorithmic non-commutative class field theory, XV: the parahoric factorization of the geodesic edge flow, the unitarity of the remainder,
and the chamber companion}\label{chap:XV}
\chaptermark{XV}
\begin{chapterabstract}
[Ch.~\ref{chap:XI}, Cor.~3.5] found numerically, on one $\Atwo$ complex, that the characteristic polynomial of
the geodesic edge flow $\LE$ factors as $\det(I-u\LE)(1-u^{3})(1-q^{3}u^{3})=\det P(u)\,G(u)$,
with $P$ the cubic Hecke pencil and $G$ a polynomial of degree $4n+6$ all of whose roots lie
on $|u|=q^{-1/2}$, and left the proof and the identification of $G$ open. This paper proves
the identity for every $\Atwo$ complex of order $q$, identifies $G$, and proves the
Riemann hypothesis for it unconditionally. The mechanism is a two-sided intertwining: the
image of $\Psi=[S^{t}\,T^{t}\,M^{t}]$ is invariant under $\LE$ \emph{and} $\LE^{t}$, so its
orthogonal complement $W=\ker S\cap\ker T\cap\ker M$ reduces $\LE$, and the local identities
$\LE\LE^{t}=qI+(q^{2}-q)T^{t}T$, $\LE^{t}\LE=qI+(q^{2}-q)S^{t}S$ --- two points of a projective
plane lie on one line --- make $R:=\LE|_W$ satisfy $RR^{t}=R^{t}R=qI$. Hence $R/\sqrt q$ is
orthogonal, $G(u)=\det(I-uR)$, and every root of $G$ has $|u|=q^{-1/2}$ for every complex,
Ramanujan or not; the Ramanujan property of the complex lives entirely in the vertex pencil.
The same calculus on pointed chambers gives the Kang--Li operator $\LB$ a companion
$B$ on $\ell^{2}(E)^{3}$, with $\LB\Phi=\Phi B$ and
$\det(I+uB)=\det\bigl((1+qu^{3})I+uK+u^{2}K^{t}\bigr)
=(1-u^{3})^{n}\det(I-u\LE)\det(I-u^{2}\LE)/\det P(u)$: the Kang--Li zeta identity in companion
form, with the parahoric block $G(u)G(u^{2})$ explaining the eigenvalues of modulus
$q^{1/4}$ of the chamber operator as square roots of the unitary remainder. A correction to
[Ch.~\ref{chap:XI}, Def.~3.1] is recorded (the geodesic condition is ``not in a common chamber''). Verified
exactly on $Y_{21},Y_{24},Y_{42},T_6$ and structurally on $Y_{168}$.
\end{chapterabstract}

\section*{Status of the results}

Every numbered statement is proved; the proofs use only the link structure of an $\Atwo$
complex of order $q$ ([Ch.~\ref{chap:X}, Def.~1.1(i)--(iii)], vertex links the incidence graph of
$\PG(2,q)$), linear algebra over $\Q$, and the results of [Ch.~\ref{chap:XI}], [Ch.~\ref{chap:XII}] that are re-derived where
used. Kang--Li \cite{KL} and Kang--Li--Wang \cite{KLW} are cited for context and for the
one statement we do not reprove (the relation between $\LB$ and its companion for $q\ne2$,
Corollary~\ref{XV:cor:KL}); Lubotzky--Samuels--Vishne \cite{LSV} for the notion of Ramanujan
complex. The flag condition [Ch.~\ref{chap:X}, Def.~1.1(iv)] is not assumed anywhere.

\section{Setting, and a correction}\label{XV:sec:setting}

$Y$ is a finite $\Atwo$ complex of order $q$: a connected pure $2$-dimensional simplicial
complex with a $\Z/3$ vertex typing, $n$ vertices, typed edges $E$ ($e=(x\to y)$,
$\tau(y)=\tau(x)+1$, $|E|=n(q^{2}+q+1)=:nN$), chambers $F$ (one vertex of each type,
$|F|=(q{+}1)|E|/3$), such that the link of every vertex $y$ is the incidence graph of
$\PG(2,q)$: in-neighbours of $y$ are the points, out-neighbours the lines, and $\{x,y,z\}$ is a
chamber iff the point $x$ lies on the line $z$. $A_1,A_2=A_1^{t}$ are the Hecke operators,
$P(u):=I-A_1u+qA_2u^{2}-q^{3}u^{3}$ the pencil, $S,T,M\in M_{n\times|E|}(\Z)$ the tail, head and
third-vertex incidences, $K:=\sum_id_id_{i+1}^{t}$ the chamber continuation
($K(e,e')=\#\{$chambers in which $e'$ follows $e\}$), all as in [Ch.~\ref{chap:XI}, \S3] and [Ch.~\ref{chap:XII}, Thm.~3.1].

\begin{definition}[Geodesic edge flow, chamber form]\label{XV:def:LE}
$\LE\in M_{|E|}(\{0,1\})$, $\LE\bigl((x\to y),(y\to z)\bigr)=1$ iff $(y\to z)\in E$ and
$\{x,y,z\}\notin F$; equivalently $\LE=T^{t}S-K$ [Ch.~\ref{chap:XII}, Thm.~3.1].
\end{definition}

\begin{remark}[Correction to {\rm[Ch.~\ref{chap:XI}, Def.~3.1]}]\label{XV:rem:corrXI}
[Ch.~\ref{chap:XI}, Def.~3.1] used ``$z$ not adjacent to $x$''. The two conditions agree exactly when $Y$ is
a flag complex ([Ch.~\ref{chap:X}, Def.~1.1(iv)]), which holds for $Y_{168}$, the only complex on which [Ch.~\ref{chap:XI}]
evaluated $\LE$. On $Y_{21}$, whose $1$-skeleton is $K_{7,7,7}$, every $z$ is adjacent to $x$
and the adjacency version is the zero matrix, while Definition~\ref{XV:def:LE} gives row sums
$q^{2}$; on $Y_{24}$ and $Y_{42}$ the two versions differ as well (Table~\ref{XV:tab:battery},
line (D)). Definition~\ref{XV:def:LE} is the projection of the geodesic operator of the
building, the operator $L_E$ of \cite{KL}, and every statement of [Ch.~\ref{chap:XI}, \S3] and [Ch.~\ref{chap:XII}] holds
for it with the proofs unchanged; the quoted numerics of [Ch.~\ref{chap:XI}] concern $Y_{168}$ and are
unaffected.
\end{remark}

\begin{lemma}[Local counts]\label{XV:lem:local}
For $e=(x\to y)$, $e'=(x'\to y)$ with the same head, and $f\in\ell^{2}(V)$, $m_f:=M^{t}f$:
\begin{enumerate}
\item[(a)] $e$ has $q^{2}$ geodesic continuations and $q^{2}$ geodesic predecessors;
\item[(b)] $\sum_{z\colon e\to(y\to z)}f(z)=(A_1f)(y)-m_f(e)$,\quad
$\sum_{z}m_f(y\to z)=q\bigl((A_2f)(y)-f(x)\bigr)$;
\item[(c)] $\#\{$common geodesic continuations of $e,e'\}=q^{2}$ if $x=x'$ and $q^{2}-q$ if
$x\ne x'$; dually for two edges with the same tail;
\item[(d)] $\sum_{x\colon(x\to y)\to e''}f(x)=(A_2f)(y)-m_f(e'')$ and
$\sum_{x}m_f(x\to y)=q\bigl((A_1f)(y)-f(z)\bigr)$ for $e''=(y\to z)$.
\end{enumerate}
\end{lemma}

\begin{proof}
In the link of $y$ the continuations of $e$ are the lines $z$ not through the point $x$:
$N-(q{+}1)=q^{2}$. (b) is [Ch.~\ref{chap:XI}, Lem.~3.2(b),(c)] with ``chamber'' in place of ``adjacent'', the
proof verbatim. (c): lines through neither of two distinct points: $N-2(q{+}1)+1=q^{2}-q$.
(d) is (b) for the opposite complex (all edges reversed, types negated), whose $\LE$ is
$\LE^{t}$.
\end{proof}

\section{Two-sided intertwining and the unitary remainder}\label{XV:sec:intertwine}

Let $\Psi:=[S^{t}\ T^{t}\ M^{t}]\colon\ell^{2}(V)^{3}\to\ell^{2}(E)$ and
$W:=(\im\Psi)^{\perp}=\ker S\cap\ker T\cap\ker M$.

\begin{theorem}[Two-sided intertwining]\label{XV:thm:two}
With $I=I_n$ and block matrices over $\ell^{2}(V)^{3}$,
\begin{gather*}
\LE\Psi=\Psi C',\quad C'=\begin{pmatrix}0&0&-qI\\ q^{2}I&A_1&qA_2\\0&-I&0\end{pmatrix};\qquad
\LE^{t}\Psi=\Psi\Cpp,\quad \Cpp=\begin{pmatrix}A_2&q^{2}I&qA_1\\0&0&-qI\\-I&0&0\end{pmatrix};\\
K\Psi=\Psi\Ch,\quad \Ch=\begin{pmatrix}0&0&qI\\(q{+}1)I&0&A_2\\0&I&0\end{pmatrix};\qquad
K^{t}\Psi=\Psi\Chp,\quad \Chp=\begin{pmatrix}0&(q{+}1)I&A_1\\0&0&qI\\I&0&0\end{pmatrix};\\
\LE\LE^{t}=qI_E+(q^{2}-q)\,T^{t}T,\qquad \LE^{t}\LE=qI_E+(q^{2}-q)\,S^{t}S .
\end{gather*}
Consequently $\im\Psi$ and $W$ are invariant under $\LE,\LE^{t},K,K^{t}$, and
$\ell^{2}(E)=\im\Psi\oplus W$ is an orthogonal decomposition reducing all four.
\end{theorem}

\begin{proof}
The first identity is [Ch.~\ref{chap:XI}, Thm.~3.3], whose proof uses only Lemma~\ref{XV:lem:local}(a),(b).
For $\LE^{t}$: $(\LE^{t}S^{t}f)(y\to z)=\sum_{x}f(x)$ over geodesic predecessors, which is
$(A_2f)(y)-m_f(y\to z)$ by (d), i.e.\ $\LE^{t}S^{t}=S^{t}A_2-M^{t}$; $(\LE^{t}T^{t}f)(y\to z)=q^{2}f(y)$
by (a), i.e.\ $\LE^{t}T^{t}=q^{2}S^{t}$; $(\LE^{t}M^{t}f)(y\to z)=\sum_xm_f(x\to y)=q((A_1f)(y)-f(z))$
by (d), i.e.\ $\LE^{t}M^{t}=qS^{t}A_1-qT^{t}$. Reading the three columns gives $\Cpp$. For
$K=T^{t}S-\LE$: $T^{t}S\Psi=T^{t}[SS^{t}\ ST^{t}\ SM^{t}]=T^{t}[NI\ A_1\ (q{+}1)A_2]$ by
[Ch.~\ref{chap:XII}, Thm.~3.1], and subtracting $\LE\Psi=\Psi C'$ gives $KS^{t}=(q{+}1)T^{t}$, $KT^{t}=M^{t}$,
$KM^{t}=T^{t}A_2+qS^{t}$, i.e.\ $\Ch$; $K^{t}=S^{t}T-\LE^{t}$ likewise gives $\Chp$. The two
product identities are Lemma~\ref{XV:lem:local}(c): $(\LE\LE^{t})(e,e')$ counts common
continuations, zero unless the heads agree, $q^{2}$ on the diagonal and $q^{2}-q$ off it, and
$(T^{t}T)(e,e')=[\text{head}(e)=\text{head}(e')]$. Invariance of $\im\Psi$ under an operator
$X$ and of $W=(\im\Psi)^{\perp}$ under $X^{t}$ are the same statement; since all four
operators and their transposes preserve $\im\Psi$, both summands are invariant under all
four. Verified as exact integer identities on all five complexes (Table~\ref{XV:tab:battery},
line (D)).
\end{proof}

\begin{theorem}[The remainder is $\sqrt q$ times an orthogonal matrix]\label{XV:thm:unitary}
Let $R:=\LE|_W$. Then $R^{t}=\LE^{t}|_W$ and
\[
RR^{t}=R^{t}R=q\,I_W .
\]
Hence $R$ is invertible, every eigenvalue of $R$ has modulus $\sqrt q$, and
$G(u):=\det(I-uR)\in\Z[u]$ has all its roots on the circle $|u|=q^{-1/2}$, with
$G(u)=(-u)^{d}\det(R)\,G\bigl(1/(qu)\bigr)$, $d=\dim W$, $\det R=\pm q^{d/2}$.
\end{theorem}

\begin{proof}
$W$ is invariant under $\LE$ and $\LE^{t}$ (Theorem~\ref{XV:thm:two}), so the restriction of
$\LE^{t}$ to $W$ is the adjoint of $R$. For $v\in W$, $Tv=0$ and $Sv=0$, so the product
identities give $\LE\LE^{t}v=qv$ and $\LE^{t}\LE v=qv$. Thus $R/\sqrt q$ is orthogonal, its
eigenvalues have modulus one, and $\det(I-uR)=\det(-uR)\det(I-(uR)^{-1})
=(-u)^{d}\det R\cdot\det(I-\tfrac1{qu}R^{t})$. $G\in\Z[u]$ because $W\cap\Z^{E}$ is a
saturated $\LE$-invariant sublattice of full rank in $W$. On the five complexes the
eigenvalues of $R$ have modulus $\sqrt q$ to $10^{-12}$ (Table~\ref{XV:tab:battery}, line (U)).
\end{proof}

\begin{remark}[What this says about the Riemann hypothesis]\label{XV:rem:RH}
Kang--Li--Wang \cite[Thm.~1.5.1]{KLW} prove that $Y$ is Ramanujan iff the nontrivial zeros of
$\det P(u)$ have $|u|=q^{-1}$, iff the nontrivial zeros of $\det(I-u\LE)$ have $|u|\in\{q^{-1},
q^{-1/2}\}$. Theorem~\ref{XV:thm:unitary} shows that the second clause of the edge criterion is
automatic: the zeros of $\det(I-u\LE)$ off the pencil are always on $|u|=q^{-1/2}$. The
Riemann hypothesis for the zeta function of an $\Atwo$ complex is therefore exactly the
Ramanujan property of the \emph{vertex} spectrum, and the ``parahoric'' zeros are on the
critical line for every complex, Ramanujan or not --- for the same reason that the Hashimoto
remainder of a graph is $\pm1$: it is a local, not a global, statement.
\end{remark}

\section{The factorization theorem}\label{XV:sec:factor}

\begin{theorem}[Parahoric factorization]\label{XV:thm:factor}
For every $\Atwo$ complex of order $q$, with $K_Y(u):=\det(I-uC'|_{\ker\Psi})$,
\begin{gather*}
\det(I-u\LE)\cdot K_Y(u)\;=\;\det P(u)\cdot G(u),\\
G(u)=\det(I-u\LE|_W)\in\Z[u],\qquad \deg G=\dim W=|E|-\rk\Psi ,
\end{gather*}
$K_Y$ divides $\det P$, and all roots of $G$ lie on $|u|=q^{-1/2}$.
\end{theorem}

\begin{proof}
By Theorem~\ref{XV:thm:two}, $\LE=\LE|_{\im\Psi}\oplus R$, so $\det(I-u\LE)=\det(I-u\LE|_{\im\Psi})
\,G(u)$. $\Psi$ induces an isomorphism $\ell^{2}(V)^{3}/\ker\Psi\to\im\Psi$ intertwining the
map induced by $C'$ with $\LE|_{\im\Psi}$, and $\ker\Psi$ is $C'$-invariant ($\Psi C'v=\LE\Psi v=0$);
hence $\det(I-uC')=K_Y(u)\det(I-u\LE|_{\im\Psi})$. Finally $\det(I-uC')=\det P(u)$: the blocks
of $C'$ commute, and the formal determinant of $I-uC'$ expanded along the first block row is
$(I-uA_1+qu^{2}A_2)+qu\cdot(-q^{2}u^{2})=P(u)$ [Ch.~\ref{chap:XI}, Prop.~3.4]. Multiply out. Degree and
roots: Theorem~\ref{XV:thm:unitary}, $R$ invertible.
\end{proof}

\begin{proposition}[The kernel of $\Psi$]\label{XV:prop:kernel}
The Gram matrix
\[
\Psi^{t}\Psi=\begin{pmatrix}NI&A_1&(q{+}1)A_2\\A_2&NI&(q{+}1)A_1\\(q{+}1)A_1&(q{+}1)A_2&qNI+A_1A_2\end{pmatrix}
\]
lies in $M_3$ of the commutative algebra $\Q[A_1,A_2]$; for a joint eigenvector with
$A_1f=\lambda f$, $A_2f=\bar\lambda f$ it is the Hermitian matrix
\begin{gather*}
H(\lambda)=\begin{pmatrix}N&\lambda&(q{+}1)\bar\lambda\\\bar\lambda&N&(q{+}1)\lambda\\
(q{+}1)\lambda&(q{+}1)\bar\lambda&qN+|\lambda|^{2}\end{pmatrix},\\
\det H(\lambda)=(qN+|\lambda|^{2})(N^{2}-|\lambda|^{2})-2N(q{+}1)^{2}|\lambda|^{2}+2(q{+}1)^{2}\,\mathrm{Re}\,\lambda^{3},
\end{gather*}
and $\dim\ker\Psi=\sum_\lambda\dim\ker H(\lambda)$ over the eigenvalues of $A_1$ with
multiplicity. For the three trivial eigenvalues $\lambda=N\omega^{j}$, $H$ has rank one, so
$\dim\ker\Psi\ge6$ [Ch.~\ref{chap:XI}, Cor.~3.5(i)]. Say $Y$ is \emph{non-degenerate} (ND) if $\det H(\lambda)\ne0$
for every nontrivial eigenvalue. Under (ND): $\ker\Psi$ is the $6$-dimensional space of
type-constant triples $(f,g,h)$ with $f_i+g_{i+1}+(q{+}1)h_{i+2}=0$,
\[
K_Y(u)=(1-u^{3})(1-q^{3}u^{3}),\qquad \deg G=4n+6 ,
\]
and Theorem~\ref{XV:thm:factor} is the identity of {\rm[Ch.~\ref{chap:XI}, Cor.~3.5(ii)]}. $Y_{21},Y_{24},Y_{42},
Y_{168}$ satisfy (ND); the thin torus $T_6$ ($q=1$) does not (Table~\ref{XV:tab:battery}, line (K)).
\end{proposition}

\begin{proof}
The block entries are [Ch.~\ref{chap:XII}, Thm.~3.1]; $A_1$ is normal, so $\ell^{2}(V)$ is an orthogonal sum of
joint eigenlines and $\Psi^{t}\Psi$ is the direct sum of the $H(\lambda)$. Expanding
$\det H$ with $x=\lambda$, $y=\bar\lambda$, $r=xy$, $c=q{+}1$, $D=qN+r$:
$N(ND-c^{2}r)-x(yD-c^{2}x^{2})+cy(c\,y^{2}-cNx)=D(N^{2}-r)-2Nc^{2}r+c^{2}(x^{3}+y^{3})$. For
$\lambda=N$ the rows are $(N,N,cN)$, $(N,N,cN)$, $(cN,cN,(c-1)N+N^{2})$, and
$(c-1)N+N^{2}-c^{2}N=N(q+N-(q{+}1)^{2})=0$, so the rank is one; the $\omega$-rotations are
conjugate. Under (ND) the kernel is the $6$-dimensional type-constant space described (it is
$C'$-invariant and contained in the $9$-dimensional space $V_0$ of type-constant triples).
$C'$ on $V_0$ has characteristic polynomial $\prod_jP_{N\omega^{j}}(u)=(1-u^{3})(1-q^{3}u^{3})(1-q^{6}u^{3})$,
and on the quotient $V_0/\ker\Psi\cong\Psi(V_0)$, the type-constant edge functions, $\LE$
acts by $q^{2}$ times the type rotation, with characteristic polynomial $1-q^{6}u^{3}$; dividing
gives $K_Y$. Then $\deg G=|E|-(3n-6)=4n+6$.
\end{proof}

\section{The chamber companion and the Kang--Li identity}\label{XV:sec:chamber}

A \emph{pointed chamber} is a pair $(e,z)$ with $e=(x\to y)\in E$ and $\{x,y,z\}\in F$; write
it $(x,y;z)$. There are $|P|=3|F|=(q{+}1)|E|$ of them. The Kang--Li operator \cite{KL} is
\[
\LB\bigl((x,y;z),(y,z;w)\bigr)=1\quad\text{iff}\quad w\ne x ,
\]
row sums $q$; it is the chamber flow $L_F$ of [Ch.~\ref{chap:XI}, Rem.~3.6] and [Ch.~\ref{chap:XII}, Def.~4.1]. Let
$\sigma,\rho,\tau\colon P\to E$ send $(x,y;z)$ to $(x\to y)$, $(y\to z)$, $(z\to x)$, and
$\Phi:=[\sigma^{*}\ \rho^{*}\ \tau^{*}]\colon\ell^{2}(E)^{3}\to\ell^{2}(P)$,
$\Phi(f,g,h)(x,y;z)=f(x,y)+g(y,z)+h(z,x)$.

\begin{theorem}[Chamber companion]\label{XV:thm:companion}
\[
\LB\Phi=\Phi B,\qquad B=\begin{pmatrix}0&0&-I_E\\ qI_E&K&K^{t}\\0&-I_E&0\end{pmatrix},
\qquad \det(I+uB)=\det\bigl((1+qu^{3})I_E+uK+u^{2}K^{t}\bigr).
\]
\end{theorem}

\begin{proof}
All sums run over the $q$ vertices $w\ne x$ with $\{y,z,w\}\in F$; with
$K((y,z),(z,w))=[\{y,z,w\}\in F]$,
\begin{align*}
(\LB\sigma^{*}f)(x,y;z)&=\textstyle\sum_{w}f(y,z)=q\,(\rho^{*}f)(x,y;z),\\
(\LB\rho^{*}g)(x,y;z)&=\textstyle\sum_{w}g(z,w)=(Kg)(y,z)-g(z,x)=(\rho^{*}Kg-\tau^{*}g)(x,y;z),\\
(\LB\tau^{*}h)(x,y;z)&=\textstyle\sum_{w}h(w,y)=(K^{t}h)(y,z)-h(x,y)=(\rho^{*}K^{t}h-\sigma^{*}h)(x,y;z).
\end{align*}
Reading the three columns gives $B$. For the determinant, permute the block rows and columns
of $I+uB$ to
\[
\begin{pmatrix}I&-uI&0\\0&I&-uI\\quI&uK^{t}&I+uK\end{pmatrix};
\]
the upper-left $2\times2$ block is unipotent with inverse
$\begin{psmallmatrix}I&uI\\0&I\end{psmallmatrix}$, and the Schur complement is
\[
I+uK-\begin{pmatrix}quI&uK^{t}\end{pmatrix}\begin{pmatrix}I&uI\\0&I\end{pmatrix}
\begin{pmatrix}0\\-uI\end{pmatrix}=I+uK+qu^{3}I+u^{2}K^{t}.
\]
No commutation of $K,K^{t}$ is used. Exact on all five complexes (Table~\ref{XV:tab:battery},
line (B)).
\end{proof}

\begin{theorem}[Kang--Li identity in companion form]\label{XV:thm:KLcomp}
For an $\Atwo$ complex of order $q$ satisfying {\rm(ND)},
\[
\det\bigl((1+qu^{3})I_E+uK+u^{2}K^{t}\bigr)\;=\;(1-u^{3})^{n}\,
\frac{\det(I-u\LE)\,\det(I-u^{2}\LE)}{\det P(u)} ,
\]
and the restriction of $B$ to $W^{3}$ has characteristic polynomial $G(u)\,G(u^{2})$.
\end{theorem}

\begin{proof}
$\ell^{2}(E)^{3}=(\im\Psi)^{3}\oplus W^{3}$ orthogonally, both summands $B$-invariant
(Theorem~\ref{XV:thm:two}). On $W$, $K=T^{t}S-\LE$ acts as $-R$ and $K^{t}$ as $-R^{t}$, so
$B|_{W^{3}}=\begin{psmallmatrix}0&0&-I\\qI&-R&-R^{t}\\0&-I&0\end{psmallmatrix}$ and, by the
elimination of Theorem~\ref{XV:thm:companion},
$\det(I+uB|_{W^{3}})=\det\bigl((1+qu^{3})I-uR-u^{2}R^{t}\bigr)=\det\bigl((I-uR)(I-u^{2}R^{t})\bigr)
=G(u)G(u^{2})$, using $RR^{t}=qI$ (Theorem~\ref{XV:thm:unitary}) and $\det(I-u^{2}R^{t})=G(u^{2})$.
On $(\im\Psi)^{3}$, $\Psi^{\oplus3}$ intertwines $B$ with
$\Bh=\begin{psmallmatrix}0&0&-I\\qI&\Ch&\Chp\\0&-I&0\end{psmallmatrix}$ on $\ell^{2}(V)^{9}$
modulo $(\ker\Psi)^{3}$, so $\det(I+uB|_{(\im\Psi)^{3}})=\det(I+u\Bh)/\det(I+u\Bh|_{(\ker\Psi)^{3}})$.
The same elimination gives $\det(I+u\Bh)=\det\bigl((1+qu^{3})I+u\Ch+u^{2}\Chp\bigr)$, a
$3\times3$ block matrix over $\Q[A_1,A_2]$; on a joint eigenline $(a,b)=(\lambda,\bar\lambda)$
it is
\[
D_{a,b}(u)=\det\begin{pmatrix}1+qu^{3}&(q{+}1)u^{2}&qu+au^{2}\\(q{+}1)u&1+qu^{3}&bu+qu^{2}\\u^{2}&u&1+qu^{3}\end{pmatrix}
=(1-u^{3})\bigl(1-bu^{2}+qau^{4}-q^{3}u^{6}\bigr)
\]
(a polynomial identity in $a,b,q,u$, verified by expansion), so
$\det(I+u\Bh)=(1-u^{3})^{n}\det\bigl(I-u^{2}A_2+qu^{4}A_1-q^{3}u^{6}\bigr)=(1-u^{3})^{n}\det P(u^{2})$,
the last step by transposition. Under (ND), $(\ker\Psi)^{3}\subset V_0^{3}$ with
$V_0$ the type-constant space; $\det(I+u\Bh|_{V_0^{3}})=\prod_jD_{N\omega^{j},N\omega^{-j}}(u)
=(1-u^{3})^{3}(1-u^{6})(1-q^{3}u^{6})(1-q^{6}u^{6})$, and on $\Psi(V_0)^{3}$ --- type-constant
edge functions, on which $K$ and $K^{t}$ act as $(q{+}1)$ times the type rotation and its
inverse --- the factor is $\prod_j(1+qu^{3}+(q{+}1)u\omega^{j}+(q{+}1)u^{2}\omega^{2j})
=\prod_j(1+u\omega^{j}+u^{2}\omega^{2j})(1+qu\omega^{j})=(1-u^{3})^{2}(1+q^{3}u^{3})$; the quotient is
$\det(I+u\Bh|_{(\ker\Psi)^{3}})=(1-u^{3})(1-q^{3}u^{3})(1-u^{6})(1-q^{3}u^{6})=K_Y(u)K_Y(u^{2})$.
Assembling,
\[
\det(I+uB)=(1-u^{3})^{n}\,\frac{\det P(u^{2})\,G(u)\,G(u^{2})}{K_Y(u)\,K_Y(u^{2})},
\]
and Theorem~\ref{XV:thm:factor} at $u$ and at $u^{2}$ turns the right side into the statement.
Certified modulo two $31$-bit primes at $3|E|+1$ points on $Y_{21},Y_{24},Y_{42},T_6$ (a
polynomial identity of degree $\le3|E|$) and at random points on $Y_{168}$
(Table~\ref{XV:tab:battery}, line (B)).
\end{proof}

\begin{corollary}[Kang--Li's zeta identity]\label{XV:cor:KL}
$\det(I+u\LB)\,K_\Phi(u)=\det(I+uB)\,G_\Phi(u)$ with $K_\Phi,G_\Phi$ the characteristic
polynomials of $B$ on $\ker\Phi$ and of $\LB$ on $\ell^{2}(P)/\im\Phi$. For $q=2$, where
$|P|=3|E|$, $K_\Phi=G_\Phi$ on the four thick complexes, and hence, with $\chi(Y)=n$,
\[
\frac{(1-u^{3})^{\chi(Y)}}{\det P(u)}=\frac{\det(I+u\LB)}{\det(I-u\LE)\,\det(I-u^{2}\LE^{t})},
\]
which is \cite[Thm.~C]{KL}, \cite[eq.~(1.6)]{KLW}, for the zeta function
$Z(Y,u)=\det(I-u\LE)^{-1}$. For general $q$, \cite{KL} gives
$G_\Phi/K_\Phi=(1-u^{3})^{\chi(Y)-n}$; verified on $T_6$, where $\chi=0$ and $n=36$. The
eigenvalues of $\LB$ off the pencil are the roots of $G(u)G(u^{2})$: moduli $q^{1/2}$ and
$q^{1/4}$, the latter being square roots of the unitary remainder --- the $q^{-1/4}$ zeros
of \cite[Thm.~1.5.1(4)]{KLW}.
\end{corollary}

\begin{proof}
$\Phi$ intertwines, $\ker\Phi$ is $B$-invariant and $\im\Phi$ is $\LB$-invariant, so both
characteristic polynomials factor through $\ell^{2}(E)^{3}/\ker\Phi\cong\im\Phi$. The rest is
Theorem~\ref{XV:thm:KLcomp}, $\det(I-u^{2}\LE)=\det(I-u^{2}\LE^{t})$, and the computations of
Table~\ref{XV:tab:battery}; the equality $K_\Phi=G_\Phi$ and the general ratio are not proved
here and are the one point where \cite{KL} is used.
\end{proof}

\section{The verification battery}\label{XV:sec:battery}

\begin{table}[ht]
\centering\scriptsize
\renewcommand{\arraystretch}{1.2}
\resizebox{\textwidth}{!}{%
\begin{tabular}{lcl}
\toprule
check & complexes & result\\
\midrule
(D) $\LE\LE^{t}=qI+(q^{2}-q)T^{t}T$, $\LE^{t}\LE=qI+(q^{2}-q)S^{t}S$; $\LE\Psi=\Psi C'$, $\LE^{t}\Psi=\Psi\Cpp$, $K\Psi=\Psi\Ch$, $K^{t}\Psi=\Psi\Chp$ & all five & exact\\
\quad adjacency version of [Ch.~\ref{chap:XI}, Def.~3.1] equals the chamber version & $Y_{168},T_6$ yes; $Y_{21},Y_{24},Y_{42}$ no & exact\\
(K) $\rk\Psi=3n-6$ (mod two primes $+$ explicit rational kernel); $K_Y=(1-u^{3})(1-q^{3}u^{3})$; $\dim W=4n+6$ & $Y_{21},Y_{24},Y_{42},Y_{168}$ & exact\\
\quad $T_6$: $\rk\Psi=93$, $\dim\ker\Psi=15$, (ND) fails, $\dim W=15$, $K_Y$ of degree $15$ & $T_6$ & exact\\
(U) $RR^{t}=R^{t}R=qI$ on $W$; $|\spec R|=\sqrt q$ to $10^{-12}$ & all five & numeric\\
(E) $\det(I-u\LE)K_Y=\det P\cdot G$ in $\Z[u]$; $\deg G=\dim W$; functional equation; roots on $|u|=q^{-1/2}$ & $Y_{21},Y_{24},Y_{42},T_6$ & exact (roots numeric)\\
\quad $Y_{168}$: six missing pencil roots $\{1,\omega,\omega^{2},q,q\omega,q\omega^{2}\}$, $678$ remainder eigenvalues of modulus $\sqrt q$ & $Y_{168}$ & numeric\\
(B) $\LB\Phi=\Phi B$; Theorem~\ref{XV:thm:KLcomp} at $3|E|+1$ points mod two primes; $\det(I+u\LB)=\det(I+uB)$ ($q=2$); ratio $(1-u^{3})^{\chi-n}$ on $T_6$ & all five ($Y_{168}$: random points) & exact / modular\\
(V) nontrivial pencil roots of modulus $q$ (vertex Ramanujan property) & all five & numeric\\
\bottomrule
\end{tabular}}
\medskip
\caption{Verification (\texttt{pgl3\_parahoric.py}). Complexes: $Y_{21},Y_{24},Y_{42},Y_{168}$
($q=2$) and $T_6$ ($q=1$) of [Ch.~\ref{chap:X}].}
\label{XV:tab:battery}
\end{table}

\section{Scope, residue, corrections}\label{XV:sec:scope}

\begin{remark}[Solved, open, impossible]\label{XV:rem:scope}
\begin{sloppypar}
\emph{Solved:} [Ch.~\ref{chap:XI}, Rem.~5.1(b)] in full --- the determinant identity for all $\Atwo$ complexes
(Theorem~\ref{XV:thm:factor}, with the exact form of the trivial factor under (ND),
Proposition~\ref{XV:prop:kernel}) and the identification of $G$ with Kang--Li's chamber factor
(Theorem~\ref{XV:thm:KLcomp}: $G(u)G(u^{2})$ is the parahoric block of the chamber companion);
the Riemann hypothesis for the parahoric factor, unconditionally (Theorem~\ref{XV:thm:unitary});
the Kang--Li identity itself in companion form, with a proof by the degeneracy calculus of
[Ch.~\ref{chap:XII}] (Theorem~\ref{XV:thm:KLcomp}). \emph{Open:} (a) a direct proof that $G_\Phi/K_\Phi=
(1-u^{3})^{\chi-n}$ (Corollary~\ref{XV:cor:KL}), which would make the Kang--Li identity
independent of \cite{KL}; (b) an exact (rather than numeric) certificate of the vertex
Ramanujan property of $Y_{21},\dots,Y_{168}$, i.e.\ of the Riemann hypothesis for their zeta
functions; (c) the exact polynomial identity on $Y_{168}$, verified here only structurally
and at random points modulo primes; (d) the $\Atwo\to\widetilde A_d$ generalization, where
the parahoric remainders should again be $\sqrt{q^{k}}$ times orthogonal matrices.
\emph{Impossible or false as stated:} see Remark~\ref{XV:rem:corrections}.
\end{sloppypar}
\end{remark}

\begin{remark}[Corrections to the task specification]\label{XV:rem:corrections}
\begin{sloppypar}
(1) The identity $\det(I-u\LE)(1-u^{3})(1-q^{3}u^{3})=\det P(u)\det(I-uL_F^{\mathrm{par}})$ with
$L_F^{\mathrm{par}}$ an operator on the $3|F|$ pointed chambers is false as stated: the right
factor would have degree up to $3|F|=(q{+}1)|E|$, not $4n+6$. The remainder $G$ is the
characteristic polynomial of $\LE$ on $W\subset\ell^{2}(E)$, and it enters the chamber
operator as the block $G(u)G(u^{2})$. (2) ``$L_F^{\mathrm{par}}\cong\im(\Psi)^{\perp}$'' is
false: $\dim W=4n+6$ while $\LB$ acts on $(q{+}1)|E|$ coordinates; the correct relation is
$\LB\Phi=\Phi B$ with $B$ on $\ell^{2}(E)^{3}$. (3) The ``Riemann hypothesis for $\Atwo$
complexes'' for the parahoric poles is true but needs no Ramanujan hypothesis: it holds for
every complex (Theorem~\ref{XV:thm:unitary}); the Ramanujan property is exactly the vertex
condition (Remark~\ref{XV:rem:RH}). (4) The formula requires (ND); on $T_6$ the trivial factor
has degree $15$ and $\deg G=15\ne4n+6$. (5) The ``block triangularization'' is in fact a
block diagonalization: $\im\Psi$ is invariant under $\LE^{t}$ too. (6) Our own [Ch.~\ref{chap:XI}, Def.~3.1] is corrected (Remark~\ref{XV:rem:corrXI}). (7) The representation-theoretic reading of
the remainder as ``the subdominant Iwahori--Hecke spectrum on a parahoric'' is \cite{KLW}'s;
here the remainder is characterized combinatorially and its unitarity proved, which is more
than the representation theory gives without a temperedness input.
\end{sloppypar}
\end{remark}

\section*{Acknowledgement of provenance and caveats}

Unrefereed preprint. All numbered statements are proved from the link structure of an
$\Atwo$ complex and linear algebra; the one external input is the relation between $\LB$
and its companion for $q\ne2$ in Corollary~\ref{XV:cor:KL}, quoted from \cite{KL}. The
modular certificates of Theorem~\ref{XV:thm:KLcomp} are identities in $\F_p[u]$ for two
primes, not in $\Z[u]$; the theorem itself is proved. Bibliographic data of
\cite{KL,KLW,LSV} were checked against publisher pages.

% generated from Paper 18/anccft18.tex
\chapter{Algorithmic non-commutative class field theory, XVIII: the unitary remainder in rank $d$}\label{chap:XVIII}
\chaptermark{XVIII}
\begin{chapterabstract}
Paper~XV proved that for an $\Atilde2$ complex of order $q$ the geodesic edge flow
satisfies $\LE\LE^{\mathsf T}=qI+(q^2-q)T^{\mathsf T}T$, so that its remainder is
$\sqrt q$ times an orthogonal matrix and the parahoric zeros of the edge zeta lie on
$|u|=q^{-1/2}$ whether or not the complex is Ramanujan; whether this survives in higher
rank was left open. It does, with the exponent $(d-1)/2$, and the proof is one line of
projective geometry. In an $\Atilde d$ complex the types force the tails of the edges into
a vertex $y$ to be the \emph{hyperplanes} of the link $\PG(d,q)$ and the successors to be
its \emph{points}, and $\{x,y,z\}$ is a simplex exactly when the point $z$ lies in the
hyperplane $x$. Since two distinct hyperplanes of $\PG(d,q)$ always meet in codimension
two --- for every $d$ --- inclusion--exclusion gives a count that does not depend on the
pair, and
\[
\LE\LE^{\mathsf T}=q^{d-1}I+q^{d-1}(q-1)\,T^{\mathsf T}T,\qquad
\LE^{\mathsf T}\LE=q^{d-1}I+q^{d-1}(q-1)\,S^{\mathsf T}S,
\]
which is Paper~XV's identity at $d=2$. Two unconditional consequences follow with no
hypothesis on the complex: \emph{exactly} $|E|-n$ singular values of $\LE$ equal
$q^{(d-1)/2}$, and $\LE$ restricts to a bijective $q^{(d-1)/2}$-isometry from $\ker S$ onto
$\ker T$ (the mixed identities $T\LE=q^{d}S$ and $S\LE^{\mathsf T}=q^{d}T$). On any
subspace of $\ker S\cap\ker T$ invariant under $\LE$ and $\LE^{\mathsf T}$ the remainder
is therefore $q^{(d-1)/2}$ times an orthogonal matrix and its zeros lie on
$|u|=q^{-(d-1)/2}$. The typing is not a convention one may vary: run the same construction
on lines against planes in $\PG(3,q)$ and the count takes two values, because two lines
meet or are skew. Four checks, exhaustive in $\PG(d,q)$ for $d\le5$ and reproducing
Paper~XV on $Y_{21}$, $Y_{24}$, $Y_{42}$. What is \emph{not} done: the invariant subspace
itself for $d\ge3$, and the worked $\Atilde3$ example, for want of a thick $\Atilde3$
complex.
\end{chapterabstract}

\section{The link, and what the types force}\label{XVIII:sec:link}

Let $X$ be an $\Atilde d$ complex of order $q$ in the sense of [Ch.~\ref{chap:X}, Def.~1.1] generalised to
rank $d$: a pure $d$-dimensional simplicial complex with vertices typed by $\Z/(d+1)$,
one vertex of each type in every chamber, and the link of every vertex isomorphic to the
flag complex of $\PG(d,q)$. Write $n$ for the number of vertices, $E$ for the set of
type-increasing directed edges ($\mathrm{type}(y)=\mathrm{type}(x)+1$ for $x\to y$), and
$S,T\in\{0,1\}^{\,n\times E}$ for the tail and head incidence matrices. The geodesic edge
flow is
\[
\LE\bigl((x\to y),(y\to z)\bigr)=1\iff\{x,y,z\}\ \text{is not a simplex of }X,
\]
the rank-$d$ form of [Ch.~\ref{chap:XV}, Def.~1.1]; for $d=2$ a $2$-simplex is a chamber and this is
[Ch.~\ref{chap:XV}, Def.~1.1] verbatim, including the correction recorded there.

\begin{lemma}[What the typing forces]\label{XVIII:lem:typing}
Fix a vertex $y$ and identify its link with the flag complex of $\PG(d,q)$, a neighbour of
type $\mathrm{type}(y)+i$ corresponding to a subspace of vector dimension $i$
$(1\le i\le d)$. Then the tails of the edges into $y$ are exactly the \emph{hyperplanes}
of $\PG(d,q)$, the heads of the edges out of $y$ are exactly its \emph{points}, and for
such $x$ and $z$ the set $\{x,y,z\}$ is a simplex of $X$ if and only if the point $z$ lies
in the hyperplane $x$.
\end{lemma}

\begin{proof}
A tail $x$ has $\mathrm{type}(x)=\mathrm{type}(y)-1\equiv\mathrm{type}(y)+d$, so $i=d$ and
$x$ is a subspace of vector dimension $d$, a hyperplane; a successor $z$ has
$\mathrm{type}(z)=\mathrm{type}(y)+1$, so $i=1$ and $z$ is a point. Two vertices of the
link span a simplex with $y$ precisely when they are incident in the link, which for a
point and a hyperplane means containment.
\end{proof}

That one paragraph is the whole difference between rank two and rank $d$, and it is what
makes the rank-two count survive: of all the pairings of types that one might have had to
handle, the typing selects the single pairing whose intersection behaviour is uniform.

\section{The identity}\label{XVIII:sec:identity}

Write $\theta_j=\theta_j(q)=(q^{\,j+1}-1)/(q-1)$ for the number of points of $\PG(j,q)$.

\begin{theorem}[Unitarity identity in rank $d$]\label{XVIII:thm:main}
For every $\Atilde d$ complex of order $q$, $d\ge2$,
\[
\LE\LE^{\mathsf T}=q^{\,d-1}I+q^{\,d-1}(q-1)\,T^{\mathsf T}T,
\qquad
\LE^{\mathsf T}\LE=q^{\,d-1}I+q^{\,d-1}(q-1)\,S^{\mathsf T}S .
\]
\end{theorem}

\begin{proof}
$(\LE\LE^{\mathsf T})(e,e')$ counts the common successors of $e$ and $e'$, which is zero
unless $e$ and $e'$ have the same head $y$; so the matrix is block diagonal by head, and
$T^{\mathsf T}T$ is the indicator of ``same head''. Fix $y$ and use
Lemma~\ref{XVIII:lem:typing}: the block is $BB^{\mathsf T}$ with $B[x][z]=1$ iff the point $z$
is not in the hyperplane $x$. On the diagonal, the number of points off a hyperplane is
$\theta_d-\theta_{d-1}=q^{\,d}$. Off the diagonal, two \emph{distinct} hyperplanes of
$\PG(d,q)$ meet in a subspace of codimension two, which carries $\theta_{d-2}$ points, so
by inclusion--exclusion the number of points on neither is
\[
\theta_d-2\theta_{d-1}+\theta_{d-2}
=(\theta_d-\theta_{d-1})-(\theta_{d-1}-\theta_{d-2})
=q^{\,d}-q^{\,d-1}=q^{\,d-1}(q-1),
\]
independent of the pair and of $d$. Hence each block is
$q^{\,d}I+q^{\,d-1}(q-1)(J-I)=q^{\,d-1}I+q^{\,d-1}(q-1)J$, which is the first identity.
The second is the same computation in the dual plane: two distinct points of $\PG(d,q)$
span a line, the hyperplanes containing both are the $\theta_{d-2}$ hyperplanes containing
that line, and the count of hyperplanes missing both points is again $q^{\,d-1}(q-1)$.
\end{proof}

\begin{remark}[Reduction to Paper XV]\label{XVIII:rem:d2}
At $d=2$ the constants are $q^{\,d-1}=q$ and $q^{\,d-1}(q-1)=q^2-q$, so
Theorem~\ref{XVIII:thm:main} is [Ch.~\ref{chap:XV}, Thm.~2.1] and its proof is [Ch.~\ref{chap:XV}, Lem.~1.3] with the
counting done once instead of case by case. Check~(A) verifies the two identities, with
these constants, on the complexes $Y_{21}$, $Y_{24}$ and $Y_{42}$ of [Ch.~\ref{chap:X}, Constr.~1.4].
\end{remark}

\begin{proposition}[The mixed identities]\label{XVIII:prop:mixed}
$T\LE=q^{\,d}S$ and $S\LE^{\mathsf T}=q^{\,d}T$.
\end{proposition}

\begin{proof}
$(T\LE)(v,f)$ counts the edges $e$ with head $v$ and $e\to f$; since $e\to f$ forces
$\mathrm{head}(e)=\mathrm{tail}(f)$, the entry vanishes unless $v=\mathrm{tail}(f)$, and
then it counts the hyperplanes of the link of $v$ not containing the point
$\mathrm{head}(f)$, of which there are $\theta_d-\theta_{d-1}=q^{\,d}$. The second
identity is dual.
\end{proof}

\begin{corollary}[Unconditional consequences]\label{XVIII:cor:uncond}
For every $\Atilde d$ complex of order $q$:
\begin{enumerate}
\item[(i)] exactly $|E|-n$ singular values of $\LE$ equal $q^{(d-1)/2}$;
\item[(ii)] $\LE$ maps $\ker S$ into $\ker T$, bijectively, and
$\|\LE w\|^2=q^{\,d-1}\|w\|^2$ there; symmetrically $\LE^{\mathsf T}$ maps $\ker T$ onto
$\ker S$ as a $q^{(d-1)/2}$-isometry.
\end{enumerate}
\end{corollary}

\begin{proof}
(i) $\LE\LE^{\mathsf T}-q^{\,d-1}I=q^{\,d-1}(q-1)T^{\mathsf T}T$ has rank
$\rank T=n$, every vertex being the head of some edge; so $q^{\,d-1}$ is an eigenvalue of
$\LE\LE^{\mathsf T}$ of multiplicity exactly $|E|-n$. (ii) For $w\in\ker S$,
$T\LE w=q^{\,d}Sw=0$ by Proposition~\ref{XVIII:prop:mixed}, and
$\|\LE w\|^2=\langle\LE^{\mathsf T}\LE w,w\rangle=q^{\,d-1}\|w\|^2$ by
Theorem~\ref{XVIII:thm:main}; both spaces have dimension $|E|-n$, so the isometry is onto.
\end{proof}

\begin{corollary}[The critical circle in rank $d$]\label{XVIII:cor:circle}
Let $W\subseteq\ker S\cap\ker T$ be invariant under $\LE$ and $\LE^{\mathsf T}$ and put
$R=\LE|_W$. Then $RR^{\mathsf T}=R^{\mathsf T}R=q^{\,d-1}I_W$, so $q^{-(d-1)/2}R$ is
orthogonal and every root of $G(u)=\det(I-uR)$ lies on $|u|=q^{-(d-1)/2}$. No hypothesis
on the complex --- in particular no Ramanujan hypothesis --- is used.
\end{corollary}

\begin{proof}
Immediate from Theorem~\ref{XVIII:thm:main}, since $T^{\mathsf T}T$ and $S^{\mathsf T}S$ vanish
on $W$; an orthogonal matrix has all eigenvalues of modulus one.
\end{proof}

\begin{remark}[The typing is not a convention]\label{XVIII:rem:typing}
The uniformity in Theorem~\ref{XVIII:thm:main} is a property of the hyperplane/point pairing, not
of $\PG(d,q)$ in general. Run the same construction on lines against planes in $\PG(3,q)$:
two distinct lines either meet or are skew, so the planes containing both number $1$ or
$0$, and the off-diagonal of $BB^{\mathsf T}$ takes two values differing by one --- $9$ and
$10$ for $q=2$, $32$ and $33$ for $q=3$, $144$ and $145$ for $q=5$ (check~(C)). Had the
types of an $\Atilde d$ complex put intermediate dimensions against each other, the
correction term would have been a sum over the classes of the Grassmann association scheme
instead of a multiple of $T^{\mathsf T}T$, and the remainder would have carried several
moduli rather than one. Lemma~\ref{XVIII:lem:typing} is what rules this out.
\end{remark}

\section{Verification}\label{XVIII:sec:battery}

The script \texttt{ad\_unitarity.py} runs the four checks of Table~\ref{XVIII:tab:battery} in
about three seconds; its output is archived as \texttt{ad\_unitarity\_output.txt}.

\begin{table}[ht]
\centering\small
\begin{tabular}{clc}
\toprule
key & check & mode\\
\midrule
(H) & $BB^{\mathsf T}=B^{\mathsf T}B=q^{d-1}I+q^{d-1}(q-1)J$ in $\PG(d,q)$, $14$ pairs $(d,q)$ & exact\\
(C) & lines against planes in $\PG(3,q)$: the off-diagonal takes two values & exact\\
(A) & the identities of \S\ref{XVIII:sec:identity} on $Y_{21},Y_{24},Y_{42}$, with $d=2$ & exact\\
(S) & exactly $|E|-n$ singular values of $\LE$ equal $q^{(d-1)/2}$ & numeric\\
\bottomrule
\end{tabular}
\caption{The verification battery, $4/4$.}
\label{XVIII:tab:battery}
\end{table}

Check~(H) is exhaustive: for $d=2,3,4,5$ and $q=2,3,5,7,11$ with $\theta_d\le1300$ it
forms the full point/hyperplane non-incidence matrix of $\PG(d,q)$ and both of its Gram
matrices, and finds the diagonal constantly $q^{\,d}$ and the off-diagonal constantly
$q^{\,d-1}(q-1)$ --- $14$ instances, including $\PG(5,3)$ on $364$ points and $\PG(4,5)$ on
$781$. Check~(A) builds the thick $\Atilde2$ complexes of [Ch.~\ref{chap:X}] and verifies
Theorem~\ref{XVIII:thm:main}, Proposition~\ref{XVIII:prop:mixed} and
Corollary~\ref{XVIII:cor:uncond}(ii) on each; check~(S) finds $126$, $144$ and $252$ singular
values equal to $\sqrt2$ out of $147$, $168$ and $294$, which are exactly $|E|-n$.

\section{Status, scope, corrections}\label{XVIII:sec:scope}

\emph{Proved.} Lemma~\ref{XVIII:lem:typing}, Theorem~\ref{XVIII:thm:main},
Proposition~\ref{XVIII:prop:mixed} and Corollary~\ref{XVIII:cor:uncond} for every $\Atilde d$ complex
of order $q$, $d\ge2$; Corollary~\ref{XVIII:cor:circle} for every such complex admitting an
invariant $W\subseteq\ker S\cap\ker T$.

\emph{Not proved, and this is the gap.} The existence of that $W$. For $d=2$ it is
[Ch.~\ref{chap:XV}, Thm.~2.1], obtained from the intertwining $\LE\Psi=\Psi C'$ with
$\Psi=[S^{\mathsf T}\ T^{\mathsf T}\ M^{\mathsf T}]$ and the cubic Hecke pencil.
Corollary~\ref{XVIII:cor:uncond} is what survives unconditionally: $\LE$ is already
$q^{(d-1)/2}$ times an isometry from $\ker S$ onto $\ker T$, and only the passage from
that pair of spaces to a single invariant one is missing.

\begin{remark}[The candidate intertwiner in rank $d$]\label{XVIII:rem:psi}
A chamber of an $\Atilde d$ complex has $d+1$ vertices, so an edge $e=(x\to y)$ lying in a
chamber determines $d-1$ further vertices, one of each type
$\mathrm{type}(y)+1,\dots,\mathrm{type}(y)+d-1$. Let $M_i(v,e)$ be the number of chambers
containing $e$ whose vertex of type $\mathrm{type}(y)+i$ is $v$, and put
\[
\Psi_d:=\bigl[\,S^{\mathsf T}\ \ T^{\mathsf T}\ \ M_1^{\mathsf T}\ \cdots\
M_{d-1}^{\mathsf T}\,\bigr]\in\Z^{\,|E|\times(d+1)n},
\]
so that $\Psi_2=[S^{\mathsf T}\ T^{\mathsf T}\ M^{\mathsf T}]$ is the intertwiner of
[Ch.~\ref{chap:XV}, Thm.~2.1] and the block count $d+1$ matches the degree of the Hecke pencil. We record
this as the natural candidate and do not claim it works: what must be shown is
$\LE\Psi_d=\Psi_dC^{(d+1)}$ and $\LE^{\mathsf T}\Psi_d=\Psi_dC^{(d+1)\prime}$ for
$(d+1)\times(d+1)$ block matrices, after which $W=(\operatorname{im}\Psi_d)^{\perp}$ would
satisfy the hypothesis of Corollary~\ref{XVIII:cor:circle} and complete the argument. Both
identities are entrywise \emph{local} --- each entry is a count inside the star of a single
vertex --- so they can in principle be settled in $\PG(d,q)$ without a global complex, as
Theorem~\ref{XVIII:thm:main} was.

\emph{Added in proof.} That has now been carried out, in Paper XX. The local identity is a
single flag count: in $\PG(d,q)$ the number of complete flags $F$ with $F_{i+1}=V$ and $F_1$
inside a hyperplane $H$ takes exactly two values, according as $V\subseteq H$ or not, and
their difference $K_i=f(d-i)f(i+1)q^{\,i}/[i+1]_q$ depends on neither $V$ nor $H$; with
$M_0=cT$ and $M_d=cS$ this gives
$\LE M_i^{\mathsf T}=K_i(T^{\mathsf T}A_{i+1}-c_{i+1}^{-1}M_{i+1}^{\mathsf T})$ and, by the
self-duality of $\PG(d,q)$, the mirror identity for $\LE^{\mathsf T}$. So $\Psi_d$ does work,
Corollary~\ref{XVIII:cor:circle} is unconditional for every $d\ge2$, and at $d=2$ the
characteristic polynomial of $C^{(3)}$ is the cubic Hecke pencil $x^3-A_1x^2+qA_2x-q^3$.
\end{remark}

\begin{remark}[Solved, open, impossible]\label{XVIII:rem:soi}
\emph{Solved:} the counting half of [Ch.~\ref{chap:XV}, Rem.~6.1(d)] --- the exact decomposition of
$\LE\LE^{\mathsf T}$ in rank $d$ --- with the exponent $(d-1)/2$ as anticipated.
\emph{Open:} (a) \emph{closed in Paper XX} --- the invariant subspace
$W_d=(\operatorname{im}\Psi_d)^{\perp}$ exists for every $d$, so
Corollary~\ref{XVIII:cor:circle} is unconditional; what remains open there is the exact rank of
$\Psi_d$, hence $\dim W_d$;
(b) a thick $\Atilde3$ complex small enough to compute with, and with it the degree-four
block companion and a worked example; (c) the exact certificate of the vertex Ramanujan
property, the other half of [Ch.~\ref{chap:XV}, Rem.~6.1(d)], which this note does not touch.
\emph{Impossible:} nothing new. In particular the mechanism does \emph{not} degenerate in
higher rank, which was the outcome one might have expected from
Remark~\ref{XVIII:rem:typing}.
\end{remark}

\begin{remark}[Corrections to the task specification]\label{XVIII:rem:corrections}
Three. (1) The specification proposed building the identity from ``the typical counting
formula for the intersection of $k$-dimensional with $(d-k)$-dimensional subspaces of
$\PG(d,q)$''. No such formula is needed and none would serve: the typing of an $\Atilde d$
complex puts hyperplanes against points (Lemma~\ref{XVIII:lem:typing}), and the only geometric
input is that two distinct hyperplanes meet in codimension two. With intermediate
dimensions the count is not constant (Remark~\ref{XVIII:rem:typing}), so an identity of the
stated shape would be false. (2) The specification asked for
$W_d=\bigcap_{i=0}^{d}\ker S_i$. Only two operators appear in Theorem~\ref{XVIII:thm:main}, and
$\ker S\cap\ker T$ already carries $RR^{\mathsf T}=q^{\,d-1}I$; whether a smaller
intersection is needed is a question about \emph{invariance}, not about the constant.
(3) The predicted $R_dR_d^{\mathsf T}=q^{\,d-1}I$ and the circle $|u|=q^{-(d-1)/2}$ are
correct, and are Corollary~\ref{XVIII:cor:circle}. The third task --- the explicit $\Atilde3$
block companion and a verified example --- is not done: no thick $\Atilde3$ complex is at
hand. The rank-two test complexes came from an exhaustive search for triangle
presentations over $\PG(2,2)$ [Ch.~\ref{chap:X}, Constr.~1.4]; there is no comparably small search in
rank three, and the known constructions of finite $\Atilde d$ quotients go through
arithmetic lattices in $\mathrm{PGL}_{d+1}$ over a local field. That is a separate problem
and we do not pretend to have solved it.
\end{remark}

\section*{Provenance of this chapter}

Unrefereed preprint. Theorem~\ref{XVIII:thm:main} and everything derived from it are proved in
full; the computations verify the geometric input exhaustively in $\PG(d,q)$ for $d\le5$
and confirm that the rank-$d$ statement specialises at $d=2$ to the published identity on
three complexes constructed independently. The one-line proof was found by asking which
pair of types the $\Z/(d+1)$ grading actually selects, and the negative half of
Remark~\ref{XVIII:rem:typing} was computed before the positive half was proved, to check that
the question had teeth.

\input{chapters/XX}
\part{Applications to physics, networks and control}
% generated from Phys 1/bloch_mass.tex
\chapter{Effective mass from spanning trees, and the absence of Bloch phases on rank-two hyperbolic lattices}\label{chap:Phys1}
\chaptermark{Phys1}
\begin{chapterabstract}
Every $\Z^{d}$-periodic tight-binding model with a finite unit cell is an abelian cover of a
finite graph, the winding numbers of the hoppings playing the role of a cocycle. We show
that for such a lattice the inverse effective-mass tensor at the bottom of the band is,
exactly, the Albanese metric of the unit cell: $(M^{-1})_{ab}=(2t/n)\langle
h_a,h_b\rangle$, where $h_a$ is the harmonic projection of the $a$-th winding cochain and
$n$ is the number of sites per cell. Three consequences follow. (i) When the flux threads a
single link $e$, the effective mass has a closed combinatorial form,
$M^{*}=n\,\kappa_0/[2t\,\tau(Y_0\!-\!e)]$, where $\kappa_0$ is the number of spanning trees
of the unit cell and $\tau(Y_0\!-\!e)$ the number of those avoiding $e$. (ii) The tensor
$\kappa_0\,(n/2t)M^{-1}$ has integer entries: \emph{effective masses on such lattices are
quantized}, in units fixed by a spanning-tree count. (iii) The diffusion tensor of the
classical walk on the same lattice obeys $\sigma^{2}=M^{-1}/[t(q{+}1)]$ with $q+1$ the
coordination number --- an Einstein relation with no other non-universal factor. We then
prove a no-go theorem in the opposite direction. For a lattice whose unit cell is a free
quotient of a thick $\Atwo$ building --- the rank-two analogue of the trees and $\{p,q\}$
tessellations used in circuit-QED hyperbolic lattices --- the first Betti number vanishes,
so no $\Z^{d}$ cover exists at all: there is no Brillouin zone, no crystal momentum, and no
band structure, and by property (T) no continuous family of unitary characters deforms the
trivial one either. The threadable fluxes form a \emph{finite} group, computed exactly as
$(\Z/2)^{6}$, $(\Z/2)^{4}\oplus\Z/14$ and $(\Z/2)^{3}\oplus(\Z/4)^{2}$ on complexes with
$21$, $24$ and $42$ sites; on the first of these only $\pi$-fluxes exist. Hyperbolic band
theory does not survive the passage from rank one to rank two. All statements are verified
by machine: $90$ checks, exact except where labelled.
\end{chapterabstract}

\section{Introduction}\label{Phys1:sec:intro}

Band theory is a statement about $\Z^{d}$ symmetry. A tight-binding model on a lattice with
a finite unit cell and a $\Z^{d}$ translation group decomposes over a Brillouin zone
$\T^{d}$, and the low-energy physics is controlled by the curvature of the lowest band at
its minimum --- the effective mass. Two recent developments make it worth asking how far
this picture is combinatorial rather than geometric.

The first is the experimental realization of \emph{hyperbolic lattices} in circuit quantum
electrodynamics~\cite{KFH}, where a coplanar-waveguide network is laid out as a $\{p,q\}$
tessellation of the hyperbolic plane, and the subsequent construction of a hyperbolic band
theory~\cite{MR1,MR2} in which the role of the Brillouin zone is played by a
higher-dimensional Jacobian and by a character variety of the (non-abelian) translation
group. The platform has since broadened to electrical
circuits~\cite{Leng22}, integrated photonics~\cite{Huang24}, and scalable superconducting
architectures with the unit cell on a genus-three surface~\cite{Xu25}. The unit cell of such
a lattice is a finite graph, and what is threaded through it is flux.

The second is the observation, standard in the mathematics of crystal lattices since Kotani
and Sunada~\cite{KS1,KS2}, that a $\Z^{d}$-periodic graph carries a canonical flat metric
on $H^{1}$ of its unit cell --- the Albanese metric --- which controls the long-time heat
kernel and the central limit theorem of the random walk, and that this metric is bound up
with the \emph{complexity} of the unit cell, i.e.\ its spanning-tree count~\cite{KS3};
explicit Kirchhoff expressions for it were given by Baker and Faber~\cite{BF}. What does not
seem to have been said is that this metric \emph{is} the inverse effective-mass tensor of
the tight-binding band bottom, exactly and without approximation, and that reading it so
turns those Kirchhoff expressions into a constraint on a measurable quantity. That is the
content of Secs.~\ref{Phys1:sec:mass}--\ref{Phys1:sec:trees}, whose sharpest consequence is a
quantization statement (Corollary~\ref{Phys1:cor:quant}): the entries of $M^{-1}$, measured in the
natural unit $2t/(n\kappa_0)$, are integers.

The paper then turns the question around. Trees and $\{p,q\}$ tessellations are rank-one
objects: the tree is the Bruhat--Tits building of $\mathrm{GL}_2$ over a local field, and
its finite quotients are ordinary finite graphs, which always have $b_1\ge2$ and therefore
always admit $\Z^{d}$ covers. The natural rank-two generalization is a two-dimensional
building of type $\Atwo$ --- a simplicial complex of sites, links and plaquettes in which
every vertex link is the incidence graph of a projective plane --- whose finite quotients
are the Ramanujan complexes~\cite{LSV:Phys1}. These are the higher-rank hyperbolic lattices. We
show in Sec.~\ref{Phys1:sec:nogo} that for them band theory does not merely become harder, it
ceases to exist: $b_1=0$ by Garland's vanishing theorem~\cite{Ga} as extended by Ballmann
and \'Swi\k{a}tkowski~\cite{BSw}, so there is no $\Z^{d}$ cover, no crystal momentum, and no
continuous flux parameter. The Aharonov--Bohm phases that a rank-two hyperbolic lattice
admits form a finite abelian group, which we compute exactly on three complexes.

The two halves meet at a single graph. The Heawood graph --- the incidence graph of the
Fano plane --- is a perfectly good rank-one unit cell: it has $b_1=8$, $50421=3\cdot7^{5}$
spanning trees, and a single-link flux gives it the effective mass $M^{*}=147/(8t)$
(Table~\ref{Phys1:tab:trees}). It is also exactly the vertex link of a thick $\Atwo$ building of
order two. Assembling Heawood links into a rank-two complex destroys $b_1$ completely
(Table~\ref{Phys1:tab:ranktwo}). The obstruction is not local.

\subsection*{Relation to previous work}

Theorem~\ref{Phys1:thm:mass} lies close to known mathematics. Once the Albanese metric is
identified with the covariance of the crystal-lattice central limit theorem~\cite{KS1} and
with a Kirchhoff ratio~\cite{KS3,BF}, Eq.~\eqref{Phys1:eq:band} follows from a second-order
perturbation argument, which we give in full for completeness rather than for novelty. What
we claim as new is: the tensor form of Theorem~\ref{Phys1:thm:mass} together with the explicit
statement that this object is the physical effective-mass tensor; the quantization
Corollary~\ref{Phys1:cor:quant} and its reading as a rigid constraint on measurable band
curvature; the deletion--contraction closed form \eqref{Phys1:eq:mstar}; the loop-flux
Corollary~\ref{Phys1:cor:loop}; and the whole of Sec.~\ref{Phys1:sec:nogo}. The closest prior result in
the physics literature that we are aware of is that of Korotyaev and Saburova~\cite{KorS},
who bound and estimate effective masses at band edges of periodic discrete and metric
graphs from geometric parameters, without spanning trees or the Albanese metric. An
``effective-mass approximation'' of a different kind --- the long-distance reduction of the
hyperbolic tight-binding Hamiltonian to a Laplace--Beltrami operator --- appears
in Ref.~\onlinecite{MR1}.

For Sec.~\ref{Phys1:sec:nogo} we have found no prior work: we are not aware of any physics
literature that places a tight-binding model on an affine building of rank two or on a
Ramanujan complex, let alone one that asks whether it has a band structure. A literature
search cannot establish a negative, and we state this only as the limit of our knowledge.

\section{Lattices from abelian covers}\label{Phys1:sec:setup}

Let $Y_0$ be a finite connected $(q{+}1)$-regular graph with $n$ vertices and
$m=(q{+}1)n/2$ edges, $E^{+}$ a choice of orientation, and $r=m-n+1=b_1(Y_0)$ its cycle
rank. Write $\partial\colon\Z^{E^{+}}\to\Z^{V}$ for the boundary map, $\partial
e=t(e)-o(e)$, and $\Delta_0=\partial\partial^{t}=(q{+}1)I-A$ for the graph Laplacian.
Let $\kappa_0$ denote the number of spanning trees of $Y_0$.

\begin{definition}[The lattice]\label{Phys1:def:lattice}
A \emph{winding cochain} is a $d$-tuple $\bm\alpha=(\alpha_1,\dots,\alpha_d)$ with
$\alpha_a\in\Z^{E^{+}}$. The associated lattice $\Yt$ has sites $V\times\Z^{d}$ and a link
from $(o(e),\bm j)$ to $(t(e),\bm j+\bm\alpha(e))$ for each $e\in E^{+}$ and each
$\bm j\in\Z^{d}$. We assume throughout that the period map
$H_1(Y_0;\Z)\to\Z^{d}$, $c\mapsto(\langle\alpha_a,c\rangle)_a$, is surjective, which is
exactly the condition that $\Yt$ be connected; in particular $d\le r$.
\end{definition}

This is not a special construction: \emph{every} $\Z^{d}$-periodic tight-binding lattice
with $n$ sites per cell and uniform hopping is of this form, with $Y_0$ the unit cell and
$\bm\alpha(e)$ recording how many cells the hopping $e$ crosses. Changing which
representative of a site one calls ``in cell $\bm 0$'' replaces $\bm\alpha$ by
$\bm\alpha+\partial^{t}\bm\varphi$; the class $[\bm\alpha]\in H^{1}(Y_0;\Z)^{d}$ is the
gauge-invariant datum.

The Hamiltonian is $H=t\,\widetilde\Delta$, with $\widetilde\Delta$ the Laplacian of $\Yt$:
uniform hopping amplitude $-t$ on every link and on-site energy $(q{+}1)t$. Floquet
transform in the $\Z^{d}$ direction gives
$\ell^{2}(\Yt)\cong\int^{\oplus}_{\T^{d}}\C^{n}\,d^{d}k/(2\pi)^{d}$ and
$\widetilde\Delta\cong\int^{\oplus}\Delta_{\bm k}$ with the Peierls-substituted Bloch
Laplacian
\begin{equation}\label{Phys1:eq:bloch}
(\Delta_{\bm k})_{uv}=(q{+}1)\delta_{uv}
-\!\!\sum_{e\colon u\to v}\!\! e^{i\bm k\cdot\bm\alpha(e)}
-\!\!\sum_{e\colon v\to u}\!\! e^{-i\bm k\cdot\bm\alpha(e)} ,
\end{equation}
a Hermitian positive-semidefinite $n\times n$ matrix, and bands
$E_j(\bm k)=t\lambda_j(\bm k)$, $\lambda_0\le\dots\le\lambda_{n-1}$. We set $\hbar=1$ and
the lattice constant to unity, so that
\begin{equation}\label{Phys1:eq:massdef}
(M^{-1})_{ab}:=\frac{\partial^{2}E_0}{\partial k_a\partial k_b}\Big|_{\bm k=0}
=t\,\frac{\partial^{2}\lambda_0}{\partial k_a\partial k_b}\Big|_{\bm k=0}.
\end{equation}
Since $\Delta_0$ is a Laplacian of a connected graph its kernel is the constant line, simple,
with gap $\lambda_1(0)>0$; the band bottom sits at $\bm k=0$ and $\lambda_0$ is analytic
near it.

\section{The effective-mass tensor is the Albanese metric}\label{Phys1:sec:mass}

Let $\Pi=I-\partial^{t}\Delta_0^{+}\partial$ be the orthogonal projection of
$\R^{E^{+}}$ onto the cycle space $\ker\partial=H_1(Y_0;\R)$, and let
$h_a:=\Pi\alpha_a$ be the harmonic representative of the $a$-th winding cochain. Define the
\emph{Albanese matrix}
\begin{equation}\label{Phys1:eq:H}
H_{ab}:=\langle h_a,h_b\rangle=\alpha_a^{t}\,\Pi\,\alpha_b .
\end{equation}

\begin{theorem}[Effective mass]\label{Phys1:thm:mass}
Under the hypothesis of Definition~\ref{Phys1:def:lattice}, $\lambda_0$ is even in $\bm k$,
simple and analytic near the origin, and
\begin{equation}\label{Phys1:eq:band}
\lambda_0(\bm k)=\frac1n\,\bm k^{t}H\bm k+O(|\bm k|^{4}),
\qquad\text{so}\qquad
M^{-1}=\frac{2t}{n}\,H .
\end{equation}
$H$ is positive definite, and it is the Gram matrix of the classes $[\alpha_a]$ in the flat
metric of the Jacobian torus $H^{1}(Y_0;\R)/H^{1}(Y_0;\Z)$, dual to the Albanese torus of
Kotani--Sunada~\cite{KS2}: if $C$ is the matrix of a $\Z$-basis of $H_1(Y_0;\Z)$,
$G=C^{t}C$ its Gram matrix and $\bm a_a=C^{t}\alpha_a$ the vector of periods, then
\begin{equation}\label{Phys1:eq:albanese}
H_{ab}=\bm a_a^{t}\,G^{-1}\bm a_b .
\end{equation}
Equivalently, in terms of the Floquet determinant $\Theta(\bm k)=\det\Delta_{\bm k}$,
\begin{equation}\label{Phys1:eq:floquet}
\frac{\partial^{2}\Theta}{\partial k_a\partial k_b}\Big|_{\bm k=0}=2\,\kappa_0\,H_{ab}.
\end{equation}
\end{theorem}

\begin{proof}
$\Delta_{-\bm k}=\overline{\Delta_{\bm k}}=\Delta_{\bm k}^{t}$, so the spectrum is even in
$\bm k$ and all odd Taylor coefficients of $\lambda_0$ vanish. Expanding
\eqref{Phys1:eq:bloch} about $\bm k=0$ gives $\Delta_{\bm k}=\Delta_0+\sum_a k_aD^{(1)}_a
+\tfrac12\sum_{ab}k_ak_bD^{(2)}_{ab}+O(|\bm k|^{3})$ with
$(D^{(1)}_a)_{uv}=-i\sum_{e\colon u\to v}\alpha_a(e)+i\sum_{e\colon v\to u}\alpha_a(e)$ and
$(D^{(2)}_{ab})_{uv}=\sum_{e\colon u\to v}\alpha_a(e)\alpha_b(e)+(u\!\leftrightarrow\!v)$.
On the constant vector $\bm1$ one has $D^{(1)}_a\bm1=i\,\partial\alpha_a$,
$\langle\bm1,D^{(1)}_a\bm1\rangle=0$ and
$\langle\bm1,D^{(2)}_{ab}\bm1\rangle=2\langle\alpha_a,\alpha_b\rangle$. Second-order
Rayleigh--Schr\"odinger perturbation theory at the simple eigenvalue $0$ therefore gives
\[
\lambda_0(\bm k)=\frac1n\sum_{ab}k_ak_b\Bigl(\langle\alpha_a,\alpha_b\rangle
-(\partial\alpha_a)^{t}\Delta_0^{+}(\partial\alpha_b)\Bigr)+O(|\bm k|^{4}),
\]
and the bracket is $\alpha_a^{t}(I-\partial^{t}\Delta_0^{+}\partial)\alpha_b=H_{ab}$. For
\eqref{Phys1:eq:albanese}, $h_a$ is the orthogonal projection of $\alpha_a$ onto the column space
of $C$, so $h_a=C(C^{t}C)^{-1}C^{t}\alpha_a$ and
$\langle h_a,h_b\rangle=\alpha_a^{t}CG^{-1}C^{t}\alpha_b$. Positive definiteness is the
surjectivity hypothesis: $\bm v^{t}H\bm v=0$ forces $\sum_av_a[\alpha_a]=0$ in
$H^{1}(Y_0;\Q)$, impossible for $\bm v\neq0$ when the periods span. Finally
$\Theta(\bm k)=\lambda_0(\bm k)\prod_{j\ge1}\lambda_j(\bm k)$, and by the matrix--tree
theorem $\prod_{j\ge1}\lambda_j(0)=n\kappa_0$, which with \eqref{Phys1:eq:band} gives
\eqref{Phys1:eq:floquet}.
\end{proof}

For $d=1$ Eq.~\eqref{Phys1:eq:band} is the effective-mass identity of
Ref.~\onlinecite{ANCFT14}, where $-n\lambda_0''(0)/2$ was identified with an
$\mathcal L$-invariant and with a diffusion constant. The content of
\eqref{Phys1:eq:albanese} for $d=1$ --- that the Albanese metric evaluates the harmonic energy ---
is Kotani--Sunada~\cite{KS1,KS3}; what is added here is the $d>1$ tensor, the identification
of the whole object with $M^{-1}$, and the consequences drawn in
Secs.~\ref{Phys1:sec:trees}--\ref{Phys1:sec:einstein}.

\begin{remark}
Equation~\eqref{Phys1:eq:floquet} is the practically useful form: $\Theta(\bm k)$ is a Laurent
polynomial in $e^{ik_1},\dots,e^{ik_d}$ with integer coefficients, computable exactly, and
its Hessian at the origin delivers $\kappa_0H$ without any eigenvalue computation. This is
how the battery of Sec.~\ref{Phys1:sec:battery} verifies Theorem~\ref{Phys1:thm:mass} in exact
arithmetic.
\end{remark}

\section{Spanning trees and the quantization of effective mass}\label{Phys1:sec:trees}

\begin{theorem}[Single-link flux]\label{Phys1:thm:trees}
Let $d=1$ and let the winding cochain be the indicator of a single link $e$, so that the
lattice is the unit cell repeated along one direction with exactly one hopping crossing the
cell boundary. Then
\begin{equation}\label{Phys1:eq:trees}
H=1-\Reff(e)=\frac{\tau(Y_0-e)}{\kappa_0},
\end{equation}
where $\Reff(e)$ is the effective resistance across $e$ when every link carries unit
resistance, $\kappa_0=\tau(Y_0)$ is the number of spanning trees of the unit cell, and
$\tau(Y_0-e)$ is the number of spanning trees avoiding $e$. Consequently
\begin{equation}\label{Phys1:eq:mstar}
M^{*}=\frac{n}{2t}\cdot\frac{\kappa_0}{\tau(Y_0-e)} .
\end{equation}
\end{theorem}

\begin{proof}
$H=\Pi_{ee}=1-(\partial^{t}\Delta_0^{+}\partial)_{ee}=1-\Reff(e)$ by the standard
resistance formula. Kirchhoff's theorem gives $\kappa_0\Reff(e)=\tau(Y_0/e)$, the number of
spanning trees containing $e$, and deletion--contraction
$\kappa_0=\tau(Y_0-e)+\tau(Y_0/e)$ finishes it. Equation~\eqref{Phys1:eq:mstar} is
\eqref{Phys1:eq:band}.
\end{proof}

\begin{corollary}[Quantization of effective mass]\label{Phys1:cor:quant}
For any winding cochain, $\kappa_0H$ is a matrix of \emph{integers}. Hence every entry of
the inverse effective-mass tensor is an integer multiple of $2t/(n\kappa_0)$,
\begin{equation}\label{Phys1:eq:quant}
\frac{n\kappa_0}{2t}\,(M^{-1})_{ab}\in\Z ,
\qquad
\det\Bigl[\frac{n\kappa_0}{2t}M^{-1}\Bigr]\in\Z ,
\end{equation}
and the admissible effective masses of a lattice with unit cell $Y_0$ form a discrete set
determined by the single integer $\kappa_0$.
\end{corollary}

\begin{proof}
By Kirchhoff's theorem in its cycle form, the Gram matrix $G$ of any $\Z$-basis of
$H_1(Y_0;\Z)$ has $\det G=\kappa_0$; since $\mathrm{adj}\,G$ is an integer matrix,
$G^{-1}\in\kappa_0^{-1}M_r(\Z)$. The period vectors $\bm a_a=C^{t}\alpha_a$ are integral, so
$\kappa_0H_{ab}=\bm a_a^{t}(\kappa_0G^{-1})\bm a_b\in\Z$ by \eqref{Phys1:eq:albanese}.
\end{proof}

This is the most directly falsifiable statement in the paper. The band curvature at the
$\Gamma$ point of such a lattice cannot be an arbitrary real number: measured in units of
$2t/n$, it is a rational with denominator dividing $\kappa_0$. Table~\ref{Phys1:tab:trees}
lists the values for six unit cells.

The algebra behind Corollary~\ref{Phys1:cor:quant} is elementary, and in the equivalent form
``the Albanese lattice of $Y_0$ has covolume $\kappa_0$'' it is classical~\cite{KS3,BF}.
What is at stake here is only what is being quantized: not a lattice covolume but the
curvature of a band, which an experiment can measure.

\begin{table}[t]
\centering
\begin{ruledtabular}
\begin{tabular}{lccccccc}
$Y_0$ & $n$ & $m$ & $q$ & $\kappa_0$ & $\Reff(e)$ & $\tau(Y_0-e)$ & $tM^{*}$\\
\colrule
$K_4$        & $4$  & $6$  & $2$ & $16$    & $1/2$   & $8$     & $4$\\
$K_{3,3}$    & $6$  & $9$  & $2$ & $81$    & $5/9$   & $36$    & $27/4$\\
$Q_3$        & $8$  & $12$ & $2$ & $384$   & $7/12$  & $160$   & $48/5$\\
$K_5$        & $5$  & $10$ & $3$ & $125$   & $2/5$   & $75$    & $25/6$\\
Petersen     & $10$ & $15$ & $2$ & $2000$  & $3/5$   & $800$   & $25/2$\\
Heawood      & $14$ & $21$ & $2$ & $50421$ & $13/21$ & $19208$ & $147/8$\\
\end{tabular}
\end{ruledtabular}
\caption{Single-link flux. $\kappa_0$ is the number of spanning trees of the unit cell,
$\tau(Y_0-e)$ the number avoiding the flux link, and $tM^{*}=n\kappa_0/[2\tau(Y_0-e)]$ by
Eq.~\eqref{Phys1:eq:mstar}. For the Heawood graph $\kappa_0=3\cdot7^{5}$ and
$\tau(Y_0-e)=8\cdot7^{4}$. All entries exact.}
\label{Phys1:tab:trees}
\end{table}

\begin{corollary}[Loop flux]\label{Phys1:cor:loop}
If the winding cochain is the indicator of an oriented cycle $C\subset Y_0$ of length
$\ell$, then $h=\alpha$, $H=\ell$ exactly, and
\begin{equation}
M^{*}=\frac{n}{2t\ell} .
\end{equation}
The effective mass depends only on the length of the flux loop and on the number of sites
per cell --- not on the rest of the unit cell.
\end{corollary}

\begin{proof}
A cycle lies in $\ker\partial$, so $\Pi\alpha=\alpha$ and $H=\|\alpha\|^{2}=\ell$.
\end{proof}

Corollary~\ref{Phys1:cor:loop} is the sharp opposite of Theorem~\ref{Phys1:thm:trees}: a single link is
maximally sensitive to the ambient cell (through $\Reff$), a closed loop is completely
insensitive to it. Both are verified in Sec.~\ref{Phys1:sec:battery} on $K_5$, where the
cyclically oriented triangle gives $H=3$ exactly while the same three links with a
non-cyclic orientation give $H=7/5$.

\begin{table}[t]
\centering
\begin{ruledtabular}
\begin{tabular}{lcccc}
$Y_0$ & $\kappa_0$ & $\kappa_0H$ & $\det(\kappa_0H)$\\
\colrule
$K_4$    & $16$    & $\left(\begin{smallmatrix}8&-4\\-4&8\end{smallmatrix}\right)$ & $48$\\[2pt]
$Q_3$    & $384$   & $\left(\begin{smallmatrix}160&-32\\-32&160\end{smallmatrix}\right)$ & $24576$\\[2pt]
$K_5$    & $125$   & $\left(\begin{smallmatrix}75&-25\\-25&75\end{smallmatrix}\right)$ & $5000$\\[2pt]
Heawood  & $50421$ & $\left(\begin{smallmatrix}19208&-9604\\-9604&19208\end{smallmatrix}\right)$ & $276710448$\\
\end{tabular}
\end{ruledtabular}
\caption{The effective-mass tensor for $d=2$, two single-link fluxes in independent
directions. $\kappa_0H=(n\kappa_0/2t)M^{-1}$ is an integer matrix
(Corollary~\ref{Phys1:cor:quant}); the off-diagonal entries are the mass anisotropy. All entries
exact.}
\label{Phys1:tab:tensor}
\end{table}

\section{An Einstein relation}\label{Phys1:sec:einstein}

The same matrix $H$ governs the classical transport problem on $\Yt$. Let $(X_i)$ be the
simple random walk on $Y_0$ and $\bm S_N=\sum_{i<N}\bm\alpha(X_i,X_{i+1})\in\Z^{d}$ the
accumulated winding --- the cell coordinate of the lifted walk on $\Yt$.

\begin{theorem}[Diffusion tensor]\label{Phys1:thm:einstein}
$\bm S_N/\sqrt N$ converges in law to a centred Gaussian with covariance
$\sigma^{2}_{ab}=H_{ab}/m$, and therefore
\begin{equation}\label{Phys1:eq:einstein}
\sigma^{2}=\frac{1}{t\,(q{+}1)}\;M^{-1}.
\end{equation}
\end{theorem}

\begin{proof}
The walk is reversible with uniform stationary law and transition matrix $A/(q{+}1)$; the
drift is $-\partial\alpha_a/(q{+}1)$ and the Poisson equation $\Delta_0g_a=-\partial\alpha_a$
has solution $g_a=-\Delta_0^{+}\partial\alpha_a$. Then
$\alpha_a+\partial^{t}g_a=\Pi\alpha_a=h_a$ are the martingale increments, and the
martingale central limit theorem gives
$\sigma^{2}_{ab}=\frac{2}{n(q+1)}\langle h_a,h_b\rangle=H_{ab}/m$. Substituting
$H=(n/2t)M^{-1}$ and $m=n(q{+}1)/2$ gives \eqref{Phys1:eq:einstein}.
\end{proof}

Equation~\eqref{Phys1:eq:einstein} says that the quantum band curvature and the classical
diffusion tensor of a lattice of this type differ by the coordination number alone. In
particular the mass anisotropy and the diffusion anisotropy are the \emph{same} tensor up
to scale, and both are quantized by Corollary~\ref{Phys1:cor:quant}. The $d=1$ case is the
crystal-lattice central limit theorem of Kotani and Sunada~\cite{KS1}, in which the
covariance is the Albanese metric; Theorem~\ref{Phys1:thm:mass} is what identifies that covariance
with a band curvature.

\section{Rank two: no Brillouin zone}\label{Phys1:sec:nogo}

Everything above needs $b_1(Y_0)\ge1$. For an ordinary finite graph this is automatic and
harmless: a connected $(q{+}1)$-regular graph with $q\ge2$ has $r=m-n+1\ge2$, so $\Z^{d}$
covers exist for every $d\le r$. The hyperbolic lattices of Refs.~\onlinecite{KFH,MR1} are
of this type --- their unit cells are finite graphs, and the difficulty there is that the
\emph{deck} group is non-abelian, not that abelian directions are missing.

We now consider the rank-two analogue. A thick $\Atwo$ building of order $q$ is the
two-dimensional simplicial complex on which $\mathrm{PGL}_3$ of a local field with residue
field of size $q$ acts; it is the rank-two counterpart of the $(q{+}1)$-regular tree, and
every vertex link in it is the incidence graph of the projective plane over $\mathbb F_q$
(for $q=2$, the Heawood graph). Its finite quotients are lattices of sites, links and
plaquettes; the Ramanujan complexes of Lubotzky, Samuels and Vishne~\cite{LSV:Phys1} are among
them, and they are built from groups acting simply transitively on the vertices in the sense
of Cartwright, Mantero, Steger and Zappa~\cite{CMSZ}. These are the natural candidates for a
higher-rank hyperbolic lattice.

\begin{theorem}[No band structure in rank two]\label{Phys1:thm:nogo}
Let $Y=\Gamma\backslash B$ be a finite quotient of a thick $\Atwo$ building of order
$q\ge2$ by a group $\Gamma$ acting freely. Then
\begin{enumerate}
\item[(i)] $b_1(Y)=0$ and $H_1(Y;\Z)$ is finite;
\item[(ii)] $Y$ admits no connected $\Z^{d}$ cover for any $d\ge1$. There is no
$\Z^{d}$-periodic lattice with unit cell $Y$, hence no Brillouin zone, no crystal momentum,
and no band structure;
\item[(iii)] the abelian covers of $Y$ are exactly those with deck group a quotient of the
finite group $H_1(Y;\Z)$, so the threadable flux configurations form the finite group
$\Hom(H_1(Y;\Z),U(1))\cong H_1(Y;\Z)$;
\item[(iv)] $\Gamma$ has Kazhdan's property (T); the trivial representation is isolated in
its unitary dual, so no continuous family of unitary characters --- abelian or not ---
deforms the trivial one.
\end{enumerate}
\end{theorem}

\begin{proof}
$\Gamma$ acts properly and cocompactly on $B$, so by Garland's vanishing
theorem~\cite{Ga} in the form extended to all $q\ge2$ by Ballmann and
\'Swi\k{a}tkowski~\cite{BSw}, $H^{1}(\Gamma;\R)=0$ and $\Gamma$ has property (T). $Y$ is
aspherical with $\pi_1(Y)=\Gamma$, so $b_1(Y)=0$ and $H_1(Y;\Z)=\Gamma^{\mathrm{ab}}$ is
finite. A cover with deck group $\Z^{d}$ requires a surjection $\Gamma\to\Z^{d}$, hence a
surjection $\Gamma^{\mathrm{ab}}\to\Z^{d}$, impossible for $d\ge1$ from a finite group; this
is (ii), and (iii) is the classification of abelian covers by
$H^{1}(Y;\,\cdot\,)=\Hom(H_1(Y;\Z),\,\cdot\,)$. Part (iv) is property (T).
\end{proof}

\begin{table}[t]
\centering
\begin{ruledtabular}
\begin{tabular}{lccccc}
$Y$ & sites & links & plaquettes & $b_1$ & $H_1(Y;\Z)$\\
\colrule
$Y_{21}$ & $21$ & $147$ & $147$ & $0$ & $(\Z/2)^{6}$\\
$Y_{24}$ & $24$ & $168$ & $168$ & $0$ & $(\Z/2)^{4}\oplus\Z/14$\\
$Y_{42}$ & $42$ & $294$ & $294$ & $0$ & $(\Z/2)^{3}\oplus(\Z/4)^{2}$\\
\end{tabular}
\end{ruledtabular}
\caption{Three finite quotients of a thick $\Atwo$ building of order $q=2$, built as
abelian covers of the three-site type quotient of a triangle presentation over
$\mathrm{PG}(2,2)$. Every vertex link is a Heawood graph. $b_1=0$ in every case, so none of
them is the unit cell of any periodic lattice; the entire flux group is the finite group in
the last column. On $Y_{21}$ the only admissible fluxes are $0$ and $\pi$. All entries
exact ($p$-adic Smith normal forms; see Sec.~\ref{Phys1:sec:battery}).}
\label{Phys1:tab:ranktwo}
\end{table}

Table~\ref{Phys1:tab:ranktwo} makes the statement concrete. On $Y_{21}$ the flux group is
$(\Z/2)^{6}$: there are exactly $64$ distinct flux configurations, each link pattern
carrying flux $0$ or $\pi$, and nothing in between. There is no knob to turn continuously,
and no momentum to plot a band against.

\begin{remark}[What this does \emph{not} say]\label{Phys1:rem:notsay}
Theorem~\ref{Phys1:thm:nogo} is a statement about translational symmetry, not about spectra. The
Laplacian and the Hecke operators of $Y$ have perfectly good spectra --- indeed the
Ramanujan complexes of Ref.~\onlinecite{LSV} have optimal spectral gaps, which is precisely
what makes them attractive as synthetic lattices. What fails is the specific device of
labelling eigenstates by a crystal momentum in a torus. Any transport theory on a rank-two
hyperbolic lattice must be built without a Brillouin zone, and by
Theorem~\ref{Phys1:thm:nogo}(iv) not even an infinitesimal one is available near the trivial
character. The effective mass of Sec.~\ref{Phys1:sec:mass} is simply not defined there: it is
the curvature of a band that does not exist.
\end{remark}

\begin{remark}[The obstruction is global]\label{Phys1:rem:global}
The Heawood graph appears on both sides of this paper. As a unit cell in its own right it
is unremarkable: $b_1=8$, $\kappa_0=3\cdot7^{5}$, and a single-link flux gives
$tM^{*}=147/8$ (Table~\ref{Phys1:tab:trees}). As the \emph{link} of a rank-two complex it is part
of a structure with $b_1=0$ (Table~\ref{Phys1:tab:ranktwo}). No local inspection of a rank-two
hyperbolic lattice reveals the obstruction; it is a property (T) phenomenon, visible only
at the level of the fundamental group.
\end{remark}

\section{Numerical verification}\label{Phys1:sec:battery}

All statements were checked by machine. The script \texttt{bloch\_mass.py} accompanying this
paper runs $90$ checks in $13$ s and reports every one as passing; its output is archived as
\texttt{bloch\_mass\_output.txt}. Exact arithmetic is used throughout except where a row is
labelled \emph{numeric}.

\begin{table}[t]
\centering
\begin{ruledtabular}
\begin{tabular}{clcc}
key & check & $\#$ & mode\\
\colrule
(S) & $H=1-\Reff(e)$ and $\kappa_0H=\tau(Y_0-e)$, Eq.~\eqref{Phys1:eq:trees} & $6/6$ & exact\\
(M) & $\partial^2_{ab}\Theta|_0=2\kappa_0H$, Eq.~\eqref{Phys1:eq:floquet} & $15/15$ & exact\\
(M) & $\lambda_0(\bm k)=\bm k^{t}H\bm k/n+O(|\bm k|^{4})$ & $15/15$ & numeric\\
(Q) & $\kappa_0H$ integral, Eq.~\eqref{Phys1:eq:quant} & $15/15$ & exact\\
(E) & $\sigma^{2}=H/m$ by Green--Kubo, Eq.~\eqref{Phys1:eq:einstein} & $15/15$ & exact\\
(G) & gap $\lambda_1(0)>0$ and $\lambda_0(0)=0$ & $15/15$ & numeric\\
(X) & $H$ reproduces Refs.~\onlinecite{ANCFT9,ANCFT14} & $4/4$ & exact\\
(C) & loop flux gives $H=\ell$, Corollary~\ref{Phys1:cor:loop} & $1/1$ & exact\\
(B) & $b_1=0$ and $H_1(Y;\Z)$, Table~\ref{Phys1:tab:ranktwo} & $4/4$ & exact\\
\colrule
 & total & $90/90$ & \\
\end{tabular}
\end{ruledtabular}
\caption{The verification battery (\texttt{bloch\_mass.py}). Unit cells: $K_4$,
$K_{3,3}$, $Q_3$, $K_5$, Petersen, Heawood, with single-link fluxes ($d=1$), two-link
fluxes ($d=2$), and on $K_4$ and $K_5$ the multi-chord and triangle patterns of
Refs.~\onlinecite{ANCFT9,ANCFT14}.}
\label{Phys1:tab:battery}
\end{table}

A word on the two \emph{numeric} rows. The band identity \eqref{Phys1:eq:band} is checked by
evaluating $\lambda_0$ at $|\bm k|=10^{-6}$ and $5\times10^{-7}$ in $40$-digit arithmetic
and taking one Richardson step, which removes the $O(|\bm k|^{2})$ truncation and leaves a
residual below $10^{-24}$; the observed relative deviations from $\bm k^{t}H\bm k/n$ range
from $2\times10^{-27}$ to $2\times10^{-26}$. The exact row (M) above it, which verifies the
Floquet-determinant form \eqref{Phys1:eq:floquet} by interpolating $\Theta$ as a Laurent
polynomial with rational coefficients and differentiating it symbolically, is the
mathematically complete check; the numeric row is a redundant confirmation that the object
being differentiated really is the bottom band.

The rank-two rows use the complex builder and the $p$-adic Smith-normal-form routine of
Ref.~\onlinecite{ANCFT10}: the complexes of Table~\ref{Phys1:tab:ranktwo} are abelian covers of
the three-site type quotient of a triangle presentation over $\mathrm{PG}(2,2)$ compatible
with a Singer cycle, the presentation axioms are re-verified at start-up, and every vertex
link is confirmed to be a Heawood graph. The Betti numbers are exact ranks; the torsion
subgroups are exact $p$-adic Smith forms with the prime set certified by the greatest common
divisor of the maximal minors.

\section{Discussion}\label{Phys1:sec:disc}

\subsection{What is testable}

Corollary~\ref{Phys1:cor:quant} is a constraint on a measurable quantity. For a synthetic lattice
assembled from a known unit cell --- a coplanar-waveguide network, a photonic-crystal
supercell, a cold-atom superlattice --- the number $\kappa_0$ of spanning trees of the cell
is a finite combinatorial datum, and the band curvature at $\Gamma$ must be an integer
multiple of $2t/(n\kappa_0)$. Departures measure either the failure of the uniform-hopping
idealization or the presence of longer-range terms; the statement is rigid enough to be
diagnostic.

Theorem~\ref{Phys1:thm:trees} says more: for a single flux link, the mass is set by the ratio of
two spanning-tree counts, so it can be engineered by choosing where the flux link sits, with
no computation beyond counting trees. Corollary~\ref{Phys1:cor:loop} gives the complementary
design rule --- route the flux around a closed loop of length $\ell$ and the mass becomes
$n/2t\ell$, insensitive to everything else in the cell.

Equation~\eqref{Phys1:eq:einstein} is a consistency check available without any quantum
measurement: the classical diffusion tensor on the same network fixes the effective-mass
tensor up to the coordination number.

\subsection{What is open}

(a) The non-abelian case. Hyperbolic band theory~\cite{MR1,MR2} replaces $\T^{d}$ by a
character variety of a non-abelian surface group. Theorem~\ref{Phys1:thm:mass} is a statement
about the abelian directions only; whether a ``non-abelian effective mass'' exists at a
higher-genus analogue of the $\Gamma$ point, and whether it is still a spanning-tree
quantity, we do not know.

(b) The quartic term. Band non-parabolicity is the next Taylor coefficient of
$\lambda_0$. It is a rational function of the same combinatorial data, but we have no
closed form for it and no quantization statement.

(c) Rank two without a Brillouin zone. Theorem~\ref{Phys1:thm:nogo} removes the tool but not the
problem: finite quotients of $\Atwo$ buildings are excellent expanders with explicit
spectra, and the question of what replaces band theory for transport on them --- given that
property (T) blocks even an infinitesimal character deformation --- seems to us the most
interesting thing this paper leaves behind.

\subsection{What is impossible}

A $\Z^{d}$-periodic lattice whose unit cell is a free quotient of a thick $\Atwo$ building,
for any $d\ge1$ and any $q\ge2$ (Theorem~\ref{Phys1:thm:nogo}). Any construction claiming a
Brillouin zone for such a lattice is in error, and the error will be that the quotient is
not free, or the building is not thick, or the complex is not a quotient of a building at
all. This is a theorem and does not depend on the state of the literature; the remark in
Sec.~\ref{Phys1:sec:intro}, that no one appears to have asked the question, does.

\section*{Provenance and caveats}

This is an unrefereed preprint. Every numbered statement is proved from finite linear
algebra together with three classical inputs that are quoted, not reproved: Kirchhoff's
matrix--tree theorem in its cycle form, the martingale central limit theorem for additive
functionals of a finite reversible chain, and the Garland--Ballmann--\'Swi\k{a}tkowski
vanishing theorem~\cite{Ga,BSw}. Kotani--Sunada~\cite{KS1,KS2,KS3} and Baker--Faber~\cite{BF}
supply the Albanese metric, its relation to the spanning-tree count, and the
crystal-lattice central limit theorem; Secs.~\ref{Phys1:sec:mass}--\ref{Phys1:sec:einstein} are a
physical reading of that material and not a replacement for it.

Attribution, statement by statement. The $d=1$ case of Theorem~\ref{Phys1:thm:mass} is
Ref.~\onlinecite{ANCFT14}, Theorem~4.2; Eq.~\eqref{Phys1:eq:trees} for $d=1$ is
Ref.~\onlinecite{ANCFT9}, Theorem~1.2(iv); Theorem~\ref{Phys1:thm:nogo}(i) is
Ref.~\onlinecite{ANCFT12}, Theorem~1.1, which in turn quotes Refs.~\onlinecite{Ga,BSw}. The
identification of the Albanese metric with a Kirchhoff ratio, and hence the mathematical
substance of Corollary~\ref{Phys1:cor:quant}, is in Refs.~\onlinecite{KS3,BF}. What the present
paper adds is the $d>1$ tensor form, the identification of that tensor with $M^{-1}$, the
deletion--contraction form \eqref{Phys1:eq:mstar}, Corollary~\ref{Phys1:cor:loop}, the Einstein relation
in the form \eqref{Phys1:eq:einstein}, and Sec.~\ref{Phys1:sec:nogo} --- the reading of the rank-two
vanishing theorem as a no-go theorem for band structure, which is the part we believe to be
without precedent. We have not found prior physics work on tight-binding models over
rank-two affine buildings; we cannot prove that none exists.

\begin{acknowledgments}
The complex builder and the exact $p$-adic Smith normal form used in Sec.~\ref{Phys1:sec:battery}
are those of Ref.~\onlinecite{ANCFT10}.
\end{acknowledgments}

% generated from Holo 1/boundary_blind.tex
\chapter{What the boundary algebra of $p$-adic AdS/CFT does not see}\label{chap:Holo1}
\chaptermark{Holo1}
\begin{chapterabstract}
\noindent
In $p$-adic AdS/CFT the bulk is the Bruhat--Tits tree of $\mathrm{PGL}_2$ over a local
field and the boundary is $\PP^1(F)$. The boundary carries two standard operator algebras,
the crossed products $C(\PP^1(F))\rtimes\Gamma$ and $L^\infty(\PP^1(F))\rtimes\Gamma$, and
both have been completely determined: the first is a Cuntz--Krieger algebra with
$K_0=\Z^{r}\oplus\Z/(r-1)$ depending only on the rank of $\Gamma$, and the second is the
hyperfinite factor of type $\III_\lambda$ with $\lambda=q^{-2}$ or $q^{-1}$ according as
the quotient graph is bipartite or not (Robertson, 2005). Neither result has ever been
cited in the holography literature, and no $p$-adic AdS/CFT paper has used an
operator-algebraic invariant at all. We import them, and add one arithmetic layer. The
bulk Hecke--Laplacian channel has $K_0=\Z\oplus\Jac(Y)$, the critical group of the
quotient graph, whose order is a product of Frobenius traces
$\frac1n\prod_f(p+1-a_p(f))$. We exhibit two cubic graphs on eight vertices whose boundary
algebras agree in $K$-theory \emph{and} in von Neumann type, and whose bulk critical groups
have coprime orders $256=2^8$ and $363=3\cdot11^2$: the boundary data do not determine the
bulk arithmetic, not even which primes divide it. This is strictly stronger than
Robertson's own counterexample, which the type does separate. We do not conclude that
holography fails: Cornelissen and Marcolli proved that the boundary quantum statistical
mechanical system, with its Busemann time evolution, reconstructs the graph completely, so
what we exhibit is a precise gap in a hierarchy whose top is a reconstruction theorem.
Finally we observe that $\lambda\ne1$ is itself a diagnostic. Continuum holography is
forced to type $\III_1$ by theorems (mixing; half-sided modular inclusion); $p$-adic
holography is type $\III_\lambda$ with $\lambda<1$, so its modular flow is periodic and its
boundary algebra is not a Haag--Kastler algebra. All computations are exact and
machine-verified.
\end{chapterabstract}

\section{Introduction}\label{Holo1:sec:intro}

$p$-adic AdS/CFT \cite{GKPSW,HMSS} replaces anti-de Sitter space by the Bruhat--Tits tree
$X$ of $\mathrm{PGL}_2(F)$, $F$ a non-archimedean local field with residue field of size
$q$, and the conformal boundary by $\PP^1(F)$, the space of ends of $X$. A discrete
torsion-free cocompact lattice $\Gamma\le\mathrm{PGL}_2(F)$ --- free of finite rank $r$ ---
plays the role of a quotient identification, exactly as a Schottky group does for BTZ
\cite{HuangJepsen}, and $Y=\Gamma\backslash X$ is a finite $(q{+}1)$-regular graph on $n$
vertices. The field has matured on correlators, tensor networks, Wilson lines and finite
temperature \cite{GubserParikh,Marcolli20,HuangJepsen,Mondal}, but it has remained
kinematic in one specific respect: no invariant from operator algebras has ever been
computed in it.

That is unfortunate, because the relevant invariants were computed twenty years ago, in a
literature the holography community does not read. Robertson \cite{Rob05} determined both
boundary algebras of exactly this situation, and Cornelissen and Marcolli \cite{CM13}
determined precisely how much of the bulk each one remembers. The purpose of this note is
to import those results, to add one layer they do not contain, and to draw one physical
consequence.

\paragraph{What is imported.}
For $\Gamma$ a free uniform tree lattice of rank $r$, Robertson proved
\cite[Thm.~1]{Rob05} that $C(\partial X)\rtimes\Gamma$ is a Cuntz--Krieger algebra with
\begin{equation}\label{Holo1:eq:K}
K_0\bigl(C(\partial X)\rtimes\Gamma\bigr)\cong\Z^{r}\oplus\Z/(r-1)\Z,
\qquad K_1\cong\Z^{r},
\end{equation}
$[1]$ generating the torsion summand --- so by Kirchberg--Phillips the algebra
\emph{depends only on $r$}. (The same $K$-groups were obtained independently in
\cite{CLM08}.) He proved further \cite[Thm.~2, Cor.~1]{Rob05} that
\begin{equation}\label{Holo1:eq:type}
L^\infty(\partial X)\rtimes\Gamma\ \cong\ \text{hyperfinite }\III_\lambda,\qquad
\lambda=\begin{cases}q^{-2}&\Gamma\subset\mathrm{PSL}_2(F)\ (\,Y\ \text{bipartite}),\\
q^{-1}&\text{otherwise},\end{cases}
\end{equation}
the ratio set being generated by the Busemann cocycle; the rank-one case goes back to
Ramagge and Robertson \cite{RR97} and the general hyperbolic-group statement to
\cite{INO08,Bowen14}. Cornelissen and Marcolli \cite{CM13} then showed that the
$C^*$-level invariants ``reconstruct only topological information'' --- nine natural
equivalences all reduce to equality of the first Betti number --- while the boundary
\emph{quantum statistical mechanical} system, $C(\partial X)\rtimes\Gamma$ equipped with
the Busemann time evolution, reconstructs $Y$ up to isomorphism.

\paragraph{What is added.}
The bulk of this geometry has an invariant that none of these papers considers. The
Hecke--Laplacian $\Delta_p=(q{+}1)I-T_p$ of $Y$ has
\begin{equation}\label{Holo1:eq:bulk}
\operatorname{coker}_\Z\Delta_p\cong\Z\oplus\Jac(Y),
\qquad
|\Jac(Y)|=\frac1n\prod_f\bigl(p+1-a_p(f)\bigr),
\end{equation}
where $\Jac(Y)$ is the critical (sandpile) group \cite{Biggs,BHN} and the product runs over
the weight-two newforms attached to $\Gamma$, so that $|\Jac(Y)|$ is a product of Frobenius
traces \cite[I]{ANCFT}. The critical group has appeared in physics --- as the global
symmetry group of fracton models on graphs \cite{GLS22,GLSS23} --- but never in holography.
Our contribution is Theorem~\ref{Holo1:thm:main}: the boundary data of \eqref{Holo1:eq:K} and
\eqref{Holo1:eq:type} \emph{together} fail to determine \eqref{Holo1:eq:bulk}, and fail badly enough
that two bulks in the same boundary class can have coprime arithmetic.

\paragraph{What we do not claim.}
Not that holography fails. \cite{CM13} is a reconstruction theorem: keep the time
evolution and the bulk comes back. What Theorem~\ref{Holo1:thm:main} locates is a gap between two
rungs of a ladder whose top rung is a positive result. Nor do we claim \eqref{Holo1:eq:K} or
\eqref{Holo1:eq:type}; both are \cite{Rob05}.

\section{Setting}\label{Holo1:sec:setting}

$Y=\Gamma\backslash X$ is a finite connected $(q{+}1)$-regular graph, $n=|V_Y|$,
$m=(q{+}1)n/2$, and
\[
r=\operatorname{rank}\Gamma=m-n+1,\qquad r-1=m-n=\tfrac{(q-1)n}{2}.
\]
The \emph{boundary channel} is the crossed product by the $\Gamma$-action on $\partial X$;
concretely $C(\partial X)\rtimes\Gamma\cong\mathcal O_{A^{\mathrm{nb}}}$ with
$A^{\mathrm{nb}}$ the non-backtracking matrix on the $2m$ oriented edges of $Y$
\cite{Rob05}, and $K_0=\operatorname{coker}(I-(A^{\mathrm{nb}})^{t})$,
$K_1=\ker(I-(A^{\mathrm{nb}})^{t})$ by \cite{Cuntz81}. The \emph{bulk channel} is
$\Delta_p$ on the $n$ vertices, with \eqref{Holo1:eq:bulk}.

\begin{remark}[Two relations, two types]\label{Holo1:rem:tworel}
Equation \eqref{Holo1:eq:type} concerns the orbit relation of $\Gamma$ on $\partial X$. It must
be distinguished from the tail equivalence relation on the Cantor limit
$\varprojlim_kV_{Y_k}$ of a tower of covers, whose ratio set is all of $q^{\Z}\cup\{0\}$
and which therefore gives type $\III_{1/q}$ unconditionally \cite[III]{ANCFT}. The orbit
relation of a type-preserving lattice realises only $q^{2\Z}$, because the Busemann cocycle
of such an element has even value; this is the source of the dichotomy in \eqref{Holo1:eq:type}
and it is easy to conflate the two.
\end{remark}

\begin{proposition}\label{Holo1:prop:coarse}
The boundary $K$-theory \eqref{Holo1:eq:K} is a function of $(q,n)$ alone, and the von Neumann
type \eqref{Holo1:eq:type} is a function of $q$ and the bipartiteness of $Y$ alone. In
particular all $(q{+}1)$-regular graphs on $n$ vertices with the same bipartiteness have
the same boundary data.
\end{proposition}

\begin{proof}
$r-1=(q-1)n/2$ and $r=(q-1)n/2+1$ depend only on $(q,n)$; \eqref{Holo1:eq:type} is as stated.
\end{proof}

\section{The bulk is not determined}\label{Holo1:sec:main}

\begin{theorem}\label{Holo1:thm:main}
There are two cubic graphs $Y_A,Y_B$ on eight vertices with
\begin{itemize}
\item the same boundary $K$-theory, $K_0=\Z^{5}\oplus\Z/4$ and $K_1=\Z^{5}$;
\item the same von Neumann type, both being non-bipartite and hence $\III_{1/q}$;
\item coprime bulk critical groups,
\[
\Jac(Y_A)=\Z/4\oplus\Z/4\oplus\Z/16,\quad |\Jac(Y_A)|=256=2^{8},
\]
\[
\Jac(Y_B)=\Z/33\oplus\Z/11,\quad\ \ |\Jac(Y_B)|=363=3\cdot11^{2}.
\]
\end{itemize}
Explicitly, with vertices $0,\dots,7$,
\[
\begin{aligned}
E(Y_A)&=\{03,06,07,12,15,16,24,25,34,37,45,67\},\\
E(Y_B)&=\{02,03,07,14,16,17,24,25,34,35,56,67\}.
\end{aligned}
\]
Consequently the boundary $K$-theory and the boundary von Neumann type, taken together, do
not determine $|\Jac(Y)|$, nor the set of primes dividing it, nor --- by \eqref{Holo1:eq:bulk}
--- the Frobenius traces $a_p(f)$ of the associated automorphic forms.
\end{theorem}

\begin{proof}
Direct computation of the three invariants in exact integer arithmetic; this is check~(P)
of Section~\ref{Holo1:sec:battery}. Both graphs contain an odd cycle, hence are non-bipartite.
\end{proof}

\begin{remark}[Comparison with \cite{Rob05}]\label{Holo1:rem:compare}
Robertson gives his own counterexample to boundary rigidity \cite[Ex.~1.5]{Rob05}: the
theta graph and the dumbbell graph, both with $\pi_1=F_2$, non-isomorphic, with the same
$C(\partial X)\rtimes\Gamma$, realised concretely inside $\mathrm{PGL}_2(\mathbb Q_2)$. His
example is separated by the von Neumann type, $\III_{1/4}$ against $\III_{1/2}$ --- indeed
that is his point, that the measure-theoretic invariant is finer. Theorem~\ref{Holo1:thm:main}
defeats both invariants at once, and does so in a way that is arithmetically meaningful
rather than merely graph-theoretic.
\end{remark}

Table~\ref{Holo1:tab:named} shows the two channels on a list of standard quotients, and
Table~\ref{Holo1:tab:spread} shows how much bulk arithmetic a single boundary class holds.

\begin{table}[ht]
\centering\small
\begin{tabular}{lrrrllr}
\toprule
$Y$ & $q$ & $n$ & $r$ & $K_0$(boundary) & vN type & $|\Jac(Y)|$\\
\midrule
$K_4$            & 2 &  4 &  3 & $\Z^3\oplus\Z/2$    & $\III_{1/q}$   & 16\\
$K_{3,3}$        & 2 &  6 &  4 & $\Z^4\oplus\Z/3$    & $\III_{1/q^2}$ & 81\\
prism $C_3\times K_2$ & 2 & 6 & 4 & $\Z^4\oplus\Z/3$ & $\III_{1/q}$   & 75\\
prism $C_4\times K_2$ & 2 & 8 & 5 & $\Z^5\oplus\Z/4$ & $\III_{1/q^2}$ & 384\\
M\"obius $M_8$   & 2 &  8 &  5 & $\Z^5\oplus\Z/4$    & $\III_{1/q}$   & 392\\
M\"obius $M_{10}$& 2 & 10 &  6 & $\Z^6\oplus\Z/5$    & $\III_{1/q^2}$ & 1815\\
prism $C_5\times K_2$ & 2 & 10 & 6 & $\Z^6\oplus\Z/5$ & $\III_{1/q}$  & 1805\\
Petersen         & 2 & 10 &  6 & $\Z^6\oplus\Z/5$    & $\III_{1/q}$   & 2000\\
Heawood          & 2 & 14 &  8 & $\Z^8\oplus\Z/7$    & $\III_{1/q^2}$ & 50421\\
$K_5$            & 3 &  5 &  6 & $\Z^6\oplus\Z/5$    & $\III_{1/q}$   & 125\\
$K_6$            & 4 &  6 & 10 & $\Z^{10}\oplus\Z/9$ & $\III_{1/q}$   & 1296\\
\bottomrule
\end{tabular}
\caption{The two channels. $K_1=\Z^r$ throughout. Within a block of equal $(q,n)$ the
boundary $K$-theory is constant; the type refines it by bipartiteness; $|\Jac(Y)|$ is
constrained by neither. Compare prism $C_5\times K_2$ and Petersen: same $(q,n)$, same
type, $|\Jac|=1805$ against $2000$. All entries exact.}
\label{Holo1:tab:named}
\end{table}

\begin{table}[ht]
\centering\small
\begin{tabular}{rrlrr}
\toprule
$n$ & $r$ & $K_0$(boundary) & distinct $|\Jac|$ found & range\\
\midrule
 8 & 5 & $\Z^{5}\oplus\Z/4$ &  5 & $256$ -- $392$\\
10 & 6 & $\Z^{6}\oplus\Z/5$ & 16 & $960$ -- $2000$\\
12 & 7 & $\Z^{7}\oplus\Z/6$ & 38 & $2808$ -- $9747$\\
14 & 8 & $\Z^{8}\oplus\Z/7$ & 54 & $12100$ -- $46998$\\
\bottomrule
\end{tabular}
\caption{Bulk critical groups inside a single boundary $K$-theory class, from $60$ random
cubic quotients at each $n$. Sampled, not exhaustive: the true counts are at least these.}
\label{Holo1:tab:spread}
\end{table}

\section{The type as a physical diagnostic}\label{Holo1:sec:type}

The recent algebraic program in holography \cite{LL1,LL2,Witten22,CPW23,CLPW23} turns on
the fact that the large-$N$ algebra of a holographic CFT, and the algebra of a black hole
exterior, are of type $\III_1$, and that crossed products by the modular flow convert them
to semifinite type $\mathrm{II}$, which is what makes generalized entropy finite. Type
$\III_1$ is not an accident there: it is forced. If the correlation functions decay --- if
the algebra is generated by mixing operators in a KMS state --- the algebra must be
$\III_1$ \cite{FLMO23}; and the half-sided modular inclusion structure that
\cite{LL1,LL2} use forces $\III_1$ by Wiesbrock's theorem \cite{Wiesbrock93}. Local
algebras of relativistic quantum field theory are the hyperfinite $\III_1$ factor
\cite{BDF87}.

Equation \eqref{Holo1:eq:type} says that $p$-adic holography is not in that class. Its boundary
algebra is $\III_\lambda$ with $\lambda=q^{-1}$ or $q^{-2}$, a Powers factor. The physical
content of $\lambda<1$ is that the modular flow is \emph{periodic}: the flow of weights is
the translation flow on $\mathbb R/T\mathbb Z$ with
\[
T=\frac{2\pi}{\log\lambda^{-1}}=\frac{2\pi}{\kappa\log q},\qquad\kappa\in\{1,2\},
\]
so that the modular Hamiltonian has spectrum in $(\log q^{\kappa})\Z$ rather than a
continuum. A $\III_\lambda$ algebra has a nontrivial discrete invariant --- Connes' $S$
invariant is $\lambda^{\Z}\cup\{0\}$ --- where a $\III_1$ algebra has none.

We draw one conclusion and refrain from a second. The conclusion: the type is a sharp,
computable obstruction to any claim that a $p$-adic holographic boundary theory is a local
quantum field theory in the Haag--Kastler sense, and it identifies precisely what would
have to change --- the residue field would have to be made continuous, $q\to1$ --- for the
$\III_1$ of \cite{LL1,Witten22} to be recovered. What we refrain from: we do not propose a
$p$-adic analogue of the crossed-product construction of \cite{Witten22}. For
$\III_\lambda$ with $\lambda<1$ the crossed product by the modular flow is a type
$\mathrm{II}_\infty$ algebra as well, but the relevant flow is periodic, and whether the
resulting trace has anything to do with an entropy is a question we have not investigated.

\section{Verification}\label{Holo1:sec:battery}

The script \texttt{boundary\_blind.py} accompanying this note runs the checks of
Table~\ref{Holo1:tab:battery} in $0.3$\,s, all passing; its output is archived as
\texttt{boundary\_blind\_output.txt}. Arithmetic is exact integer throughout except the two
rows marked \emph{sampled}.

\begin{table}[ht]
\centering\small
\begin{tabular}{clc}
\toprule
key & check & mode\\
\midrule
(K) & \eqref{Holo1:eq:K} holds on 11 standard quotients, including $q=3,4$ & exact\\
(B) & $r-1=(q-1)n/2$: the boundary $K$-theory is a function of $(q,n)$ & exact\\
(J) & witness families at $(q,n)=(2,6),(2,8),(2,10)$: same boundary, different $\Jac$ & exact\\
(P) & Theorem~\ref{Holo1:thm:main}: same $K$-theory, same type, coprime $|\Jac|$ & exact\\
(S) & Table~\ref{Holo1:tab:spread}: the spread inside one boundary class & sampled\\
(R) & the Busemann cocycle is $q^{\Z}$-valued (normalisation check only) & exact\\
\bottomrule
\end{tabular}
\caption{The verification battery (\texttt{boundary\_blind.py}), $7/7$ passing.}
\label{Holo1:tab:battery}
\end{table}

Row (R) deserves a caveat we state rather than bury: it verifies only that the cylinder
measures $\nu(\Omega_v)=q^{1-d(O,v)}$ give a $q^{\Z}$-valued Radon--Nikodym cocycle. That
the \emph{essential range} of that cocycle is attained --- which is what Krieger's theorem
\cite{Krieger76} needs --- is quoted from \cite{RR97,Rob05}, not verified here; no finite
computation could verify it.

\section{Scope}\label{Holo1:sec:scope}

\emph{Quoted, not proved here.} Equations \eqref{Holo1:eq:K} and \eqref{Holo1:eq:type} and the
rigidity observation \cite{Rob05}; the reconstruction theorem \cite{CM13}; Krieger's
classification \cite{Krieger76}; the Cuntz machine \cite{Cuntz81}; the forcing of
$\III_1$ in the continuum \cite{FLMO23,Wiesbrock93,BDF87}.

\emph{New here.} Theorem~\ref{Holo1:thm:main} and Tables~\ref{Holo1:tab:named}--\ref{Holo1:tab:spread}: the
arithmetic layer, and the observation that $K$-theory and type together are still blind to
it. The import of \cite{Rob05,CM13} into $p$-adic holography, where no operator-algebraic
invariant has previously been computed. The reading of $\lambda\ne1$ in
Section~\ref{Holo1:sec:type} as a diagnostic separating $p$-adic from continuum holography.

\emph{Open.} (a) Whether the Cornelissen--Marcolli reconstruction \cite{CM13}, which
recovers $Y$ and hence $\Jac(Y)$, can be made to recover the arithmetic \emph{directly} ---
an explicit formula for $|\Jac(Y)|$ from the boundary KMS data would be a genuine
holographic dictionary entry for Frobenius traces. (b) The higher-rank case: for
$\widetilde A_n$ buildings the boundary type is $\III_{1/q}$ or $\III_{1/q^2}$ according to
the parity of $n$ \cite{CR03}, and holographic constructions on buildings exist
\cite{GMP22,Marcolli20}, but none is arithmetic; the bulk invariant there is the rank-two
critical group of \cite[X]{ANCFT}, and we have not computed the corresponding witnesses.
(c) Whether the periodicity of the modular flow has an entropic meaning.

\section*{Provenance of this chapter}

Unrefereed preprint. The three exact computations behind Theorem~\ref{Holo1:thm:main} --- the
Cuntz--Krieger $K$-groups, the bipartiteness, and the Smith normal form of the
Hecke--Laplacian --- are performed by the accompanying script and are reproducible in
seconds. Everything else in Sections~\ref{Holo1:sec:intro}--\ref{Holo1:sec:type} is quoted, with the
attribution set out in Section~\ref{Holo1:sec:scope}. We searched the physics literature for any
prior use of an operator-algebraic invariant in $p$-adic AdS/CFT, and for any occurrence of
type $\III_\lambda$ with $\lambda\ne1$ in holography, and found none; a search cannot
establish a negative, and we state this as the limit of our knowledge.

% generated from Net 1/tower_lambda.tex
\chapter{The sign of $\lambda$: an algebraic certificate for how a graph tower is generated}\label{chap:Net1}
\chaptermark{Net1}
\begin{chapterabstract}
\noindent
For a tower of finite graphs $Y_0\leftarrow Y_1\leftarrow\cdots$ the $\ell$-adic size of
the sandpile group often obeys a three-parameter law
$\ord_\ell|\Jac(Y_k)|=\mu\ell^{k}+\lambda k+\nu$. In the established theory of abelian
$\ell$-towers this law is a theorem and $\mu,\lambda$ are non-negative by construction,
being the Iwasawa invariants of a torsion $\Zl[[T]]$-module. We observe that the tower of
iterated line digraphs --- the standard design iteration for de Bruijn and Kautz
interconnection networks --- obeys the same law with $\lambda=-1$, on every base graph we
tested, with $\mu=n(q{+}1)/q$ exactly. Since $\lambda<0$ cannot be the $\lambda$-invariant
of any torsion $\Zl[[T]]$-module, a measured $\lambda<0$ certifies that the family is not an
abelian $\ell$-tower. We also find that the certificate detects \emph{normality}, not
commutativity: dihedral (non-abelian but Galois) towers give $\lambda=3>0$, and only the
non-normal line-digraph tower goes negative. We are careful about what is new. The abelian
law is due to Vall\`ieres and collaborators; the line-digraph tower is classical in network
design (Fiol--Alegre--Yebra, 1983) and its spanning-tree arithmetic is Knuth's, in the
sandpile form of Levine; indeed we show that $\lambda=-1$ telescopes directly from that
classical recursion, and verify the telescoping on five bases. What appears to be new is
organising the iteration as a tower, extracting $(\mu,\lambda,\nu)$ from it, obtaining a
negative $\lambda$ at all, and reading the $-1$ as an Euler characteristic --- the
one-dimensional Eisenstein kernel removed at each level. Since the spanning-tree count is a
standard network reliability measure, the law is a scaling law for a reliability measure
whose linear coefficient records the design method. All $27$ towers are computed in exact
integer arithmetic.
\end{chapterabstract}

\section{Introduction}\label{Net1:sec:intro}

Given a family of networks that grows by refinement --- a sequence
$Y_0\leftarrow Y_1\leftarrow\cdots$ in which each term covers or refines the last --- one
may ask how the family was generated. Existing answers in network science are statistical
(likelihood, cross-validation, graphlet summaries) or geometric: the renormalisation
exponents of Song, Havlin and Makse \cite{SHM05,SHM06} infer a growth mechanism from how
box-covering statistics scale. This note offers an \emph{algebraic} answer, in a narrow
but exactly computable setting.

The invariant is the sandpile (critical) group $\Jac(Y)$ \cite{Biggs}, whose order is the
number of spanning trees --- itself a classical measure of network reliability
\cite{Boesch86,Colbourn87,BH09}. Along a tower, $\ord_\ell|\Jac(Y_k)|$ frequently obeys
\begin{equation}\label{Net1:eq:law}
\ord_\ell\bigl|\Jac(Y_k)\bigr|\;=\;\mu\,\ell^{k}\;+\;\lambda\,k\;+\;\nu
\qquad(k\ \text{large}).
\end{equation}

\paragraph{What is established.}
For abelian $\ell$-towers, \eqref{Net1:eq:law} is a theorem. Vall\`ieres \cite{Val21} proved it
for a family of towers over bouquets; McGown and Vall\`ieres \cite{MV24} proved it over an
arbitrary connected multigraph; Gonet \cite{Gonet22} obtained the Sylow structure from the
$\La=\Zl[[T]]$-module decomposition, and Kleine and M\"uller \cite{KM24} and others have
since extended the picture to $\Zl^{d}$-towers, ramified towers, weighted graphs and
$p$-adic Lie voltage groups. In every one of these constructions the tower is a
\emph{voltage} (derived-graph) tower, hence Galois by construction, and
$\mu,\lambda\ge0$ because they are the Iwasawa invariants of a finitely generated torsion
$\La$-module \cite{Washington}. The non-negativity is stated flatly throughout that
literature.

Separately, the tower of iterated line digraphs is classical. Fiol, Alegre and Yebra
\cite{FAY83,FYA84} introduced line-digraph iteration as \emph{the} design method for dense
interconnection networks, which is how de Bruijn and Kautz networks are built; and the
spanning-tree arithmetic of one iteration is Knuth's \cite{Knuth67}, refined by Levine
\cite{Levine11} to the group statement $\Jac(G)\cong\Jac(LG)/(\text{$q$-torsion})$, with the
de Bruijn and Kautz cases worked out in \cite{BK11,CHP15}.

\paragraph{What is here.}
These two bodies of work have not met. We put the line-digraph iteration into the form
\eqref{Net1:eq:law} and find
\begin{equation}\label{Net1:eq:minus1}
\mu=\frac{n(q+1)}{q},\qquad \lambda=-1,\qquad \nu=\text{const},
\end{equation}
on every base graph tested (Section~\ref{Net1:sec:measure}). Since $\lambda<0$ is impossible for
a torsion $\La$-module, this is a certificate (Section~\ref{Net1:sec:cert}). We show in
Section~\ref{Net1:sec:where} that \eqref{Net1:eq:minus1} telescopes from the classical recursion, so
the arithmetic is not new --- only its reading is --- and we give the Euler-characteristic
interpretation of the $-1$. Section~\ref{Net1:sec:net} draws the network consequence.

\paragraph{What we do not claim.}
Not \eqref{Net1:eq:law} in the abelian case \cite{Val21,MV24,Gonet22}; not the line-digraph
tower as an object \cite{FAY83,FYA84}; not its spanning-tree counts
\cite{Knuth67,Levine11,BK11,CHP15}. A referee will observe that \eqref{Net1:eq:minus1} follows
from \cite{Knuth67,Levine11}, and will be right; we say so in Section~\ref{Net1:sec:where} and
verify it.

\section{Setting}\label{Net1:sec:setting}

A \emph{tower} is a sequence of finite connected graphs with maps
$Y_{k+1}\to Y_k$. We consider four generating mechanisms.
\begin{enumerate}
\item[(a)] \emph{Abelian}: $Y_k$ is the derived graph of a voltage assignment in
$\Z/\ell^{k}$, so the tower is Galois with group $\Z/\ell^{k}$.
\item[(b)] \emph{Rank two}: voltages in $(\Z/\ell)^{k}$.
\item[(c)] \emph{Non-abelian but normal}: voltages in the dihedral group $D_{\ell^{k}}$.
\item[(d)] \emph{Non-normal}: $Y_{k+1}=L(Y_k)$, the line digraph, starting from the
non-backtracking digraph of a base graph. Level-$k$ vertices are the non-backtracking
$k$-paths; this is the Fiol--Alegre--Yebra iteration \cite{FAY83}.
\end{enumerate}
For (a)--(c) $|\Jac(Y_k)|$ is the number of spanning trees; for (d), which is a digraph
tower, it is the number of spanning arborescences. Both are computed exactly as integer
determinants.

\section{Measurements}\label{Net1:sec:measure}

Table~\ref{Net1:tab:main} reports $(\mu,\lambda,\nu)$ fitted from the last three levels of each
tower --- fitting from the tail is deliberate, since the defects of \eqref{Net1:eq:law} are
finitely supported [Ch.~\ref{chap:IX}] and the first levels may be transient --- together with
the number of levels on which the fit holds.

\begin{table}[ht]
\centering\small
\begin{tabular}{llrr}
\toprule
mechanism & instances & $\lambda$ observed & law holds?\\
\midrule
(a) abelian $\Z/2^{k}$ & 15 towers, 6 bases & $1,3,5,9$ (all $>0$) & yes\\
(b) rank two $(\Z/2)^{k}$ & 3 towers & --- & \textbf{no} (as required)\\
(c) dihedral $D_{2^{k}}$ & 4 towers & $3$ & yes\\
(d) line digraph & 5 bases & $\mathbf{-1}$ (every base) & yes\\
\bottomrule
\end{tabular}
\caption{The four mechanisms. For (b) the one-variable law \eqref{Net1:eq:law} fails on all
three towers, as it must: a $\Zl^{2}$-extension needs the two-variable theory of
\cite{CM81,DV22,KM24}. All entries exact.}
\label{Net1:tab:main}
\end{table}

\begin{table}[ht]
\centering\small
\begin{tabular}{lrrrrl}
\toprule
base $Y_0$ & $n$ & $q$ & $\mu$ measured & $n(q{+}1)/q$ & $\ord_q|\Jac(Y_k)|$, $k=1,2,\dots$\\
\midrule
$K_4$          &  4 & 2 &  6 &  6 & $7,18,41,88,183,374$\\
$K_{3,3}$      &  6 & 2 &  9 &  9 & $5,22,57,128,271$\\
$Q_3$ (cube)   &  8 & 2 & 12 & 12 & $14,37,84,179$\\
Petersen       & 10 & 2 & 15 & 15 & $13,42,101,220$\\
$C_4\times K_2$&  8 & 2 & 12 & 12 & $14,37,84,179$\\
\bottomrule
\end{tabular}
\caption{The non-normal (line-digraph) towers. In every case $\lambda=-1$ and $\mu$ equals
$n(q{+}1)/q$ exactly, and the fitted law reproduces every computed level. All entries
exact.}
\label{Net1:tab:ld}
\end{table}

The contrast in Table~\ref{Net1:tab:main} is the point, and it is not the contrast we expected.
Negative $\lambda$ does not track non-commutativity: the dihedral towers are non-abelian and
give $\lambda=3>0$. It tracks \emph{normality}. Every tower in the graph-Iwasawa literature
is a voltage tower and hence Galois; the line-digraph tower is not, and it is the only one
that goes negative.

\section{The certificate}\label{Net1:sec:cert}

\begin{theorem}\label{Net1:thm:cert}
Suppose a tower satisfies \eqref{Net1:eq:law} with $\lambda<0$. Then the tower is not an
abelian $\ell$-tower: there is no voltage assignment in $\Zl$ whose derived graphs are the
$Y_k$.
\end{theorem}

\begin{proof}
For an abelian $\ell$-tower the inverse limit
$X_\infty=\varprojlim_k\Jac(Y_k)\otimes\Zl$ is a finitely generated torsion
$\La=\Zl[[T]]$-module [Ch.~\ref{chap:IV}], \cite{Gonet22,MV24}, and its Iwasawa invariants
satisfy $\mu=\sum_ik_i\ge0$ and $\lambda=\sum_jm_j\deg g_j\ge0$ by the structure theorem
\cite[Thm.~13.12]{Washington}; these are the coefficients in \eqref{Net1:eq:law}.
\end{proof}

\begin{remark}[What the certificate does not say]\label{Net1:rem:caveat}
Theorem~\ref{Net1:thm:cert} rules out the characteristic ideal of a \emph{single} torsion
$\La$-module. It does not rule out a virtual object --- a formal difference of two
$\La$-modules, or the cohomology of a complex --- and indeed the interpretation in
Section~\ref{Net1:sec:where} is exactly of that kind. It is therefore wrong to say that
$\lambda=-1$ contradicts Iwasawa theory; the correct statement is that it is the invariant
of a difference, not of a module. Two further limitations: the certificate is
one-directional (Table~\ref{Net1:tab:main}(c) shows non-abelian towers with $\lambda>0$), and it
presupposes that \eqref{Net1:eq:law} holds at all, which for rank-two towers it does not.
\end{remark}

\begin{remark}[A parity coincidence]\label{Net1:rem:parity}
Kataoka \cite{Kataoka26} shows that every pair $(\lambda,\mu)$ with $\lambda$ \emph{odd} is
realised by an unramified $\Zl$-covering of a bouquet. Our $\lambda=-1$ is odd. We have no
reason to think this is more than a coincidence, but it is the kind of thing a reader
notices, so we record it.
\end{remark}

\section{Where the $-1$ comes from}\label{Net1:sec:where}

\begin{proposition}\label{Net1:prop:telescope}
Let $Y_1$ be a $q$-regular Eulerian digraph on $n_1$ vertices and $Y_{k+1}=L(Y_k)$. The
classical line-digraph tree formula \cite{Knuth67}, in the form used in
{\rm[Ch.~\ref{chap:III}]} and equivalent to Levine's $\Jac(G)\cong\Jac(LG)/(\text{$q$-torsion})$
{\rm\cite{Levine11}}, gives
\[
\ord_q\bigl|\Jac(Y_{k+1})\bigr|-\ord_q\bigl|\Jac(Y_k)\bigr|=n_k(q-1)-1 .
\]
With $n_k=n(q{+}1)q^{k-1}$ this telescopes to \eqref{Net1:eq:minus1}.
\end{proposition}

\begin{proof}
Summing the increment from level $1$ to level $k$ gives
$n(q{+}1)(q^{k-1}-1)-(k-1)+\ord_q|\Jac(Y_1)|$, which is
$\frac{n(q+1)}{q}q^{k}-k+\text{const}$.
\end{proof}

So $\lambda=-1$ is a consequence of a $1967$ formula. We verify the telescoping numerically
on all five bases (check (K) of Section~\ref{Net1:sec:battery}): the measured sequences of
Table~\ref{Net1:tab:ld} agree term by term with the telescoped prediction.

\paragraph{The Euler characteristic.}
The $-1$ in the increment is not a normalisation artefact. The one-step increment is
$n_k(q-1)-\chi_k$ with $\chi_k=\dim\ker\Delta^{\rm dir}_k=1$: at every level a
one-dimensional kernel --- the constant, or Eisenstein, direction --- is removed
[Ch.~\ref{chap:III}]. Summing a constant Euler characteristic over $k$ levels is what produces
a term linear in $k$ with negative coefficient, and it is why no single $\La$-module can
have this as its characteristic ideal: the object is a difference, one term of which is the
Eisenstein line. Check (E) verifies the increment identity level by level.

\section{The network reading}\label{Net1:sec:net}

The two mechanisms of Table~\ref{Net1:tab:main} are not artificial. Torus, circulant and
chordal-ring interconnects are voltage covers of a small base, hence abelian towers; de
Bruijn and Kautz interconnects are built by line-digraph iteration \cite{FAY83,FYA84,FL92},
hence non-normal towers. The quantity being measured, $|\Jac(Y_k)|$, is the spanning-tree
count, a standard reliability measure \cite{Boesch86,Colbourn87,BH09}. Equation
\eqref{Net1:eq:law} is therefore a three-parameter scaling law for a reliability measure along a
design family, in which
\begin{itemize}
\item $\mu$ is the bulk growth rate --- for the line-digraph family, $n(q{+}1)/q$, i.e.\
proportional to the level-one size;
\item $\lambda$ records the design method: positive for covering constructions, $-1$ for
line-digraph iteration;
\item $\nu$ is an offset with no evident structural meaning.
\end{itemize}
We are not aware of prior use of an algebraic growth exponent for this purpose; the closest
precedent is the use of renormalisation exponents to infer a growth mechanism
\cite{SHM06}. Whether the law survives the perturbations a real network family suffers ---
irregularity, missing edges, coarse-grainings that are not covers --- we have not tested,
and Section~\ref{Net1:sec:scope} lists it first among the open questions.

\section{Verification}\label{Net1:sec:battery}

The script \texttt{tower\_lambda.py} accompanying this note builds and measures $27$ towers
and runs the checks of Table~\ref{Net1:tab:battery} in about $7$ minutes, all passing; its output
is archived as \texttt{tower\_lambda\_output.txt}. All spanning-tree and arborescence counts
are exact integer determinants.

\begin{table}[ht]
\centering\small
\begin{tabular}{clc}
\toprule
key & check & mode\\
\midrule
(A) & \eqref{Net1:eq:law} holds with $\lambda>0$ on 15 abelian towers over 6 bases & exact\\
(L) & $\lambda=-1$ and $\mu=n(q{+}1)/q$ on 5 line-digraph towers (Table~\ref{Net1:tab:ld}) & exact\\
(E) & the increment is $n_k(q-1)-\chi$ with $\chi=1$, level by level & exact\\
(K) & Prop.~\ref{Net1:prop:telescope}: the measured sequence equals the telescoped recursion & exact\\
(R) & the one-variable law fails on all 3 rank-two $(\Z/2)^{k}$ towers & exact\\
(D) & dihedral towers: $\lambda=3>0$, so the certificate detects normality & exact\\
(C) & the certificate assembled over the whole battery & exact\\
\bottomrule
\end{tabular}
\caption{The verification battery (\texttt{tower\_lambda.py}), $7/7$ passing.}
\label{Net1:tab:battery}
\end{table}

\section{Scope}\label{Net1:sec:scope}

\emph{Quoted.} The abelian law \cite{Val21,MV24,Gonet22,KM24}; non-negativity of Iwasawa
invariants \cite{Washington}; the line-digraph iteration \cite{FAY83,FYA84}; its
tree arithmetic \cite{Knuth67,Levine11,BK11,CHP15}; spanning trees as a reliability measure
\cite{Boesch86,Colbourn87,BH09}.

\emph{New here, as far as we can determine.} Putting the line-digraph iteration in the form
\eqref{Net1:eq:law}; the value $\lambda=-1$, there being no negative $\lambda$ anywhere in the
graph-Iwasawa literature; the Euler-characteristic reading of it; the certificate
(Theorem~\ref{Net1:thm:cert}) with its caveats (Remark~\ref{Net1:rem:caveat}); and the observation
that the certificate separates normal from non-normal rather than abelian from non-abelian.

\emph{Open.} (a) Robustness. Every tower here is exact; real network families are not, and
we do not know how \eqref{Net1:eq:law} degrades under edge noise or under refinements that are
not covers. This is the question that decides whether any of this is usable. (b) Whether
$\lambda<0$ occurs for any tower other than the line-digraph family --- we found none, but
tested only four mechanisms. (c) Whether the structural refinement of $\lambda$ as a count
of persistently growing cyclic factors is already contained in \cite{Gonet22} or in
Kataoka's work on Fitting ideals; we were not able to settle this. (d) The rank-two
non-normal case, where [Chs.~\ref{chap:X}--\ref{chap:XII}] shows no abelian sector exists at all.

\section*{Provenance of this chapter}

Unrefereed preprint. Theorem~\ref{Net1:thm:cert} is a restatement of the structure theorem for
$\La$-modules and is proved above modulo that theorem. Proposition~\ref{Net1:prop:telescope} is
a summation of a classical formula. Everything numerical is produced by the accompanying
script in exact integer arithmetic. We searched for a negative $\lambda$, for a non-normal
graph tower with a growth law, and for an Euler-characteristic reading of such a law, and
found none; a search cannot establish a negative, and our coverage of engineering
proceedings in particular is incomplete.

% generated from Ctrl 1/hidden_eigs.tex
\chapter{Power-sum lower bounds on hidden eigenvalues in positive realizations}\label{chap:Ctrl1}
\chaptermark{Ctrl1}
\begin{chapterabstract}
\noindent
A positive realization of a transfer function of order $n$ is a nonnegative matrix of some
size $N\ge n$ carrying the $n$ poles among its eigenvalues; the $s=N-n$ remaining
eigenvalues are the \emph{hidden} eigenvalues, whose characterization Benvenuti and Farina
list as an open problem. We give a lower bound on $s$ from the power sums of the poles: if
$p_k$ denotes the $k$-th power sum and $\rho$ the spectral radius, then
$s\ge\max_{k\ge1}\lceil[-p_k]_+/\rho^k\rceil$. At $k=1$ this reproduces the quantity
$\zeta$ of Benvenuti's lower bound; for $k\ge2$ it is, as far as we can determine, new, and
it is strictly stronger. On the family $\M_m(\theta)=\{1\}\cup m\cdot\{\pm i\theta\}$ the
published lower bound is vacuous ($N\ge n$) while ours gives $N\ge 2m(1+\theta^{2})$; when
$\theta^{2}>1-1/2m$ this meets Benvenuti's own upper bound $4m$ and therefore determines
the minimal size exactly, where previously only $n\le N\le 4m$ was known. The bound is
attained: on $\{1\}\cup m\cdot\{\pm i\}$ it gives $N\ge4m$, and a direct sum of $m$ copies
of the $4$-cycle permutation matrix realizes $N=4m$. We also record two limitations. The
method has a hard ceiling --- it can never prove more than $N\le2n$ --- so it is
complementary to, not competitive with, the digraph bounds that give unbounded ratios; and
the ratio $2$ itself is not new territory, being exactly what the Benvenuti--Farina family
of 1999 achieves. Finally we explain why no such bound can exist in the
Boyle--Handelman/Kim--Ormes--Roush setting: there one appends only zeros, and zeros change
no power sum. All claims are machine-verified in exact rational arithmetic.
\end{chapterabstract}

\section{Introduction}\label{Ctrl1:sec:intro}

Let $H(z)$ be a rational transfer function of order $n$ with pole multiset
$\Lambda=\{\lambda_1,\dots,\lambda_n\}$ and spectral radius
$\rho=\max_i|\lambda_i|$. A \emph{positive realization} of $H$ is a triple
$(A,b,c)$ with $A\in\R_{\ge0}^{N\times N}$, $b,c\ge0$ and
$H(z)=c(zI-A)^{-1}b$. Existence is settled \cite{Farina96,ADFB96,OMK84}, and a clean
criterion is available in terms of an invariant polyhedral cone \cite{BF04}. What is not
settled is \emph{size}. The poles of $H$ are a subset of the eigenvalues of any positive
realization, and the remaining $s=N-n$ eigenvalues --- the \emph{hidden eigenvalues} --- are
the reason $N$ may exceed $n$. Benvenuti and Farina put it as follows
\cite[\S IX]{BF04}: ``An important issue related to minimality of positive systems is the
study of `hidden' eigenvalues, i.e.\ of the eigenvalues which are not poles and that
possibly one has to add in order to obtain a minimal positive realization. A full
characterization of this property may lead to a deeper and valuable insight into the
problem.'' The minimality problem remains open; a recent statement of its status is
\cite{TS25}.

This note contributes one inequality. It is elementary, and its only merit is that it
appears not to have been written down beyond its first instance.

\paragraph{The bound.} Write $p_k=p_k(\Lambda)=\sum_{i}\lambda_i^{k}$ for the $k$-th power
sum of the poles and $[x]_+=\max(x,0)$. Theorem~\\ref{Ctrl1:thm:main} states
\begin{equation}\label{Ctrl1:eq:bound}
s\;\ge\;\Bigl\lceil\;\max_{k\ge1}\;\frac{[-p_k]_+}{\rho^{k}}\;\Bigr\rceil .
\end{equation}

\paragraph{What is already known.} The case $k=1$ of \\eqref{Ctrl1:eq:bound} is published.
Benvenuti \cite[Thm.~4.7]{Ben20ELA} proves $N\ge N_D+\sum_{|s|<1}m(s)+\zeta$ with
$\zeta=\lceil-m_{\mathcal R_1}-\sum_{|s|<1}s\,m(s)\rceil$ when that quantity is positive,
and the proof is exactly the $k=1$ argument: ``since the trace of a nonnegative matrix is
nonnegative \dots\ since each one of these additional eigenvalues has to be not greater
than $1$, then their number is not lesser than $\zeta$''
\cite[pp.~379--380]{Ben20ELA}. The companion paper \cite{Ben20SCL} transfers the same
quantity to the realization dimension. We have found no power sum beyond the first
anywhere in this literature: not in \cite{Ben20ELA,Ben20SCL}, not in the digraph bound of
Nagy and Matolcsi \cite{NM03}, not in \cite{BF04,TS25}. Section~\\ref{Ctrl1:sec:improve} shows
that $k=2$ already does strictly better.

\paragraph{What we do not claim.} Two disclaimers, stated here rather than buried.
\begin{enumerate}
\item[(a)] The ratio $N/n=2$ is \emph{not} new territory. The family
$H(z,m)=1/\bigl((z-1)(z^{2^{m}}+1)\bigr)$ of Benvenuti and Farina \cite{BF99} has
$n=2^{m}+1$ and minimal positive realization of dimension $N=2^{m+1}=2n-2$, so ratio $2$ is
exactly what they already achieve; and Nagy and Matolcsi \cite{NM03} obtain \emph{unbounded}
ratios, with $n=4$ fixed and $N\ge m+1$ for arbitrary $m$.
\item[(b)] Worse, ratio $2$ is a hard ceiling for the present method.
Since $p_k\ge-\sum_i|\lambda_i|^{k}\ge-n\rho^{k}$, the right-hand side of \\eqref{Ctrl1:eq:bound}
never exceeds $n$, so \\eqref{Ctrl1:eq:bound} can never prove more than $N\le 2n$
(Proposition~\\ref{Ctrl1:prop:ceiling}). The bound is therefore complementary to the combinatorial
bounds of \cite{NM03}, not a competitor to them.
\end{enumerate}
What we do claim is (i) the inequality for $k\ge2$, (ii) that it strictly improves the
published lower bound, and (iii) that on an explicit family it closes the published gap
completely.

\section{The bound}\label{Ctrl1:sec:bound}

\begin{theorem}\label{Ctrl1:thm:main}
Let $A\in\R_{\ge0}^{N\times N}$ and let $\Lambda$ be a sub-multiset of $\operatorname{spec}A$
of size $n$ containing an eigenvalue of modulus $\rho(A)$; write $\rho=\rho(A)$ and
$s=N-n$. Then for every $k\ge1$
\[
s\;\ge\;\frac{-p_k(\Lambda)}{\rho^{k}},
\qquad\text{hence}\qquad
s\;\ge\;\Bigl\lceil\max_{k\ge1}\frac{[-p_k(\Lambda)]_+}{\rho^{k}}\Bigr\rceil .
\]
\end{theorem}

\begin{proof}
Let $\M$ be the complementary multiset, $|\M|=s$, so that
$\operatorname{spec}A=\Lambda\uplus\M$. Every eigenvalue of $A$ has modulus at most
$\rho(A)=\rho$, so $\bigl|\sum_{\mu\in\M}\mu^{k}\bigr|\le\sum_{\mu\in\M}|\mu|^{k}\le
s\rho^{k}$. Since $A$ is nonnegative, $\Tr(A^{k})\ge0$: it is a sum of products of
nonnegative entries. Therefore
\[
0\;\le\;\Tr(A^{k})\;=\;p_k(\Lambda)+\sum_{\mu\in\M}\mu^{k}\;\le\;p_k(\Lambda)+s\rho^{k},
\]
which rearranges to the claim; $s$ is an integer, so the ceiling may be taken.
\end{proof}

The hypothesis that $\Lambda$ contains a dominant eigenvalue is what makes $\rho$ computable
from the poles alone; in the realization setting it holds because $\rho(H)$ is a pole
\cite[Lem.~7]{BF04}.

\begin{proposition}[$k=1$ is Benvenuti's $\zeta$]\label{Ctrl1:prop:k1}
At $k=1$, \\eqref{Ctrl1:eq:bound} reads $s\ge\lceil-p_1(\Lambda)/\rho\rceil$, which after
normalising $\rho=1$ and splitting $\Lambda$ into its dominant and non-dominant parts is
the quantity $\zeta$ of {\rm\cite[Thm.~4.7]{Ben20ELA}}.
\end{proposition}

\begin{proof}
Normalise $\rho=1$. The dominant part of $\Lambda$ contributes $m_{\mathcal R_1}$ to
$p_1$, since the $h$-th roots of unity sum to zero for $h>1$
\cite[proof of Thm.~4.7]{Ben20ELA}, and the non-dominant part contributes
$\sum_{|s|<1}s\,m(s)$. Hence $-p_1=-m_{\mathcal R_1}-\sum_{|s|<1}s\,m(s)$, and the ceiling
of its positive part is $\zeta$.
\end{proof}

\begin{proposition}[Ceiling of the method]\label{Ctrl1:prop:ceiling}
The right-hand side of \\eqref{Ctrl1:eq:bound} is at most $n$; consequently
Theorem~\\ref{Ctrl1:thm:main} can never certify $N>2n$.
\end{proposition}

\begin{proof}
$-p_k\le\sum_i|\lambda_i|^{k}\le n\rho^{k}$.
\end{proof}

\section{Strict improvement, and an exact determination}\label{Ctrl1:sec:improve}

Fix $0<\theta<1$ and $m\ge1$, and let
\[
\M_m(\theta)=\{1\}\;\cup\;m\ \text{copies of}\ \{\,i\theta,\,-i\theta\,\},
\qquad n=2m+1,\ \rho=1 .
\]
Only $1$ is dominant, so in the notation of \cite{Ben20ELA} we have $r=1$,
$N_D=1$ and $m_{\mathcal R_1}=1$.

\begin{theorem}\label{Ctrl1:thm:improve}
For $\M_m(\theta)$:
\begin{enumerate}
\item[(i)] the lower bound of {\rm\cite[Thm.~4.7]{Ben20ELA}} gives $N\ge 2m+1=n$, i.e.\ it is
vacuous, because $\zeta=0$;
\item[(ii)] Theorem~\\ref{Ctrl1:thm:main} at $k=2$ gives
$s\ge\lceil 2m\theta^{2}-1\rceil$, hence $N\ge 2m+1+\lceil2m\theta^{2}-1\rceil$, which
exceeds $n$ whenever $\theta^{2}>1/2m$;
\item[(iii)] the upper bound of {\rm\cite[Thm.~4.4]{Ben20ELA}} gives $N\le 4m$;
\item[(iv)] if $\theta^{2}>1-\dfrac{1}{2m}$ then the bounds in (ii) and (iii) coincide and
\[
N_{\min}=4m .
\]
\end{enumerate}
\end{theorem}

\begin{proof}
(i) $p_1(\M_m(\theta))=1$ and the non-dominant numbers cancel in conjugate pairs, so
$m_{\mathcal R_1}+\sum_{|s|<1}s\,m(s)=1\ge0$ and $\zeta=0$; the bound reads
$N\ge N_D+2m+0=2m+1$. (ii) $p_2=1+m\bigl((i\theta)^{2}+(-i\theta)^{2}\bigr)=1-2m\theta^{2}$,
so $[-p_2]_+/\rho^{2}=2m\theta^{2}-1$ when positive. (iii) is \cite[Thm.~4.4]{Ben20ELA}
with $\kappa(i\theta)=4$, valid since $i\theta\in\Pi_4\setminus\Pi_3$ for
$\theta>1/\sqrt3$. (iv) The bound in (ii) equals $4m$ as soon as
$\lceil2m\theta^{2}-1\rceil\ge2m-1$, i.e.\ $2m\theta^{2}-1>2m-2$, i.e.\
$\theta^{2}>1-1/2m$.
\end{proof}

Table~\\ref{Ctrl1:tab:improve} shows the three quantities. The criterion $\theta^{2}>1-1/2m$ of
Theorem~\\ref{Ctrl1:thm:improve}(iv) predicts exactly which rows close: it holds for
$\theta=9/10$ only at $m=1$ ($0.81>0.5$) and fails at $m=4$ ($0.81<0.875$), where the bound
indeed falls one short of the upper bound; it holds throughout for $\theta=99/100$ from
$m=2$ on and for $\theta=999/1000$, and there the gap closes and $N$ is determined.

\begin{table}[ht]
\centering\small
\begin{tabular}{rrrrrr}
\toprule
$\theta$ & $m$ & $n$ & \cite[Thm.~4.7]{Ben20ELA} & Theorem~\\ref{Ctrl1:thm:main} & \cite[Thm.~4.4]{Ben20ELA}\\
 & & & lower & lower & upper\\
\midrule
$9/10$      &  1 &  3 &  3 &  \textbf{4} &  4\\
$9/10$      &  4 &  9 &  9 & 15 & 16\\
$9/10$      & 16 & 33 & 33 & 58 & 64\\
\midrule
$99/100$    &  2 &  5 &  5 &  \textbf{8} &  8\\
$99/100$    &  4 &  9 &  9 & \textbf{16} & 16\\
$99/100$    &  8 & 17 & 17 & \textbf{32} & 32\\
$99/100$    & 16 & 33 & 33 & \textbf{64} & 64\\
\midrule
$999/1000$  &  4 &  9 &  9 & \textbf{16} & 16\\
$999/1000$  & 16 & 33 & 33 & \textbf{64} & 64\\
\bottomrule
\end{tabular}
\caption{Lower and upper bounds on the minimal size $N$ for $\M_m(\theta)$. Bold: the new
lower bound meets the published upper bound, so $N$ is determined exactly. All entries
exact.}
\label{Ctrl1:tab:improve}
\end{table}

\section{Sharpness}\label{Ctrl1:sec:sharp}

Let $\M_m=\{1\}\cup m\cdot\{i,-i\}$, all of modulus one, $n=2m+1$.

\begin{theorem}\label{Ctrl1:thm:sharp}
$p_2(\M_m)=1-2m$, so Theorem~\\ref{Ctrl1:thm:main} gives $s\ge2m-1$ and $N\ge4m$. The bound is
attained: if $C_4$ is the $4$-cycle permutation matrix, with characteristic polynomial
$x^{4}-1$ and spectrum $\{1,i,-1,-i\}$, then $C_4^{\oplus m}$ is nonnegative of size $4m$
and contains $\M_m$ in its spectrum. Hence $N_{\min}=4m=2n-2$.
\end{theorem}

\begin{proof}
The power sum is immediate. For the construction, $\operatorname{spec}C_4^{\oplus m}$ is
$m$ copies of $\{1,i,-1,-i\}$, which contains one $1$ and $m$ copies of each of $\pm i$.
\end{proof}

\begin{remark}\label{Ctrl1:rem:caveat}
Two caveats. First, for $m\ge2$ the multiset $\M_m$ violates the standing hypothesis of
\cite[Thms.~3.2, 4.1]{Ben20ELA} that the spectral radius occur with dominant multiplicity:
here $m(1)=1$ while the dominant multiplicity is $m$. The formula of
\cite[Thm.~4.1]{Ben20ELA} nevertheless evaluates to $N_D=4m$ on $\M_m$, so the conclusion
$N_{\min}=4m$ is also reachable by that route, by an argument through roots of unity rather
than through traces; we note this so that Theorem~\\ref{Ctrl1:thm:sharp} is not mistaken for a new
determination. It is a sharpness check on \\eqref{Ctrl1:eq:bound}, not a new value of $N$. Second,
$4m=2n-2$ is the Benvenuti--Farina ratio of \cite{BF99}, as recalled in
Section~\\ref{Ctrl1:sec:intro}.
\end{remark}

\begin{remark}[Tutorial Example 8]\label{Ctrl1:rem:ex8}
For the poles $\{1,\pm0.9i\}$ of \cite[Ex.~8]{BF04}, $p_2=-31/50$ and \\eqref{Ctrl1:eq:bound}
gives $s\ge1$, i.e.\ $N\ge4$. This is attained by the nonnegative circulant with first row
$(0,\tfrac{19}{20},0,\tfrac1{20})$, whose characteristic polynomial is
$(x-1)(x+1)(x^{2}+\tfrac{81}{100})$: the hidden eigenvalue is $-1$. The example in
\cite{BF04} is $4$-dimensional, so the bound is tight there as well.
\end{remark}

\section{Why there is no analogue for the Boyle--Handelman problem}\label{Ctrl1:sec:bh}

It is worth saying why the same traces behave completely differently one category over.
In the inverse eigenvalue problem solved by Boyle and Handelman over $\R$ \cite{BH91} and by
Kim, Ormes and Roush over $\Z$ \cite{KOR00}, one asks whether $\Lambda$ is the \emph{nonzero}
spectrum of a nonnegative matrix; the matrix may be enlarged, but only by appending
\emph{zeros}. Zeros contribute nothing to any power sum, so the trace conditions
$\operatorname{tr}_k\ge0$ must already hold for $\Lambda$ itself: they are a genuine
obstruction, not a counting device, and no lower bound of the shape \\eqref{Ctrl1:eq:bound} on the
size of the dilation can exist. Equivalently: in that category the analogue of
Theorem~\\ref{Ctrl1:thm:main} is the statement $p_k(\Lambda)\ge0$, with no $s$ in it.

This is the exact point at which results from symbolic dynamics fail to transfer to the
positive realization problem. A negative $p_2$ closes the fixed-nonzero-spectrum problem
outright --- shift equivalence preserves the nonzero spectrum with multiplicities
\cite[\S7.4]{LM95}, so no dilation of any size helps, an argument used in this form in
\cite[II]{ANCFT} --- while in the present setting the same negative $p_2$ merely counts
hidden eigenvalues, and indeed \cite[Ex.~8]{BF04} exhibits a transfer function with
$p_2=-0.62$ and a $4$-dimensional positive realization. The two categories differ exactly by
the freedom to add nonzero eigenvalues.

\section{Verification}\label{Ctrl1:sec:battery}

The script \texttt{hidden\_eigs.py} accompanying this note runs the checks of
Table~\\ref{Ctrl1:tab:battery} in $23$\,s, all passing; its output is archived as
\texttt{hidden\_eigs\_output.txt}. Arithmetic is exact rational throughout; the only
floating-point quantity anywhere is the spectral radius in check (T), and that row is
labelled accordingly.

\begin{table}[ht]
\centering\small
\begin{tabular}{clc}
\toprule
key & check & mode\\
\midrule
(T) & Theorem~\\ref{Ctrl1:thm:main} holds on $120$ random nonnegative integer matrices, $\Lambda$ a rational factor & exact $+$ numeric $\rho$\\
(K) & at $k=1$ the bound equals $\zeta$ of \cite[Thm.~4.7]{Ben20ELA} & exact\\
(I) & $k=2$ strictly exceeds \cite[Thm.~4.7]{Ben20ELA} in $15/15$ instances (Table~\\ref{Ctrl1:tab:improve}) & exact\\
(S) & Theorem~\\ref{Ctrl1:thm:sharp}: bound $=4m$ attained by $C_4^{\oplus m}$, $m\le32$ & exact\\
(E) & Remark~\\ref{Ctrl1:rem:ex8}: bound $=1$, attained by an explicit rational circulant & exact\\
(C) & Proposition~\\ref{Ctrl1:prop:ceiling}: the bound never exceeds $n$, on $300$ random pole sets & exact\\
(Z) & appending zeros leaves every power sum unchanged (Section~\\ref{Ctrl1:sec:bh}) & exact\\
\bottomrule
\end{tabular}
\caption{The verification battery (\texttt{hidden\_eigs.py}), $7/7$ passing.}
\label{Ctrl1:tab:battery}
\end{table}

In check (T) the spectral radius is replaced by a rational \emph{lower} bound, which makes
the tested inequality strictly stronger than the theorem; no violation occurred. The bound
was attained in none of those $120$ random instances, which is the expected behaviour: like
the $k=1$ bound it is weak on generic spectra and bites only when the poles are strongly
oscillatory.

\section{Scope}\label{Ctrl1:sec:scope}

\emph{Established here.} Theorem~\\ref{Ctrl1:thm:main} for all $k$; its coincidence with
\cite[Thm.~4.7]{Ben20ELA} at $k=1$; strict improvement and an exact determination of
$N_{\min}$ on $\M_m(\theta)$ for $\theta^{2}>1-1/2m$; sharpness on $\M_m$ and on
\cite[Ex.~8]{BF04}; the ceiling $N\le2n$; the structural reason the bound has no analogue
in \cite{BH91,KOR00}.

\emph{Open.} (a) Whether power sums can be combined with the cone criterion \cite{BF04} or
with the digraph bound \cite{NM03} to pass the ratio-$2$ ceiling. (b) Whether the
Johnson--Loewy--London inequalities \cite{LoewyLondon78}, which mix power sums of different
orders, give a bound not implied by \\eqref{Ctrl1:eq:bound}; we have not investigated this and it
is the obvious next step. (c) Whether the bound is ever tight outside the families of
Sections~\\ref{Ctrl1:sec:improve}--\\ref{Ctrl1:sec:sharp}; it was not tight in any of our $120$ random
instances.

\emph{Not claimed.} Ratio $2$, which is \cite{BF99}; unbounded ratios, which are
\cite{NM03}; and the $k=1$ case, which is \cite{Ben20ELA}.

\section*{Provenance of this chapter}

Unrefereed preprint. Theorem~\\ref{Ctrl1:thm:main} is elementary and its proof is complete above.
The bounds of \cite{Ben20ELA} quoted for comparison in Table~\\ref{Ctrl1:tab:improve} are
evaluated by us from the published statements of Theorems 4.4 and 4.7 of that paper and are
reproduced in the script; we have not reproved them. The symbolic-dynamics facts in
Section~\\ref{Ctrl1:sec:bh} are quoted from \cite{LM95,BH91,KOR00}. We searched the literature for
a power-sum bound with $k\ge2$ and found none; a search cannot establish a negative, and we
state this as the limit of our knowledge.

\part{The state of the programme}
\chapter{Solved, open, impossible}\label{chap:state}

This chapter consolidates the closing remarks of the chapters. Verdicts are quoted, not
re-argued; the chapter and statement numbers point to where the argument is.

Four of the open problems of the first edition of this list have since been settled, three
of them affirmatively and one by a no-go theorem, and a fifth has been settled in a way that
was not on the menu at all. Composite $q$ is proved (Chapter~\ref{chap:XVII}); the rank-$d$
unitarity mechanism is proved, and then made unconditional (Chapters~\ref{chap:XVIII},
\ref{chap:XX}); the search for a cyclic class computing the $L$-invariant is closed by an
obstruction (Chapter~\ref{chap:XIX}); and the bidirected sandpile theory asked for by the
assembly chapters turns out not to exist, for a reason that is itself a theorem
(Chapter~\ref{chap:Bio7}).

\section{Corrections made by later chapters to earlier ones}

\begin{itemize}
\item The untwisted boundary type is $\mathrm{III}_{1/q}$ outright, not a
$\beta$-family: Chapter~\ref{chap:III}, Remark~1.6 corrects Chapter~\ref{chap:I},
Theorem~2.4(iii). The applied note of Chapter~\ref{chap:Holo1} sharpens this further:
$\mathrm{III}_{1/q^2}$ when the graph is bipartite.
\item The line-digraph sandpile recursion is $|\Jac(Y_{k+1})|=q^{\,n_k(q-1)-1}|\Jac(Y_k)|$:
Chapter~\ref{chap:III}, Theorem~2.1 corrects the exponent $2^{n_k-1}$ of
Chapter~\ref{chap:II}, Remark~5.2 (the two agree for $q=2$). This recursion should be
compared with Levine's theorem on sandpile groups of line digraphs (J.\ Combin.\ Theory
Ser.\ A 118 (2011)); the series flags the comparison and does not claim priority.
\item The naive $\Z_p[[T]]$ conjecture for the congruence tower
(Chapter~\ref{chap:II}, Conjecture~3.6(ii-b)) is withdrawn in the same chapter and in
Chapter~\ref{chap:III}; three independent certificates against it are $\lambda=-1$
(Ch.~\ref{chap:III}), $A_kA_k^{\mathsf T}\ne A_k^{\mathsf T}A_k$ (Ch.~\ref{chap:V}), and
the torsion-freeness of the $p$-part of the pro-module with $\mu>0$
(Ch.~\ref{chap:VIII}).
\item The shadow of a $q$-regular level needs $q-1$ return arcs, not $q$:
Chapter~\ref{chap:VIII}, Remark~1.4.
\item The geodesic edge flow must be defined by ``not in a common chamber'', not ``not
adjacent'': Chapter~\ref{chap:XV}, Remark~1.2 corrects Chapter~\ref{chap:XI},
Definition~3.1; the two agree only on flag complexes, and of the five test complexes
only $Y_{168}$ is a flag complex.
\item The $p$-primary hypothesis in the tail-norm theorem was an artefact of the proof, not
of the statement: Chapter~\ref{chap:XVII}, Theorem~\ref{XVII:thm:main} removes it for every
$q\ge2$, because the rank of a line-digraph adjacency matrix is $n$ over \emph{every} field
(Ch.~\ref{chap:XVII}, Lemma~\ref{XVII:lem:rank}), which is what the counting argument needs.
\item The invariant subspace left open at Chapter~\ref{chap:XVIII}, Remark~4.2(a) does
exist: Chapter~\ref{chap:XX}, Corollary~\ref{XX:cor:W}. The candidate intertwiner recorded
there was the right one, and the identity it needs is a single flag count.
\item The expectation, stated in the assembly chapters, that double-strandedness would be
handled by a sandpile theory on a quotient graph is withdrawn:
Chapter~\ref{chap:Bio7}, Theorem~\ref{Bio7:thm:anti}. Chapter~\ref{chap:Bio6} had already
shown that no new lemma is \emph{needed}, by passing to the double cover instead; Chapter
\ref{chap:Bio7} shows that the quotient route is not merely unnecessary but unavailable.
\item Cross-references between the original papers drifted as sections were
renumbered (for instance ``[V, Prob.~5.1]'' is Problem~4.1 of Chapter~\ref{chap:V};
``[I, Prop.~1.7(iii)]'' is Proposition~1.9(iii) of Chapter~\ref{chap:I}). The
concordance and the chapter-wise numbering of this book are authoritative.
\item The type III framing of ``holographic failure'' and the claim that the control
positive-realization problem inherits the trace obstruction were both wrong and are
retracted in the applied chapters themselves (Chapters~\ref{chap:Holo1},
\ref{chap:Ctrl1}).
\end{itemize}

\section{Proved}

\begin{itemize}
\item The boundary channel is arithmetically blind: $K_0(C(\partial X)\rtimes\Gamma)
\cong\Z^r\oplus\Z/(r-1)$ (Ch.~\ref{chap:I}, Thm.~3.7).
\item The trace obstruction: no non-negative matrix, of any size, is shift equivalent to
the Ihara companion (Ch.~\ref{chap:II}, Thms.~1.2, 1.5).
\item The shadow realization: $K_0(\mathcal O_{B(Y)})\cong\Z\oplus\Jac(Y)$, $K_1\cong\Z$,
functorially and at every level of the tower (Ch.~\ref{chap:VI}, Thm.~2.1;
Ch.~\ref{chap:VIII}, Thm.~1.3); the signed-shadow lemma for any integer matrix with
positive diagonal (Ch.~\ref{chap:X}, Lem.~5.1).
\item The abelian limit module $\Lambda^n/D(1+T)$ with exact control
(Ch.~\ref{chap:IV}, Thm.~2.1); the exceptional zero of order two and the formula
$\tfrac12\Theta''(0)=-\kappa(Y_0)\|h_\alpha\|^2$ (Ch.~\ref{chap:IX}, Thm.~1.2); the five
faces of the $L$-invariant (Ch.~\ref{chap:XIV}, Cor.~7.2).
\item The degeneracy calculus and the Euler rank of the degeneracy complex
(Ch.~\ref{chap:V}, Prop.~1.2, Thm.~2.1); the tail norm and its pro-module
(Ch.~\ref{chap:V}, Thm.~3.1); vanishing of the upper Hecke operator on the norm kernel
and the snake factorization (Ch.~\ref{chap:VII}, Lem.~1.1, Thm.~2.1); the Eisenstein
boundary sequence $0\to\Z/n_k\to\coker\Delta_k^0\to\Jac(Y_k)\to0$ and its
non-splitting on $K_4$ (Ch.~\ref{chap:XIII}, Thm.~3.2).
\item \textbf{For every $q\ge2$}, $\ker(\tau_*|\Jac(Y_{k+1}))=\Jac(Y_{k+1})[q]\cong
(\Z/q)^{r_k}$, and every invariant factor divisible by a prime $p\mid q$ is divisible by
$p^{v_p(q)}$ (Ch.~\ref{chap:XVII}, Thm.~\ref{XVII:thm:main}); hence the pro-module has no
$q$-torsion (Ch.~\ref{chap:XVII}, Cor.~\ref{XVII:cor:notorsion}) and the sequence
$0\to\mathfrak X_\infty^\vee\to K_0(\mathcal O_\infty)\to\Z[1/q]\to0$ splits for every $q$
(Ch.~\ref{chap:XVII}, Thm.~\ref{XVII:thm:split}). The $p$-prime case is
Ch.~\ref{chap:VIII}, Thm.~2.8; the inductive limit is Kirchberg with
$\operatorname{Tor}K_0$ dual to the pro-module (Ch.~\ref{chap:VIII}, Thms.~3.2, 3.3); the
head correspondence with $K_0$-index $\eta_*$ and the Hecke round trips
(Ch.~\ref{chap:XVI}, Thms.~2.3, 3.1, 4.1).
\item In rank two: no cubic trace obstruction (Ch.~\ref{chap:X}, Thm.~4.2); $b_1=0$ for
thick quotients, hence no abelian towers (Ch.~\ref{chap:XII}, Thm.~1.1); the covering
norm theorem and the incidence calculus (Ch.~\ref{chap:XII}, Thms.~2.1, 3.1); the
edge-flow intertwiner $L_E\Psi=\Psi C'$ (Ch.~\ref{chap:XI}, Thm.~3.3); the two-sided
reduction $\ell^2(E)=\operatorname{im}\Psi\oplus W$ and the unitarity of the remainder,
$RR^{\mathsf T}=qI$ (Ch.~\ref{chap:XV}, Thms.~2.1, 2.2, 3.1).
\item \textbf{In every rank $d\ge2$}: the unitarity identity
$L_EL_E^{\mathsf T}=q^{\,d-1}I+q^{\,d-1}(q-1)T^{\mathsf T}T$ and its dual
(Ch.~\ref{chap:XVIII}, Thm.~\ref{XVIII:thm:main}), whose whole content is that two
hyperplanes of $\mathrm{PG}(d,q)$ always meet in codimension two; exactly $|E|-n$ singular
values equal $q^{(d-1)/2}$ and $L_E:\ker S\to\ker T$ is a bijective isometry
(Ch.~\ref{chap:XVIII}, Cor.~\ref{XVIII:cor:uncond}). The invariant subspace that makes the
critical circle unconditional is supplied in Ch.~\ref{chap:XX}: the number of complete flags
of $\mathrm{PG}(d,q)$ with $F_{i+1}=V$ and $F_1$ inside a hyperplane $H$ takes exactly two
values according as $V\subseteq H$ or not (Thm.~\ref{XX:thm:flag}), which \emph{is} the
intertwining $L_EN_i^{\mathsf T}=q^{\,i}(N_0^{\mathsf T}A_{i+1}-N_{i+1}^{\mathsf T})$
(Thm.~\ref{XX:thm:inter}); hence $W_d=(\operatorname{im}\Psi_d)^\perp$ is two-sided
invariant and every eigenvalue of $L_E|_{W_d}$ has modulus $q^{(d-1)/2}$ with no hypothesis
whatever (Cor.~\ref{XX:cor:W}). The companion matrix is explicit in every rank and its
characteristic polynomial is the local Hecke polynomial of $\mathrm{PGL}_{d+1}$
(Thm.~\ref{XX:thm:hecke}).
\item In the applications: the bundle-chain splitting theorem
$\K(G)\cong\bigoplus(\Z/r_c)^{\ell_c-1}\oplus\K(\mathrm{Sk})$ and the canonicity of the
skeleton (Ch.~\ref{chap:Bio3}); the two-family theorem $\K(D_w)=\Z/c$
(Ch.~\ref{chap:Bio4}); the power-sum bound on hidden eigenvalues
(Ch.~\ref{chap:Ctrl1}); the sign of $\lambda$ as a normality certificate
(Ch.~\ref{chap:Net1}).
\item In double-stranded assembly: the reverse-complement double cover is balanced, so the
whole single-strand apparatus applies to it verbatim and the five plus-strand and two
minus-strand ribosomal operons of \emph{E.\ coli} merge into one seven-copy family
(Ch.~\ref{chap:Bio6}, Thm.~\ref{Bio6:thm:merge}, Cor.~\ref{Bio6:cor:rrna}); the quotient is
a signed graph exactly when no repeated $k$-mer is its own reverse complement, which is
automatic for odd $k$ (Ch.~\ref{chap:Bio7}, Prop.~\ref{Bio7:prop:parity}); it is unbalanced
exactly when the sequence carries an inverted repeat of length at least $k$
(Thm.~\ref{Bio7:thm:balance}); and the Reiner--Tseng covering sequence holds on it, with a
non-split $2$-primary part arising from a real bacteriophage (Thm.~\ref{Bio7:thm:rt}).
\end{itemize}

\section{Verified computationally, not proved}

The parahoric factorization on $Y_{168}$ (Ch.~\ref{chap:XI}, Cor.~3.5) was numeric
before Ch.~\ref{chap:XV} proved it in general; the Ramanujan property of the five
$\widetilde A_2$ test complexes is numeric to $10^{-6}$ (Ch.~\ref{chap:X}, Thm.~4.5(i));
the character mechanism for kernel orders along finite $\widetilde A_2$ chains is
verified to $10^{-6}$ (Ch.~\ref{chap:XII}, Thm.~2.1(iii)); finiteness of
$K_0(\mathcal O_{L_E})$ rests on $1\notin\operatorname{spec}L_E$ on $Y_{168}$
(Ch.~\ref{chap:XI}, Thm.~3.7); the non-commutativity certificate
$A_kA_k^{\mathsf T}\ne A_k^{\mathsf T}A_k$ is checked on $K_4$ only (Ch.~\ref{chap:V},
Prop.~3.4); the exhaustive statements about chord diagrams hold for $m\le5$
(Ch.~\ref{chap:Bio3}, Rem.~9) and are known to fail at $m=6$.

Two items are new. The corank of $\Psi_d$ is at least $d(d+1)$ by
Ch.~\ref{chap:XX}, Theorem~\ref{XX:thm:corank}, and equals it on all three thick
$\widetilde A_2$ complexes tested; equality for $d\ge3$ is unproved, and with it the exact
value of $\dim W_d$. And a signed $r$-bundle chain of $s$ internal vertices raises the
$r$-rank by exactly $s$ in every case tested (Ch.~\ref{chap:Bio7},
Thm.~\ref{Bio7:thm:series}); only the inequality is proved.

\section{Open}

\begin{enumerate}
\item \textbf{The main conjecture for the degeneracy algebra.} Construct, from
$L$-values alone, an element $L_H\in K_1(H_\Sigma)$ equal to the characteristic class
$[\Theta_H]$ modulo the image of $K_1(H)$ (Ch.~\ref{chap:XIII}, Rem.~5.1); before that,
determine the image of $K_1(H)$, i.e.\ how non-unique the class is, and whether
$H\to H_\Sigma$ is stably flat (Ch.~\ref{chap:XIII}, Rem.~5.2). This is Problem~4.1 of
Ch.~\ref{chap:V} in its final form, and it is now the oldest unsettled question in the
programme.
\item \textbf{Rank two realization.} A non-negative integer matrix shift equivalent
over $\Z$ to the graded companion $C^{(3)}$ --- an instance of the weak Generalized
Spectral Conjecture of Boyle--Handelman (Ch.~\ref{chap:XI}, Rem.~2.4).
\item \textbf{Congruence towers in rank two.} The growth law of $\Jac^H$ along a genuine
infinite congruence tower of $\widetilde A_2$ complexes; no second level has been
computed (Ch.~\ref{chap:XII}, Rem.~6.1).
\item \textbf{The size of the unitary part.} Whether
$\operatorname{corank}\Psi_d=d(d+1)$ for $d\ge3$, and whether $W_d$ is the \emph{largest}
two-sided invariant subspace inside $\ker S\cap\ker T$ (Ch.~\ref{chap:XX},
Rem.~\ref{XX:rem:soi}). Also a thick $\widetilde A_3$ complex small enough to compute with,
and an exact certificate of the vertex Ramanujan property of the test complexes
(Ch.~\ref{chap:XV}, Rem.~6.1(c)).
\item \textbf{Counting double-stranded reconstructions.} The reverse-complement involution
induces a perfect \emph{symmetric} linking form on $\K(D_k)$
(Ch.~\ref{chap:Bio7}, Rem.~\ref{Bio7:rem:pairing}). At tree level it is a bijection from
the arborescences converging on a root to those diverging from the image root, not an
action, so there are no orbits to sum over; and for $\phi$X174 the group order is not a
square, so there is no metabolizer either. What a double-stranded count is built from is
open.
\item \textbf{Non-pairwise arrangement invariants.} Any presentation of the arrangement
group for repeats with three or more copies that is not built from pairwise data
(Ch.~\ref{chap:Bio4}, Open~(a)); and the even-$k$ case of the quotient, which is a signed
graph with half-edges branched over the palindromic repeated $k$-mers
(Ch.~\ref{chap:Bio7}, Rem.~\ref{Bio7:rem:halfedge}).
\end{enumerate}

\section{Impossible}

No non-negative matrix or dilation realizes the Ihara companion; no Hecke-dynamical
$C^*$-algebra carries $\Z\oplus\Jac(Y)$ as $K_0$ with the Ihara flow as gauge dynamics;
no separable $C^*$-algebra has $K_0\supseteq\mathfrak X_\infty$; no $\Z_p[[T]]$-module
produces $\lambda=-1$; no unital $*$-homomorphism $\mathcal O_{k+1}\to\mathcal O_k$
induces $\eta_*$; no cyclic cocycle on the commutative Hecke--Dirac triple computes the
$L$-invariant; no $L$-invariant exists for the Iwahori tower; no $\Z_p$-tower of thick
$\widetilde A_2$ complexes exists; no cubic trace obstruction exists; no invariant of the
unsigned chord-crossing graph computes the arrangement group of two-copy repeats, and no
pairwise invariant at all computes it for three-copy repeats; no coset structure
underlies long-read constraints in the reconstruction torsor.

Two entries are added here, and both were open problems on the previous list.

\emph{No odd cyclic class computes the graph $L$-invariant by a linear construction.}
The crossed product proposed for the job is Morita-trivial --- the deck action on a
connected $\Z$-cover is free, so $C_0(\widetilde Y_\alpha)\rtimes\Z\cong\bigoplus_V\mathcal
K$ and $K_1=0$ (Ch.~\ref{chap:XIX}, Prop.~\ref{XIX:prop:compacts}). In the Floquet
commutant, where $K_1$ is not zero, the natural element does not exist, because the pencil
degenerates exactly at the point one wants to evaluate at; what the odd pairing reads there
is the \emph{order} of the exceptional zero, which is two for every voltage class, whereas
$\kappa_0\|h_\alpha\|^2$ is its leading \emph{coefficient}
(Ch.~\ref{chap:XIX}, Prop.~\ref{XIX:prop:notinv}, Cor.~\ref{XIX:cor:winding}). And in any
algebra whose $K_1$ has rank one, a pairing linear in each slot is a product of two linear
forms and therefore vanishes on a subspace, while $\kappa_0\|h_\alpha\|^2$ is positive
definite; so no such construction can work once $b_1\ge2$
(Ch.~\ref{chap:XIX}, Thm.~\ref{XIX:thm:obstruction}). What is \emph{not} excluded is an
algebra with $\operatorname{rank}K_1\ge b_1$: the reduced group $C^*$-algebra of
$\pi_1(Y)$ is free of rank $b_1$ and sits exactly at that threshold.

\emph{The directed sandpile group of a reverse-complement de Bruijn graph has no covering
theory, and its signed quotient does not contract.} With $P$ the permutation matrix of the
involution, $P\Delta P=\Delta^{\mathsf T}$ and never $\Delta$: the deck involution is an
anti-automorphism of the Laplacian, so the $\pm1$ eigenlattices are not $\Delta$-stable and
the classical covering factorization is unavailable
(Ch.~\ref{chap:Bio7}, Thm.~\ref{Bio7:thm:anti}). The Reiner--Tseng sequence that does hold
computes an invariant of the undirected shadow, which by the matrix-tree theorem counts
spanning trees; the quantity of interest counts Eulerian circuits and lives on the digraph.
Separately, a signed $r$-bundle chain obeys the series law --- the presentation column at
the chain's endpoint differs from the contracted graph by $s\,r\,d_0$ --- so
$\K(\Sigma)\not\cong(\Z/r)^s\oplus\K(\Sigma/C)$, and the chain contraction that makes a
bacterial genome computable on one strand has no bidirected analogue
(Ch.~\ref{chap:Bio7}, Thm.~\ref{Bio7:thm:series}, Cor.~\ref{Bio7:cor:nocompute}).

\backmatter
\begin{thebibliography}{999}
\bibitem{Lyons2005} R.~Lyons, \emph{Asymptotic enumeration of spanning trees}, Combin.\ Probab.\ Comput.\ \textbf{14} (2005), no.~4, 491--522. [Primary-source verified: Cambridge UP, DOI 10.1017/S096354830500684X.]

\bibitem{Kesten1959} H.~Kesten, \emph{Symmetric random walks on groups}, Trans.\ Amer.\ Math.\ Soc.\ \textbf{92} (1959), no.~2, 336--354. [Crossref-verified: DOI 10.1090/S0002-9947-1959-0109367-6.]

\bibitem{McKay1981} B.~D.~McKay, \emph{The expected eigenvalue distribution of a large regular graph}, Linear Algebra Appl.\ \textbf{40} (1981), 203--216. [Crossref-verified: DOI 10.1016/0024-3795(81)90150-6.]

\bibitem{Williams} R.~F.~Williams, \emph{Classification of subshifts of finite type}, Ann.\ of Math.\ (2) \textbf{98} (1973), no.~1, 120--153; errata ibid.\ \textbf{99} (1974), no.~2, 380--381. [Main paper Crossref-verified (DOI 10.2307/1970908; Crossref records the start page only, the end page 153 is the standard record); errata primary-source verified: JSTOR, DOI 10.2307/1970903.]

\bibitem{Pimsner} M.~V.~Pimsner, \emph{A class of $C^*$-algebras generalizing both Cuntz--Krieger algebras and crossed products by $\Z$}, in: Free probability theory (Waterloo, ON, 1995), Fields Inst.\ Commun.\ \textbf{12}, Amer.\ Math.\ Soc., Providence, RI, 1997, pp.~189--212. [Crossref-verified: DOI 10.1090/fic/012/08, pp.~189--212; Crossref dates the chapter December 1996, the volume is conventionally cited 1997.]

\bibitem{Sunada} T.~Sunada, \emph{Unitary representations of fundamental groups and the spectrum of twisted Laplacians}, Topology \textbf{28} (1989), no.~2, 125--132. [Crossref-verified: DOI 10.1016/0040-9383(89)90015-3.]

\bibitem{Levine} L.~Levine, \emph{Sandpile groups and spanning trees of directed line graphs}, J.\ Combin.\ Theory Ser.\ A \textbf{118} (2011), no.~2, 350--364. [Crossref-verified: DOI 10.1016/j.jcta.2010.04.001.]

\bibitem{Vallieres} D.~Vallières, \emph{On abelian $\ell$-towers of multigraphs}, Ann.\ Math.\ Qu\'ebec \textbf{45} (2021), no.~2, 433--452. [Crossref-verified: DOI 10.1007/s40316-020-00152-4. Both earlier out-of-band page ranges (335--359 and 133--152) were incorrect.]

\bibitem{McGownVallieres} K.~J.~McGown and D.~Vallières, \emph{On abelian $\ell$-towers of multigraphs II}, Ann.\ Math.\ Qu\'ebec \textbf{47}, no.~2, 461--473. [Crossref-verified: DOI 10.1007/s40316-021-00183-5; Crossref carries an online date of November 2021 and an issue date of October 2023 for this record. Both earlier out-of-band venues (J.\ Number Theory \textbf{238} and Ann.\ Math.\ Qu\'ebec \textbf{46}) were incorrect. A third part in the series exists and is not verified here.]

\bibitem{Gonet} S.~R.~Gonet, \emph{Iwasawa theory of Jacobians of graphs}, Algebraic Combinatorics \textbf{5} (2022), no.~5, 827--848. [Crossref-verified: DOI 10.5802/alco.225. This supersedes the earlier entry entirely: the previously listed title (``Jacobians of finite and infinite voltage covers of graphs'') and venue (J.\ Algebraic Combin.) were incorrect in \emph{both} out-of-band passes --- note that Algebraic Combinatorics and J.\ Algebraic Combin.\ are distinct journals.]

\bibitem{CK} J.~Cuntz and W.~Krieger, \emph{A class of $C^*$-algebras and topological Markov chains}, Invent.\ Math.\ \textbf{56} (1980), 251--268; together with the standard graph-algebra $K$-theory and simplicity/pure-infiniteness criteria descending from it.

\bibitem{Powers} R.~T.~Powers, \emph{Representations of uniformly hyperfinite algebras and their associated von Neumann rings}, Ann.\ of Math.\ (2) \textbf{86} (1967), no.~1, 138--171. [Primary-source verified: JSTOR, DOI 10.2307/1970364.]

\bibitem{ArakiWoods} H.~Araki and E.~J.~Woods, \emph{A classification of factors}, Publ.\ Res.\ Inst.\ Math.\ Sci.\ \textbf{4} (1968), no.~1, 51--130. [Primary-source verified: EMS Press / PRIMS article page.]

\bibitem{Krieger} W.~Krieger, \emph{On ergodic flows and the isomorphism of factors}, Math.\ Ann.\ \textbf{223} (1976), 19--70. [Primary-source verified: Springer, Math.\ Ann.\ \textbf{223}, pp.~19--70, February 1976.]

\bibitem{Connes} A.~Connes, \emph{Classification of injective factors. Cases $\mathrm{II}_1$, $\mathrm{II}_\infty$, $\mathrm{III}_\lambda$, $\lambda\neq1$}, Ann.\ of Math.\ (2) \textbf{104} (1976), no.~1, 73--115. [Primary-source verified: JSTOR, DOI 10.2307/1971057.]

\bibitem{FeldmanMoore} J.~Feldman and C.~C.~Moore, \emph{Ergodic equivalence relations, cohomology, and von Neumann algebras.\ I, II}, Trans.\ Amer.\ Math.\ Soc.\ \textbf{234} (1977), 289--324 and 325--359. [Both parts verified: Part I Crossref-verified, DOI 10.1090/S0002-9947-1977-0578656-4, pp.~289--324; Part II primary-source verified, DOI 10.1090/S0002-9947-1977-0578730-2, pp.~325--359. The related 1975 Bulletin announcement, Bull.\ Amer.\ Math.\ Soc.\ \textbf{81} (1975), no.~5, 921--924, DOI 10.1090/S0002-9904-1975-13888-2, is a distinct paper and is not the reference intended here.]

\bibitem{LevineIII} L.~Levine, \emph{Sandpile groups and spanning trees of directed line graphs}, J.\ Combin.\ Theory Ser.\ A \textbf{118} (2011), no.~2, 350--364. [Crossref-verified: DOI 10.1016/j.jcta.2010.04.001.]

\bibitem{Washington} L.~C.~Washington, \emph{Introduction to Cyclotomic Fields}, 2nd ed., Grad.\ Texts in Math.\ \textbf{83}, Springer, 1997; Chapter~13 for the structure theory of $\La$-modules. [Bibliographic data primary-source verified via the Springer book page: author, title, second edition, GTM~\textbf{83}, \copyright~1997. The chapter-level pointer is a reading aid from recollection.]

\bibitem{Ba} H.~Bass, \emph{The Ihara--Selberg zeta function of a tree lattice}, Internat.\ J.\ Math.\ \textbf{3} (1992), 717--797.

\bibitem{GS} R.~Greenberg and G.~Stevens, \emph{$p$-adic $L$-functions and $p$-adic periods of modular forms}, Invent.\ Math.\ \textbf{111} (1993), 407--447. [DOI 10.1007/BF01231294.]

\bibitem{MTT} B.~Mazur, J.~Tate and J.~Teitelbaum, \emph{On $p$-adic analogues of the conjectures of Birch and Swinnerton-Dyer}, Invent.\ Math.\ \textbf{84} (1986), 1--48. [DOI 10.1007/BF01388731.]

\bibitem{ST} H.~M.~Stark and A.~A.~Terras, \emph{Zeta functions of finite graphs and coverings}, Adv.\ Math.\ \textbf{121} (1996), 124--165.

\bibitem{Wa} L.~C.~Washington, \emph{Introduction to Cyclotomic Fields}, 2nd ed., Grad.\ Texts in Math.\ \textbf{83}, Springer, 1997; Chapter~7 for Weierstrass preparation and Chapter~13 for $\La$-modules.

\bibitem{CC} A.~H.~Chamseddine and A.~Connes, \emph{The spectral action principle}, Comm.\ Math.\ Phys.\ \textbf{186} (1997), 731--750.

\bibitem{Co} A.~Connes, \emph{Noncommutative geometry}, Academic Press, San Diego, 1994.

\bibitem{JLO} A.~Jaffe, A.~Lesniewski and K.~Osterwalder, \emph{Quantum $K$-theory. I. The Chern character}, Comm.\ Math.\ Phys.\ \textbf{118} (1988), 1--14.

\bibitem{KS1} M.~Kotani and T.~Sunada, \emph{Albanese maps and off diagonal long time asymptotics for the heat kernel}, Comm.\ Math.\ Phys.\ \textbf{209} (2000), 633--670.

\bibitem{KS2} M.~Kotani and T.~Sunada, \emph{Standard realizations of crystal lattices via harmonic maps}, Trans.\ Amer.\ Math.\ Soc.\ \textbf{353} (2001), 1--20.

\bibitem{Lo} J.-L.~Loday, \emph{Cyclic homology}, 2nd ed., Grundlehren Math.\ Wiss.\ \textbf{301}, Springer, Berlin, 1998.

\bibitem{CFKSV} J.~Coates, T.~Fukaya, K.~Kato, R.~Sujatha and O.~Venjakob, \emph{The $\mathrm{GL}_2$ main conjecture for elliptic curves without complex multiplication}, Publ.\ Math.\ Inst.\ Hautes \'Etudes Sci.\ \textbf{101} (2005), 163--208. [Primary-source verified: the IH\'ES journal page, Volume 101 (2005), pp.~163--208, authors and title as listed.]

\bibitem{Co:XIII} P.~M.~Cohn, \emph{Free rings and their relations}, 2nd ed., London Math.\ Soc.\ Monographs \textbf{19}, Academic Press, London, 1985.

\bibitem{NR} A.~Neeman and A.~Ranicki, \emph{Noncommutative localisation in algebraic $K$-theory I}, Geom.\ Topol.\ \textbf{8} (2004), 1385--1425.

\bibitem{Sch} A.~H.~Schofield, \emph{Representations of rings over skew fields}, London Math.\ Soc.\ Lecture Note Ser.\ \textbf{92}, Cambridge Univ.\ Press, Cambridge, 1985.

\bibitem{KR} E.~Kirchberg and M.~R\o rdam, \emph{Non-simple purely infinite $C^*$-algebras}, Amer.\ J.\ Math.\ \textbf{122} (2000), 637--666.

\bibitem{Ph} N.~C.~Phillips, \emph{A classification theorem for nuclear purely infinite simple $C^*$-algebras}, Doc.\ Math.\ \textbf{5} (2000), 49--114.

\bibitem{R} I.~Raeburn, \emph{Graph algebras}, CBMS Regional Conference Series in Mathematics \textbf{103}, Amer.\ Math.\ Soc., Providence, RI, 2005.

\bibitem{Ro} M.~R\o rdam, \emph{Classification of nuclear, simple $C^*$-algebras}, in: Classification of nuclear $C^*$-algebras. Entropy in operator algebras, Encyclopaedia Math.\ Sci.\ \textbf{126}, Springer, Berlin, 2002, pp.~1--145.

\bibitem{RS} J.~Rosenberg and C.~Schochet, \emph{The K\"unneth theorem and the universal coefficient theorem for Kasparov's generalized $K$-functor}, Duke Math.\ J.\ \textbf{55} (1987), 431--474.

\bibitem{Bla} B.~Blackadar, \emph{$K$-theory for operator algebras}, 2nd ed., Math.\ Sci.\ Res.\ Inst.\ Publ.\ \textbf{5}, Cambridge Univ.\ Press, Cambridge, 1998.

\bibitem{Kas} G.~G.~Kasparov, \emph{The operator $K$-functor and extensions of $C^*$-algebras}, Math.\ USSR-Izv.\ \textbf{16} (1981), no.~3, 513--572.

\bibitem{KP} J.~Kaminker and I.~Putnam, \emph{$K$-theoretic duality for shifts of finite type}, Comm.\ Math.\ Phys.\ \textbf{187} (1997), 509--522.

\bibitem{CMSZ} D.~I.~Cartwright, A.~M.~Mantero, T.~Steger and A.~Zappa, \emph{Groups acting simply transitively on the vertices of a building of type $\widetilde A_2$, I}, Geom.\ Dedicata \textbf{47} (1993), 143--166.

\bibitem{DKM} A.~M.~Duval, C.~J.~Klivans and J.~L.~Martin, \emph{Critical groups of simplicial complexes}, Ann.\ Comb.\ \textbf{17} (2013), 53--70.

\bibitem{KL} M.-H.~Kang and W.-C.~W.~Li, \emph{Zeta functions of complexes arising from $\mathrm{PGL}(3)$}, Adv.\ Math.\ \textbf{256} (2014), 46--103.

\bibitem{KOR} K.~H.~Kim, N.~S.~Ormes and F.~W.~Roush, \emph{The spectra of nonnegative integer matrices via formal power series}, J.\ Amer.\ Math.\ Soc.\ \textbf{13} (2000), 773--806.

\bibitem{BH91} M.~Boyle and D.~Handelman, \emph{The spectra of nonnegative matrices via symbolic dynamics}, Ann.\ of Math.\ (2) \textbf{133} (1991), 249--316.

\bibitem{BH93} M.~Boyle and D.~Handelman, \emph{Algebraic shift equivalence and primitive matrices}, Trans.\ Amer.\ Math.\ Soc.\ \textbf{336} (1993), 121--149.

\bibitem{BS} M.~Boyle and S.~Schmieding, \emph{Strong shift equivalence and the generalized spectral conjecture for nonnegative matrices}, Linear Algebra Appl.\ \textbf{498} (2016), 231--243; arXiv:1501.04697.

\bibitem{BSw} W.~Ballmann and J.~\'Swi\k{a}tkowski, \emph{On $L^{2}$-cohomology and property (T) for automorphism groups of polyhedral cell complexes}, Geom.\ Funct.\ Anal.\ \textbf{7} (1997), 615--645.

\bibitem{Ga} H.~Garland, \emph{$p$-adic curvature and the cohomology of discrete subgroups of $p$-adic groups}, Ann.\ of Math.\ (2) \textbf{97} (1973), no.~3, 375--423. [Primary-source verified: JSTOR record, DOI 10.2307/1970829.]

\bibitem{KLW} M.-H.~Kang, W.-C.~W.~Li and C.-J.~Wang, \emph{The zeta functions of complexes from $\mathrm{PGL}(3)$: a representation-theoretic approach}, Israel J.\ Math.\ \textbf{177} (2010), 335--348.

\bibitem{LSV} A.~Lubotzky, B.~Samuels and U.~Vishne, \emph{Ramanujan complexes of type $\widetilde A_d$}, Israel J.\ Math.\ \textbf{149} (2005), 267--299.

\bibitem{KFH} A.~J.~Koll\'ar, M.~Fitzpatrick, and A.~A.~Houck, \emph{Hyperbolic lattices in circuit quantum electrodynamics}, Nature \textbf{571}, 45 (2019).

\bibitem{MR1} J.~Maciejko and S.~Rayan, \emph{Hyperbolic band theory}, Sci.\ Adv.\ \textbf{7}, eabe9170 (2021).

\bibitem{MR2} J.~Maciejko and S.~Rayan, \emph{Automorphic Bloch theorems for hyperbolic lattices}, Proc.\ Natl.\ Acad.\ Sci.\ USA \textbf{119}, e2116869119 (2022).

\bibitem{Leng22} P.~M.~Lenggenhager \emph{et al.}, \emph{Simulating hyperbolic space on a circuit board}, Nat.\ Commun.\ \textbf{13}, 4373 (2022).

\bibitem{Huang24} L.~Huang \emph{et al.}, \emph{Hyperbolic photonic topological insulators}, Nat.\ Commun.\ \textbf{15}, 1647 (2024).

\bibitem{Xu25} X.~Xu, A.~A.~Mahmoud, N.~Gorgichuk, R.~Thomale, S.~Rayan, and M.~Mariantoni, \emph{A scalable superconducting circuit framework for emulating physics in hyperbolic space}, arXiv:2510.23827 (2025).

\bibitem{KS3} M.~Kotani and T.~Sunada, \emph{Jacobian tori associated with a finite graph and its abelian covering graphs}, Adv.\ Appl.\ Math.\ \textbf{24}, 89 (2000).

\bibitem{BF} M.~Baker and X.~Faber, \emph{Metric properties of the tropical Abel--Jacobi map}, J.\ Algebr.\ Comb.\ \textbf{33}, 349 (2011).

\bibitem{KorS} E.~Korotyaev and N.~Saburova, \emph{Effective masses for Laplacians on periodic graphs}, J.\ Math.\ Anal.\ Appl.\ \textbf{436}, 104 (2016).

\bibitem{LSV:Phys1} A.~Lubotzky, B.~Samuels, and U.~Vishne, \emph{Ramanujan complexes of type $\widetilde A_d$}, Isr.\ J.\ Math.\ \textbf{149}, 267 (2005).

\bibitem{ANCFT9} R.~Chen, \emph{Algorithmic non-commutative class field theory, IX: exceptional zeros, the graph $\mathcal L$-invariant, and the operator-algebraic Mazur--Tate--Teitelbaum formula}, preprint (2026).

\bibitem{ANCFT10} R.~Chen, \emph{Algorithmic non-commutative class field theory, X: $\widetilde A_2$ complexes, higher simplicial Laplacians, and the fate of the trace obstruction in rank two}, preprint (2026).

\bibitem{ANCFT12} R.~Chen, \emph{Algorithmic non-commutative class field theory, XII: $\widetilde A_2$ towers, no abelian sector, the covering norm theorem, and the two digraph towers}, preprint (2026).

\bibitem{ANCFT14} R.~Chen, \emph{Algorithmic non-commutative class field theory, XIV: the graph $\mathcal L$-invariant as spectral action, Bloch effective mass and Albanese length, and what the cyclic Chern character sees}, preprint (2026).

\bibitem{BHN} R.~Bacher, P.~de~la~Harpe and T.~Nagnibeda, \emph{The lattice of integral flows and the lattice of integral cuts on a finite graph}, Bull.\ Soc.\ Math.\ France \textbf{125} (1997), 167--198.

\bibitem{Biggs} N.~Biggs, \emph{Chip-firing and the critical group of a graph}, J.\ Algebraic Combin.\ \textbf{9} (1999), 25--45.

\bibitem{Bowen14} L.~Bowen, \emph{The type and stable type of the boundary of a Gromov hyperbolic group}, Geom.\ Dedicata \textbf{172} (2014), 363--386.

\bibitem{BDF87} D.~Buchholz, C.~D'Antoni and K.~Fredenhagen, \emph{The universal structure of local algebras}, Comm.\ Math.\ Phys.\ \textbf{111} (1987), 123--135.

\bibitem{CPW23} V.~Chandrasekaran, G.~Penington and E.~Witten, \emph{Large $N$ algebras and generalized entropy}, JHEP \textbf{04} (2023), 009.

\bibitem{CLPW23} V.~Chandrasekaran, R.~Longo, G.~Penington and E.~Witten, \emph{An algebra of observables for de Sitter space}, JHEP \textbf{02} (2023), 082.

\bibitem{ANCFT} R.~Chen, \emph{Algorithmic non-commutative class field theory, I--XVI}, preprints (2026); here III (the line-digraph recursion and the Euler term), IV (the limit module), IX (the defects), X--XII (rank two).

\bibitem{CLM08} G.~Cornelissen, O.~Lorscheid and M.~Marcolli, \emph{On the $K$-theory of graph $C^*$-algebras}, Acta Appl.\ Math.\ \textbf{102} (2008), 57--69.

\bibitem{CM13} G.~Cornelissen and M.~Marcolli, \emph{Graph reconstruction and quantum statistical mechanics}, J.\ Geom.\ Phys.\ \textbf{72} (2013), 110--117.

\bibitem{Cuntz81} J.~Cuntz, \emph{A class of $C^*$-algebras and topological Markov chains II: reducible chains and the Ext-functor for $C^*$-algebras}, Invent.\ Math.\ \textbf{63} (1981), 23--50.

\bibitem{CR03} P.~Cutting and G.~Robertson, \emph{Type III actions on boundaries of $\widetilde A_n$ buildings}, J.\ Operator Theory \textbf{49} (2003), 25--44.

\bibitem{FLMO23} K.~Furuya, N.~Lashkari, M.~Moosa and S.~Ouseph, \emph{Information loss, mixing and emergent type $\mathrm{III}_1$ factors}, JHEP \textbf{08} (2023), 111.

\bibitem{GMP22} E.~Gesteau, M.~Marcolli and S.~Parikh, \emph{Holographic tensor networks from hyperbolic buildings}, JHEP \textbf{10} (2022), 169.

\bibitem{GLS22} P.~Gorantla, H.~T.~Lam and S.-H.~Shao, \emph{Fractons on graphs and complexity}, Phys.\ Rev.\ B \textbf{106} (2022), 195139.

\bibitem{GLSS23} P.~Gorantla, H.~T.~Lam, N.~Seiberg and S.-H.~Shao, \emph{Gapped lineon and fracton models on graphs}, Phys.\ Rev.\ B \textbf{107} (2023), 125121.

\bibitem{GKPSW} S.~S.~Gubser, J.~Knaute, S.~Parikh, A.~Samberg and P.~Witaszczyk, \emph{$p$-adic AdS/CFT}, Comm.\ Math.\ Phys.\ \textbf{352} (2017), 1019--1059.

\bibitem{GubserParikh} S.~S.~Gubser and S.~Parikh, \emph{Geodesic bulk diagrams on the Bruhat--Tits tree}, Phys.\ Rev.\ D \textbf{96} (2017), 066024.

\bibitem{HMSS} M.~Heydeman, M.~Marcolli, I.~Saberi and B.~Stoica, \emph{Tensor networks, $p$-adic fields, and algebraic curves: arithmetic and the $\mathrm{AdS}_3/\mathrm{CFT}_2$ correspondence}, Adv.\ Theor.\ Math.\ Phys.\ \textbf{22} (2018), 93--176.

\bibitem{HuangJepsen} A.~Huang and C.~B.~Jepsen, \emph{Finite temperature at finite places}, JHEP \textbf{03} (2025), 097.

\bibitem{INO08} M.~Izumi, S.~Neshveyev and R.~Okayasu, \emph{The ratio set of the harmonic measure of a random walk on a hyperbolic group}, Israel J.\ Math.\ \textbf{163} (2008), 285--316.

\bibitem{Krieger76} W.~Krieger, \emph{On ergodic flows and the isomorphism of factors}, Math.\ Ann.\ \textbf{223} (1976), 19--70.

\bibitem{LL1} S.~Leutheusser and H.~Liu, \emph{Causal connectability between quantum systems and the black hole interior in holographic duality}, Phys.\ Rev.\ D \textbf{108} (2023), 086019.

\bibitem{LL2} S.~Leutheusser and H.~Liu, \emph{Emergent times in holographic duality}, Phys.\ Rev.\ D \textbf{108} (2023), 086020.

\bibitem{Marcolli20} M.~Marcolli, \emph{Holographic codes on Bruhat--Tits buildings and Drinfeld symmetric spaces}, Pure Appl.\ Math.\ Q.\ \textbf{16} (2020), 1--33.

\bibitem{Mondal} A.~Mondal, S.~Parikh, P.~Pradhan and R.~Sengar, \emph{Toward holography on biregular trees}, Phys.\ Rev.\ D \textbf{112} (2025), 086011.

\bibitem{RR97} J.~Ramagge and G.~Robertson, \emph{Factors from trees}, Proc.\ Amer.\ Math.\ Soc.\ \textbf{125} (1997), 2051--2055.

\bibitem{Rob05} G.~Robertson, \emph{Boundary operator algebras for free uniform tree lattices}, Houston J.\ Math.\ \textbf{31} (2005), 913--935.

\bibitem{Wiesbrock93} H.-W.~Wiesbrock, \emph{Half-sided modular inclusions of von Neumann algebras}, Comm.\ Math.\ Phys.\ \textbf{157} (1993), 83--92.

\bibitem{Witten22} E.~Witten, \emph{Gravity and the crossed product}, JHEP \textbf{10} (2022), 008.

\bibitem{BH09} J.~S.~Baras and P.~Hovareshti, \emph{Efficient and robust communication topologies for distributed decision making in networked systems}, Proc.\ 48th IEEE Conf.\ Decision and Control (2009).

\bibitem{BK11} H.~Bidkhori and S.~Kishore, \emph{Counting the spanning trees of a directed line graph}, Combin.\ Probab.\ Comput.\ \textbf{20} (2011), 11--25.

\bibitem{Boesch86} F.~T.~Boesch, \emph{On unreliability polynomials and graph connectivity in reliable network synthesis}, J.\ Graph Theory \textbf{10} (1986), 339--352.

\bibitem{CHP15} S.~H.~Chan, H.~D.~L.~Hollmann and D.~V.~Pasechnik, \emph{Sandpile groups of generalized de Bruijn and Kautz graphs and circulant matrices over finite fields}, J.\ Algebra \textbf{421} (2015), 268--295.

\bibitem{Colbourn87} C.~J.~Colbourn, \emph{The Combinatorics of Network Reliability}, Oxford University Press, 1987.

\bibitem{CM81} A.~A.~Cuoco and P.~Monsky, \emph{Class numbers in $\Z_p^{d}$-extensions}, Math.\ Ann.\ \textbf{255} (1981), 235--258.

\bibitem{DV22} S.~DuBose and D.~Vall\`ieres, \emph{On $\Z_\ell^{d}$-towers of graphs}, preprint, arXiv:2207.01711 (2022).

\bibitem{FAY83} M.~A.~Fiol, I.~Alegre and J.~L.~A.~Yebra, \emph{Line digraph iterations and the $(d,k)$ problem for directed graphs}, Proc.\ 10th Int.\ Symp.\ Computer Architecture (1983), 174--177.

\bibitem{FYA84} M.~A.~Fiol, J.~L.~A.~Yebra and I.~Alegre, \emph{Line digraph iterations and the $(d,k)$ digraph problem}, IEEE Trans.\ Comput.\ \textbf{C-33} (1984), 400--403.

\bibitem{FL92} M.~A.~Fiol and A.~S.~Llad\'o, \emph{The partial line digraph technique in the design of large interconnection networks}, IEEE Trans.\ Comput.\ \textbf{41} (1992), 848--857.

\bibitem{Gonet22} S.~R.~Gonet, \emph{Iwasawa theory of Jacobians of graphs}, Algebraic Combinatorics \textbf{5} (2022), 827--848.

\bibitem{Kataoka26} T.~Kataoka, \emph{Construction of graph coverings with prescribed Iwasawa invariants}, preprint, arXiv:2603.23088 (2026).

\bibitem{KM24} S.~Kleine and K.~M\"uller, \emph{On the growth of the Jacobian in $\Z_p^{l}$-voltage covers of graphs}, Algebraic Combinatorics \textbf{7} (2024), no.~4.

\bibitem{Knuth67} D.~E.~Knuth, \emph{Oriented subtrees of an arc digraph}, J.\ Combin.\ Theory \textbf{3} (1967), 309--314.

\bibitem{Levine11} L.~Levine, \emph{Sandpile groups and spanning trees of directed line graphs}, J.\ Combin.\ Theory Ser.~A \textbf{118} (2011), 350--364.

\bibitem{MV24} K.~McGown and D.~Vall\`ieres, \emph{On abelian $\ell$-towers of multigraphs III}, Ann.\ Math.\ Qu\'ebec \textbf{48} (2024), 1--19.

\bibitem{SHM05} C.~Song, S.~Havlin and H.~A.~Makse, \emph{Self-similarity of complex networks}, Nature \textbf{433} (2005), 392--395.

\bibitem{SHM06} C.~Song, S.~Havlin and H.~A.~Makse, \emph{Origins of fractality in the growth of complex networks}, Nature Physics \textbf{2} (2006), 275--281.

\bibitem{Val21} D.~Vall\`ieres, \emph{On abelian $\ell$-towers of multigraphs}, Ann.\ Math.\ Qu\'ebec \textbf{45} (2021), 433--452.

\bibitem{ADFB96} B.~D.~O.~Anderson, M.~Deistler, L.~Farina and L.~Benvenuti, \emph{Nonnegative realization of a linear system with nonnegative impulse response}, IEEE Trans.\ Circuits Syst.\ I \textbf{43} (1996), 134--142.

\bibitem{Ben20ELA} L.~Benvenuti, \emph{The NIEP and the positive realization problem}, Electron.\ J.\ Linear Algebra \textbf{36} (2020), 367--384.

\bibitem{Ben20SCL} L.~Benvenuti, \emph{A lower bound on the dimension of minimal positive realizations for discrete time systems}, Systems Control Lett.\ \textbf{135} (2020), 104595.

\bibitem{BF99} L.~Benvenuti and L.~Farina, \emph{An example of how positivity may force realizations of ``large'' dimension}, Systems Control Lett.\ \textbf{36} (1999), 261--266.

\bibitem{BF04} L.~Benvenuti and L.~Farina, \emph{A tutorial on the positive realization problem}, IEEE Trans.\ Automat.\ Control \textbf{49} (2004), 651--664.

\bibitem{Farina96} L.~Farina, \emph{On the existence of a positive realization}, Systems Control Lett.\ \textbf{28} (1996), 219--226.

\bibitem{KOR00} K.~H.~Kim, N.~S.~Ormes and F.~W.~Roush, \emph{The spectra of nonnegative integer matrices via formal power series}, J.\ Amer.\ Math.\ Soc.\ \textbf{13} (2000), 773--806.

\bibitem{LM95} D.~Lind and B.~Marcus, \emph{An Introduction to Symbolic Dynamics and Coding}, Cambridge University Press, 1995.

\bibitem{LoewyLondon78} R.~Loewy and D.~London, \emph{A note on an inverse problem for nonnegative matrices}, Linear Multilinear Algebra \textbf{6} (1978), 83--90.

\bibitem{NM03} B.~Nagy and M.~Matolcsi, \emph{A lowerbound on the dimension of positive realizations}, IEEE Trans.\ Circuits Syst.\ I \textbf{50} (2003), 782--784.

\bibitem{OMK84} Y.~Ohta, H.~Maeda and S.~Kodama, \emph{Reachability, observability and realizability of continuous-time positive systems}, SIAM J.\ Control Optim.\ \textbf{22} (1984), 171--180.

\bibitem{TS25} H.~Taghavian and J.~Sj\"olund, \emph{Minimal positive Markov realizations}, preprint, arXiv:2502.21102 (2025).

\bibitem{AEdB} T.~van Aardenne-Ehrenfest and N.~G.~de Bruijn, \emph{Circuits and trees in oriented linear graphs}, Simon Stevin \textbf{28} (1951), 203--217.

\bibitem{BBT} G.~Bresler, M.~Bresler and D.~Tse, \emph{Optimal assembly for high throughput shotgun sequencing}, BMC Bioinformatics \textbf{14} (Suppl.~5) (2013), S18.

\bibitem{BidkhoriKishore} H.~Bidkhori and S.~Kishore, \emph{Counting the spanning trees of a directed line graph}, Combin.\ Probab.\ Comput.\ \textbf{20} (2011), 11--25.

\bibitem{CHP} S.~H.~Chan, H.~D.~L.~Hollmann and D.~V.~Pasechnik, \emph{Sandpile groups of generalized de Bruijn and Kautz graphs and circulant matrices over finite fields}, J.\ Algebra \textbf{421} (2015), 268--295.

\bibitem{Bio2} R.~Chen, \emph{The sandpile group as a coordinate system on genome reconstructions}, companion preprint (Bio~2), 2026.

\bibitem{KSP} C.~Kingsford, M.~C.~Schatz and M.~Pop, \emph{Assembly complexity of prokaryotic genomes using short reads}, BMC Bioinformatics \textbf{11} (2010), 21.

\bibitem{Knuth} D.~E.~Knuth, \emph{Oriented subtrees of an arc digraph}, J.\ Combin.\ Theory \textbf{3} (1967), 309--314.

\bibitem{Loukides} G.~Loukides, S.~P.~Pissis et al., \emph{Beyond the BEST theorem: fast assessment of Eulerian trails}, in: Fundamentals of Computation Theory (FCT 2021), Lecture Notes in Comput.\ Sci.\ \textbf{12867}, Springer, 2021, pp.~162--175.

\bibitem{NagarajanPop} N.~Nagarajan and M.~Pop, \emph{Parametric complexity of sequence assembly: theory and applications to next generation sequencing}, J.\ Comput.\ Biol.\ \textbf{16} (2009), 897--908.

\bibitem{ObscuraTomescu} N.~Obscura Acosta and A.~I.~Tomescu, \emph{Simplicity in Eulerian circuits: uniqueness and safety}, preprint, arXiv:2208.08522 (2022).

\bibitem{RT} V.~Reiner and D.~Tseng, \emph{Critical groups of covering, voltage and signed graphs}, Discrete Math.\ \textbf{318} (2014), 10--40.

\bibitem{SmithTutte} C.~A.~B.~Smith and W.~T.~Tutte, \emph{On unicursal paths in a network of degree 4}, Amer.\ Math.\ Monthly \textbf{48} (1941), 233--237.

\bibitem{Zaslavsky} T.~Zaslavsky, \emph{Signed graphs}, Discrete Appl.\ Math.\ \textbf{4} (1982), 47--74.

\bibitem{AE} S.~F.~Altschul and B.~W.~Erickson, \emph{Significance of nucleotide sequence alignments: a method for random sequence permutation that preserves dinucleotide and codon usage}, Mol.\ Biol.\ Evol.\ \textbf{2} (1985), 526--538.

\bibitem{ABCS} R.~Arratia, B.~Bollob\'as, D.~Coppersmith and G.~B.~Sorkin, \emph{Euler circuits and DNA sequencing by hybridization}, Discrete Appl.\ Math.\ \textbf{104} (2000), 63--96.

\bibitem{BK} H.~Bidkhori and S.~Kishore, \emph{Counting the spanning trees of a directed line graph}, Combin.\ Probab.\ Comput.\ \textbf{20} (2011), 11--25.

\bibitem{BW} G.~R.~Brightwell and P.~Winkler, \emph{Counting Eulerian circuits is $\#$P-complete}, in: Proc.\ ALENEX/ANALCO 2005, SIAM, 2005, pp.~259--262.

\bibitem{Bio1} R.~Chen, \emph{Assembly ambiguity is not a function of branch statistics: the sandpile group of a de Bruijn graph}, companion preprint (Bio~1, version~2), 2026.

\bibitem{CMN} C.~J.~Colbourn, W.~J.~Myrvold and E.~Neufeld, \emph{Two algorithms for unranking arborescences}, J.\ Algorithms \textbf{20} (1996), 268--281.

\bibitem{CC:Bio2} P.~Creed and M.~Cryan, \emph{The number of Euler tours of random directed graphs}, Electron.\ J.\ Combin.\ \textbf{20} (2013), P13.

\bibitem{HLMPPW} A.~E.~Holroyd, L.~Levine, K.~M\'esz\'aros, Y.~Peres, J.~Propp and D.~B.~Wilson, \emph{Chip-firing and rotor-routing on directed graphs}, in: In and Out of Equilibrium 2, Progress in Probability \textbf{60}, Birkh\"auser, 2008, pp.~331--364.

\bibitem{uShuffle} M.~Jiang, J.~Anderson, J.~Gillespie and M.~Mayne, \emph{uShuffle: a useful tool for shuffling biological sequences while preserving the $k$-let counts}, BMC Bioinformatics \textbf{9} (2008), 192.

\bibitem{KMUW} D.~Kandel, Y.~Matias, R.~Unger and P.~Winkler, \emph{Shuffling biological sequences}, Discrete Appl.\ Math.\ \textbf{71} (1996), 171--185.

\bibitem{KM} J.~D.~Kececioglu and E.~W.~Myers, \emph{Combinatorial algorithms for DNA sequence assembly}, Algorithmica \textbf{13} (1995), 7--51.

\bibitem{MB} P.~Medvedev and M.~Brudno, \emph{Maximum likelihood genome assembly}, J.\ Comput.\ Biol.\ \textbf{16} (2009), 1101--1116.

\bibitem{PDDK} V.~B.~Priezzhev, D.~Dhar, A.~Dhar and S.~Krishnamurthy, \emph{Eulerian walkers as a model of self-organized criticality}, Phys.\ Rev.\ Lett.\ \textbf{77} (1996), 5079--5082.

\bibitem{PW} J.~G.~Propp and D.~B.~Wilson, \emph{How to get a perfectly random sample from a generic Markov chain and generate a random spanning tree of a directed graph}, J.\ Algorithms \textbf{27} (1998), 170--217.

\bibitem{Wilson} D.~B.~Wilson, \emph{Generating random spanning trees more quickly than the cover time}, in: Proc.\ 28th ACM Symp.\ Theory of Computing (1996), pp.~296--303.

\bibitem{Alar} D.~Alar, M.~Celaya, L.~D.~Garc\'ia Puente, S.~Henson and A.~Wheeler, \emph{The sandpile group of a thick cycle graph}, Electron.\ J.\ Graph Theory Appl.\ \textbf{10} (2022), no.~2.

\bibitem{ChenSchedler} W.~Chen and T.~Schedler, \emph{Concrete and abstract structure of the sandpile group for thick trees with loops}, preprint, arXiv:math/0701381 (2007).

\bibitem{BCALM} R.~Chikhi, A.~Limasset and P.~Medvedev, \emph{Compacting de Bruijn graphs from sequencing data quickly and in low memory}, Bioinformatics \textbf{32} (2016), i201--i208.

\bibitem{Flye} M.~Kolmogorov, J.~Yuan, Y.~Lin and P.~A.~Pevzner, \emph{Assembly of long, error-prone reads using repeat graphs}, Nat.\ Biotechnol.\ \textbf{37} (2019), 540--546.

\bibitem{ABS} R.~Arratia, B.~Bollob\'as and G.~B.~Sorkin, \emph{The interlace polynomial of a graph}, J.\ Combin.\ Theory Ser.~B \textbf{92} (2004), 199--233.

\bibitem{LasVergnas} M.~Las Vergnas, \emph{Le polyn\^ome de Martin d'un graphe eul\'erien}, Ann.\ Discrete Math.\ \textbf{17} (1983), 397--411.

\bibitem{MMN} C.~Merino, I.~Moffatt and S.~Noble, \emph{The critical group of a combinatorial map}, Combinatorial Theory \textbf{5} (2025), no.~3; arXiv:2308.13342.

\bibitem{MacrisPule} N.~Macris and J.~V.~Pul\'e, \emph{An alternative formula for the number of Euler trails for a class of digraphs}, Discrete Math.\ \textbf{154} (1996), 301--305.

\bibitem{Jonsson} J.~Jonsson, \emph{On the number of Euler trails in directed graphs}, Math.\ Scand.\ \textbf{90} (2002), 191--214.

\bibitem{PSX} D.~Perkinson, N.~Salter and T.~Xu, \emph{A note on the critical group of a line graph}, Electron.\ J.\ Combin.\ \textbf{18} (2011), P124.

\bibitem{PTT} P.~A.~Pevzner, H.~Tang and G.~Tesler, \emph{De novo repeat classification and fragment assembly}, Genome Res.\ \textbf{14} (2004), 1786--1796.

\bibitem{PTW} P.~A.~Pevzner, H.~Tang and M.~S.~Waterman, \emph{An Eulerian path approach to DNA fragment assembly}, Proc.\ Natl.\ Acad.\ Sci.\ USA \textbf{98} (2001), 9748--9753.

\bibitem{Bio3} R.~Chen, \emph{The sandpile group of a de Bruijn graph splits along its repeats}, companion preprint (Bio~3), 2026.

\bibitem{Lauri} J.~Lauri, \emph{On a determinant formula for enumerating Euler trails in a class of digraphs}, Rend.\ Sem.\ Mat.\ Messina Ser.~II \textbf{5} (1999), 75--90.

\bibitem{Toth} L.~T\'othm\'er\'esz, \emph{On canonical sandpile actions of embedded graphs}, Advances in Combinatorics (2026), no.~6; arXiv:2504.05275.

\end{thebibliography}

\end{document}
